% This file was converted to LaTeX by Writer2LaTeX ver. 1.4
% see http://writer2latex.sourceforge.net for more info
\documentclass[10pt]{article}
\usepackage[utf8]{inputenc}
\usepackage[T2A,T1]{fontenc}
\usepackage[russian,ngerman]{babel}
\usepackage{amsmath}
\usepackage{amssymb,amsfonts,textcomp}
\usepackage{array}
\usepackage{supertabular}
\usepackage{hhline}
\usepackage{hyperref}
\hypersetup{colorlinks=true, linkcolor=blue, citecolor=blue, filecolor=blue, urlcolor=blue}
\usepackage{graphicx}
% footnotes configuration
\makeatletter
\renewcommand\thefootnote{\arabic{footnote}}
\renewcommand\@makefnmark{\mbox{\textstyleFootnoteanchor{\@thefnmark}}}
\makeatother
\newcommand\textsubscript[1]{\ensuremath{{}_{\text{#1}}}}
% Text styles
\newcommand\textstyleFootnoteSymbol[1]{\textsuperscript{#1}}
\newcommand\textstyleannotationreference[1]{{\fontsize{8pt}{9.6pt}\selectfont #1}}
\newcommand\textstylepostdetails[1]{#1}
\newcommand\textstyleListLabelclv[1]{#1}
\newcommand\textstyleListLabelclvi[1]{#1}
\newcommand\textstyleListLabelclvii[1]{#1}
\newcommand\textstyleListLabelclviii[1]{#1}
\newcommand\textstyleListLabellxxxii[1]{{\fontsize{11pt}{13.2pt}\selectfont \textmd{#1}}}
\newcommand\textstyleListLabellxxxiii[1]{#1}
\newcommand\textstyleListLabellxxxiv[1]{#1}
\newcommand\textstyleListLabellxxxv[1]{#1}
\newcommand\textstyleListLabelcxci[1]{{\fontsize{10pt}{12.0pt}\selectfont \textrm{#1}}}
\newcommand\textstyleListLabelcxcii[1]{#1}
\newcommand\textstyleListLabelcxciii[1]{#1}
\newcommand\textstyleListLabelcxciv[1]{#1}
\newcommand\textstyleListLabelc[1]{#1}
\newcommand\textstyleListLabelci[1]{#1}
\newcommand\textstyleListLabelcii[1]{#1}
\newcommand\textstyleListLabelciii[1]{#1}
\newcommand\textstyleInternetlink[1]{#1}
\newcommand\textstylehgkelc[1]{#1}
\newcommand\textstyleListLabelcxxvii[1]{#1}
\newcommand\textstyleListLabelcxxviii[1]{#1}
\newcommand\textstyleListLabelcxxix[1]{#1}
\newcommand\textstyleListLabelcxxx[1]{#1}
\newcommand\textstyleListLabelclxxxii[1]{#1}
\newcommand\textstyleListLabelclxxxiii[1]{#1}
\newcommand\textstyleListLabelclxxxiv[1]{#1}
\newcommand\textstyleListLabelclxxxv[1]{#1}
\newcommand\textstylest[1]{#1}
\newcommand\textstyleacopre[1]{#1}
\newcommand\textstylehighlight[1]{#1}
\newcommand\textstyleFootnoteanchor[1]{\textsuperscript{#1}}
% Headings and outline numbering
\makeatletter
\renewcommand\section{\@startsection{section}{1}{0.0cm}{0.5in}{0.3335in}{\normalfont\normalsize\fontsize{12pt}{14.4pt}\selectfont\bfseries\raggedright}}
\renewcommand\subsection{\@startsection{subsection}{2}{0.0cm}{0.25in}{0.1665in}{\normalfont\normalsize\fontsize{12pt}{14.4pt}\selectfont\bfseries\raggedright}}
\renewcommand\subsubsection{\@startsection{subsubsection}{3}{0.0cm}{0in}{0.1mm}{\normalfont\normalsize\fontsize{12pt}{14.4pt}\selectfont}}
\renewcommand\paragraph{\@startsection{paragraph}{4}{0.0cm}{0in}{0.1mm}{\normalfont\normalsize\fontsize{12pt}{14.4pt}\selectfont}}
\renewcommand\@seccntformat[1]{\csname @textstyle#1\endcsname{\csname the#1\endcsname}\csname @distance#1\endcsname}
\setcounter{secnumdepth}{0}
\newcommand\@distancesection{}
\newcommand\@textstylesection[1]{#1}
\newcommand\@distancesubsection{}
\newcommand\@textstylesubsection[1]{#1}
\newcommand\@distancesubsubsection{}
\newcommand\@textstylesubsubsection[1]{#1}
\newcommand\@distanceparagraph{}
\newcommand\@textstyleparagraph[1]{#1}
\makeatother
\makeatletter
\newcommand\arraybslash{\let\\\@arraycr}
\makeatother
\raggedbottom
% Paragraph styles
\renewcommand\familydefault{\rmdefault}
\newenvironment{styleStandard}{\renewcommand\baselinestretch{1.25}\setlength\leftskip{0cm}\setlength\rightskip{0cm}\setlength\parindent{0cm}\setlength\parfillskip{0pt plus 1fil}\setlength\parskip{0in plus 1pt}\writerlistparindent\writerlistleftskip\leavevmode\normalfont\normalsize\fontsize{12pt}{14.4pt}\selectfont\writerlistlabel\ignorespaces}{\unskip\vspace{0in plus 1pt}\par}
\newenvironment{styleContentsHeading}{\renewcommand\baselinestretch{1.0}\setlength\leftskip{0cm}\setlength\rightskip{0cm plus 1fil}\setlength\parindent{0cm}\setlength\parfillskip{0pt plus 1fil}\setlength\parskip{0.1252in plus 0.012520001in}\writerlistparindent\writerlistleftskip\leavevmode\normalfont\normalsize\fontsize{12pt}{14.4pt}\selectfont\bfseries\writerlistlabel\ignorespaces}{\unskip\vspace{0in plus 1pt}\par}
\newenvironment{styleelenconumeratoii}{\renewcommand\baselinestretch{1.25}\setlength\leftskip{0.5in}\setlength\rightskip{0in plus 1fil}\setlength\parindent{-0.25in}\setlength\parfillskip{0pt plus 1fil}\setlength\parskip{0.1252in plus 0.012520001in}\writerlistparindent\writerlistleftskip\leavevmode\normalfont\normalsize\fontsize{12pt}{14.4pt}\selectfont\writerlistlabel\ignorespaces}{\unskip\vspace{0in plus 1pt}\par}
\newenvironment{styletablecaption}{\renewcommand\baselinestretch{1.0}\setlength\leftskip{0cm plus 1fil}\setlength\rightskip{0cm plus 1fil}\setlength\parindent{0cm}\setlength\parfillskip{0pt}\setlength\parskip{0.0835in plus 0.00835in}\writerlistparindent\writerlistleftskip\leavevmode\normalfont\normalsize\writerlistlabel\ignorespaces}{\unskip\vspace{0.0835in plus 0.00835in}\par}
\newenvironment{stylelsSourceline}{\setlength\leftskip{0.748in}\setlength\rightskip{0in plus 1fil}\setlength\parindent{0in}\setlength\parfillskip{0pt plus 1fil}\setlength\parskip{0in plus 1pt}\writerlistparindent\writerlistleftskip\leavevmode\normalfont\normalsize\fontsize{12pt}{14.4pt}\selectfont\itshape\writerlistlabel\ignorespaces}{\unskip\vspace{0in plus 1pt}\par}
\newenvironment{styleFootnote}{\renewcommand\baselinestretch{1.0}\setlength\leftskip{0cm}\setlength\rightskip{0cm}\setlength\parindent{0cm}\setlength\parfillskip{0pt plus 1fil}\setlength\parskip{0in plus 1pt}\writerlistparindent\writerlistleftskip\leavevmode\normalfont\normalsize\fontsize{12pt}{14.4pt}\selectfont\writerlistlabel\ignorespaces}{\unskip\vspace{0in plus 1pt}\par}
\newenvironment{styleTabellenInhalt}{\renewcommand\baselinestretch{1.25}\setlength\leftskip{0cm}\setlength\rightskip{0cm}\setlength\parindent{0cm}\setlength\parfillskip{0pt plus 1fil}\setlength\parskip{0in plus 1pt}\writerlistparindent\writerlistleftskip\leavevmode\normalfont\normalsize\fontsize{12pt}{14.4pt}\selectfont\writerlistlabel\ignorespaces}{\unskip\vspace{0in plus 1pt}\par}
\newenvironment{styleListParagraph}{\renewcommand\baselinestretch{1.0}\setlength\leftskip{0.5in}\setlength\rightskip{0in plus 1fil}\setlength\parindent{0in}\setlength\parfillskip{0pt plus 1fil}\setlength\parskip{0in plus 1pt}\writerlistparindent\writerlistleftskip\leavevmode\normalfont\normalsize\fontsize{11pt}{13.2pt}\selectfont\bfseries\writerlistlabel\ignorespaces}{\unskip\vspace{0in plus 1pt}\par}
\newenvironment{stylecaption}{\renewcommand\baselinestretch{1.0}\setlength\leftskip{0cm}\setlength\rightskip{0cm plus 1fil}\setlength\parindent{0cm}\setlength\parfillskip{0pt plus 1fil}\setlength\parskip{0.0835in plus 0.00835in}\writerlistparindent\writerlistleftskip\leavevmode\normalfont\normalsize\fontsize{11pt}{13.2pt}\selectfont\bfseries\writerlistlabel\ignorespaces}{\unskip\vspace{0.0835in plus 0.00835in}\par}
\newenvironment{styleAuthor}{\renewcommand\baselinestretch{1.0}\setlength\leftskip{0.5in}\setlength\rightskip{0in}\setlength\parindent{0in}\setlength\parfillskip{0pt plus 1fil}\setlength\parskip{0in plus 1pt}\writerlistparindent\writerlistleftskip\leavevmode\normalfont\normalsize\fontsize{12pt}{14.4pt}\selectfont\mdseries\writerlistlabel\ignorespaces}{\unskip\vspace{0in plus 1pt}\par}
\newenvironment{styleDLiteratur}{\renewcommand\baselinestretch{1.0}\setlength\leftskip{0.3937in}\setlength\rightskip{0in plus 1fil}\setlength\parindent{-0.3937in}\setlength\parfillskip{0pt plus 1fil}\setlength\parskip{0.0417in plus 0.00417in}\writerlistparindent\writerlistleftskip\leavevmode\normalfont\normalsize\fontsize{9pt}{10.8pt}\selectfont\writerlistlabel\ignorespaces}{\unskip\vspace{0in plus 1pt}\par}
\newenvironment{styleLiteraturStandard}{\setlength\leftskip{0.1972in}\setlength\rightskip{0in}\setlength\parindent{-0.1972in}\setlength\parfillskip{0pt plus 1fil}\setlength\parskip{0in plus 1pt}\writerlistparindent\writerlistleftskip\leavevmode\normalfont\normalsize\fontsize{8.5pt}{10.2pt}\selectfont\writerlistlabel\ignorespaces}{\unskip\vspace{0in plus 1pt}\par}
\newenvironment{styleannotationtext}{\renewcommand\baselinestretch{1.25}\setlength\leftskip{0cm}\setlength\rightskip{0cm}\setlength\parindent{0cm}\setlength\parfillskip{0pt plus 1fil}\setlength\parskip{0in plus 1pt}\writerlistparindent\writerlistleftskip\leavevmode\normalfont\normalsize\writerlistlabel\ignorespaces}{\unskip\vspace{0in plus 1pt}\par}
% List styles
\newcounter{saveenum}
\newcommand\writerlistleftskip{}
\newcommand\writerlistparindent{}
\newcommand\writerlistlabel{}
\newcommand\writerlistremovelabel{\aftergroup\let\aftergroup\writerlistparindent\aftergroup\relax\aftergroup\let\aftergroup\writerlistlabel\aftergroup\relax}
\newcounter{listWWNumxixleveli}
\newcounter{listWWNumxixlevelii}[listWWNumxixleveli]
\newcounter{listWWNumxixleveliii}[listWWNumxixlevelii]
\newcounter{listWWNumxixleveliv}[listWWNumxixleveliii]
\renewcommand\thelistWWNumxixleveli{\roman{listWWNumxixleveli}}
\renewcommand\thelistWWNumxixlevelii{\alph{listWWNumxixlevelii}}
\renewcommand\thelistWWNumxixleveliii{\roman{listWWNumxixleveliii}}
\renewcommand\thelistWWNumxixleveliv{\arabic{listWWNumxixleveliv}}
\newcommand\labellistWWNumxixleveli{\textstyleListLabelclv{\thelistWWNumxixleveli.}}
\newcommand\labellistWWNumxixlevelii{\textstyleListLabelclvi{\thelistWWNumxixlevelii.}}
\newcommand\labellistWWNumxixleveliii{\textstyleListLabelclvii{\thelistWWNumxixleveliii.}}
\newcommand\labellistWWNumxixleveliv{\textstyleListLabelclviii{\thelistWWNumxixleveliv.}}
\newenvironment{listWWNumxixleveli}{\def\writerlistleftskip{\setlength\leftskip{0.75in}}\def\writerlistparindent{}\def\writerlistlabel{}\def\item{\def\writerlistparindent{\setlength\parindent{-0.5in}}\def\writerlistlabel{\stepcounter{listWWNumxixleveli}\makebox[0cm][l]{\labellistWWNumxixleveli}\hspace{-0.635cm}\writerlistremovelabel}}}{}
\newenvironment{listWWNumxixlevelii}{\def\writerlistleftskip{\setlength\leftskip{1in}}\def\writerlistparindent{}\def\writerlistlabel{}\def\item{\def\writerlistparindent{\setlength\parindent{-0.25in}}\def\writerlistlabel{\stepcounter{listWWNumxixlevelii}\makebox[0cm][l]{\labellistWWNumxixlevelii}\hspace{-1.905cm}\writerlistremovelabel}}}{}
\newenvironment{listWWNumxixleveliii}{\def\writerlistleftskip{\setlength\leftskip{1.5in}}\def\writerlistparindent{}\def\writerlistlabel{}\def\item{\def\writerlistparindent{\setlength\parindent{-0.1252in}}\def\writerlistlabel{\stepcounter{listWWNumxixleveliii}\makebox[0cm][r]{\labellistWWNumxixleveliii}\hspace{-3.4919918cm}\writerlistremovelabel}}}{}
\newenvironment{listWWNumxixleveliv}{\def\writerlistleftskip{\setlength\leftskip{2in}}\def\writerlistparindent{}\def\writerlistlabel{}\def\item{\def\writerlistparindent{\setlength\parindent{-0.25in}}\def\writerlistlabel{\stepcounter{listWWNumxixleveliv}\makebox[0cm][l]{\labellistWWNumxixleveliv}\hspace{-4.4449997cm}\writerlistremovelabel}}}{}
\newcounter{listWWNumxleveli}
\newcounter{listWWNumxlevelii}[listWWNumxleveli]
\newcounter{listWWNumxleveliii}[listWWNumxlevelii]
\newcounter{listWWNumxleveliv}[listWWNumxleveliii]
\renewcommand\thelistWWNumxleveli{\arabic{listWWNumxleveli}}
\renewcommand\thelistWWNumxlevelii{\alph{listWWNumxlevelii}}
\renewcommand\thelistWWNumxleveliii{\arabic{listWWNumxleveliii}}
\renewcommand\thelistWWNumxleveliv{\arabic{listWWNumxleveliv}}
\newcommand\labellistWWNumxleveli{\textstyleListLabellxxxii{(\thelistWWNumxleveli)}}
\newcommand\labellistWWNumxlevelii{\textstyleListLabellxxxiii{\thelistWWNumxlevelii.}}
\newcommand\labellistWWNumxleveliii{\textstyleListLabellxxxiv{\thelistWWNumxleveliii.}}
\newcommand\labellistWWNumxleveliv{\textstyleListLabellxxxv{\thelistWWNumxleveliv.}}
\newenvironment{listWWNumxleveli}{\def\writerlistleftskip{\setlength\leftskip{0.3937in}}\def\writerlistparindent{}\def\writerlistlabel{}\def\item{\def\writerlistparindent{\setlength\parindent{-0.3543in}}\def\writerlistlabel{\stepcounter{listWWNumxleveli}\makebox[0cm][l]{\labellistWWNumxleveli}\hspace{-0.10007603cm}\writerlistremovelabel}}}{}
\newenvironment{listWWNumxlevelii}{\def\writerlistleftskip{\setlength\leftskip{1in}}\def\writerlistparindent{}\def\writerlistlabel{}\def\item{\def\writerlistparindent{\setlength\parindent{-0.25in}}\def\writerlistlabel{\stepcounter{listWWNumxlevelii}\makebox[0cm][l]{\labellistWWNumxlevelii}\hspace{-1.905cm}\writerlistremovelabel}}}{}
\newenvironment{listWWNumxleveliii}{\def\writerlistleftskip{\setlength\leftskip{1.6252in}}\def\writerlistparindent{}\def\writerlistlabel{}\def\item{\def\writerlistparindent{\setlength\parindent{-0.25in}}\def\writerlistlabel{\stepcounter{listWWNumxleveliii}\makebox[0cm][l]{\labellistWWNumxleveliii}\hspace{-3.4930081cm}\writerlistremovelabel}}}{}
\newenvironment{listWWNumxleveliv}{\def\writerlistleftskip{\setlength\leftskip{2in}}\def\writerlistparindent{}\def\writerlistlabel{}\def\item{\def\writerlistparindent{\setlength\parindent{-0.25in}}\def\writerlistlabel{\stepcounter{listWWNumxleveliv}\makebox[0cm][l]{\labellistWWNumxleveliv}\hspace{-4.4449997cm}\writerlistremovelabel}}}{}
\newcounter{listWWNumxxiiileveli}
\newcounter{listWWNumxxiiilevelii}[listWWNumxxiiileveli]
\newcounter{listWWNumxxiiileveliii}[listWWNumxxiiilevelii]
\newcounter{listWWNumxxiiileveliv}[listWWNumxxiiileveliii]
\renewcommand\thelistWWNumxxiiileveli{\roman{listWWNumxxiiileveli}}
\renewcommand\thelistWWNumxxiiilevelii{\alph{listWWNumxxiiilevelii}}
\renewcommand\thelistWWNumxxiiileveliii{\roman{listWWNumxxiiileveliii}}
\renewcommand\thelistWWNumxxiiileveliv{\arabic{listWWNumxxiiileveliv}}
\newcommand\labellistWWNumxxiiileveli{\textstyleListLabelcxci{\thelistWWNumxxiiileveli.}}
\newcommand\labellistWWNumxxiiilevelii{\textstyleListLabelcxcii{\thelistWWNumxxiiilevelii.}}
\newcommand\labellistWWNumxxiiileveliii{\textstyleListLabelcxciii{\thelistWWNumxxiiileveliii.}}
\newcommand\labellistWWNumxxiiileveliv{\textstyleListLabelcxciv{\thelistWWNumxxiiileveliv.}}
\newenvironment{listWWNumxxiiileveli}{\def\writerlistleftskip{\setlength\leftskip{0.75in}}\def\writerlistparindent{}\def\writerlistlabel{}\def\item{\def\writerlistparindent{\setlength\parindent{-0.5in}}\def\writerlistlabel{\stepcounter{listWWNumxxiiileveli}\makebox[0cm][l]{\labellistWWNumxxiiileveli}\hspace{-0.635cm}\writerlistremovelabel}}}{}
\newenvironment{listWWNumxxiiilevelii}{\def\writerlistleftskip{\setlength\leftskip{1in}}\def\writerlistparindent{}\def\writerlistlabel{}\def\item{\def\writerlistparindent{\setlength\parindent{-0.25in}}\def\writerlistlabel{\stepcounter{listWWNumxxiiilevelii}\makebox[0cm][l]{\labellistWWNumxxiiilevelii}\hspace{-1.905cm}\writerlistremovelabel}}}{}
\newenvironment{listWWNumxxiiileveliii}{\def\writerlistleftskip{\setlength\leftskip{1.5in}}\def\writerlistparindent{}\def\writerlistlabel{}\def\item{\def\writerlistparindent{\setlength\parindent{-0.1252in}}\def\writerlistlabel{\stepcounter{listWWNumxxiiileveliii}\makebox[0cm][r]{\labellistWWNumxxiiileveliii}\hspace{-3.4919918cm}\writerlistremovelabel}}}{}
\newenvironment{listWWNumxxiiileveliv}{\def\writerlistleftskip{\setlength\leftskip{2in}}\def\writerlistparindent{}\def\writerlistlabel{}\def\item{\def\writerlistparindent{\setlength\parindent{-0.25in}}\def\writerlistlabel{\stepcounter{listWWNumxxiiileveliv}\makebox[0cm][l]{\labellistWWNumxxiiileveliv}\hspace{-4.4449997cm}\writerlistremovelabel}}}{}
\newcounter{listWWNumxiileveli}
\newcounter{listWWNumxiilevelii}[listWWNumxiileveli]
\newcounter{listWWNumxiileveliii}[listWWNumxiilevelii]
\newcounter{listWWNumxiileveliv}[listWWNumxiileveliii]
\renewcommand\thelistWWNumxiileveli{\roman{listWWNumxiileveli}}
\renewcommand\thelistWWNumxiilevelii{\alph{listWWNumxiilevelii}}
\renewcommand\thelistWWNumxiileveliii{\roman{listWWNumxiileveliii}}
\renewcommand\thelistWWNumxiileveliv{\arabic{listWWNumxiileveliv}}
\newcommand\labellistWWNumxiileveli{\textstyleListLabelc{\thelistWWNumxiileveli.}}
\newcommand\labellistWWNumxiilevelii{\textstyleListLabelci{\thelistWWNumxiilevelii.}}
\newcommand\labellistWWNumxiileveliii{\textstyleListLabelcii{\thelistWWNumxiileveliii.}}
\newcommand\labellistWWNumxiileveliv{\textstyleListLabelciii{\thelistWWNumxiileveliv.}}
\newenvironment{listWWNumxiileveli}{\def\writerlistleftskip{\setlength\leftskip{0.75in}}\def\writerlistparindent{}\def\writerlistlabel{}\def\item{\def\writerlistparindent{\setlength\parindent{-0.5in}}\def\writerlistlabel{\stepcounter{listWWNumxiileveli}\makebox[0cm][l]{\labellistWWNumxiileveli}\hspace{-0.635cm}\writerlistremovelabel}}}{}
\newenvironment{listWWNumxiilevelii}{\def\writerlistleftskip{\setlength\leftskip{1in}}\def\writerlistparindent{}\def\writerlistlabel{}\def\item{\def\writerlistparindent{\setlength\parindent{-0.25in}}\def\writerlistlabel{\stepcounter{listWWNumxiilevelii}\makebox[0cm][l]{\labellistWWNumxiilevelii}\hspace{-1.905cm}\writerlistremovelabel}}}{}
\newenvironment{listWWNumxiileveliii}{\def\writerlistleftskip{\setlength\leftskip{1.5in}}\def\writerlistparindent{}\def\writerlistlabel{}\def\item{\def\writerlistparindent{\setlength\parindent{-0.1252in}}\def\writerlistlabel{\stepcounter{listWWNumxiileveliii}\makebox[0cm][r]{\labellistWWNumxiileveliii}\hspace{-3.4919918cm}\writerlistremovelabel}}}{}
\newenvironment{listWWNumxiileveliv}{\def\writerlistleftskip{\setlength\leftskip{2in}}\def\writerlistparindent{}\def\writerlistlabel{}\def\item{\def\writerlistparindent{\setlength\parindent{-0.25in}}\def\writerlistlabel{\stepcounter{listWWNumxiileveliv}\makebox[0cm][l]{\labellistWWNumxiileveliv}\hspace{-4.4449997cm}\writerlistremovelabel}}}{}
\newcommand\liststyleWWNumx{%
\renewcommand\theenumi{\arabic{enumi}}
\renewcommand\theenumii{\alph{enumii}}
\renewcommand\theenumiii{\arabic{enumiii}}
\renewcommand\theenumiv{\arabic{enumiv}}
\renewcommand\labelenumi{(\theenumi)}
\renewcommand\labelenumii{\theenumii.}
\renewcommand\labelenumiii{\theenumiii.}
\renewcommand\labelenumiv{\theenumiv.}
}
\newcommand\labellistWWNumxvleveli{\textstyleListLabelcxxvii{{}-}}
\newcommand\labellistWWNumxvlevelii{\textstyleListLabelcxxviii{o}}
\newcommand\labellistWWNumxvleveliii{\textstyleListLabelcxxix{[F0A7?]}}
\newcommand\labellistWWNumxvleveliv{\textstyleListLabelcxxx{[F0B7?]}}
\newenvironment{listWWNumxvleveli}{\def\writerlistleftskip{\setlength\leftskip{0.5in}}\def\writerlistparindent{}\def\writerlistlabel{}\def\item{\def\writerlistparindent{\setlength\parindent{-0.25in}}\def\writerlistlabel{\makebox[0cm][l]{\labellistWWNumxvleveli}\hspace{-0.635cm}\writerlistremovelabel}}}{}
\newenvironment{listWWNumxvlevelii}{\def\writerlistleftskip{\setlength\leftskip{1in}}\def\writerlistparindent{}\def\writerlistlabel{}\def\item{\def\writerlistparindent{\setlength\parindent{-0.25in}}\def\writerlistlabel{\makebox[0cm][l]{\labellistWWNumxvlevelii}\hspace{-1.905cm}\writerlistremovelabel}}}{}
\newenvironment{listWWNumxvleveliii}{\def\writerlistleftskip{\setlength\leftskip{1.5in}}\def\writerlistparindent{}\def\writerlistlabel{}\def\item{\def\writerlistparindent{\setlength\parindent{-0.25in}}\def\writerlistlabel{\makebox[0cm][l]{\labellistWWNumxvleveliii}\hspace{-3.175cm}\writerlistremovelabel}}}{}
\newenvironment{listWWNumxvleveliv}{\def\writerlistleftskip{\setlength\leftskip{2in}}\def\writerlistparindent{}\def\writerlistlabel{}\def\item{\def\writerlistparindent{\setlength\parindent{-0.25in}}\def\writerlistlabel{\makebox[0cm][l]{\labellistWWNumxvleveliv}\hspace{-4.4449997cm}\writerlistremovelabel}}}{}
\newcounter{listWWNumxxiileveli}
\newcounter{listWWNumxxiilevelii}[listWWNumxxiileveli]
\newcounter{listWWNumxxiileveliii}[listWWNumxxiilevelii]
\newcounter{listWWNumxxiileveliv}[listWWNumxxiileveliii]
\renewcommand\thelistWWNumxxiileveli{\roman{listWWNumxxiileveli}}
\renewcommand\thelistWWNumxxiilevelii{\alph{listWWNumxxiilevelii}}
\renewcommand\thelistWWNumxxiileveliii{\roman{listWWNumxxiileveliii}}
\renewcommand\thelistWWNumxxiileveliv{\arabic{listWWNumxxiileveliv}}
\newcommand\labellistWWNumxxiileveli{\textstyleListLabelclxxxii{\thelistWWNumxxiileveli.}}
\newcommand\labellistWWNumxxiilevelii{\textstyleListLabelclxxxiii{\thelistWWNumxxiilevelii.}}
\newcommand\labellistWWNumxxiileveliii{\textstyleListLabelclxxxiv{\thelistWWNumxxiileveliii.}}
\newcommand\labellistWWNumxxiileveliv{\textstyleListLabelclxxxv{\thelistWWNumxxiileveliv.}}
\newenvironment{listWWNumxxiileveli}{\def\writerlistleftskip{\setlength\leftskip{0.75in}}\def\writerlistparindent{}\def\writerlistlabel{}\def\item{\def\writerlistparindent{\setlength\parindent{-0.5in}}\def\writerlistlabel{\stepcounter{listWWNumxxiileveli}\makebox[0cm][l]{\labellistWWNumxxiileveli}\hspace{-0.635cm}\writerlistremovelabel}}}{}
\newenvironment{listWWNumxxiilevelii}{\def\writerlistleftskip{\setlength\leftskip{1in}}\def\writerlistparindent{}\def\writerlistlabel{}\def\item{\def\writerlistparindent{\setlength\parindent{-0.25in}}\def\writerlistlabel{\stepcounter{listWWNumxxiilevelii}\makebox[0cm][l]{\labellistWWNumxxiilevelii}\hspace{-1.905cm}\writerlistremovelabel}}}{}
\newenvironment{listWWNumxxiileveliii}{\def\writerlistleftskip{\setlength\leftskip{1.5in}}\def\writerlistparindent{}\def\writerlistlabel{}\def\item{\def\writerlistparindent{\setlength\parindent{-0.1252in}}\def\writerlistlabel{\stepcounter{listWWNumxxiileveliii}\makebox[0cm][r]{\labellistWWNumxxiileveliii}\hspace{-3.4919918cm}\writerlistremovelabel}}}{}
\newenvironment{listWWNumxxiileveliv}{\def\writerlistleftskip{\setlength\leftskip{2in}}\def\writerlistparindent{}\def\writerlistlabel{}\def\item{\def\writerlistparindent{\setlength\parindent{-0.25in}}\def\writerlistlabel{\stepcounter{listWWNumxxiileveliv}\makebox[0cm][l]{\labellistWWNumxxiileveliv}\hspace{-4.4449997cm}\writerlistremovelabel}}}{}
\setlength\tabcolsep{1mm}
\renewcommand\arraystretch{1.3}
\newcounter{Table}
\renewcommand\theTable{\arabic{Table}}
\newcounter{Figure}
\renewcommand\theFigure{\arabic{Figure}}
\title{Die Syntax von Vergleichskonstruktionen: formale Ansätze}
\author{}
\date{2024-06-28}
\begin{document}
\clearpage\setcounter{page}{1}\begin{styleStandard}
\textbf{The (in)transparency of meaning change and variation: a study of the indefinite }\textbf{\textit{cualquiera}}\textbf{ in European and Argentinian Spanish}
\end{styleStandard}

\begin{styleStandard}
Olga Kellert
\end{styleStandard}

\begin{styleStandard}
This book is a modified version of my habilitation thesis submitted to the University of Göttingen in 2021
\end{styleStandard}

\begin{styleStandard}
(University of Göttingen)
\end{styleStandard}

\clearpage\begin{styleContentsHeading}
Tables and Figures
\end{styleContentsHeading}

\begin{styleelenconumeratoii}
\textbf{Tables}
\end{styleelenconumeratoii}

\clearpage\begin{styleelenconumeratoii}
\textbf{Figures}
\end{styleelenconumeratoii}

\clearpage\begin{styleContentsHeading}
List of abbreviations
\end{styleContentsHeading}

\begin{styleStandard}
A\&M \ \ Alonso-Ovalle \& Menéndez-Benito
\end{styleStandard}

\begin{styleStandard}
ArgSp\ \ Argentinian Spanish
\end{styleStandard}

\begin{styleStandard}
CdE \ \ \ \ Corpus del Espanol
\end{styleStandard}

\begin{styleStandard}
CdE\_G/H\ \ \textit{Corpus del Espanol\_Genre/Historical}
\end{styleStandard}

\begin{styleStandard}
CdE\_web\textit{ \ \ Corpus del Español web dialects}
\end{styleStandard}

\begin{styleStandard}
CdE\_now\_EurSp.\ \ \textit{Corpus del Español now European Spanish}
\end{styleStandard}

\begin{styleStandard}
COM \ \ common
\end{styleStandard}

\begin{styleStandard}
CORDIAM \ \ \textit{Corpus Diacrónico y Diatópico del Español de América}
\end{styleStandard}

\begin{styleStandard}
CORDE \ \ \textit{Corpus Diacrónico del Español}
\end{styleStandard}

\begin{styleStandard}
C\&P\ \ \ \ Company Company \& Pozas-Loyo
\end{styleStandard}

\begin{styleStandard}
Decl\ \ \ \ Declarative
\end{styleStandard}

\begin{styleStandard}
Def\ \ \ \ Definite article
\end{styleStandard}

\begin{styleStandard}
DER \ \ \ \ deragotory
\end{styleStandard}

\begin{styleStandard}
DILE \ \ \textit{Diccionario Latinoamericano de la Lengua Española}
\end{styleStandard}

\begin{styleStandard}
DLE\ \ \ \ \textit{Diccionario de la Lengua Española}
\end{styleStandard}

\begin{styleStandard}
EVAL \ \ evaluative meanings: e.g. ‘ordinary’, ‘nonsense’
\end{styleStandard}

\begin{styleStandard}
EVAL 1 \ \ evaluative meaning ‘ordinary’
\end{styleStandard}

\begin{styleStandard}
EVAL 2 \ \ evaluative meaning ‘nonsense’
\end{styleStandard}

\begin{styleStandard}
EXH \ \ \ \ exhaustivity operator
\end{styleStandard}

\begin{styleStandard}
EurSp.\ \ European Spanish
\end{styleStandard}

\begin{styleStandard}
FCI \ \ \ \ Free Choice Interpretation
\end{styleStandard}

\begin{styleStandard}
Fr.\ \ \ \ French
\end{styleStandard}

\begin{styleStandard}
IGN\ \ \ \ ignorance
\end{styleStandard}

\begin{styleStandard}
IND\ \ \ \ indiscriminative ‘not just any’
\end{styleStandard}

\begin{styleStandard}
Ind\ \ \ \ Indicative
\end{styleStandard}

\begin{styleStandard}
It.\ \ \ \ Italian
\end{styleStandard}

\begin{styleStandard}
LF\ \ \ \ Logical Form
\end{styleStandard}

\begin{styleStandard}
Occ.\ \ \ \ Occurence
\end{styleStandard}

\begin{styleStandard}
RCI \ \ \ \ \ Random Choice Interpretation
\end{styleStandard}

\begin{styleStandard}
P* \ \ \ \ pragmatically determined property 
\end{styleStandard}

\begin{styleStandard}
P,Q,R,S \ \ Predicates
\end{styleStandard}

\begin{styleStandard}
\textsc{qualquier \ \ }all forms from Old Spanish to Modern Spanish 
\end{styleStandard}

\begin{styleStandard}
SHIFT \ \ shifts of alternatives from individuals to propositions
\end{styleStandard}

\begin{styleStandard}
SIM \ \ \ \ similarity
\end{styleStandard}

\begin{styleStandard}
Sp.\ \ \ \ Spanish
\end{styleStandard}

\begin{styleStandard}
Subj\ \ \ \ Subject
\end{styleStandard}

\begin{styleStandard}
Sbjv\ \ \ \ Subjunctive
\end{styleStandard}

\begin{styleStandard}
Subord.\ \ Subordinate
\end{styleStandard}

\begin{styleStandard}
s.v.\ \ \ \ sub voce
\end{styleStandard}

\begin{styleStandard}
XP \ \ \ \ Phrase (e.g. NP, PP, VP, PrP, KindP, ModifP)
\end{styleStandard}

\begin{styleStandard}
x, y, z \ \ Individual Variables 
\end{styleStandard}

\begin{styleStandard}
${\wedge}$,${\vee}$ \ \ \ \ \ (and, or) 
\end{styleStandard}

\begin{styleStandard}
${\exists}$,${\forall}$ \ \ \ \ (exists, all) Quantifiers
\end{styleStandard}

\begin{styleStandard}
$\lambda $ \ \ \ \ Lambda operator
\end{styleStandard}

\begin{styleStandard}
{\textless}e,t{\textgreater} \ \ \ \ predicate type
\end{styleStandard}

\begin{styleStandard}
{\textless}e{\textgreater} \ \ \ \ individual type
\end{styleStandard}

\begin{styleStandard}
{\textless}s,t{\textgreater} \ \ \ \ proposition type
\end{styleStandard}

\begin{styleStandard}
{\textless}t{\textgreater} \ \ \ \ truth value type
\end{styleStandard}

\begin{styleStandard}
[23CA?]\ \ \ \ meaning cannot apply
\end{styleStandard}

\begin{styleStandard}
$\iota $ \ \ \ \ iota-Operator
\end{styleStandard}

\clearpage\clearpage\setcounter{page}{1}\section[1 Introduction]{1 Introduction}
\begin{styleStandard}
Meaning variation and change of words or complex linguistic expressions are an integral part of language (change) (Blank 1997, Traugott 1989, Eckardt 2006, among others). However, the question as to how meaning change and variation exactly work is still an open question in linguistic research. Describing the meaning change of words in time and space and the mechanism underlying it is a challenging task as the historical and/or the cognitive or psychological conditions under which meaning change happens are not always given or known. This is probably one of the reasons why there are different views on the underlying mechanisms of meaning change and variation. Some linguists take a more algorithmic and logical perspective on describing this mechanism assuming that meaning change is based on a compositional reanalysis of the prior meaning into a new meaning, that is, the reanalysis can be made fully transparent and explainable by formal semantic principles as argued by formal diachronic semanticists (Eckardt 2006, Bar-Asher Siegal 2015, Condoravdi 2015, Aloni 2021, among many others). According to these authors, it is possible to describe the reanalysis involved in meaning change on the level of morphosyntactic form (F) and on the level of meaning (M) (see also Detgers et al. 2021). The important claim of formal diachronic semantics on meaning change is that at both points of time (t1 and t2), F suits M in a compositional manner—this is necessary for a reanalysis to take place (see Eckardt 2006, Bar-Asher Siegal 2015, Aloni 2021). The compositionality principle can be summarized as follows. A complex linguistic expression E such as \textit{the boy is sleeping} is compositional if the meaning of E depends on, and only on, (i) the (morpho)syntactic structure of E and (ii) the meanings of E’s simple linguistic units (Frege 1997). In other words, the meaning of E depends on the meaning of the definite article \textit{the} and the noun \textit{boy}, which combined together syntactically provide a more complex meaning of ‘a particular boy mentioned previously in the discourse’, the meaning of the verbal aspect morphology expressed in \textit{is +Verb+ing} and the meaning of the verbal lemma \textit{sleep}, which combined together provide the meaning of ‘is sleeping’ and finally the combined meaning of complex expressions has the sentence meaning ‘the boy is sleeping’. The general idea of diachronic formal semantics is that meaning variation and change take place in a transparent way and every diachronic stage or step in the reanalysis process involves compositionality that provides input to reanalysis.
\end{styleStandard}

\begin{styleStandard}
\ \ Other linguists take a different perspective assuming that meaning change does not need to be fully compositional or transparent such as meaning change that happens with associations or analogies triggered in a sociohistoric or situational context (Croft \& Cruse 2004, De Smet \& Fischer 2017, among many others). 
\end{styleStandard}

\begin{styleStandard}
\ \ This book has the general goal of testing the compositionality and the transparency principle in formal diachronic semantics on meaning change and variation of the modal indefinite \textit{cualquiera} ‘any’ in Argentinian and European Spanish in synchrony and diachrony. More precisely, my main goal is to investigate the question how the modal meaning of the indefinite \textit{cualquiera} ‘any’ changed into evaluative meanings ‘common/unremarkable’ in European and Argentinian Spanish and into the evaluative meaning ‘nonsense’ in Argentinian Spanish (empirical research question) and how this change fits models of linguistic reanalysis in the frameworks of compositional analysis of meaning change in formal diachronic semantics (theoretical research question).
\end{styleStandard}

\subsection[1.1 Structure of the thesis]{1.1 Structure of the thesis}
\begin{styleStandard}
In Sections 1.2 and 1.3, I provide a brief introduction into the meanings of \textit{cualquiera}, that will provide the scope of my study. In section 1.4, I give a brief overview of current state of the art on meaning change of \textit{cualqueira} and in 1.5, I present very briefly the general methods that will be used to approach my research questions (empirical and theoretical questions). 
\end{styleStandard}

\begin{styleStandard}
\ \ Chapter 2 offers a detailed review of the data and analyses discussed in the literature concerning \textit{cualquiera} in Modern European and Latin American Spanish in synchrony and diachrony. We will see that most synchronic and diachronic approaches do not investigate the variety of meanings of \textit{cualquiera} in synchrony and diachrony and how these new meanings are related to the modal meaning ‘any’ and what might have triggered this meaning change.
\end{styleStandard}

\begin{styleStandard}
\ \ In Chapter 3, I will present new data and a new analysis of special readings of \textit{cualquiera} in modern European Spanish by extracting new corpus data and testing the corpus analysis by speakers judgements from questionnaires and analysing them.
\end{styleStandard}

\begin{styleStandard}
\ \ In chapter 4, I will focus on \textit{cualquiera} in Argentinian Spanish in synchrony and diachrony and suggest a synchronic and diachronic analysis for \textit{cualquiera} in this variety. 
\end{styleStandard}

\begin{styleStandard}
\ \ Chapter 5 presents my research on the diachrony of \textit{cualquiera }in European Spanish. In this chapter, I will focus on the distribution of \textit{cualquiera} in various linguistic contexts or variables such as verbal position and verbal properties of the sentence in which \textit{cualquiera} occurs across time in more detail. I will raise the question of whether these variables play any role in the meaning change of \textit{cualquiera} across time using a quantitative/statistical analysis. In this chapter, I will provide a formal semantic analysis of different steps that led to meaning change of \textit{cualquiera} and its development in the diachrony.
\end{styleStandard}

\begin{styleStandard}
\ \ Finally, Chapter 6 summarises the most important findings of this thesis and their implications for theoretical research on meaning change.
\end{styleStandard}

\subsection[1.2 Modal reading of Spanish cualquiera]{1.2 Modal reading of Spanish \textit{cualquiera}}
\begin{styleStandard}
The main interest here is the indefinite \textit{cualquiera}, which can roughly be translated with ‘any’ and is usually described as a free choice (FC) item (FCI) in (certain varieties of) Spanish (see Brucart 1999, Sánchez López 1999, Alonso-Ovalle \& Menéndez-Benito 2011/2018, henceforth A\&M 2011/2018, Rivero 2011, Eberenz 2000). To introduce this item and before I turn to its semantics (which is the main focus of this book), let us have a brief look at its syntax and morphology:
\end{styleStandard}

\begin{styleStandard}
\textbf{\ \ }a.\ \ cualquier\ \ libro\ \ b.\ \ una\ \ camisa\ \ cualquiera
\end{styleStandard}

\begin{styleStandard}
\ \ \ \ \textsc{cualquiera}\textstyleFootnoteSymbol{ }\footnote{\textrm{\ Since }\textrm{\textit{cualquiera}}\textrm{ has different interpretations, as will be shown in this chapter and elaborated on throughout the thesis, I do not translate it in the glosses.}}\ \ book\ \ \ \ a\ \ Shirt\ \  \textsc{cualquiera}
\end{styleStandard}

\begin{styleStandard}
\ \ \ \ ‘any book’\ \ \ \ ‘any shirt’
\end{styleStandard}

\begin{styleStandard}
\ \ \ \ \ \ \ \  (Aloni et al. 2011: 2 / DLE, s.v. \textit{cualquiera})
\end{styleStandard}

\begin{styleStandard}
\textbf{\ \ }\ \ Esto\ \ lo\ \ puede\ \ hacer\ \ cualquiera.\ \ 
\end{styleStandard}

\begin{styleStandard}
\ \ \ \ this\ \ it\ \ can\ \ do\ \ \textsc{cualquiera}\ \ 
\end{styleStandard}

\begin{styleStandard}
\ \ \ \ ‘Anybody can do this.’
\end{styleStandard}

\begin{styleStandard}
\ \ \ \ (@ 1\footnote{\textrm{\ The symbol @ indicates that the source of this example is from the Internet. The number represents the index under which the example can be found in the Appendix, where all sources are listed.}})
\end{styleStandard}

\begin{styleStandard}
In (1), \textit{cualquiera} is an invariable indefinite element, which can occur pre- or postnominally. In its prenominal position in (1a), it usually drops the \textit{{}-a }and is not preceded by any article, in contrast to the postnominal position, as in (1b), where it is preceded by an indefinite article and cannot drop the \textit{{}-a}. In (2), it is an indefinite pronoun. In this function, the final \textit{{}-a} is always preserved. It does not generally vary for number and gender, its morphological plural form (\textit{cualesquiera}) being almost obsolete today (de Bruyne 2011: 236).
\end{styleStandard}

\begin{styleStandard}
FCIs are indefinites that come along with a free choice implicature or inference (i.e. each individual is a permissible option) (see Arregui 2006, Aloni \& Port 2013, Aloni 2021, A\&M 2011/2018, among many others): 
\end{styleStandard}

\begin{styleStandard}
\textbf{\ \ }\ \ \textit{Puedes}\ \ \textit{traer}\ \ \textit{cualquier}\ \ \textit{libro.}\ \ \ \ 
\end{styleStandard}

\begin{styleStandard}
\ \ \ \ can\ \ bring\ \ \textsc{cualquier}\ \ book\ \ \ \ 
\end{styleStandard}

\begin{styleStandard}
\ \ \ \ ‘You can bring any book.’
\end{styleStandard}

\begin{styleStandard}
\ \ \ \ Conventional meaning: ‘You can bring me a book and each book is a possible option.’
\end{styleStandard}

\begin{styleStandard}
\ \ \ \ (Aloni et al. 2011: 2)
\end{styleStandard}

\begin{styleStandard}
\textit{Cualquier(a)} is described in the literature as a lexically encoded free choice indefinite because its free choice inference seems to be integrated into its semantic content in sentences like (3) (see Aloni \& Port 2013, Aloni 2021, A\&M 2011/2018), whereas the free choice inference is possible but not necessarily available with plain indefinites such as \textit{uno} in (4) (see Aloni 2021):
\end{styleStandard}

\begin{styleStandard}
\textbf{\ \ }\ \ \textit{Puedes}\ \ \textit{traer}\ \ \textit{un}\ \ \textit{libro.}\ \ \ \ 
\end{styleStandard}

\begin{styleStandard}
\ \ \ \ can\ \ bring\ \ a\ \ book\ \ \ \ 
\end{styleStandard}

\begin{styleStandard}
\ \ \ \ ‘You can bring a book.’
\end{styleStandard}

\begin{styleStandard}
\ \ \ \ a.\ \ Conventional meaning: ‘You can bring a book.’
\end{styleStandard}

\begin{styleStandard}
\ \ \ \ b.\ \ Free choice inference: ‘Each book is a possible option.’
\end{styleStandard}

\begin{styleStandard}
\ \ \ \ \ \ (Aloni et al. 2011: 2)
\end{styleStandard}

\begin{styleStandard}
FCIs are known to appear in modal contexts such as with modal verbs—followed by \textit{que} ‘that’—to introduce free relative clauses with a subjunctive verb, as in (5) (see Quer 2000, Giannakidou 2001, among many others). This also applies to \textit{cualquiera}:
\end{styleStandard}

\begin{styleStandard}
\textbf{\ \ }\ \ Cualquiera que sea tu caso, vamos a intentar que en este artículo encuentres la respuesta.
\end{styleStandard}

\begin{styleStandard}
\ \ \ \ ‘Whatever your case is, we will try to make sure you find the answer in this article.’
\end{styleStandard}

\begin{styleStandard}
\ \ \ \ FC: ‘Each case is a possible option.’\ \ 
\end{styleStandard}

\begin{styleStandard}
\ \ \ \ \ \ (@ 2)
\end{styleStandard}

\begin{styleStandard}
Importantly, it has been noted in the literature that \textit{cualquiera} as well as comparable items in other Romance languages can have a series of different other interpretations that diverge from the free choice implicature (cf. A\&M 2011: 29, Rivero 2011 for Spanish; Fălăuş 2015 for \textit{oarecare }in Romanian; and Chierchia 2006: 543 for \textit{qualunque} in Italian). As will be shown, this is also the case in Spanish, and the main aim of this book is: (i) to investigate the interpretations of \textit{cualquiera} from a synchronic and diachronic perspective with the aim to probe compositional meaning change, and (ii) to find a common source for all, or at least a subset, of these interpretations. In order to probe a compositional analysis of meaning change, I will investigate whether the variation of semantic interpretation correlates with morphosyntactic and semantic/pragmatic properties.
\end{styleStandard}

\begin{styleStandard}
\ \ In what follows, I present interpretations other than the free choice interpretation triggered in modal contexts, such as in (4) or (5).
\end{styleStandard}

\subsection[1.3 Special readings of Spanish cualquiera]{1.3 Special readings of Spanish \textit{cualquiera}}
\begin{styleStandard}
First, the indefinite \textit{cualquiera} can have a random choice interpretation (RCI) in postnominal position with volitional verbs with perfective aspect (see 6) (cf. Choi \& Romero 2008, A\&M 2011: 29 and Rivero 2011 for Spanish; Fălăuş 2015 for \textit{oarecare }in Romanian; Chierchia 2006: 543, Kellert 2021 for \textit{qualunque/qualsiasi} in Italian; and Vlachou 2007, 2012 for \textit{quelconque }in French):
\end{styleStandard}

\begin{styleStandard}
\textbf{\ \ }\ \ Juan cogió\ \ una\ \ carta\ \ cualquiera.
\end{styleStandard}

\begin{styleStandard}
\ \ \ \ Juan grabbed\ \ a\ \ card\ \ \textsc{cualquiera}
\end{styleStandard}

\begin{styleStandard}
\ \ \ \ ‘Juan picked a random card.’
\end{styleStandard}

\begin{styleStandard}
\ \ \ \ Random choice interpretation (RCI): ‘Juan picked a card, and he chose it indiscriminately— he could have picked any other.’
\end{styleStandard}

\begin{styleStandard}
\ \ \ \ (A\&M 2018: 13)
\end{styleStandard}

\begin{styleStandard}
RCI is also known under the term ‘indiscriminative function’, i.e. it conveys the agent’s indifference with respect to his/her choice from a given set of alternatives (for Spanish, see A\&M 2015: 14, Horn 2000; for English; Vlachou 2007, 2012 for French). The indiscriminative function has been also observed for \textit{wh}{}-\textit{ever} relatives in English; for further illustration of this function, see the sentence in (7), where \textit{whatever} indicates that the agent’s (in the case the speaker’s) choice of tool was indiscriminate (von Fintel 2000):
\end{styleStandard}

\begin{styleStandard}
\textbf{\ \ }\ \ I grabbed whatever tool was handy. \ \ 
\end{styleStandard}

\begin{styleStandard}
\ \ \ \ \ \ (von Fintel 2000: 32)
\end{styleStandard}

\begin{styleStandard}
Second, postnominal \textit{cualquiera} can also have an evaluative interpretation (EVAL). It describes the speaker’s or agent’s (negative) attitude towards some entity, for example that he/she finds Juan ordinary\footnote{\ \textrm{Note that EVAL is a descriptive term that refers to different attitude predicates that have different flavours or readings such as ‘ordinary’, ‘common’, or ‘unremarkable’. These different flavours highly depend on the context (see Section 3).}} (for Spanish, see A\&M 2011: 29 and Rivero 2011; for Romanian, Fălăuş 2015; for Italian, Chierchia 2006: 543, Kellert 2021, for French Vlachou 2007, 2012; Šimík 2008 for Czech):\textstyleFootnoteSymbol{ }
\end{styleStandard}

\clearpage\begin{styleStandard}
\textbf{\ \ }\ \ Juan\ \ es\ \ un\ \ hombre\ \ \textit{cualquiera.}\ \ (Spanish)
\end{styleStandard}

\begin{styleStandard}
\ \ \ \ Juan\ \ is\ \ a\ \ man\ \ \textsc{cualquiera}\ \ 
\end{styleStandard}

\begin{styleStandard}
\ \ \ \ ‘Juan is an ordinary man’ (EVAL)
\end{styleStandard}

\begin{styleStandard}
\ \ \ \ (@ 3)
\end{styleStandard}

\begin{styleStandard}
Third, the postnominal \textit{cualquiera} can also have a depreciatory meaning of ‘worthless’ or refer to low quality (see Rivero 2011, A\&M 2011 for Spanish; Fălăuş 2015 for Romanian \textit{oarecare}; Vlachou 2007, 2012 on French and Greek FCIs). I consider this meaning to be a subset of the evaluative interpretation (EVAL):
\end{styleStandard}

\clearpage\begin{styleStandard}
\textbf{\ \ }\ \ Es\ \ un\ \ hombre\ \ \textit{cualquiera.}\ \ \ \ (Spanish)
\end{styleStandard}

\begin{styleStandard}
\ \ \ \ is\ \ a\ \ man\ \ \textsc{cualquiera}\ \ \ \ 
\end{styleStandard}

\begin{styleStandard}
\ \ \ \ ‘He is a man without great qualities.’ (EVAL)
\end{styleStandard}

\begin{styleStandard}
\ \ \ \ (@ 4)
\end{styleStandard}

\begin{styleStandard}
The evaluative reading is particularly interesting, because it leads to further processes of recategorisation and lexicalisation. For example, the French FCI \textit{quelconque} ‘which-ever’ shows a shift from a free choice indefinite to a gradable adjective:
\end{styleStandard}

\clearpage\begin{styleStandard}
\textbf{\ \ }\ \ Elle\ \ m’a\ \ paru\ \ assez\ \ quelconque.\ \ (French)
\end{styleStandard}

\begin{styleStandard}
\ \ \ \ she\ \ me has\ \ looked\ \ quite\ \ \textsc{quelconque \ \ }
\end{styleStandard}

\begin{styleStandard}
\ \ \ \ ‘She looked quite common to me.’ \ \ 
\end{styleStandard}

\begin{styleStandard}
\ \ \ \ (Vlachou 2012: 1530)
\end{styleStandard}

\begin{styleStandard}
Similarly, the French FCI \textit{n’importe quoi} (‘whatever’, lit. ‘does not matter what’) can also be used to express EVAL. In contrast to the adjectival use of \textit{quelconque} in (10), \textit{n’importe} \textit{quoi} is a noun that can be modified by an adjective and has the meaning ‘nonsense/bullshit’:
\end{styleStandard}

\clearpage\begin{styleStandard}
\textbf{\ \ }\ \ Il\ \ raconte\ \ du\ \ grand/vrai\ \ n’importe quoi.\ \ (French)
\end{styleStandard}

\begin{styleStandard}
\ \ \ \ he\ \ tells\ \ of-the\ \ big/true\ \ \textsc{n’importe quoi \ \ }
\end{styleStandard}

\begin{styleStandard}
\ \ \ \ ‘He talks really bullshit.’\ \ 
\end{styleStandard}

\begin{styleStandard}
\ \ \ \ (@ 138)
\end{styleStandard}

\begin{styleStandard}
In the course of this book, I will show similar effects for Argentinian Spanish (ArgSp), where \textit{cualquiera} has lexicalised the meaning EVAL and can nowadays be used as a gradable adjective modifiable by degree adverbs or the degree modifier \textit{muy} or \textit{re} (Rizzo Salierno 2013: 27):
\end{styleStandard}

\clearpage\begin{styleStandard}
\textbf{\ \ }\ \ La\ \ peli\ \ es\ \ (re\footnote{\textrm{\ The degree modifier }\textrm{\textit{re}}\textrm{ is used in Argentinian Spanish to reinforce the meaning of linguistic expressions. When }\textrm{\textit{re}}\textrm{ is combined with gradable adjectives, it has the meaning of ‘very’ (see Rizzo Salierno 2013 and Section 4 and, in particular, footnote 34, for further discussion).}}/muy)\ \ cualquiera.\ \ (ArgSp)
\end{styleStandard}

\begin{styleStandard}
\ \ \ \ the\ \ movie\ \ is\ \ very\ \ \textsc{cualquiera} \ \ 
\end{styleStandard}

\begin{styleStandard}
\ \ \ \ ‘The movie is (very) unremarkable/bad/not of great interest.’
\end{styleStandard}

\begin{styleStandard}
\ \ \ \ (@ 5 / @ 6)
\end{styleStandard}

\begin{styleStandard}
\textit{Cualquiera} in ArgSp can also be used as a noun under transitive verbs, similar to \textit{n’importe quoi}, in (11) with the meaning ‘nonsense/bullshit’:
\end{styleStandard}

\clearpage\begin{styleStandard}
\textbf{\ \ }\ \ Me\ \ dijiste\ \ cualquiera.\ \ (ArgSp)
\end{styleStandard}

\begin{styleStandard}
\ \ \ \ me\ \ said\ \ \textsc{cualquiera} \ \ 
\end{styleStandard}

\begin{styleStandard}
\ \ \ \ ‘You told me bullshit/nonsense.’
\end{styleStandard}

\begin{styleStandard}
\ \ \ \ (@ 7)
\end{styleStandard}

\begin{styleStandard}
In this book, I distinguish between \textit{cualquiera} with the interpretation of ‘unremarkable, bad, not of great interest’, which is mostly found as a predicative complement as in (12) (henceforth: EVAL 1) and \textit{cualquiera} embedded under transitive verbs with the meaning ‘bullshit/nonsense’ as in (13) (henceforth: EVAL 2).
\end{styleStandard}

\begin{styleStandard}
\ \ We find the same lexicalisation pattern in Mexican Spanish for the lexical item \textit{equis} (originally meaning ‘the letter x’), which can be used as an indefinite with a free choice interpretation ‘any’ in conditional clauses as in (14) as shown in Kellert (2020a: 137), but can also function as a gradable adjective that expresses EVAL 1 as shown in (15) (see Kellert 2020a: 139):
\end{styleStandard}

\clearpage\begin{styleStandard}
\textbf{\ \ }\ \ Si por equis razón no vas a venir, avisame\ \ (MexSp)
\end{styleStandard}

\begin{styleStandard}
\ \ \ \ ‘For whatever reason you cannot come, inform me.’\ \ 
\end{styleStandard}

\begin{styleStandard}
\ \ \ \ (Kellert 2020a: 137)
\end{styleStandard}

\clearpage\begin{styleStandard}
\textbf{\ \ }\ \ un hombre bastante/demasiado/muy equis\ \ (MexSp)
\end{styleStandard}

\begin{styleStandard}
\ \ \ \ ‘a very ordinary man (with no special qualities)’\ \ 
\end{styleStandard}

\begin{styleStandard}
\ \ \ \ (Kellert 2020a: 138)
\end{styleStandard}

\begin{styleStandard}
Finally, \textit{cualquiera} can be interpreted as Negative Polarity Item (NPI) in European Spanish, i.e. as an existential indefinite in scope of negation with the interpretation of ‘not… anything’ (see De Bruyne 2011, §497, Arregui 2006, Chierchia 2006, among others):
\end{styleStandard}

\clearpage\begin{styleStandard}
\textbf{\ \ }\ \ Estos últimos, además de no comer carne y pescado, no comen cualquier cosa que sea de origen animal, como la miel, la leche y los huevos.
\end{styleStandard}

\begin{styleStandard}
\ \ \ \ ‘The latter ones, apart from not eating meat and fish, don’t eat anything of animal origin such as honey, milk, and eggs.’
\end{styleStandard}

\begin{styleStandard}
\ \ \ \ (@ 8)
\end{styleStandard}

\begin{styleStandard}
However, the NPI interpretation of \textit{cualquiera} is restricted to relative clause modification, because without relative clause modification, \textit{cualquiera} has the meaning ‘everything’ or a ‘not just any’ interpretation under negation (see Arregui 2006, among others; for the general phenomonen, see Chierchia 2006):
\end{styleStandard}

\clearpage\begin{styleStandard}
\textbf{\ \ }\ \ No\ \ comen\ \ cualquier\ \ cosa.
\end{styleStandard}

\begin{styleStandard}
\ \ \ \ not\ \ eat\ \ \textsc{cualquier}\ \ thing
\end{styleStandard}

\begin{styleStandard}
\ \ \ \ ‘They don’t eat everything.’ or They don’t eat just anything. They eat something special.’
\end{styleStandard}

\begin{styleStandard}
\ \ \ \ \# ‘They don’t eat anything. Not even a piece of bread.’
\end{styleStandard}

\begin{styleStandard}
\ \ \ \ (@ 9)
\end{styleStandard}

\begin{styleStandard}
Given all the meanings in the examples from (1) to (17), the question arises as to whether these meanings represent lexical ambiguity, that is, where each reading of \textit{cualquiera} has a different denotation (as with the word \textit{bank} in English, which denotes either a financial institution or the land alongside a river or lake), or rather that \textit{cualquiera} has only one denotation and the different meanings described above are the result of some pragmatic inference on top of the semantic denotation:
\end{styleStandard}

\begin{styleStandard}
\textbf{\ \ }\ \ Hypothesis 1: \textit{Cualquiera} has different denotations (a case of polysemy).
\end{styleStandard}

\begin{styleStandard}
\textbf{\ \ }\ \ Hypothesis 2: \textit{Cualquiera} has one denotation and several meanings follow from this denotation.
\end{styleStandard}

\begin{styleStandard}
In this book, I will argue that the latter case is true for European Spanish and that \textit{cualquiera} has only one denotation there, namely the denotation of an existential quantifier with the meaning ‘some (or other)’ (see Fălăuş 2014, Aloni 2021, among others for analysing FCIs as existential quantifiers), whereas it is lexically ambiguous in Argentinian Spanish (‘some or other’ and ‘bad’).
\end{styleStandard}

\begin{styleStandard}
\ \ Despite some progress in the research on free choice indefinites, there is still no agreement about the semantic relation between RCI and EVAL 1 (A\&M 2015: 4), that is, on whether these interpretations or at least some subset of them have a common semantic source (i.e. semantic denotation) and, if so, how one interpretation is related to another. A\&M (2011) assume as a working hypothesis that \textit{cualquiera} is ambiguous between the random choice interpretation (RCI) and the evaluative interpretation as ‘unremarkable’ (EVAL 1) in Modern Spanish. However, as they mention themselves, this working hypothesis should be studied in detail (A\&M 2011: 41). In this book, I will test this hypothesis. Moreover, I will also study \textit{cualquiera} in Argentinian Spanish, which allows EVAL 2, in contrast to European Spanish, as will be shown in Section 1.2. The main question of this book is how EVAL 1 and EVAL 2 are related to FCI or RCI. 
\end{styleStandard}

\subsection[1.4 Brief overview of the current state of research and open issues]{1.4 Brief overview of the current state of research and open issues}
\begin{styleStandard}
Most synchronic accounts of \textit{cualquiera} focus on one of its particular interpretations such as the random choice interpretation (see Choi \& Romero 2008, A\&M 2011) or the free choice interpretation (see Menéndez-Benito 2010, Aloni \& Port 2013, Fălăuş 2015, among others). There is no systematic investigation of the distribution and restrictions of the evaluative meaning of \textit{cualquiera}. Moreover, previous accounts do not try to unify the meanings that we are concerned with in this investigation or derive one meaning from the other. It thus remains to be seen whether the evaluative meaning can be derived from the free choice function of indefinites.
\end{styleStandard}

\begin{styleStandard}
\ \ The main aim of this book is to fill this gap by offering a new analysis that can unify different meanings of \textit{cualquiera} that look unrelated in synchrony, namely the FCI and RCI meanings ‘any’ and ‘random’ and the EVAL meanings ‘ordinary, bad, common, nonsense’. In order to achieve this goal, it will be crucial to incorporate a diachronic perspective on meaning variation. In fact, as I will demonstrate in this investigation for Spanish \textit{cualquiera}, it turns out that at some diachronic stage, EVAL1 was derived from FCI through a certain syntactic and/or semantic reanalysis of FCIs. By looking at the history of the meanings and functions of \textit{cualquiera} we will be able to identify how one particular meaning evolved from the other and not the other way round. In sum, a diachronic perspective can help us to reconstruct the relation between different meanings in a compositional way.
\end{styleStandard}

\begin{styleStandard}
\ \ This perspective of combining synchrony and diachrony has not been discussed so far in the literature on the semantic evolution of \textit{cualquiera}. The few existing diachronic studies concentrate on the question of how the relative pronoun \textit{qual} and the verb form \textit{quiera} (3\textsc{sg.pres.subjv.}) became grammaticalised into one complex compound (see Palomo 1934, Lombard 1938–39/47–48, Rivero 1988, Company Company \& Pozas-Loyo 2009, henceforth C\&P 2009, Mackenzie 2019, among others). These accounts, however, do not discuss the question of what consequences this grammaticalisation process has for the interpretational component of \textit{cualquiera}. Furthermore, it is still unclear when and how the different meanings associated with this form have emerged.
\end{styleStandard}

\begin{styleStandard}
\ \ Moreover, there is no systematic quantitative analysis of diachronic data on \textit{cualquiera} that includes a statistical significance test giving thus empirical support for a change. While C\&P (2009) and Aloni \& Port (2013) have presented a quantitative description of \textit{cualquiera}, they did not use any statistical test to verify their hypotheses. Also, the classification of the data in the literature is not satisfactory. Aloni \& Port (2013) look at a number of syntactic and semantic properties of \textit{cualquiera}, but they do not attempt to establish a correlation between the different factors considered. From their results we do not learn whether different functions correlate with different syntactic positions of \textit{cualquiera} (e.g. +/– prenominal position, +/– preverbal, +/– ap-pearance in episodic context or copular verbs). 
\end{styleStandard}

\begin{styleStandard}
\ \ In view of the shortcomings of previous studies, it seems relevant to use a more detailed syntactic classification for a corpus analysis of synchronic and diachronic data on \textit{cualquiera}. In particular, I will look at the morphosyntactic and semantic properties of the verb forms that co-occur in the sentences in which \textit{cualquiera} appears and describe them with respect to their modality, lexical specification, tense, and aspect. This will allow us to observe correlations between the different meanings of \textit{cualquiera} and the aforementioned verbal properties. As for the diachrony of \textit{cualquiera}, I will investigate the first occurrences of its different meanings in order to draw conclusions about the sequence in which these meanings appear over time as well as their semantic source.
\end{styleStandard}

\begin{styleStandard}
\ \ As a result of all of the above, my book is the first monograph to contain a full analysis of \textit{cualquiera} taking into account synchronic and diachronic perspectives combined with a solid empirical basis grounded on corpus data and quantitative analysis, as well as a combination of formal and functional theoretical approaches.
\end{styleStandard}

\subsection[1.5 Data and methodology]{1.5 Data and methodology}
\begin{styleStandard}
In order to describe the evaluative properties of \textit{cualquiera} in Modern European Spanish, I have extracted data with a noun +\textit{ cualquiera} (henceforth (\textsc{un)} N \textit{cualquiera}) from the \textit{Corpus del Español web dialects} (CdE\_web), which is annotated for different regional varieties of Spanish. The data on the postnominal in Argentinian Spanish as well as on the pronoun \textit{cualquiera} in Argentinian Spanish and in other Latin American varieties including European Spanish was extracted from the same corpus. I also consulted native speakers of European and Argentinian Spanish personally and digitally on Facebook.
\end{styleStandard}

\begin{styleStandard}
\ \ In order to respond to my diachronic research question, I have chosen 150 randomised occurences per century from the 13th to the 20th centuries from the \textit{Corpus del Español\_hist-gen} (henceforth CdE\_hist-gen) and have annotated them according to 30 syntactic and semantic properties. As some of the texts contained in CdE\_hist-gen are taken from editions that have low philological reliability, as reported in Octavio de Toledo \& Rodríguez Molina (2017), I selected only those texts that these authors classify as suitable for linguistic analysis. For the diachronic research in Argentinian Spanish, I used the diachronic corpus CORDIAM (\textit{Corpus Diacrónico y Diatópico del Español}, CORDIAM, www.cordiam.org, which contains a total of 10,889,176 words from 19 Latin-American countries from 1494 to 1905, with documents from archives, journals, and literature. It is annotated for geographical origins. The corpus CORDIAM does not contain many instances of \textit{cualquiera} in Argentinian Spanish in comparison to CdE\_hist-gen (around 150 in total). I therefore also used a larger diachronic corpus, CORDE (\textit{Corpus Diacrónico del Español}), which is a diachronic corpus with written texts spanning from Old Spanish to 1974 with annotation of geographic origin (country), which contains 544 instances of \textit{cualquiera} in total in ArgSp in CORDE. I analysed all diachronic data from ArgSp from CORDIAM and CORDE as the data collection was relatively small. 
\end{styleStandard}

\begin{styleStandard}
\ \ The data will be published on GitHub upon the acceptance of this publication.
\end{styleStandard}

\begin{styleStandard}
\ \ The quantitative analysis was based on statistical methods such as the chi-square test and the Linear Regression Model (see Baayen 2006). I have used the R software programme for the statistical analysis of the data. The generation of the plots and histograms was carried out in Microsoft Excel. 
\end{styleStandard}

\begin{styleStandard}
\ \ For the semantic analysis of synchronic and diachronic data, I used formal semantic frameworks of FCIs (Chierchia 2006, Aloni 2021, Aloni \& Port 2013, A\&M 2011/2018, among others) and studied models of semantic reanalysis within formal semantics (Eckardt 2006, Aloni \& Port 2013, among others).
\end{styleStandard}

\clearpage\section[2 Current state of research on cualquiera]{\textbf{2 Current state of research on }\textbf{\textit{cualquiera}}}
\begin{styleStandard}
In this chapter, I review the current state of research on modern uses of \textit{cualquiera} in different varieties of Spanish. Many empirical and theoretical studies of \textit{cualquiera} are based on European Spanish \textit{cualquiera}, as will be shown in this chapter. There are few descriptive studies of \textit{cualquiera} in Argentinian Spanish (see 2.1.4). However, these studies do not relate the description of \textit{cualquiera }to the theoretical studies discussed in Sections 2.2 through 2.5. I will fill this gap in Chapters 3 and 4 by analysing the variation between Modern European and Modern Argentinian Spanish in theoretical terms and by looking at the diachrony of \textit{cualquiera} in Chapters 4 and 5, which will show the syntactic and semantic conditions under which the evaluative meaning can occur in diachrony. 
\end{styleStandard}

\subsection[2.1 Informal description of syntactic and semantic properties of modern cualquiera]{\textbf{2.1 Informal description of syntactic and semantic properties of modern }\textbf{\textit{cualquiera}}}
\begin{styleStandard}
In European Spanish, \textit{cualquiera} can be used as a pronoun, determiner, noun, and postnominal modifier. As a pronoun and determiner, it has a restricted distribution with respect to verb mood and modality. However, as I will show, this restricted distribution is not present with \textit{cualquiera} as a noun or postnominal modifier. In modern Argentinian Spanish, the word can also be used as a gradable adjective without the verb mood restrictions found in European Spanish. One aim of the book is to understand the variation in Spanish and the variation across syntactic categories. As many uses of \textit{cualquiera} are the same in both varieties (Argentinian and European Spanish), besides its use as an evaluative noun with the meaning of ‘nonsense’ and as an adjective in the sense of ‘bad/ordinary’ in Argentinian Spanish but not in European Spanish, I only explicitly mark the language variety Argentinian Spanish, when these special uses are important for the discussion.
\end{styleStandard}

\subsubsection[2.1.1 The pronoun cualquiera and the determiner cualquier]{\textbf{2.1.1 The pronoun }\textbf{\textit{cualquiera }}\textbf{and the determiner }\textbf{\textit{cualquier}}}
\begin{styleStandard}
\textit{Cualquiera} can be used as an indefinite pronoun in modern Spanish to express generic statements, i.e. statements about every person or every entity, especially when it is used as a subject in declarative sentences (Brucart 1999, Sánchez López 1999, Eberenz 2000, DLE, s.v. \textit{cualquiera}):
\end{styleStandard}

\begin{styleStandard}
\textbf{\ \ }Esto lo hace cualquiera. 
\end{styleStandard}

\begin{styleStandard}
\textbf{\ \ }‘Anybody does this.’
\end{styleStandard}

\begin{styleStandard}
\ \ (@ 12)
\end{styleStandard}

\begin{styleStandard}
\textbf{\ \ }Cualquiera lo hubiera hecho mejor. 
\end{styleStandard}

\begin{styleStandard}
\textbf{\ \ }‘Anybody would have done it better.’
\end{styleStandard}

\begin{styleStandard}
\textbf{\ \ }(@ 13)
\end{styleStandard}

\begin{styleStandard}
The indefinite pronoun \textit{cualquiera} can be followed by a prepositional phrase [\textit{de }DP] ‘of DP’ (Sánchez López 1999: 1041, DLE, s.v. \textit{cualquiera}) in which case the DP restricts the quantificational domain of \textit{cualquiera}, e.g. in examples \ and \ the quantification applies to a particular set of\textit{ listed }reasons and translators, respectively (for the quantificational property of \textit{cualquiera}, see Section 2.2):
\end{styleStandard}

\begin{styleStandard}
\label{bkm:Ref72400832}\textbf{\ \ }cualquiera de las razones enumeradas
\end{styleStandard}

\begin{styleStandard}
\ \ ‘any of the listed reasons’
\end{styleStandard}

\begin{styleStandard}
\ \ (@ 14a)
\end{styleStandard}

\begin{styleStandard}
\label{bkm:Ref533953098}\textbf{\ \ }cualquiera de esos traductores
\end{styleStandard}

\begin{styleStandard}
\ \ ‘any of those translators’
\end{styleStandard}

\begin{styleStandard}
\ \ (@ 14b)
\end{styleStandard}

\begin{styleStandard}
\textit{Cualquiera} can also be used as a (quantificational) determiner\footnote{\textrm{\ In this book, I will use the terminology of Zamparelli (2000) to refer to quantifiers in prenominal position such as }\textrm{\textit{some/any/every N}}\textrm{ and Spanish counterparts as (quantificational) determiners.}} similar to\textit{ cada} ‘each’ in declarative sentences to express generic statements (Sánchez López 1999: 1034, 1036; DLE, s.v. \textit{cualquiera}; Menéndez-Benito 2010; Rivero 2011):
\end{styleStandard}

\begin{styleStandard}
\textbf{\ \ }Cualquier persona tiene derecho a decidir sobre su cuerpo.
\end{styleStandard}

\begin{styleStandard}
\ \ ‘Any person has the right to decide about her/his body.’
\end{styleStandard}

\begin{styleStandard}
\ \ (@ 15)
\end{styleStandard}

\begin{styleStandard}
\textbf{\ \ }Cualquier ciudadano tiene derecho a elegir. 
\end{styleStandard}

\begin{styleStandard}
\ \ ‘Any citizen has the right to choose.’
\end{styleStandard}

\begin{styleStandard}
\ \ (@ 14c)
\end{styleStandard}

\begin{styleStandard}
The similar interpretation of \textit{cualquiera} and a universal quantifiers such as \textit{cada }‘each’\textit{ }and \textit{todo} ‘all’ can be illustrated by the adverbial modification of both universal quantifiers and \textit{cualquiera} by \textit{casi} ‘almost’ and other adverbs such as \textit{excepto} ‘except/but’ that restrict the domain of quantification from every possible element to a smaller part (see Arregui 2006, among others):
\end{styleStandard}

\begin{styleStandard}
\textbf{\ \ }Casi todos saben hablar inglés. 
\end{styleStandard}

\begin{styleStandard}
\textbf{\ \ }‘Almost everbody knows how to speak English.’
\end{styleStandard}

\begin{styleStandard}
\textbf{\ \ }(@ 17)
\end{styleStandard}

\begin{styleStandard}
\textbf{\ \ }Casi cualquiera puede decirte eso.
\end{styleStandard}

\begin{styleStandard}
\textbf{\ \ }‘Almost anybody can tell you that.’
\end{styleStandard}

\begin{styleStandard}
\textbf{\ \ }(Arregui 2006, e.g. 65 a.)
\end{styleStandard}

\begin{styleStandard}
\textit{Cualquiera} can be used as a determiner inside a Determiner Phrase (DP) that specifies a relative clause with a subjunctive as shown in . This relative clause has an indifference interpretation, ‘I don’t care about the flaws he has’ (see DLE, s.v. \textit{cualquiera}):
\end{styleStandard}

\begin{styleStandard}
\label{bkm:Ref512708}\textbf{\ \ }Cualquier defecto que tenga, yo lo voy a querer igual.
\end{styleStandard}

\begin{styleStandard}
\ \ ‘Whatever flaws he/it has, I will love him/it anyway.’
\end{styleStandard}

\begin{styleStandard}
\ \ (Rizzo Salierno 2013: 24, similar to @ 18)
\end{styleStandard}

\begin{styleStandard}
Subjunctive relative clauses like that in \ are interpreted similar to conditional clauses as in \ (see Quer 2000):
\end{styleStandard}

\begin{styleStandard}
\label{bkm:Ref513122}\textbf{\ \ }Aunque tenga defectos, yo lo voy a querer igual.
\end{styleStandard}

\begin{styleStandard}
\ \ ‘Although/Even if he/it has flaws, I will love him/it anyway.’
\end{styleStandard}

\begin{styleStandard}
\ \ (similar to @ 19)
\end{styleStandard}

\begin{styleStandard}
Both types of conditionals require a subjunctive. The subjunctive includes worlds other than the actual world. Thus, the example in \ not only takes into account the actual world where the person under discussion in might have only a specific flaw but all other flaws in all possible worlds that the speaker can think of (see Quer 2000, Chierchia 2006, Becker 2014 on the role of subjunctive). The use of the subjunctive in the relative clause produces the effect that the assertion ‘I will love her’ is pragmatically strengthened, because it is evaluated with respect to worlds in which the person under discussion might have some or all kinds of flaws. The assertion is considered as true without any restriction. 
\end{styleStandard}

\begin{styleStandard}
\ \ Let us now turn to \textit{cualquiera} as a pronoun, which can have different syntactic functions, such as that of a subject pronoun, as shown in , or it can be used as a pronoun having a function of a copular predicate inside a conditional clause with a subjunctive, as in :
\end{styleStandard}

\begin{styleStandard}
\label{bkm:Ref533761017}\textbf{\ \ }Cualquiera sea el defecto que tiene, yo lo voy a querer igual.
\end{styleStandard}

\begin{styleStandard}
\ \ ‘Whatever the flaw she has may be, I will love him/it anyway.’
\end{styleStandard}

\begin{styleStandard}
\ \ (similar to @ 20)
\end{styleStandard}

\begin{styleStandard}
\label{bkm:Ref533761034}\textbf{\ \ }Sea cualquiera el defecto que tiene, yo lo voy a querer igual.
\end{styleStandard}

\begin{styleStandard}
\ \ ‘Whatever be the flaw he has, I will love him/it anyway.’
\end{styleStandard}

\begin{styleStandard}
\ \ (similar to @ 21)
\end{styleStandard}

\begin{styleStandard}
It is a very common observation in the literature that free choice indefinites have a restricted distribution (Quer 2000, Menéndez-Benito 2010, A\&M 2011/2018, among many others). They are usually licensed in modal contexts such as with modal and generic verbs or in the subjunctive (see Chapter 1) and ruled out in episodic\footnote{\textrm{\ I define a context as episodic if the sentence at issue involves an event variable e, and e is existentially closed: Jean is eating fruits. ${\exists}$e [eating fruits (e) \& Agent (e, Jean)].}} sentences, which are usually marked by perfective aspect such as in \ (see Menéndez-Benito 2010, Aloni et al. 2010, Dayal 2013, among others):
\end{styleStandard}

\begin{styleStandard}
\label{bkm:Ref41139913}\ \ a.\ \ *Juan cogió cualquiera de las cartas de esta baraja.
\end{styleStandard}

\begin{styleStandard}
\ \ \ \ ‘Juan took-\textsc{pfv} any of the cards from this deck.’
\end{styleStandard}

\begin{styleStandard}
\ \ b.\ \ *John took any of the cards from this deck. 
\end{styleStandard}

\begin{styleStandard}
\ \ \ \ (Menéndez-Benito 2010: 34)
\end{styleStandard}

\begin{styleStandard}
However, as has been noted by Quer (2000), Rivero (2011), and Menéndez-Benito (2010), among others, it is possible to find the determiner \textit{cualquiera} in episodic contexts only if the noun phrase is modified by a relative clause (so-called ‘subtrigging’ examples):
\end{styleStandard}

\begin{styleStandard}
\label{bkm:Ref41139369}\ \ a.\ \ *Juan compró cualquier revista.
\end{styleStandard}

\begin{styleStandard}
\ \ \ \ ‘Juan bought any magazine.’
\end{styleStandard}

\begin{styleStandard}
\ \ b.\ \ Juan compró cualquier revista que tuviera fotos de actualidad.
\end{styleStandard}

\begin{styleStandard}
\ \ \ \ ‘Juan bought whatever magazine had photos of hot news.’ 
\end{styleStandard}

\begin{styleStandard}
\ \ \ \ (Rivero 2011: 67)
\end{styleStandard}

\begin{styleStandard}
Quer (2000) argues that free choice indefinites (FCIs) such as \textit{cualquiera} are licensed by the presence of a relative clause in examples such as , which brings in a modal context in which it is tied to the quantificational interpretation of characterising sentences in the past, whereby a habitual or generic operator quantifies over worlds or situations (see Krifka et al. 1995). The sentence in \ has a conditional interpretation associated with generic/characterising statements: if there was a magazine with hot news, John would buy it, or every time there was a magazine with photos of hot news, John would buy that magazine. The typical cases of licensing of FCIs under subtrigging involve modal context such as a conditional-like interpretation. To sum up, given its restriction to modal context, \textit{cualquiera} as a determiner and as a pronoun is ungrammatical as the object of a present or past tense indicative verb that refers to a specific event or situation as in . 
\end{styleStandard}

\begin{styleStandard}
\ \ What has not been investigated in the literature so far is the use of determiner or pronoun \textit{cualquiera} under indicative copular verbs that predicate over individuals or entities. In the following context, the property of being a man is predicated over the person Juan. In this context, only the indefinite article \textit{un} is judged as acceptable, whereas the determiner \textit{cualquier} and the pronoun \textit{cualquiera} are not\textit{,} according to speaker judgements:
\end{styleStandard}

\begin{styleStandard}
\textbf{\ \ }Juan\ \ es\ \ un\ \ hombre\ \ (de\ \ 30\ \ años).
\end{styleStandard}

\begin{styleStandard}
\ \ Juan\ \ is\ \ a\ \ man\ \ of\ \ 30\ \ years
\end{styleStandard}

\begin{styleStandard}
\ \ ‘Juan is a (30-year-old) man.’
\end{styleStandard}

\begin{styleStandard}
\ \ (@ 22)
\end{styleStandard}

\begin{styleStandard}
\label{bkm:Ref513974}\textbf{\ \ }*Juan es cualquier hombre (de 30 años)/*Juan es cualquiera. 
\end{styleStandard}

\begin{styleStandard}
\ \ Juan is any 30-year-old man /Juan is anyone.
\end{styleStandard}

\begin{styleStandard}
One possible explanation of the ungrammaticality in \ is that \textit{cualquier/cualquiera} are not used in a modal or generic context necessary for its licensing as shown in (25) through (36). However, if \textit{cualquiera} is used in a licensing context such as a partitive phrase as in \ or a subtrigging clause as in b), the sentences are grammatical as expected:
\end{styleStandard}

\begin{styleStandard}
\textbf{\ \ }\label{bkm:Ref57364726}a.\ \ Si tu hospedaje es cualquiera de estos hoteles a continuación ¡nos lo vas a agradecer!
\end{styleStandard}

\begin{styleStandard}
\textbf{\ \ }\ \ ‘If you are staying at any of these hotels mentioned below, you will thank us!’
\end{styleStandard}

\begin{styleStandard}
\textbf{\ \ }\ \ (@ 23)
\end{styleStandard}

\begin{styleStandard}
\textbf{\ \ }b.\ \ Sabio es cualquiera que sepa seguir los movimientos, las fases de la vida para mantener el y\=\in yáng en armonía en su interior.
\end{styleStandard}

\begin{styleStandard}
\ \ \ \ ‘Wise is anyone who knows how to follow the movements, the life cycles…’ 
\end{styleStandard}

\begin{styleStandard}
\ \ \ \ (Genís Sol, \textit{Flujo del espíritu humano}, 35)
\end{styleStandard}

\begin{styleStandard}
To sum up, \textit{cualquiera} as a pronoun or determiner is not licensed in episodic or affirmative predicative contexts without an overt modal. 
\end{styleStandard}

\begin{styleStandard}
\ \ One rescuing context of the determiner or the pronoun \textit{cualquiera} in predicative contexts without a modal is negation, which shows that \textit{cualquiera} is polarity-sensitive (see Chierchia 2006, Sections 2.2.2 and 2.4 on FCIs under negation):
\end{styleStandard}

\begin{styleStandard}
\textbf{\ \ }El Sr. Miguel Ángel no es cualquier empleado de CC.OO, es director del gabinete de estudios de CC.OO
\end{styleStandard}

\begin{styleStandard}
\ \ ‘Mr Miguel Angel is not just any worker of CC.OO, he is the director of CC.OO’s research bureau’
\end{styleStandard}

\begin{styleStandard}
\ \ (@ 24)
\end{styleStandard}

\begin{styleStandard}
\textbf{\ \ }La curva que sale no es cualquier curva. Es una parabola
\end{styleStandard}

\begin{styleStandard}
\ \ ‘The resulting curve is not just any curve. It is a parabola’
\end{styleStandard}

\begin{styleStandard}
\ \ (@ 25)
\end{styleStandard}

\begin{styleStandard}
To sum up, \textit{cualquiera} as a pronoun and determiner is not permitted with indicative mood, which refers to actual situations or property descriptions in the actual world without subtrigging or negation.
\end{styleStandard}

\subsubsection[2.1.2 Postnominal cualquiera and cualquiera as a noun]{\textbf{2.1.2 Postnominal }\textbf{\textit{cualquiera}}\textbf{\textit{ }}\textbf{and }\textbf{\textit{cualquiera}}\textbf{ as a noun}}
\begin{styleStandard}
As will be shown in this section, postnominal \textit{cualquiera} has a status different from that of the pronominal and the determiner \textit{cualquiera}. The first difference is that postnominal \textit{cualquiera} is usually used with indefinite nouns, as in \ (i.e. the structure \textsc{un} N \textit{cualquiera}). In contrast to the determiner and pronoun \textit{cualquiera}, \textsc{un} N \textit{cualquiera} can be used in contexts that refer to actual situations or properties of actual entities/individuals. The question is thus whether \textsc{un} N \textit{cualquiera} still expresses modality in these contexts.
\end{styleStandard}

\begin{styleStandard}
\ \ In , it is asserted that Juan took some card. The semantic contribution of \textit{cualquiera} is to express that this card was taken randomly, which is the so-called random choice interpretation (RCI). It is commonly assumed in the literature that this interpretation correlates with the appearance of volitional verbs, which imply some agent that has desires or goals, such as \textit{coger }‘to take/to pick’ as in . The RCI implies a modal meaning expressing that the agent Juan could have picked any other card (see Menéndez-Benito 2005, Choi \& Romero 2008, A\&M 2011/2018, Sánchez López 1999: 1041; Arregui 2006; Rivero 2011):
\end{styleStandard}

\begin{styleStandard}
\label{bkm:Ref533761157}\textbf{\ \ }Juan cogió una carta cualquiera
\end{styleStandard}

\begin{styleStandard}
\ \ Random Choice (RC): Juan picked a card and he chose it indiscriminately—he could have picked any other.
\end{styleStandard}

\begin{styleStandard}
\ \ (Menéndez-Benito 2005: 2)
\end{styleStandard}

\begin{styleStandard}
However, in affirmative contexts such as in –, with indicative copular\footnote{\textrm{\ The copula }\textrm{\textit{be}}\textrm{ has different uses (see Higgins 1979, Zamparelli 2000, Author 2015, and references therein). It can be used predicationally such as }\textrm{\textit{Today is a common day (for me)}}\textrm{, or it can be used identificationally/ specificationally }\textrm{\textit{Today is the 21st of July 2010 }}\textrm{or}\textrm{\textit{ Today is Monday}}\textrm{. The predicative use usually comes with an indefinite article such as }\textrm{\textit{Today is a/*the nice day}}\textrm{. Another test to tease the functions apart is to look at the type of the subject the copula takes. In the predicational use, the subject is an individual type, e.g}\textrm{\textit{. Jean is a man}}\textrm{. In the identificational use, the subject is expressed by a demonstrative or deictic pronoun, e.g. }\textrm{\textit{This is Jean }}\textrm{or}\textrm{\textit{ It’s Jean.}}\textrm{ According to this typology of copular verbs, it is easy to see that }\textsc{un}\textrm{ N}\textrm{\textit{ cualquiera}}\textrm{ in }\textrm{\textit{Juan es un hombre cualquiera}}\textrm{ ‘Juan is an ordinary man’ or in }\textrm{\textit{Es un hombre cualquiera}}\textrm{ should be interpreted predicationally, i.e. }\textsc{un}\textrm{ N}\textrm{\textit{ cualquiera}}\textrm{ is a copular complement that describes some property about Juan. From this point on, I will use the term predicative }\textrm{\textit{cualquiera}}\textrm{ to refer to }\textsc{un}\textrm{ N }\textrm{\textit{cualquiera}}\textrm{ as a predicative complement of predicative verbs such as }\textrm{\textit{ser}}\textrm{ ‘be’ and similar verbs (}\textrm{\textit{parecer}}\textrm{ ‘resemble’, etc.). }} verbs as in \ or with perception verbs such as \textit{mirar }‘to look (at)’ or \textit{ver} ‘to see’ like in , \textit{cualquiera} has the evaluative interpretation ‘ordinary/common/unremarkable’ (see A\&M 2011/2018, DLE, s.v. \textit{cualquiera}):
\end{styleStandard}

\begin{styleStandard}
\label{bkm:Ref533761081}\textbf{\ \ }Juan es un estudiante cualquiera. 
\end{styleStandard}

\begin{styleStandard}
‘Juan is an unremarkable student.’ (EVAL)
\end{styleStandard}

\begin{styleStandard}
\textbf{\ \ }(A\&M 2018: 3)
\end{styleStandard}

\begin{styleStandard}
\textbf{\ \ }Eran dos mujeres cualquiera. 
\end{styleStandard}

\begin{styleStandard}
\textbf{\ \ }‘They were two ordinary women.’ (EVAL)
\end{styleStandard}

\begin{styleStandard}
\textbf{\ \ }(Rizzo Salierno 2013: 24, similar to @ 26)
\end{styleStandard}

\begin{styleStandard}
\label{bkm:Ref533953228}\textbf{\ \ }Te miró y vio un hombre cualquiera, no vio al diablo 
\end{styleStandard}

\begin{styleStandard}
\ \ ‘(S)he saw you and (s)he saw an ordinary man, not the devil’ (EVAL)
\end{styleStandard}

\begin{styleStandard}
\ \ (@ 14d)
\end{styleStandard}

\begin{styleStandard}
The evaluative function can also be expressed in a sentence that refers to an episodic context, in which \textit{cualquiera} modifies a noun that must refer to some entity or individual that the speaker knows and evaluates negatively’ see the interpretation in (i) (note that the native speakers I consulted also admit the so-called ignorance interpretation (IGN), as in (ii), where the speaker simply does not know the exact nature of the N):
\end{styleStandard}

\begin{styleStandard}
\label{bkm:Ref54686931}\textbf{\ \ }La carta de ajuste, en la pantalla, se complementaba con la música ramplona de una zarzuela cualquiera 
\end{styleStandard}

\begin{styleStandard}
\ \ i. \ \ ‘The testcard, in the screen, was complemented with the vulgar sound of some ordinary zarzuela’ \ \ (EVAL)
\end{styleStandard}

\begin{styleStandard}
\ \ ii.\ \ ‘…with the vulgar sound of some zarzuela that I don’t know’ (IGN)
\end{styleStandard}

\begin{styleStandard}
\ \ (@ 14e)
\end{styleStandard}

\begin{styleStandard}
Another type of the evaluative meaning, namely the pejorative or derogatory interpretation, can be observed with \textit{cualquiera} as a noun in \textsc{un} \textit{cualquiera},\footnote{\textrm{\ One could also assume that }\textrm{\textit{cualquiera}}\textrm{ in }\textrm{\textit{un(a) cualquiera}}\textrm{ is an adjective and that the noun has been elided (e.g. ‘person’). “Nominal gapping is permitted with PPs and postposed restricted adjectives in Spanish, as in }\textrm{\textit{El Ø rojo}}\textrm{ ‘the red one’” (Bosque 2001: 34).}} see the examples in \ and . The morphological plural form \textit{cualesquiera,} which was used in earlier periods of Spanish (see Section 2.5), is no longer accepted in Modern (European) Spanish, as shown in \ (DLE, s.v. \textit{cualquiera}), although it is still marginally found in literary language:\footnote{\textrm{\textsuperscript{\ }}\textrm{In this work, I abstract away from singular and plural differences of }\textrm{\textit{cualquiera}}\textrm{ with respect to their interpretation. }}
\end{styleStandard}

\begin{styleStandard}
\label{bkm:Ref150778470}\textbf{\ \ }Juan es un cualquiera.
\end{styleStandard}

\begin{styleStandard}
\ \ ‘Juan is a nobody’ 
\end{styleStandard}

\begin{styleStandard}
\ \ (Martínez 2005: 36)
\end{styleStandard}

\begin{styleStandard}
\label{bkm:Ref150778472}\textbf{\ \ }María es una cualquiera
\end{styleStandard}

\begin{styleStandard}
\ \ ‘María is a prostitute’
\end{styleStandard}

\begin{styleStandard}
\ \ (@ 27)
\end{styleStandard}

\begin{styleStandard}
\label{bkm:Ref533761216}\textbf{\ \ }Son unos/unas cualquiera/*cualesquiera.
\end{styleStandard}

\begin{styleStandard}
\ \ ‘They are nobodies/prostitutes’
\end{styleStandard}

\begin{styleStandard}
\ \ (@ 28/@ 29)
\end{styleStandard}

\begin{styleStandard}
It seems that no modality is involved in contexts where \textsc{un} (N) \textit{cualquiera} has an evaluative meaning, such as under copular verbs. The question is thus what status \textit{cualquiera} has in these predicative contexts with affirmative sentence mood and whether it is still used as an indefinite element there. It is often assumed in the literature that postnominal \textit{cualquiera} acts as a determiner (see 2.1.3 on data and discussion of this assumption). Note that the postnominal \textit{cualquiera} cannot be analysed as a (degree) adjective in European Spanish, because it cannot be modified by degree adverbs (see 2.1.4 for ArgSp):
\end{styleStandard}

\begin{styleStandard}
\textbf{\ \ }Un hombre (*muy) cualquiera.
\end{styleStandard}

\begin{styleStandard}
\ \ ‘A man (*very) cualquiera’
\end{styleStandard}

\begin{styleStandard}
\ \ (@ 4)
\end{styleStandard}

\begin{styleStandard}
This conclusion is confirmed by the absence of morphological agreement. Adjectives usually inflect for person, number, and gender in Modern Spanish in contrast to \textit{cualquiera} which is invariable (Seco 1986, s.v \textit{cualquiera}, Rainer 1993, s.v. \textit{cualquiera}). In its marginal use in written Spanish, the plural is realised on \textit{cual} (originally a realtive pronoun) and not on the ending -\textit{a} in \textit{cualquiera}, as we would expect if \textit{cualquiera} were lexicalised as an adjective:
\end{styleStandard}

\begin{styleStandard}
\textbf{\ \ }Ordinario,.\ \ ordinaria,\ \ ordinaries,\ \ ordinarias
\end{styleStandard}

\begin{styleStandard}
\textbf{\ \ }ordinary.\textsc{m.sg.}\ \ ordinary.\textsc{f.sg}.\ \ ordinary.\textsc{m.pl}.\ \ ordinary.\textsc{f.pl.}
\end{styleStandard}

\begin{styleStandard}
\ \ ‘Ordinary’
\end{styleStandard}

\begin{styleStandard}
\textbf{\ \ }cualquiera, \ \ cual(es)quiera,\ \ *cualquiero,\ \ *cual(es)quieros,\ \ *cual(es)quieras
\end{styleStandard}

\begin{styleStandard}
\textbf{\ \ }cualquiera.\textsc{sg}.\ \ cualquiera.\textsc{pl}.\ \ cualquiera.\textsc{m.sg}.\ \ cualquiera.\textsc{m.pl.}\ \ cualquiera.\textsc{f.pl.}
\end{styleStandard}

\begin{styleStandard}
\ \ ‘any’
\end{styleStandard}

\begin{styleStandard}
These restrictions can be accounted for diachronically as \textit{cualquiera} is a morphological compound of the relative pronoun \textit{cual} and the verb \textit{querer} (see Chapter 5). This is why no gender or plural agreement on the ending \textit{{}-a} (originally probably a subjunctive form) is possible in written Spanish (see Chapter 5).
\end{styleStandard}

\begin{styleStandard}
\ \ However, we do find the plural form -\textit{s} on the ending -\textit{a} in some Latin American varieties of Spanish (e.g. “\textit{cualquieras que sean….”}, see Rainer 1993: 298) and according to C\&P (2009) few cases of spoken European Spanish (e.g. \textit{unas cualquieras}, see C\&P 2009: 1125), which can be interpreted as a loss of this item’s transparency as a morphological compound of the relative pronoun \textit{cual} and the verb \textit{querer} in these varieties. As will be shown in 2.1.4, \textit{cualquiera} can be analysed as an adjective in these varieties because of degree modification (see 2.1.4 for details):
\end{styleStandard}

\begin{styleStandard}
\textbf{\ \ }La peli es (muy/re) cualquiera. 
\end{styleStandard}

\begin{styleStandard}
‘The movie is (very) ordinary/mediocre’\ \ (okay in ArgSp/ *EurSp)
\end{styleStandard}

\begin{styleStandard}
\textbf{\ \ }(@ 5 / @ 6)
\end{styleStandard}

\begin{styleStandard}
The variation in Spanish can be interpreted as a sign that the reanalysis of \textit{cualquiera} from a relative clause into a degree predicate is more advanced in ArgSp than in EurSp. 
\end{styleStandard}

\subsubsection[2.1.3 Syntactic analyses of postnominal cualquiera]{\textbf{2.1.3 Syntactic analyses of postnominal }\textbf{\textit{cualquiera}}}
\begin{styleStandard}
The syntactic and semantic status of postnominal \textit{cualquiera} is still unclear in the literature. Lapesa (2000 [1974]) interprets the definiteness restriction of \textit{un/*el hombre cualquiera} as evidence that \textit{cualquiera} is a quantificational determiner like \textit{cualquier} in \textit{cualquier hombre}. The ‘determiner’ \textit{cualquiera} can co-occur with \textit{un} and not with \textit{el} because \textit{un} does not have any semantic content according to Lapesa (2000 [1974]). Determiners that do not have any semantic meaning are analysed as weak determiners in the literature, whereas determiners that contribute semantic content to the overall meaning are called strong determiners (see Zamparelli 2000). 
\end{styleStandard}

\begin{styleStandard}
\ \ The same observation and conclusion about the determiner status of \textit{cualquiera} has been made by Rivero (2011), who has shown that postnominal \textit{cualquiera} cannot combine with quantified or definite nouns, but only with bare nouns or indefinite nouns (\textit{un, una,} \textit{unos, unas} N-sg/N-pl) or nouns modified by numerals (\textit{dos, tres,} etc.). The following examples in \ and \ are from Rivero (2011: 72):
\end{styleStandard}

\begin{styleStandard}
\label{bkm:Ref63458156}\ \ a.\ \ * Puedes comprar \{muchos/algunos/todos/los/estos\} libros cualesquiera.
\end{styleStandard}

\begin{styleStandard}
\ \ \ \ ‘You can buy \{many/some/all/the/those\} random books’
\end{styleStandard}

\begin{styleStandard}
\ \ b.\ \ * Puedes comprar \{este, ese, aquel, el\} libro cualquiera.
\end{styleStandard}

\begin{styleStandard}
\ \ \ \ ‘You can buy \{this, that, the\} random book.’
\end{styleStandard}

\begin{styleStandard}
\label{bkm:Ref533762162}\ \ a.\ \ Compra (otras)\{dos, tres, cuatro, etc.\} revistas cualesquiera.
\end{styleStandard}

\begin{styleStandard}
\ \ \ \ ‘Buy any \{two, three, four, etc.\} (more) journals’
\end{styleStandard}

\begin{styleStandard}
\ \ b.\ \ Compra unas revistas cualesquiera.
\end{styleStandard}

\begin{styleStandard}
\ \ \ \ ‘Buy any journals.’
\end{styleStandard}

\begin{styleStandard}
To conclude, it is assumed in the literature on postnominal \textit{cualquiera} that it is a strong (quantificational) determiner because it cannot co-occur with other determiners of the same kind. As will be shown in Section 3.6 on corpus data from Modern Spanish, postnominal \textit{cualquiera} can co-occur with other strong determiners. Similar data will be shown on diachronic corpus data in Section 6.2.2 I will conclude, that postnominal \textit{cualquiera} is not a determiner. 
\end{styleStandard}

\begin{styleStandard}
\ \ In order to explain the licensing of postnominal \textit{cualquiera} in episodic contexts, Rivero (2011) assumes that postnominal \textit{cualquiera} encodes lexical modality; that is, it does not need a licensor, as it licenses itself. According to Rivero (2011), the item \textit{quiera} inside \textit{cualquiera} is interpreted as a modal operator that creates a modal context and licenses \textit{cualquiera}. The following structure is a representation of the Logical From (LF), according to which the item \textit{quiera} is interpreted as a modal operator that scopes over the sentence with a perfective verb:
\end{styleStandard}

\begin{styleStandard}
\textbf{\ \ }LF [\textsubscript{Sentence} [\textsubscript{Modal} \textbf{quiera} ] [\textsubscript{Sentence} Juan compró [\textsubscript{NP} [\textsubscript{NumP} una ] revista cual [\textsubscript{Modal} quiera ]]]]
\end{styleStandard}

\begin{styleStandard}
Although Rivero (2011) did not consider predicative contexts in affirmative sentences in which \textsc{un} N \textit{cualquiera} is also lincensed, one could assume that her explanation covers these cases as well:
\end{styleStandard}

\begin{styleStandard}
\textbf{\ \ }LF [\textsubscript{Sentence} [\textsubscript{Modal} \textbf{quiera} ] [\textsubscript{Sentence} Juan es [\textsubscript{NP} [\textsubscript{NumP} un ] hombre cual [\textsubscript{Modal} quiera ]]]]
\end{styleStandard}

\begin{styleStandard}
However, Rivero (2011) is silent about in what sense \textit{cualquiera} expresses modality in unlicensed contexts such as episodic or predicative contexts.
\end{styleStandard}

\subsubsection[2.1.4 Cualquiera in Argentinian Spanish]{\textbf{2.1.4 }\textbf{\textit{Cualquiera}}\textbf{ in Argentinian Spanish}}
\begin{styleStandard}
A special use of \textit{cualquiera} is found in Argentinian Spanish. According to Conde (2011) it has a function of an indefinite pronoun with the meaning similar to \textit{cualquier cosa} ‘anything’ and meanings derived thereof in expressions like \textit{decir cualquiera} or \textit{mandar cualquiera} ‘say something out of place, inconvenient, tell a lie’ or in expressions like \textit{estar en cualquiera} ‘to be in a bad situation, be in trouble’ (literally ‘to be in any [place]’; see also Company Company 2009 for \textit{estar en cualquiera}). This use is documented for Lunfardo:\footnote{\textrm{\ }\textrm{\textit{Lunfardo}}\textrm{ is a variety of Argentinian Spanish, which originated from a massive immigration from Europe between 1881 and 1914. Italians were the largest group of immigrants in this period, comprising 48\% (see Devoto 2002, Kellert \& Francia in press, among others). }\textrm{\textit{Lunfardo}}\textrm{ was originally spoken only in the River Plate region, with a vocabulary enriched from several immigrant languages including Italian (Conde 2011). Sociolinguistic properties of Lunfardo still reflect some connection to the Italian immigration (see Kellert \& Francia in press).}}
\end{styleStandard}

\begin{styleStandard}
\ \ Lexical entry of \textit{cualquiera} (Conde 2011, \textit{Diccionario etimológico del lunfardo}, s.v.)
\end{styleStandard}

\begin{styleStandard}
\ \ \textit{cualquiera}. pron. indef. En las exprs. ss.: DECIR o MANDAR CUALQUIERA: decir algo fuera de lugar, inconveniente o disparatado; inventar una mentira. {\textbar} 2. ESTAR EN CUALQUIERA: pertenecer a ambientes delictivos o relacionados con la drogadicción; estar distraído. (Del esp. \textit{cualquiera}: persona intedeterminada, que pasa a significar aquí “cualquier cosa”.) 
\end{styleStandard}

\begin{styleStandard}
‘\textit{cualquiera.} indef. pronoun. In the expressions such as: DECIR o MANDAR CUALQUIERA: say something out of place, inconvenient, or unreflected; to invent a lie. {\textbar} 2. 2. ESTAR EN CUALQUIERA: to belong to criminal environment or to environments related to drug addiction; to be distracted. (From Sp. \textit{cualquiera}: indeterminate person, which comes to mean here “anything”)’
\end{styleStandard}

\begin{styleStandard}
However, this use of \textit{cualquiera} is not restricted to this variety. According to a note by Di Tullio (2015) in the framework of DILE (\textit{Diccionario Latinoamericano de la Lengua Española}), under verbs such as \textit{decir} or \textit{mandar} ‘to say’, \textit{cualquiera} means ‘nonsense/incorrect/false’:
\end{styleStandard}

\begin{styleStandard}
\ \ A: \ \ Me\ \ dijiste\ \ que\ \ venías\textit{\ \ }tremprano.\ \ B: \ \ \textit{Cualquiera}.
\end{styleStandard}

\begin{styleStandard}
\ \ \ \ me\ \ told\ \ that\ \ would.come\ \ early\ \  \ \ \textsc{cualquiera}
\end{styleStandard}

\begin{styleStandard}
\ \ A: ‘You said that you would come early.’ B: ‘Nonsense’
\end{styleStandard}

\begin{styleStandard}
\ \ \ \ (Di Tullio 2015; DILE, s.v.))
\end{styleStandard}

\begin{styleStandard}
\ \ mandaste\ \ \textit{cualquiera}.\ \ 
\end{styleStandard}

\begin{styleStandard}
\ \ said\ \ \textsc{cualquiera}\ \ 
\end{styleStandard}

\begin{styleStandard}
\ \ ‘You said nonsense.’
\end{styleStandard}

\begin{styleStandard}
(Di Tullio 2015; DILE, s.v.)
\end{styleStandard}

\begin{styleStandard}
There is a similar pejorative use of the free choice indefinite pronoun \textit{n’importe quoi }(lit. ‘it doesn’t matter what’)\textit{ }under verbs of saying and copular verbs in Modern French, where it has been nominalised (see Paillard 1997, Pescarini 2009, Section 3 for details):
\end{styleStandard}

\begin{styleStandard}
\ \ dire/c’est \ \ (du grand)\ \ n’importe quoi.
\end{styleStandard}

\begin{styleStandard}
\ \ say/it is\ \ (of-the big)\ \ \textsc{n’importe quoi}
\end{styleStandard}

\begin{styleStandard}
\ \ ‘to say/it’s (total) bullshit.’
\end{styleStandard}

\begin{styleStandard}
\ \ (@ 30 / @ 31)
\end{styleStandard}

\begin{styleStandard}
In this use, \textit{cualquiera} can be modified by a number of degree modifiers or affixes. According to Rizzo Salierno (2013), all of them refer to extreme or high degrees (see Rizzo Salierno 2013: 27).
\end{styleStandard}

\begin{styleStandard}
\ \ \textit{re-}: prefix\footnote{\ \textrm{I disagree with the analysis of }\textrm{\textit{re}}\textrm{ as a prefix as it is always written as a separate word in social media and can be reinforced by repetition }\textrm{\textit{re re cualquiera}}\textrm{ ‘very very bad’.} } that indicates an extreme degree of a gradable predicate 
\end{styleStandard}

\begin{styleStandard}
\ \ a. \textit{muy}, \textit{mal }
\end{styleStandard}

\begin{styleStandard}
\ \ b.\textit{ demasiado}, \textit{absulutamente}, \textit{totalmente }‘too, absolutely, totally’
\end{styleStandard}

\begin{styleStandard}
\ \ c. -\textit{ito}, -\textit{azo}: affective suffixes 
\end{styleStandard}

\begin{styleStandard}
\ \ d. -\textit{ísimo}: suffix indicating extreme degree, mostly used with feminine nouns
\end{styleStandard}

\begin{styleStandard}
\ \ \textit{Es\ \ demasiado\ \ cualquiera,}\ \ casi\ \ me\ \ hace\ \ estrellar\ \ mi\ \ cabeza
\end{styleStandard}

\begin{styleStandard}
\ \ is\ \ too\ \ \textsc{cualquiera}\ \ almost\textit{\ \ }me\textit{\ \ }makes\textit{\ \ }slam\ \ my\ \ head
\end{styleStandard}

\begin{styleStandard}
\ \ contra\ \ el\ \ monitor.\textsc{\ \ }\ \ \ \ \ \ \ \ \ \ 
\end{styleStandard}

\begin{styleStandard}
\ \ against\ \ the\ \ monitor\textsc{\ \ }\ \ \ \ \ \ \ \ \ \ 
\end{styleStandard}

\begin{styleStandard}
\ \ ‘It’s too \textit{bad/ordinary}, it almost makes me slam my head against the monitor.’
\end{styleStandard}

\begin{styleStandard}
\ \ (Rizzo Salierno 2013: 27; ArgSp)
\end{styleStandard}

\begin{styleStandard}
\ \ Si\ \ es\ \ una\ \ mina,\ \ \textit{re\ \ cualquiera.}
\end{styleStandard}

\begin{styleStandard}
\ \ yes\ \ is\ \ a\ \ girl\ \ very\ \ \textsc{cualquiera}
\end{styleStandard}

\begin{styleStandard}
\ \ ‘Yes, she is a very ordinary girl.’
\end{styleStandard}

\begin{styleStandard}
\ \ (Rizzo Salierno 2013: 26; ArgSp)
\end{styleStandard}

\begin{styleStandard}
\ \ Bueno \ \ vieja,\ \ pero\ \ entonces\ \ cambiale\ \ el\ \ título\ \ al\ \ thread
\end{styleStandard}

\begin{styleStandard}
\ \ well\ \ old lady\ \ but\ \ then\ \ change\ \ the\ \ title\ \ to-the\ \ thread
\end{styleStandard}

\begin{styleStandard}
\ \ Porque\ \ si\ \ no\ \ e\textit{s\ \ cualquiera\ \ mal.}\footnote{\textrm{\ If }\textrm{\textit{mal}}\textrm{ has a status of a degree modifier modifying }\textrm{\textit{cualquiera}}\textrm{ then the syntactic order is inverse in }\textrm{\textit{cualquiera mal/cualunque mal}}\textrm{ as compared to }\textrm{\textit{re cualquiera}}\textrm{, where the degree modifier is placed before }\textrm{\textit{cualquiera}}\textrm{ rather than after it. In the future, I will test the hypothesis as to whether }\textrm{\textit{cualquiera}}\textrm{ can be analysed as a degree modifier itself in }\textrm{\textit{cualquiera mal}}\textrm{. }}
\end{styleStandard}

\begin{styleStandard}
\ \ because\ \ otherwise\ \ is\ \ \textsc{cualquiera}\textit{\ \ }bad
\end{styleStandard}

\begin{styleStandard}
\ \ ‘Well, old lady, but then change the title of the thread, otherwise it’s really \textit{bad}.’
\end{styleStandard}

\begin{styleStandard}
\ \ (Rizzo Salierno 2013: 27; ArgSp)
\end{styleStandard}

\begin{styleStandard}
\ \ Si\ \ es\ \ un\ \ ranking\ \ de\ \ todos\ \ los\ \ tiempos\ \ este\ \ que\ \ pusiste\ \ es\ \ \textit{cualquier-ita}!!!
\end{styleStandard}

\begin{styleStandard}
\ \ if\ \ is\ \ a\ \ ranking\ \ of\ \ all\ \ the\ \ times\ \ this\ \ that\ \ put\ \ is\ \ \textsc{cualquiera-ita}
\end{styleStandard}

\begin{styleStandard}
\ \ ‘If this is a ranking of every time, this one you posted is \textit{bad}!!!’
\end{styleStandard}

\begin{styleStandard}
\ \ (Rizzo Salierno 2013: 27; ArgSp)
\end{styleStandard}

\begin{styleStandard}
\ \ \textit{Cualquieraza}\ \ ese\ \ video.
\end{styleStandard}

\begin{styleStandard}
\ \ \textsc{cualquiera}.\textsc{affective}\ \ that\ \ video
\end{styleStandard}

\begin{styleStandard}
\ \ ‘That video is \textit{ordinary/bad}.’
\end{styleStandard}

\begin{styleStandard}
\ \ (Rizzo Salierno 2013: 27; ArgSp)
\end{styleStandard}

\begin{styleStandard}
\ \ Es\ \ \textit{cualquierísima\ \ }tu\ \ planteo.
\end{styleStandard}

\begin{styleStandard}
\ \ is\ \ \textsc{cualquiera}.\textsc{elative}\ \ your\ \ proposal
\end{styleStandard}

\begin{styleStandard}
\ \ ‘Your proposal is extremely \textit{bad}.’
\end{styleStandard}

\begin{styleStandard}
\ \ (Rizzo Salierno 2013: 27; ArgSp)
\end{styleStandard}

\begin{styleStandard}
\ \ Es\ \ \textit{absolutamente\ \ cualquiera.\ \ }
\end{styleStandard}

\begin{styleStandard}
\ \ is\ \ absolutely\ \ \textsc{cualquiera}\ \ 
\end{styleStandard}

\begin{styleStandard}
\ \ ‘It is absolutely \textit{bad}.’
\end{styleStandard}

\begin{styleStandard}
\ \ (Rizzo Salierno 2013: 28; ArgSp)
\end{styleStandard}

\begin{styleStandard}
\ \ Es\ \ \textit{totalmente\ \ cualquiera\ \ }eso\ \ de\ \ haber\ \ puesto\ \ a\ \ Roger\ \ con\ \ la voz de Syd.
\end{styleStandard}

\begin{styleStandard}
\ \ is\ \ totally\ \ \textsc{cualquiera}\textit{\ \ }that\ \ of\ \ having\ \ put\ \ to\ \ Roger\ \ with\ \ the voice of Syd
\end{styleStandard}

\begin{styleStandard}
\ \ ‘Putting Roger with Syd’s voice is totally \textit{bad}.’
\end{styleStandard}

\begin{styleStandard}
\ \ (Rizzo Salierno 2013: 28; ArgSp)
\end{styleStandard}

\begin{styleStandard}
\ \ Eso\ \ de\ \ la\ \ tanga\ \ es\ \ \textit{completamente\ \ cualquiera.}
\end{styleStandard}

\begin{styleStandard}
\ \ that\ \ of\ \ the\ \ thong\ \ is\ \ completely\ \ \textsc{cualquiera}
\end{styleStandard}

\begin{styleStandard}
\ \ ‘That thong thing is completely \textit{bad}.’
\end{styleStandard}

\begin{styleStandard}
\ \ (Rizzo Salierno 2013: 28; ArgSp)
\end{styleStandard}

\begin{styleStandard}
\textit{Cualquiera} can also have an attributive use in a postnominal position as in European Spanish (see 2.1.3):
\end{styleStandard}

\begin{styleStandard}
\ \ Esto\ \ es\ \ \textit{un\ \ video\ \ cualquiera.}
\end{styleStandard}

\begin{styleStandard}
\ \ this\ \ is\ \ a\ \ video\ \ \textsc{cualquiera}
\end{styleStandard}

\begin{styleStandard}
\ \ ‘This is an ordinary video.’
\end{styleStandard}

\begin{styleStandard}
\ \ (Rizzo Salierno 2013: 26; ArgSp)
\end{styleStandard}

\begin{styleStandard}
\ \ Acá\ \ les\ \ dejo\ \ \textit{las\ \ fotos\ \ cualquiera\ \ }de\ \ hoy.
\end{styleStandard}

\begin{styleStandard}
\ \ here\ \ you\ \ leave\ \ the\ \ photos\ \ \textsc{cualquiera}\ \ of\ \ today
\end{styleStandard}

\begin{styleStandard}
\ \ ‘Here you have \textit{ordinary} photos of today.’
\end{styleStandard}

\begin{styleStandard}
\ \ (Rizzo Salierno 2013: 26; ArgSp)
\end{styleStandard}

\begin{styleStandard}
According to Rizzo Salierno (2013), \textit{cualquiera} can also be used under transitive verbs and be interpreted in a pejorative or derogatory way in Argentinian Spanish. However, she does not say anything about its different syntactic function in this context:
\end{styleStandard}

\begin{styleStandard}
\ \ \textit{Me\ \ contestó\ \ cualquiera\ \ }o\ \ ni\ \ me\ \ respondía.
\end{styleStandard}

\begin{styleStandard}
\ \ me\ \ answered\ \ \textsc{cualquiera}\textit{\ \ }or\ \ not even\ \ me\ \ responded
\end{styleStandard}

\begin{styleStandard}
\ \ ‘(S)he would either answer me \textit{bullshit} or not even answer.’
\end{styleStandard}

\begin{styleStandard}
\ \ (Rizzo Salierno 2013: 26; ArgSp)
\end{styleStandard}

\begin{styleStandard}
\ \ Se\ \ ve\ \ que\ \ \textit{me\ \ dijo\ \ cualquiera\ \ }la\ \ mina.
\end{styleStandard}

\begin{styleStandard}
\ \ it\ \ looks\ \ that\ \ me\ \ told\ \ \textsc{cualquiera}\ \ the\ \ girl
\end{styleStandard}

\begin{styleStandard}
\ \ ‘Apparently she told me bullshit, that girl.’
\end{styleStandard}

\begin{styleStandard}
\ \ (Rizzo Salierno 2013: 26; ArgSp)
\end{styleStandard}

\begin{styleStandard}
To sum up, \textit{cualquiera} can be used as an adjective with degree affixes and has an evaluative meaning of ‘ordinary/bad’ in ArgSp (Rizzo Salierno 2013). It can also be used under transitive verbs such as \textit{decir/mandar} \textit{cualquiera} or \textit{estar en cualquiera} as a direct object comparable to \textit{cualquier cosa} with a deragotary meaning of ‘bullshit’ (Conde 2011, Di Tullio 2013, Company Company 2009):
\end{styleStandard}

\begin{center}
\tablefirsthead{}
\tablehead{}
\tabletail{}
\tablelasttail{}
\begin{supertabular}{|m{1.5913599in}|m{1.9886599in}|m{1.94906in}|}
\hhline{~--}
\multicolumn{1}{m{1.5913599in}|}{} &
{\centering \textbf{Copular} \textbf{predicate}\par}

{\centering \textit{es cualquiera}\par}

\centering \textit{is}\textsc{ cualquiera} &
{\centering \textbf{Transitive verbs}\par}

{\centering \textit{decir cualquiera}\par}

\centering\arraybslash \textit{to say}\textsc{ cualquiera}\\\hline
\textbf{Evaluative meaning} &
\centering + (EVAL ‘ordinary, bad’) &
\centering\arraybslash + (EVAL ‘bullshit’)\\\hline
\textbf{Degree modification} &
\centering + &
\centering\arraybslash ?\\\hline
\end{supertabular}
\end{center}
\begin{styletablecaption}
Table \stepcounter{Table}{\theTable}:\textit{ Cualquiera} with evaluative meanings in ArgSp.
\end{styletablecaption}

\begin{styleStandard}
In Chapters 3 and 4, I will present new data and offer a semantic analysis of \textit{cualquiera} with EVAL.
\end{styleStandard}

\subsubsection[2.1.5 Summary ]{\textbf{2.1.5 Summary}\textbf{ }}
\begin{styleStandard}
The following table sums up the current state of research with respect to different syntactic functions of \textit{cualquiera} and their interpretations. The main observation is that the indicative mood is correlated with different interpretations—Eval, RC, and ignorance—in both European and Argentinian Spanish. The difference between the two varieties is that in Argentinian Spanish, \textit{cualquiera} with evaluative interpretation is used less restrictively across different syntactic categories, whereas in European Spanish it is only \textsc{un} (N) \textit{cualquiera} that allows the evaluative reading:
\end{styleStandard}

\begin{center}
\tablefirsthead{\hhline{~------}
\multicolumn{1}{m{0.71155983in}|}{} &
\multicolumn{2}{m{1.6712599in}|}{\centering \textbf{Syntactic Context}} &
\multicolumn{4}{m{3.38866in}|}{\centering \textbf{Semantic Interpretation}}\\\hline
\textbf{Syntactic Category} &
\centering \textbf{Degree modification} &
\centering \textbf{Indicative mood in episodic or predicative context} &
\centering \textbf{Evaluative interpretation (EVAL)} &
\centering \textbf{Free choice (FC)} &
{\centering \textbf{Random choice}\par}

\centering \textbf{(RC)} &
\centering\arraybslash \textbf{Ignorance interpretation}\\\hline
\textbf{Pronoun} &
{\centering no\par}

 &
{\centering no (EurSp)\par}

\centering yes (ArgSp) &
{\centering no (EurSp)\par}

\centering yes (ArgSp) &
\centering yes &
\centering no &
\centering\arraybslash no\\\hline
\textbf{\textit{cualquier }}\textbf{N} &
\centering no &
{\centering no (EurSp)\par}

\centering yes (ArgSp) &
{\centering no (EurSp)\par}

\centering yes (ArgSp) &
\centering yes &
\centering no &
\centering\arraybslash no\\\hline
\textbf{\textit{un/una}}\textbf{ (N) }\textbf{\textit{cualquiera}} &
{\centering no (EurSp)\par}

\centering yes (ArgSp) &
\centering yes &
\centering yes &
\centering yes (with artitives and under modals) &
{\centering yes\par}

{\centering (episodic tense,\par}

\centering volitional verbs) &
\centering\arraybslash yes (episodic tense)\\}
\tablehead{\hhline{~------}
\multicolumn{1}{m{0.71155983in}|}{} &
\multicolumn{2}{m{1.6712599in}|}{\centering \textbf{Syntactic Context}} &
\multicolumn{4}{m{3.38866in}|}{\centering \textbf{Semantic Interpretation}}\\\hline
\textbf{Syntactic Category} &
\centering \textbf{Degree modification} &
\centering \textbf{Indicative mood in episodic or predicative context} &
\centering \textbf{Evaluative interpretation (EVAL)} &
\centering \textbf{Free choice (FC)} &
{\centering \textbf{Random choice}\par}

\centering \textbf{(RC)} &
\centering\arraybslash \textbf{Ignorance interpretation}\\\hline
\textbf{Pronoun} &
{\centering no\par}

 &
{\centering no (EurSp)\par}

\centering yes (ArgSp) &
{\centering no (EurSp)\par}

\centering yes (ArgSp) &
\centering yes &
\centering no &
\centering\arraybslash no\\\hline
\textbf{\textit{cualquier }}\textbf{N} &
\centering no &
{\centering no (EurSp)\par}

\centering yes (ArgSp) &
{\centering no (EurSp)\par}

\centering yes (ArgSp) &
\centering yes &
\centering no &
\centering\arraybslash no\\\hline
\textbf{\textit{un/una}}\textbf{ (N) }\textbf{\textit{cualquiera}} &
{\centering no (EurSp)\par}

\centering yes (ArgSp) &
\centering yes &
\centering yes &
\centering yes (with artitives and under modals) &
{\centering yes\par}

{\centering (episodic tense,\par}

\centering volitional verbs) &
\centering\arraybslash yes (episodic tense)\\}
\tabletail{}
\tablelasttail{}
\begin{supertabular}{|m{0.71155983in}|m{0.69975984in}|m{0.8927598in}|m{0.83515984in}|m{0.7955598in}|m{0.7143598in}|m{0.8073598in}|}

\end{supertabular}
\end{center}
\begin{styletablecaption}
Table \stepcounter{Table}{\theTable}: Syntactic and semantic functions of \textit{cualquiera} in Modern European and Argentinian Spanish.
\end{styletablecaption}

\begin{styleStandard}
In the following section, I will look at the distribution of different interpretations of \textit{cualquiera} as defined in the literature more closely.
\end{styleStandard}

\subsection[2.2 Formal semantic analysis of cualquiera]{\textbf{2.2 Formal semantic analysis of }\textbf{\textit{cualquiera}}}
\begin{styleStandard}
In this section, I will give a brief introduction of semantic analysis of \textit{cualquiera} from the literature which is necessary to understand the formal semantic analysis in this section. I will first introduce into the notion of total variation and present formal analyses of \textit{cualquiera}. 
\end{styleStandard}

\subsubsection[2.2.1 Total variation of cualquiera and partial variation of algún ]{\textbf{2.2.1 Total variation of }\textbf{\textit{cualquiera}}\textbf{ and partial variation of }\textbf{\textit{algún}}\textbf{ }}
\begin{styleStandard}
\textit{Algún} is an epistemic indefinite that signals the speaker’s ignorance (IGN) with respect to the identity of the individual or object denoted by the indefinite. This makes it incompatible with continuations that specify the referent of the \textit{algún} phrase, as illustrated in \ (see A\&M 2008 for \textit{algún} phrases, Kratzer and Shimoyama 2002 for German \textit{irgendein}, Aloni \& Port 2013 for Italian \textit{un qualche, }Chierchia 2006 for \textit{un }NP\textit{ qualunque}, Jayez \& Tovena 2007, and others for French \textit{quelque} and \textit{un }NP\textit{ quelconque},\textit{ }Fălăuş 2014 for Romanian \textit{un }NP\textit{ oarecare}):
\end{styleStandard}

\begin{styleStandard}
\label{bkm:Ref533761877}\ \ Maria se casó con algún medico. \# En concreto con el Dr. Smith.
\end{styleStandard}

\begin{styleStandard}
\ \ ‘Maria married some doctor. Concretely, Dr. Smith.’
\end{styleStandard}

\begin{stylelsSourceline}
\textup{\ \ (A\&M 2013: 36) }
\end{stylelsSourceline}

\begin{styleStandard}
The impossibility of such types of continuation is captured by the “anti-singleton condition”, which expresses the idea captured by A\&M (2011) and Giannakidou \& Quer (2013) that \textit{algún} can only be used if there is more than one alternative known to the speaker (e.g\textit{. algún estudiante} ‘some or other student’). This condition expresses \textit{epistemic indeterminacy} or \textit{anti-specificity} (see also Giannakidou and Quer 2013, Etxeberria \& Giannakidou 2017). Specificity is given when “a speaker uses an indefinite noun phrase and intends to refer to a particular referent” (von Heusinger 2011: 10). Giannakidou \& Quer (2013) suggest that the phenomenon of referential vagueness is the absence of specificity, i.e., the absence of referential anchoring and the absence of referential intent. To sum up, \textit{algún} is anti-specific and expresses referential vagueness. It signals the lack of referential intent that characterises specificity (von Heusinger 2002, 2007, 2011). Referential vagueness consists of the minimal variation in possible values, and in uncertainty about which one is the actual value. 
\end{styleStandard}

\begin{styleStandard}
\ \ One crucial dimension along which epistemic indefinites vary concerns the extent of variation (‘freedom of choice’) imposed on the domain of quantification, which can be \textit{total }or \textit{partial} (see Fălăuş 2014, among others, and 2.2.3). It has been argued that certain epistemic indefinites (e.g. \textit{irgendein, un }NP \textit{qualunque, un }NP\textit{ oarecare}) can sustain TOTAL variation, requiring that \textit{all }relevant alternatives in the domain of quantification qualify as \textit{equally} possible options. For example, the sentence in , with the Romanian epistemic indefinite \textit{un }NP\textit{ oarecare}, states that Maria has the obligation to work with a colleague and that this could be any colleague:
\end{styleStandard}

\begin{styleStandard}
\label{bkm:Ref533761905}\ \ Maria\ \ \textit{trebuie}\ \ să\ \ lucreze\ \ cu\ \ \textit{un}\ \ coleg\ \ \textbf{\textit{oarecare.}}
\end{styleStandard}

\begin{styleStandard}
\ \ Maria\ \ must\ \ \textsc{subj}\ \ work\ \ with\ \ a\ \ colleague\ \ \textsc{oarecare}
\end{styleStandard}

\begin{styleStandard}
\ \ ‘Maria must work with a colleague/ any colleague.’
\end{styleStandard}

\begin{styleStandard}
\ \ (Fălăuş 2014: 32)
\end{styleStandard}

\begin{styleStandard}
In contrast to this total freedom of choice, there are epistemic indefinites that trigger a weaker inference—partial variation—\textit{some, but not necessarily all }alternatives in the relevant domain are epistemic possibilities. As such, they are compatible with the exclusion of some of the possible options. The difference can be most readily observed in the scenario in , as was shown first by Alonso-Ovalle and Menéndez-Benito (2010) and restated by Fălăuş 2014:
\end{styleStandard}

\begin{styleStandard}
\label{bkm:Ref533761939}\ \ Maria, Juan, and Pedro are playing hide-and-seek in their country house. Juan is hiding. 
\end{styleStandard}

\begin{styleStandard}
Pedro believes that Juan is inside the house, rather than in the bathroom or in the kitchen.
\end{styleStandard}

\begin{styleStandard}
\ \ a.\ \ Juan \textit{tiene} que estar en \textbf{\textit{alguna}} habitación de la casa.
\end{styleStandard}

\begin{styleStandard}
\ \ b.\ \ Juan \textit{trebuie} să fie în \textbf{\textit{vreo}} cameră din casă.
\end{styleStandard}

\begin{styleStandard}
\ \ \ \ ‘Juan must be in a room of the house.’
\end{styleStandard}

\begin{styleStandard}
\ \ \ \ (Fălăuş 2014: 32)
\end{styleStandard}

\begin{styleStandard}
The context makes it clear that not all rooms of the house are possible choices, that is, the variation associated with the indefinite is limited to a subset of elements in the relevant domain. While the use of \textit{algún }is perfectly acceptable in this scenario, a total variation item like \textit{un }NP\textit{ oarecare/qualunque/cualquiera} would be deviant (see Fălăuş 2014).
\end{styleStandard}

\begin{styleStandard}
\ \ In contrast with \textit{cualquiera}, \textit{algún} requires non-exhaustive variation. The oddness in the example in \ illustrates a presupposition conflict. The two anti-specific indefinites \textit{algún} and \textit{cualquiera} require opposite things: the former requires non-exhaustive variation, while the latter (the free choice indefinite) requires exhaustive variation (but see Section 2.3 for grammatical instances of \textit{algún N cualquiera} with evaluative interpretation of \textit{cualquiera}):
\end{styleStandard}

\begin{styleStandard}
\label{bkm:Ref533761963}\ \ \# Juan ha hablado con algún estudiante cualquiera.
\end{styleStandard}

\begin{styleStandard}
\ \ Juan has talked with some student any 
\end{styleStandard}

\begin{styleStandard}
\ \ (Etxeberria \& Giannakidou 2017: 19)
\end{styleStandard}

\begin{styleStandard}
N \textit{cualquiera} is anti-specific and referentially vague because it is not used referentially and the speaker does not refer to a specific individual (see Giannakidou \& Quer 2013, among others). Menéndez-Benito (2010) assumes that \textit{cualquiera} occurs with nouns that do not denote singletons. The anti-singleton restriction follows from the FCI reading of \textit{cualquiera}, which compares at least two alternatives and conveys that all alternatives count as possible or that the agent does not discriminate between at least two alternatives for his choice. This explanation can explain why \textit{cualquiera} as an FCI does not occur with definite nouns that denote unique individuals like proper nouns do. 
\end{styleStandard}

\begin{styleStandard}
\ \ Summary of \textit{cualquiera/algún}: both are anti-specific indefinites and both express referential vagueness. They signal the lack of referential intent that characterises specificity. The indefinite \textit{cualquiera} expresses total variation, meaning that all options are evaluated epistemically, whereas in case of \textit{algún} only a subset of possibilities are considered as candidates. 
\end{styleStandard}

\subsubsection[2.2.2 Formal analysis of free choice]{\textbf{2.2.2 Formal analysis of free choice}}
\begin{styleStandard}
This subsection introduces the formal analysis of free choice Indefinites within the Alternative Semantic Framework.
\end{styleStandard}

\paragraph[2.2.2.1 Preliminaries: An introduction to alternative semantics]{\textit{2.2.2.1 Preliminaries: An introduction to alternative semantics}}
\begin{styleStandard}
In order to show how \textit{cualquiera} can be analysed in alternative semantics, I will give a short introduction to the latter based on Fălăuş (2014).
\end{styleStandard}

\begin{styleStandard}
\ \ On this approach, the semantics of FCIs involves considering alternative semantic values. The theory was developed in a framework of bidimensional semantics (cf. Rooth 1985), where alternatives are introduced and computed separately from regular semantic values. The starting point is the interpretation of an exchange such as the following:
\end{styleStandard}

\begin{styleStandard}
\ \ A: Who came to the party? 
\end{styleStandard}

\begin{styleStandard}
B: Paul and Sue.
\end{styleStandard}

\begin{styleStandard}
This utterance conveys that Paul and Sue are the only individuals in the context for which a certain property holds (here: having come to the party), a meaning that goes beyond what is explicitly uttered. There are two components leading to this enriched, non-literal interpretation: (i) a set of alternatives (e.g. contextually relevant individuals) and (ii) the exclusion of the (non-entailed) alternatives. This enriched or strengthened meaning, the one obtained via the consideration and exclusion of alternatives, is often triggered by focus.
\end{styleStandard}

\begin{styleStandard}
\ \ Free choice items also activate alternatives. A phrase such as \textit{any student} activates domain alternatives, that is, subsets of (contextually relevant) students (e.g \{Mario\}, \{Marco\}, etc.). Free choice items are indefinites (i.e. existentially quantified elements) that systematically activate sets of alternatives that expand (compositionally) until they find a (quantificational) operator that selects them (e.g. universal quantifier, existential quantifier, negation, question operator) as shown in (83). Alternatives can be of different types (e.g. individual, propositional) and as a consequence are accessible to both sentential and non-sentential (generalised) quantificational operators, which eventually determine the interpretation of the indefinite. Individual alternatives can expand and give rise to alternatives of a higher type, e.g. propositional alternatives, which then get caught by a sentential operator ${\exists}$p or ${\forall}$p, such as the ones in \ (see Kratzer \& Shimoyama 2002, Chierchia 2006, Aloni 2007, Alonso-Ovalle \& Menéndez-Benito 2011, among others):
\end{styleStandard}

\begin{styleStandard}
\label{bkm:Ref533761448}\ \ Indefinites in alternative semantics after Kratzer and Shimoyama (2002: 8)
\end{styleStandard}

\begin{styleStandard}
\ \ For [[$\alpha $]]\textsuperscript{w,g }${\subseteq}$ D\textsubscript{{\textless}s,t{\textgreater}}:
\end{styleStandard}

\begin{styleStandard}
\ \ a.\ \ [[$\alpha $ ]]\textsuperscript{w,g}=\{$\lambda $w[2032?].${\exists}$p [p${\in}$[[$\alpha $]]\textsuperscript{w,g} and p(w[2032?]) = 1]\}
\end{styleStandard}

\begin{styleStandard}
\ \ b.\ \ [[$\alpha $ ]]\textsuperscript{w,g}=\{$\lambda $w[2032?].${\forall}$p [p${\in}$[[$\alpha $]]\textsuperscript{w,g} → p(w[2032?]) = 1]\}
\end{styleStandard}

\begin{styleStandard}
The contribution of the indefinite remains constant, namely, to provide a set of alternatives, its different readings being the result of associating with different operators.
\end{styleStandard}

\begin{styleStandard}
\ \ When a set of propositional alternatives combines with the existential operator as defined in , it yields the proposition that is true in case at least one alternative in the relevant set is true. In contrast, the universal quantifier in \ yields the proposition that is true in case every alternative in the relevant set is true. 
\end{styleStandard}

\begin{styleStandard}
\ \ Now, importantly, as will be shown below in Section 2.2.2.2, \textit{cualquiera} provides alternatives that always combine with the universal operator ${\forall}$p in modal contexts (i.e. under modal verbs, conditionals, subtrigging, and other modalised constructions). Let us see how this approach works for \textit{cualquier libro} ‘some book/any book’ in . The domain alternatives of \textit{cualquier libro }consists of all individual books of a contextually relevant set, for instance \{\textit{Don Quijote}, \textit{El camino}, …\}. These alternatives grow to become propositions \{You can buy \textit{Don Quijote}, You can buy \textit{El camino}, …\}:
\end{styleStandard}

\begin{styleStandard}
\label{bkm:Ref533761676}\ \ Puedes tomar cualquier libro. 
\end{styleStandard}

\begin{styleStandard}
\ \ ‘You can take any book.’
\end{styleStandard}

\begin{styleStandard}
\ \ (@ 32)
\end{styleStandard}

\begin{styleStandard}
In order to derive the free choice Meaning in , that every alternative (i.e. every book) is a possible option and not just some of them, we compute the enriched meaning of the sentence in , by interpreting the exhaustivity expressed by \textit{cualquier}, i.e. we negate all alternatives that do not match the assertion, which yields the following enriched meaning (‘All alternatives must be taken into account’). The enriched meaning does not contain the statement according to which only one alternative can be taken into account (represented by ${\perp}$ which basically means not-computable):
\end{styleStandard}

\begin{styleStandard}
\label{bkm:Ref533761693}\ \ Total exhaustivity of alternatives (a ${\vee}$ b ${\vee}$ c) =
\end{styleStandard}

\begin{styleStandard}
\ \ \ \ You can take (a ${\vee}$ b ${\vee}$ c) ${\wedge}$
\end{styleStandard}

\begin{styleStandard}
\ \ \ \ ¬ You can take (a ${\wedge}$ b) ${\wedge}$
\end{styleStandard}

\begin{styleStandard}
\ \ \ \ ¬ You can take (a ${\wedge}$ c) ${\wedge}$
\end{styleStandard}

\begin{styleStandard}
\ \ \ \ ¬ You can take (b ${\wedge}$ c) ${\wedge}$
\end{styleStandard}

\begin{styleStandard}
\ \ \ \ ¬ You can take (a ${\vee}$ b) ${\wedge}$
\end{styleStandard}

\begin{styleStandard}
\ \ \ \ ¬ You can take (b ${\vee}$ c) ${\wedge}$
\end{styleStandard}

\begin{styleStandard}
\ \ \ \ ¬ You can take (a ${\vee}$ c) ${\wedge}$
\end{styleStandard}

\begin{styleStandard}
\ \ \ \ ¬ You can take a ${\wedge}$ ¬ You can take b ${\wedge}$ ¬ You can take c = ${\perp}$
\end{styleStandard}

\begin{styleStandard}
To sum up, I have discussed so far that \textit{cualquiera} triggers alternatives and total exhaustification over the alternatives in modal contexts (see 2.2). In the following section, I will show one possible formalisation of FCIs in Alternative Semantics.
\end{styleStandard}

\paragraph[2.2.2.2 Aloni’s (2007) formalisation of FCIs]{\textit{2.2.2.2 Aloni’s (2007) formalisation of FCIs}}
\begin{styleStandard}
According to Aloni (2007), Aloni \& et al. (2013), and Aloni (2021), FCIs trigger the application of the total exhaustification over alternatives as demonstrated in \ and represented formally by an exhaustivity operator \textbf{EXH} (see also Menéndez-Benito 2010). This operator has a similar function as the focus adverb \textit{only} which associates with alternatives and expresses that only a certain alternative is true. The SHIFT\footnote{\ Shift means here type shifting, i.e. conversion of one type into another type, e.g. from a type of individuals into a type of propositions. } function expresses that the exhaustification of domain alternatives shifts to the exhaustification of propositional alternatives, that is, for \ it is possible that only Peter falls or it is possible that only Mary falls. The total exhaustification expresses that there is no alternative proposition that is left out by the exhaustification:
\end{styleStandard}

\begin{styleStandard}
\label{bkm:Ref150778689}\ \ Anyone can fall. (Aloni et al. 2011)
\end{styleStandard}

\begin{styleStandard}
\ \ a.\ \ [${\forall}$](${\lozenge}$(\textsc{SHIFT}‹\textsubscript{s,t}›\textsubscript{ }(\textbf{EXH}[anyone, fall])))
\end{styleStandard}

\begin{styleStandard}
\ \ b.\ \ ${\lozenge}$ only Peter fall\ \ ${\lozenge}$ only Mary fall\ \ ${\lozenge}$ only Bob fall\ \ ...
\end{styleStandard}

\begin{styleStandard}
\ \ \ \ (Aloni et al. 2011)
\end{styleStandard}

\begin{styleStandard}
She accounts for the distribution of FCIs under modals such as in \ and negation as in , as well as in contexts that also recover universal quantification of exhaustified propositional alternatives without always expressing an overt modal verb such as subtrigging in \ and degree comparatives as in :
\end{styleStandard}

\begin{styleStandard}
\label{bkm:Ref63544535}\ \ \textit{Canonical FCI}
\end{styleStandard}

\begin{styleStandard}
\ \ a.\ \ Sentence: You can bring any book.
\end{styleStandard}

\begin{styleStandard}
\ \ b.\ \ Logical form: [${\forall}$](${\lozenge}$\textsc{SHIFT}‹\textsubscript{s,t}› (\textbf{EXH}[any book, $\lambda $x. you bring me x])))
\end{styleStandard}

\begin{styleStandard}
\ \ c.\ \ Predicted meaning: For each book it is possible that you bring me only that book.
\end{styleStandard}

\begin{styleStandard}
\label{bkm:Ref63544537}\ \ \textit{Embedding of the FC effect}
\end{styleStandard}

\begin{styleStandard}
\ \ a.\ \ Sentence: You cannot bring any book.
\end{styleStandard}

\begin{styleStandard}
\ \ b.\ \ Logical form: [¬][${\forall}$](${\lozenge}$\textsc{SHIFT}‹\textsubscript{s,t}› (\textbf{EXH}[any book, $\lambda $x. you bring me x])))
\end{styleStandard}

\begin{styleStandard}
\ \ c.\ \ Predicted meaning: You cannot freely choose which book you bring me.
\end{styleStandard}

\begin{styleStandard}
It is not the case that for each book you can bring me that book.
\end{styleStandard}

\begin{styleStandard}
\label{bkm:Ref54857683}\ \ \textit{Licensing by subtrigging (see Aloni 2007)}
\end{styleStandard}

\begin{styleStandard}
\ \ a.\ \ Sentence: Anyone who tried to jump fell.
\end{styleStandard}

\begin{styleStandard}
\ \ b.\ \ Logical form: [${\forall}$](↓\textsc{SHIFT}\textsubscript{e}(\textbf{EXH}[anyone, who tried to jump]) fell)
\end{styleStandard}

\begin{styleStandard}
\ \ c.\ \ Predicted meaning: All persons who tried to jump fell.
\end{styleStandard}

\begin{styleStandard}
\label{bkm:Ref54857674}\ \ \textit{FCIs in Comparatives (see Aloni 2021)}
\end{styleStandard}

\begin{styleStandard}
\ \ a.\ \ Sentence: John is taller than any girl.
\end{styleStandard}

\begin{styleStandard}
\ \ b.\ \ Logical form: [${\forall}$](\textsc{SHIFT}\textsubscript{e}(\textbf{EXH}[d, $\lambda $d.tall (j,d)]) {\textgreater} \textsc{SHIFT}\textsubscript{e}(\textbf{EXH}[d, $\lambda $d.tall (any girl, d)])
\end{styleStandard}

\begin{styleStandard}
\ \ c.\ \ Predicted meaning: For all girls x, John is taller than x.
\end{styleStandard}

\begin{styleStandard}
\ \ \ \ (Aloni et al. 2011)
\end{styleStandard}

\begin{styleStandard}
Aloni (2007, 2021) and Aloni et al. (2013) also account for the ungrammaticality of FCIs such as \textit{any} in episodic contexts in \ that refer to some factual events by showing that the exhaustivity operator cannot apply there. This is represented as [23CA?] in . 
\end{styleStandard}

\begin{styleStandard}
\label{bkm:Ref517204}\ \ \textit{Ruling out FCIs in episodic contexts }
\end{styleStandard}

\begin{styleStandard}
\ \ a.\ \ Sentence: \# Anyone fell.
\end{styleStandard}

\begin{styleStandard}
\ \ b.\ \ Logical form: [${\forall}$]((\textsc{SHIFT}‹\textsubscript{s,t}› (\textbf{EXH}[anyone, fell]))
\end{styleStandard}

\begin{styleStandard}
\ \ c.\ \ Predicted meaning: [23CA?]
\end{styleStandard}

\begin{styleStandard}
\ \ \ \ (Aloni 2021)
\end{styleStandard}

\begin{styleStandard}
Aloni argues that the sentence in \ would give the following incorrect interpretation: for each person only that person fell. A past event described by the perfective aspect can either express that everyone fell or that a certain person/group of people fell. A modal verb as in \textit{Anyone can fall} rescues the interpretation, as it shifts the exaustivity to different worlds (i.e. it becomes interpretable as ‘For every person it is possible that this person fell’) and not only the actual word (see Menéndez-Benito 2010, Aloni 2021, among others). This is how these authors account for the relation between modality and free choice indefinites.
\end{styleStandard}

\begin{styleStandard}
\ \ Although Aloni makes a correct prediction for FCIs such as English \textit{any} and the determiner/pronoun\textit{ cualquiera} in European Spanish, which need a modalised context to be licensed and are not permitted in episodic contexts (see 2.1), her analysis does not account for the distribution of (postnominal) \textit{cualquiera} in episodic contexts and with copular verbs in present or past tense in European and Argentinian Spanish (see Sections 2.1 and 2.2.1). The question is how the distribution of \textit{cualquiera} should be accounted for in a system where \textit{cualquiera} triggers alternatives and total exhaustivity over these alternatives. I will deal with the current state research on postnominal \textit{cualquiera} in episodic contexts in Sections 2.3.
\end{styleStandard}

\paragraph[2.2.2.3 Summary]{\textit{2.2.2.3 Summary}}
\begin{styleStandard}
It has been shown that, according to the current state of research, FCIs such as \textit{any} or \textit{cualquiera} on its free choice meaning have a universal interpretation if embedded under a possibility modal and other modalised constructions such as subtrigging that recover a possibility modal. According to Aloni (2007), FCIs trigger total exhaustification over alternatives represented by an exhaustivity operator EXH. This operator associates with a universal quantifier that operates on propositions. In this way Aloni can account for the universal inference of FCIs, which is a conventional implicature due to the impossibility of cancelling under negation.
\end{styleStandard}

\subsubsection[2.2.3 Semantic analysis of cualquiera in episodic contexts ]{\textbf{2.2.3 Semantic analysis of }\textbf{\textit{cualquiera}}\textbf{ in episodic contexts}\textbf{ }}
\paragraph[2.2.3.1 Informal description]{\textit{2.2.3.1 Informal description}}
\begin{styleStandard}
We saw in Section 2.2.1 that \textit{cualquiera} is supposed to be excluded from episodic contexts with a free choice interpretation. However, it can express agent-related modality, such as the random choice interpretation. According to Choi (2007), Choi \& Romero (2008), and A\&M (2011/2018), postnominal \textit{cualquiera} used in episodic contexts under volitional verbs has an RC interpretation:
\end{styleStandard}

\begin{styleStandard}
\ \ Juan\ \ necesitaba\ \ un\ \ pisapapeles,\ \ de\ \ modo\ \ que\ \ cogió\ \ un\ \ libro\ \ cualquiera
\end{styleStandard}

\begin{styleStandard}
\ \ Juan\ \ needed\ \ a\ \ paperweight\ \ of\ \ way\ \ that\ \ he-took\ \ a\ \ book\ \ \textsc{cualquiera}
\end{styleStandard}

\begin{styleStandard}
\ \ de\ \ la\ \ estantería\ \ y\ \ lo\ \ puso\ \ encima\ \ de\ \ la\ \ pila.\ \ 
\end{styleStandard}

\begin{styleStandard}
\ \ from\ \ the\ \ shelf\ \ and\ \ it\ \ he-put\ \ on.top\ \ of\ \ the\ \ pile\ \ 
\end{styleStandard}

\begin{styleStandard}
\ \ ‘John needed a paperweight, so he took a random book from the shelf and put it on top of the pile’
\end{styleStandard}

\begin{styleStandard}
\ \ (Choi \& Romero 2008: 97)
\end{styleStandard}

\begin{styleStandard}
A\&M (2011, 2018) assume that postnominal \textit{cualquiera} is ambiguous between RC and EVAL (but see Choi \& Romero 2008, who do not mention the ambiguity):
\end{styleStandard}

\begin{styleStandard}
\label{bkm:Ref533953640}\ \ Juan compró un libro cualquiera. ‘Juan bought a book.’
\end{styleStandard}

\begin{styleStandard}
\ \ a.\ \ John bought a book, and he choose it indiscriminately. (Random Choice)
\end{styleStandard}

\begin{styleStandard}
\ \ b.\ \ John bought a book, and the book is not special or remarkable (according to the speaker). (Unremarkable reading, EVAL)
\end{styleStandard}

\begin{styleStandard}
\ \ \ \ (A\&M 2018: 2)
\end{styleStandard}

\begin{styleStandard}
In order to distinguish EVAL from RC, A\&M (2011, 2018) have created different contexts:
\end{styleStandard}

\begin{styleStandard}
\label{bkm:Ref54686706}\ \ Scenario. Juan went to the bookstore. He wanted to buy \textit{The Unbearable Lightness of Being}, and did so. I don’t think this book is special in any way.
\end{styleStandard}

\begin{styleStandard}
\label{bkm:Ref533953745}\ \ Scenario. Juan went to the bookstore, and bought a book at random. The book turned out to be \textit{The Unbearable Lightness of Being}. I think this book is remarkable.
\end{styleStandard}

\begin{styleStandard}
In the scenario in [], I can truthfully utter [] on its unremarkable interpretation (since I consider the book that Juan bought unremarkable) but not on its random choice interpretation (as Juan didn’t make an indiscriminate choice). In contrast, my utterance of ] in the scenario in [] would be false on the unremarkable interpretation but true on the random choice interpretation. 
\end{styleStandard}

\begin{styleStandard}
(A\&M 2018: 2f.)
\end{styleStandard}

\begin{styleStandard}
According to A\&M (2018), the noun modified by \textit{cualquiera} must represent an object that the agent of the sentence must have chosen indiscriminately on the RC interpretation. They assume that \textit{cualquiera} appears in object position of volitional verbs on the RC interpretation. A\&M (2011) assume that the RCI is restricted (e.g. not possible with unaccusative verbs), whereas EVAL is always possible:
\end{styleStandard}

\begin{styleStandard}
While the unremarkable interpretation is always possible, the random choice interpretation has a restricted distribution. First of all, for this interpretation to obtain, the sentence has to describe a volitional event (Choi \& Romero 2008) and \textit{uno cualquiera} should not be in the subject position of agentive verbs. If that condition is not met, \textit{uno cualquiera}\footnote{\textrm{\ A\&M refer to }\textsc{un}\textrm{\textit{ }}\textrm{N}\textrm{\textit{ cualquiera}}\textrm{ as }\textrm{\textit{uno cualquiera}}\textrm{.}} only has the unremarkable interpretation. Copular sentences are a clear case. The sentences in [(91)] convey that the subject is unremarkable within the class of individuals denoted by the noun phrase: (a) says that Juan is a regular clerk and (b) that Juan is a student just like the rest.
\end{styleStandard}

\begin{styleStandard}
\ \ a.\ \ Juan es un oficinista cualquiera. ‘Juan is an unremarkable clerk.’
\end{styleStandard}

\begin{styleStandard}
\ \ b.\ \ Juan es un estudiante cualquiera. ‘Juan is an unremarkable student.’
\end{styleStandard}

\begin{styleStandard}
\ \ \ \ (A\&M 2018: 3)
\end{styleStandard}

\begin{styleStandard}
The following examples only have the ‘unremarkable’ reading (=Eval) and not RC, because they do not contain volitional verbs:
\end{styleStandard}

\begin{styleStandard}
\ \ a.\ \ Ayer\ \ Juan\ \ tropezó\ \ con\ \ un\ \ objeto\ \ cualquiera.
\end{styleStandard}

\begin{styleStandard}
\ \ \ \ yesterday\ \ Juan\ \ stumbled\ \ with\ \ an\ \ object\ \ \textsc{cualquiera}
\end{styleStandard}

\begin{styleStandard}
\ \ \ \ ‘Yesterday, Juan stumbled on an unremarkable object.’ 
\end{styleStandard}

\begin{styleStandard}
\ \ \ \ (A\&M 2018: 3)
\end{styleStandard}

\begin{styleStandard}
\ \ b.\ \ La\ \ levadura\ \ destrozó\ \ un\ \ molde\ \ cualquiera.
\end{styleStandard}

\begin{styleStandard}
\ \ \ \ The\ \ yeast\ \ destroyed\ \ a\ \ baking pan\ \ \textsc{cualquiera}
\end{styleStandard}

\begin{styleStandard}
\ \ \ \ ‘The yeast destroyed an unremarkable baking pan.’
\end{styleStandard}

\begin{styleStandard}
\ \ \ \ (A\&M 2018: 3)
\end{styleStandard}

\begin{styleStandard}
Nouns modified by postnominal \textit{cualquiera} that represent subjects of volitional verbs such as \textit{hablar} ‘speak’ or objects of non-volitional verbs cannot have RC as the authors assume, but only Eval:
\end{styleStandard}

\begin{styleStandard}
\ \ Habló un estudiante cualquiera.
\end{styleStandard}

\begin{styleStandard}
\ \ ‘An unremarkable student spoke.’ (\# ‘a random student spoke’)
\end{styleStandard}

\begin{styleStandard}
\ \ (A\&M 2018: 4)
\end{styleStandard}

\begin{styleStandard}
Furthermore, A\&M (2011) assume that \textit{cualquiera} can also have a free choice interpretation under some modals, which they call harmonic interpretation (HI): 
\end{styleStandard}

\begin{styleStandard}
When \textit{uno cualquiera} is embedded under some modals, another possibility arises: \textit{uno cualquiera} introduces a distribution effect with respect to the worlds that the modal ranges over. For instance, \textit{Coge una carta cualquiera!} (‘take any card!’) can be interpreted as conveying that any card is a permitted option. 
\end{styleStandard}

\begin{styleStandard}
(A\&M 2018: 1)
\end{styleStandard}

\begin{styleStandard}
A\&M (2011) assume that \textit{cualquiera} is lexically ambiguous between RCI/HI on the one hand and EVAL on the other:
\end{styleStandard}

\begin{styleStandard}
As a working hypothesis, we will assume that \textit{uno cualquiera} is ambiguous between one form that can give rise to the random choice and harmonic interpretations (the form that we will focus on), and another that conveys the unremarkable interpretation. To determine whether this hypothesis is tenable it would be necessary to provide a detailed analysis of the unremarkable interpretation, a non-trivial task that we will not be able to undertake here.
\end{styleStandard}

\begin{styleStandard}
(A\&M 2018: 8)
\end{styleStandard}

\begin{styleStandard}
To sum up, A\&M (20011, 2018) show that RCI is restricted to objects of volitional verbs and it does not arise under non-volitional verbs. The ‘unremarkable’ reading is completely unrestricted; it can appear everywhere:
\end{styleStandard}

\begin{center}
\tablefirsthead{\hline
\textbf{Interpretation} &
\multicolumn{2}{m{3.7649598in}|}{\centering \textbf{Restrictions}}\\\hline
 &
\multicolumn{1}{m{1.8427598in}|}{\centering Available} &
\centering\arraybslash Unavailable\\\hline
Random choice &
\multicolumn{1}{m{1.8427598in}|}{object of volitional verbs} &
subject of volitional verbs

non-volitional verbs\\\hline
Unremarkable &
\multicolumn{2}{m{3.7649598in}|}{\centering unrestricted}\\}
\tablehead{\hline
\textbf{Interpretation} &
\multicolumn{2}{m{3.7649598in}|}{\centering \textbf{Restrictions}}\\\hline
 &
\multicolumn{1}{m{1.8427598in}|}{\centering Available} &
\centering\arraybslash Unavailable\\\hline
Random choice &
\multicolumn{1}{m{1.8427598in}|}{object of volitional verbs} &
subject of volitional verbs

non-volitional verbs\\\hline
Unremarkable &
\multicolumn{2}{m{3.7649598in}|}{\centering unrestricted}\\}
\tabletail{}
\tablelasttail{}
\begin{supertabular}{|m{1.8427598in}|m{1.8427598in}m{1.8434598in}|}

\end{supertabular}
\end{center}
\begin{styletablecaption}
Table \stepcounter{Table}{\theTable}: The distribution of RCI and EVAL of \textsc{un}\textit{ }N \textit{cualquiera} (A\&M 2011, 2018).
\end{styletablecaption}

\begin{styleStandard}
Under their approach, the restriction to volitional verbs follows from the modal component of \textsc{un} N \textit{cualquiera }(cf. 2.1.3). The modal component of \textsc{un} N \textit{cualquiera} expresses that every alternative of \textsc{un} N \textit{cualquiera} conforms to the agent’s decision to choose this alternative. Thus, RCI can only apply to volitional verbs, as they express the agent’s decisions to do something (e.g. picking some book).
\end{styleStandard}

\begin{styleStandard}
\ \ The random reading disappears whenever there is nothing for the agent to choose, as for example with unaccusative verbs, because unaccusative verbs do not match with the decision condition according to which the agent does something in accordance to his/her decision (see A\&M 2011, 2018).
\end{styleStandard}

\begin{styleStandard}
\ \ A\&M (2018) assume lexical ambiguity between RCI/HI and EVAL and that EVAL is unrestricted. This hypothesis needs to be tested (see Chapter 3).
\end{styleStandard}

\begin{styleStandard}
\textbf{\ \ }Chierchia (2006: 574) observes yet another interpretation of Free Choice indefinites in episodic contexts. He notes that the existential FCI \textsc{un} N \textit{qualsiasi} in Italian can be used in episodic contexts under a covert modal operator that expresses speaker’s modality, namely speaker’s ignorance (see also 2.2.1 for the ignorance reading expressed by Spanish \textit{algún}):
\end{styleStandard}

\begin{styleStandard}
\label{bkm:Ref533762757}\ \ Gianni è uscito di corsa e non sapendo che fare, ha bussato ad una porta qualsiasi.
\end{styleStandard}

\begin{styleStandard}
\ \ ‘Gianni ran out and not knowing what to do, knocked on any door.’
\end{styleStandard}

\begin{styleStandard}
Speaker’s ignorance: it is possible that for every door x Gianni knocked at the door x.
\end{styleStandard}

\begin{styleStandard}
\ \ (Chierchia 2006: 574)
\end{styleStandard}

\begin{styleStandard}
Note that free choice relatives in English can also express an ignorance reading (IGN) in episodic context (see Condoravdi 2015: 216), that is, in (95) the speaker is ignorant about the identity of the book John bought:
\end{styleStandard}

\begin{styleStandard}
\ \ Whatever book John bought was very expensive.
\end{styleStandard}

\begin{styleStandard}
According to Chierchia the sentence in \ could be interpreted as something like “it follows from what the speaker knows that Gianni knocked at a door.” The FCI implicature would then be “it is consistent with what the speaker knows that any door might have been the one knocked at” (Chierchia 2006: 574). Chierchia thus assumes that the speaker considers it a possibility that for every door x, Gianni knocked at the door x. It is important to mention that the modality is associated with the speaker and not with the agent \textit{Gianni}, according to Chierchia (2006) (see 2.3 for agent’s modality). 
\end{styleStandard}

\begin{styleStandard}
\ \ Recall from the Section 2.1 that native speakers agree on the possibility of an ignorance interpretation of \textsc{un} N \textit{cualquiera} in episodic contexts (see , repeated here in ):
\end{styleStandard}

\begin{styleStandard}
\label{bkm:Ref72400981}\textbf{\ \ }La carta de ajuste, en la pantalla, se complementaba con la música ramplona de una zarzuela cualquiera 
\end{styleStandard}

\begin{styleStandard}
\ \ i. \ \ ‘The testcard, in the screen, was complemented with the vulgar sound of some ordinary zarzuela’ \ \ (EVAL)
\end{styleStandard}

\begin{styleStandard}
\ \ ii.\ \ ‘…with the vulgar sound of some zarzuela that I don’t know’ (IGN)
\end{styleStandard}

\begin{styleStandard}
\ \ (@ 14e)
\end{styleStandard}

\begin{styleStandard}
The ignorance interpretation is not mentioned in the literature on Spanish \textit{cualquiera} (see Kellert \& Francia in press, on the ignorance interpretation on a similar item in ArgSp: \textit{cualunque}). In this book, I do not pursue the study on the ignorance reading in episodic contexts of \textit{cualquiera} either. 
\end{styleStandard}

\paragraph[2.2.3.2 Formal Analysis of the Random Choice interpretation]{\textit{2.2.3.2 Formal Analysis of the Random Choice interpretation}}
\begin{styleStandard}
The random choice interpretation, RCI, has been described very well in the literature. I will illustrate two prominent proposals, which both assume that postnominal \textit{cualquiera} has a RCI whenever \textit{cualquiera} occurs under volitional verbs (see Choi \& Romero 2008, A\&M 2018). As I will show, neither approach accounts for the evaluative function of postnominal \textit{cualquiera}.
\end{styleStandard}

\begin{styleStandard}
Choi \& Romero (2008) assume that \textit{cualquiera} has a similar denotation as von Fintel’s (2000) denotation of \textit{wh-ever} in English. Von Fintel (2000) adopts Dayal’s (1997) insight that \textit{\nobreakdash-ever} Free Relatives (FRs) introduce a layer of quantification over possible worlds and proposes that \nobreakdash-\textit{ever} in \nobreakdash-\textit{ever} FRs induces a presupposition of variation on either counterfactual worlds or epistemic worlds. Von Fintel (2000) points out that a subtype of \textit{\nobreakdash-ever} FRs expresses “indifference” on somebody’s part. Compare \ and . Both of them assert the same proposition paraphrasable using a definite description, namely, the proposition that the person who was at the top of the ballot won the election yesterday (see Choi \& Romero 2008: 88).
\end{styleStandard}

\begin{styleStandard}
\label{bkm:Ref533762240}\ \ a.\ \ In yesterday’s election, who was at the top of the ballot won.
\end{styleStandard}

\begin{styleStandard}
\ \ b.\ \ In yesterday’s election, whoever was at the top of the ballot won.
\end{styleStandard}

\begin{styleStandard}
Unlike a., b. conveys an extra meaning triggered by \textit{\nobreakdash-ever}, such that the identity of who was at the top of the ballot did not matter to winning yesterday’s election. In the sense that the identity of the denotation of \textit{\nobreakdash-ever} FRs does not matter for the general nature or outcome of the election.
\end{styleStandard}

\begin{styleStandard}
\ \ In von Fintel (2000), this essential link follows from the presupposition of variation given in , which is identified as the nature of \textit{\nobreakdash-ever}’s contribution. The presupposition of variation tells us that if the individual denoted by an \textit{\nobreakdash-ever} FR had been different, the truth value of the assertion in the actual world would still be valid in all the counterfactual worlds.
\end{styleStandard}

\begin{styleStandard}
\label{bkm:Ref533762335}\ \ Presupposition of variation for b: If the person who was at the top of the ballot had been different, the same thing would have happened: that (new) person would have won.
\end{styleStandard}

\begin{styleStandard}
That is, the presupposition triggered by “whoever was on the top of the ballot won”\textit{ }conveys that if the person who won the election had been different in all the counterfactual worlds, the truth of the proposition “the person at the top of the ballot won” would also hold in the counterfactual worlds. From this presupposition of variation, it is inferred that regardless of who was at the top of the ballot, “being at the top of the ballot” and “winning yesterday’s election” are in an essential relation.
\end{styleStandard}

\begin{styleStandard}
\ \ In von Fintel (2000), a sentence containing an \textit{\nobreakdash-ever} FR is formalised as in . In the formulae, F indicates the modal base for \textit{\nobreakdash-ever} FRs, which is a set of worlds on which the presupposition of variation operates. P refers to the denotation of the NP property contained in the \textit{\nobreakdash-ever} FR, and Q refers to the property expressed by the rest of the sentence. Sentences containing an \textit{\nobreakdash-ever} FR assert that the thing that has P is Q in the actual world, as shown in . Note that, according to von Fintel (2000), \textit{{}-ever} FRs denote definite descriptions on the assertion level. This is represented by the iota-Operator $\iota $. The presupposition triggered by \textit{\nobreakdash-ever} says that in all worlds (of the corresponding modal base) that are different from the actual world only with respect to the referent of the \textit{\nobreakdash-ever} FR, the asserted proposition has in w[2032?] whatever truth value it has in the actual world w\textsubscript{0}. (see Choi \& Romero 2008: 88):
\end{styleStandard}

\begin{styleStandard}
\label{bkm:Ref533762404}\ \ Whatever (w\textsubscript{0})(F)(P)(Q) according to von Fintel (2000)
\end{styleStandard}

\begin{styleStandard}
\ \ a.\ \ \textbf{Asserts}: Q(w\textsubscript{0})($\iota $ x.P(w\textsubscript{0}) (x))
\end{styleStandard}

\begin{styleStandard}
\ \ b.\ \ \textbf{Presupposes}: ${\forall}$ w[2032?] ${\in}$ min\textsubscript{w0} [F ${\cap}$ $\lambda $ w[2033?]. $\iota $ x.P(w[2033?])(x) ${\neq}$ $\iota $ x.P(w\textsubscript{0})(x)]: Q (w[2032?]) ($\iota $~x.P(w[2032?])(x)) = Q (w\textsubscript{0}) ($\iota $ x.P(w\textsubscript{0})(x))
\end{styleStandard}

\begin{styleStandard}
Choi \& Romero apply the formalisation\textstyleannotationreference{ }in \ to the sentence in . The modal base F is counterfactual, and thus a presupposition of counterfactual variation is conveyed, as in (b).
\end{styleStandard}

\begin{styleStandard}
\label{bkm:Ref150779072}\ \ In yesterday’s election, whoever was at the top of the ballot won.
\end{styleStandard}

\begin{styleStandard}
\ \ a.\ \ \textbf{Assertion}: $\lambda $ w\textsubscript{0}. win($\iota $ y.top-of-ballot(y,w\textsubscript{0}), w\textsubscript{0})
\end{styleStandard}

\begin{styleStandard}
\ \ b.\ \ \textbf{Presupposition}: $\lambda $ w\textsubscript{0 }${\forall}$ w[2032?] ${\in}$ min\textsubscript{w0} [F ${\cap}$ $\lambda $ w[2033?]. [$\iota $ y.top-of-ballot(y,w[2033?] ${\neq}$ $\iota $ y.top-of-ballot(y,w\textsubscript{0})]: win ($\iota $ y.top-of-ballot(y,w[2032?]), w[2032?]) = win ($\iota $ y.top-of-ballot(y,w\textsubscript{0}), w\textsubscript{0})
\end{styleStandard}

\begin{styleStandard}
In prose:
\end{styleStandard}

\begin{styleStandard}
\ \ a.\ \ \textbf{Assertion}: In w\textsubscript{0}, the person who was at the top of the ballot in w\textsubscript{0} won.
\end{styleStandard}

\begin{styleStandard}
\ \ b.\ \ \textbf{Presupposition}: In each world w[2032?], a counterfactual world of w\textsubscript{0}, if someone else had been at the top of the ballot in w[2032?], the person who was at the top of the ballot in w[2032?] won in w[2032?] iff the person who was at the the top of the ballot in w\textsubscript{0} won in w\textsubscript{0}.
\end{styleStandard}

\begin{styleStandard}
Choi (2007) and Choi \& Romero (2008) apply von Fintel’s formalisation of \textit{wh-ever} to the Spanish indefinite \textit{cualquiera} with a minor change: the ordinary meaning of Spanish indefinite \textit{cualquiera} contains an existential quantifier and not an iota operator as in von Fintel’s analysis of \textit{wh-ever}. The Spanish example in \ expresses the assertion that some book was picked and put on a pile and a presupposition of variation according to which there are other maximally similar worlds in which a book has been picked and put on a pile (see Choi 2007, Choi \& Romero 2008: 91):
\end{styleStandard}

\begin{styleStandard}
\label{bkm:Ref517736}\ \ Juan cogió un libro cualquiera y lo puso encima de la pila. ‘Juan took some book and put it on top of the pile.’ 
\end{styleStandard}

\begin{styleStandard}
\ \ a.\ \ \textbf{Assertion}: $\lambda $ w\textsubscript{0}. ${\exists}$x.book(x,w\textsubscript{0}) \& pick(j,x,w\textsubscript{0}) \& put-on-pile(j,x,w\textsubscript{0})
\end{styleStandard}

\begin{styleStandard}
\ \ b.\ \ \textbf{Presupposition}: \newline
$\lambda $ w\textsubscript{0 }${\forall}$ w[2032?] ${\in}$ min\textsubscript{w0} [F ${\cap}$ $\lambda $ w[2033?]. \{x:book(x,w[2033?])\} ${\neq}$ \{x:book(x,w\textsubscript{0})\}]:\newline
${\exists}$x.book(x,w[2032?]) \& pick(j,x,w[2032?]) \& put-on-pile(j,x,w[2032?])\newline
${\exists}$x.book(x,w\textsubscript{0}) \& pick(j,x,w\textsubscript{0}) \& put-on-pile(j,x,w\textsubscript{0})
\end{styleStandard}

\begin{styleStandard}
In prose:
\end{styleStandard}

\begin{styleStandard}
\ \ a.\ \ \textbf{Assertion}: In the actual world w\textsubscript{0}, there is some book in w\textsubscript{0} that John picked up and put on the pile in w\textsubscript{0}.
\end{styleStandard}

\begin{styleStandard}
\ \ b.\ \ \textbf{Presupposition}: In all counterfactual worlds w[2032?] minimally different from w\textsubscript{0} with respect to the identity of the set of books, there is some book in w[2032?] that John picked up and put on the pile in w[2032?] iff there is some book in w\textsubscript{0} that John picked up and put on the pile in w\textsubscript{0}.
\end{styleStandard}

\begin{styleStandard}
In contrast to Choi \& Romero 2008, A\&M (2018) assume that the random choice interpretation introduces a modal domain consisting of worlds compatible with agent’s decisions. A\&M (2018) describe RC using a modal component in b according to which the relevant cards that Juan did not choose but could have chosen are part of the cards the agent had at the time he took the decision \textit{d}. The modal component states that Juan’s decision could be fulfilled by taking any card in the domain. The following components are part of truth conditions (for the formalisation see A\&M 2018):
\end{styleStandard}

\begin{styleStandard}
\ \ \ \ Juan\ \ cogió\ \ una\ \ carta\ \ cualquiera.
\end{styleStandard}

\begin{styleStandard}
\ \ \ \ Juan\ \ took\ \ a\ \ card\textit{\ \ }\textsc{cualquiera}
\end{styleStandard}

\begin{styleStandard}
\ \ \ \ ‘Juan took a random card.’
\end{styleStandard}

\begin{styleStandard}
\label{bkm:Ref150779122}\ \ a.\ \ \textit{Existential Component}: In \textit{w}, there is an event \textit{e} of Juan taking a card \textit{x}.
\end{styleStandard}

\begin{styleStandard}
\ \ b.\ \ \textit{Modal Component}: For every (relevant) card \textit{y} in \textit{w}, there is a world \textit{w}[2032?] compatible with \textit{d}\textit{\textsubscript{e}} where there is an event \textit{e}[2032?] of Juan taking \textit{y} that fulfills \textit{d}\textit{\textsubscript{e}}.
\end{styleStandard}

\begin{styleStandard}
\ \ \ \ (A\&M 2018: 32)
\end{styleStandard}

\begin{styleStandard}
I dub this approach \textit{intentional randomness}, as the agent who picked a random card had taken the decision to do so. 
\end{styleStandard}

\paragraph[2.2.3.3 Discussion]{\textit{2.2.3.3 Discussion}}
\begin{styleStandard}
I would like to discuss three critical points concerning the previous analyses of \textit{cualquiera}. First, no study so far suggests a detailed analysis of the evaluative interpretation of \textsc{un} N \textit{cualquiera} in episodic and predicative contexts and makes a proposal as to how to analyse the relation between the evaluative meaning and RC/FC. 
\end{styleStandard}

\begin{styleStandard}
\ \ The second critical point concerns the description of the random choice interpretation. I will show in Chapter 4 that \textit{cualquiera} can express unintentional/unvoluntary randomness too, which is the result of an unreflected/unplanned action. Let us consider the following scenario:
\end{styleStandard}

\begin{styleStandard}
\ \ Scenario: Imagine Juan is a lunatic. He wanders around in his sleep. He takes different objects, among which are cards:
\end{styleStandard}

\begin{styleStandard}
\ \ Juan cogió una carta cualquiera.
\end{styleStandard}

\begin{styleStandard}
\ \ ‘Juan took a random card.’
\end{styleStandard}

\begin{styleStandard}
In this scenario, \textit{cualquiera} has RC even though the event of taking the actual card was not intentional and even though Juan did not decide to take a card. Moreover, I will show in Chapter 4 on diatopic variation of \textit{cualquiera} in ArgSp, that there are volitional/intentional verbs that have an evaluative interpretation as their only meaning when they select \textit{cualquiera} as a direct object , such as \textit{mandar cualquiera }‘to say bullshit’ (see also Fr. \textit{dire n’importe quoi }‘to say bullshit’). The following examples, also from Argentinian Spanish, demonstrates such a case.
\end{styleStandard}

\begin{styleStandard}
\ \ Perdón\ \ respondí\ \ cualquiera\ \ 
\end{styleStandard}

\begin{styleStandard}
\ \ sorry\ \ answered\ \ \textsc{cualquiera}
\end{styleStandard}

\begin{styleStandard}
\ \ ‘Sorry, I answered \textit{nonsense}. 
\end{styleStandard}

\begin{styleStandard}
\ \ (@ 33, \textstylepostdetails{Buenos Aires/neighrborhood Congreso, 25 years old }ArgSp)
\end{styleStandard}

\begin{styleStandard}
\ \ Osea respondí re cualquiera lo que me preguntaba, no sé qué carajo, habré estado muy dormida pero como que contesté las preguntas por la mitad… 
\end{styleStandard}

\begin{styleStandard}
\ \ ‘I answered with complete nonsense to what I was asked, I don’t know, what the hell, I must have been very tired because I answered the questions only partially.’
\end{styleStandard}

\begin{styleStandard}
\ \ (@ 34, female blogger, Buenos Aires, 25 years old; ArgSp)
\end{styleStandard}

\begin{styleStandard}
We thus need to find a better distinction between RCI and the evaluative interpretation that is based on more criteria than just the presence of a volitional/intentional verb. The relation between random choice and evaluative meaning needs to be explained. I address this issue in Chapter 4. 
\end{styleStandard}

\begin{styleStandard}
\ \ Another critical point of previous approaches to RCI is that they are too restrictive. Both approaches of RC in 2.3.1 and 2.3.2 suggest that there is a correlation between volitional verbs and RC of postnominal \textit{cualquiera}. However, this is too restrictive. There are examples of \textsc{un} N \textit{cualquiera} in mathematical contexts under copular verbs or in subject position with free choice interpretation. Note that \textsc{un} N \textit{cualquiera} should not have RC in subject or predicative position under either approach in 2.3. Let us consider the following examples:
\end{styleStandard}

\begin{styleStandard}
\label{bkm:Ref43031271}\ \ El punto P de coordenadas (a,b) es un punto cualquiera.
\end{styleStandard}

\begin{styleStandard}
\ \ ‘The point P of coordinates (a,b) is a random point.’
\end{styleStandard}

\begin{styleStandard}
\ \ (@ 35)
\end{styleStandard}

\begin{styleStandard}
\label{bkm:Ref43031168}\ \ una muestra cualquiera nos vale.
\end{styleStandard}

\begin{styleStandard}
\ \ ‘Any random sample is good for us.’ 
\end{styleStandard}

\begin{styleStandard}
\ \ (@ 36)
\end{styleStandard}

\begin{styleStandard}
In these examples, random choice is implicit; that is, what is random about the point or the sample is that the point P or the sample is defined in a way such that they were picked randomly or can be picked randomly. Random Choice is thus not restricted to the occurrence of volitional verbs expressed in the sentence, but to contexts in which it makes sense to recover a choice. Note that in these examples, the evaluative reading of ‘ordinary, average, common’ does not make sense at all, because there is no point in the coordination system that can be described as unremarkable or common for the speaker in a mathematical context and it does not make sense to talk about an unremarkable sample in statistics. I thus disagree with A\&M (2011, 2018), who predict an unremarkable reading of \textsc{un} N \textit{cualquiera} in affirmative sentences in predicative position such as \ and in subject position such as .
\end{styleStandard}

\subsection[2.3 Cualquiera under negation]{\textbf{2.3 }\textbf{\textit{Cualquiera}}\textbf{ under negation}}
\begin{styleStandard}
In this section, I will look into FCIs under negation more closely, because negation will play a role in the corpus data description of \textsc{un} N \textit{cualquiera} in Chapter 3.
\end{styleStandard}

\begin{styleStandard}
\ \ The use of FCIs under negation was first described by Horn (2000). The speaker in \ says that she wants to distinguish or discriminate among the possible candidates she wants to go to bed with. In other words, the speaker wants to be selective. According to Horn, \textit{not} negates the indiscriminacy. Thus, \textit{not just any} has an anti-indiscriminative interpretation (see Horn 2000, Aloni \& Port 2013, among many others):
\end{styleStandard}

\begin{styleStandard}
\label{bkm:Ref42805181}\ \ I don’t want to go to bed with just anyone anymore. I have to be attracted to them sexually.
\end{styleStandard}

\begin{styleStandard}
\ \ (Aloni et al. 2011: 8)
\end{styleStandard}

\begin{styleStandard}
Fălăuş (2013) offers a more general analysis of FCIs under negation. According to Fălăuş, Italian \textit{qualunque} in \ has the denotation of an existential quantifier like \textit{some} in English (see Fălăuş 2013 for formal representation and Section 2.1.3). Moreover, \textit{qualunque} triggers alternatives \{a, b, c…\}, such as alternative books \{\textit{Don Quijote}, \textit{War and Peace}, …\}. Postnominal \textit{qualunque} is focused, which triggers a domain restriction on the quantification of ‘some’. The sentence with a focused \textit{un libro qualunque} is interpreted as meaning that Gianni did not read a book in the initially considered domain of quantification (say D), but he read a book in some other domain (say D[2032?]), e.g. some exceptional book. The book in the domain D[2032?] is singled out pragmatically, according to Fălăuş (2013).
\end{styleStandard}

\begin{styleStandard}
\label{bkm:Ref42805226}\ \ Gianni\ \ non\ \ ha\ \ letto\ \ un\ \ libro\ \ \MakeUppercase{qualunque}
\end{styleStandard}

\begin{styleStandard}
\ \ Gianni\ \ \textsc{neg}\ \ has\ \ Read\ \ a\ \ book\ \ \textsc{qualunque}
\end{styleStandard}

\begin{styleStandard}
\ \ ‘Gianni didn’t read just any book.’
\end{styleStandard}

\begin{styleStandard}
\ \ (Italian, Fălăuş 2013: 95)
\end{styleStandard}

\begin{styleStandard}
The negation and focus interact in . If \textit{qualunque} were not focused in , negation would operate on \textit{every} domain alternative and \textit{qualunque} would be interpreted as a negative polarity item similar to English \textit{any} under negation. It would get a different interpretation, namely the plain negation of ‘Gianni did not read any book’. However, such a function of plain negation is expressed by negative words such as \textit{niente/nessun} ‘anything/any’ in Italian:
\end{styleStandard}

\begin{styleStandard}
\ \ Gianni non ha letto niente/nessun libro.\ \ (Italian)
\end{styleStandard}

\begin{styleStandard}
\ \ ‘Gianni did not ready anything/any book.’
\end{styleStandard}

\begin{styleStandard}
Instead, the negation in \ operates on a \textit{restricted} set of alternatives; that is, it negates every domain alternative out of a particular set of alternatives.
\end{styleStandard}

\begin{styleStandard}
\ \ The same facts have been observed for Spanish \textit{cualquiera}, which is stressed under negation according to Arregui (2006), represented in capital letters as CUALQUIER (see Arregui 2006). If \textit{cualquier} is not stressed, it is ungrammatical, because it would be then interpreted as ‘some’ under negation as in :
\end{styleStandard}

\begin{styleStandard}
\ \ No compró \MakeUppercase{cualquier} libro.
\end{styleStandard}

\begin{styleStandard}
\ \ ‘S/he didn't buy just any book.’
\end{styleStandard}

\begin{styleStandard}
\ \ (Arregui 2006: e.g. 2)
\end{styleStandard}

\begin{styleStandard}
\label{bkm:Ref533953565}\ \ *No compró cualquier libro.
\end{styleStandard}

\begin{styleStandard}
\ \ (intended NPI interpretation: ‘S/he didn’t buy any book.’) 
\end{styleStandard}

\begin{styleStandard}
\ \ (Arregui 2006: e.g. 2)
\end{styleStandard}

\begin{styleStandard}
There are other indefinite items that are used in the scope of negation with the NPI interpretation, namely negative indefinites such as \textit{ningún} or \textit{nadie}:
\end{styleStandard}

\begin{styleStandard}
\ \ No vió a nadie.
\end{styleStandard}

\begin{styleStandard}
\ \ ‘S/he didn’t see anyone.’
\end{styleStandard}

\begin{styleStandard}
\ \ (@ 37)
\end{styleStandard}

\begin{styleStandard}
What remains unclear under Fălăuş (2013) and Arregui’s (2006) proposal is the role of stress. It is a general rule that all FCIs, regardless of whether they are stressed or not and no matter what syntactic position and function they have, produce the anti-indiscriminative interpretation under negation (see 2.2.2). Stress is a side effect at best, not a necessary condition:
\end{styleStandard}

\begin{styleStandard}
\label{bkm:Ref42806294}\ \ Mi padre no es cualquiera.\footnote{\ If stress played an important role, we would not observe the robust correlation between the anti-discriminative effect of FCIs under negation (see Section 3) }
\end{styleStandard}

\begin{styleStandard}
\ \ ‘My dad is not just anyone.’
\end{styleStandard}

\begin{styleStandard}
\ \ (@ 38)
\end{styleStandard}

\begin{styleStandard}
\label{bkm:Ref42806303}\ \ Hoy no es cualquier día, sino ese en el que todo puede cambiar.
\end{styleStandard}

\begin{styleStandard}
\ \ ‘Today is not just any day, but that day when everything can change.’
\end{styleStandard}

\begin{styleStandard}
\ \ (@ 39)
\end{styleStandard}

\begin{styleStandard}
Another critical point is that this proposal does not say anything about the domain restriction of the indefinite \textit{qualunque}. What property do the books in the old domain D and new domain D[2032?] in \ have? How shall we account for the fact that the new domain D[2032?] must contain a particular or a distinguished book that is not like any other books?
\end{styleStandard}

\begin{styleStandard}
\ \ Gianni non ha letto un libro qualunque. \# Ha letto un libro come tanti.\ \ (Italian)
\end{styleStandard}

\begin{styleStandard}
\ \ ‘Gianni did not read just any book. \# He read a book like any other book.’\ \ 
\end{styleStandard}

\begin{styleStandard}
Another approach to FCIs under negation is offered by Aloni (2007). According to her, FCIs under negation operate on the universal inference triggered by FCIs. The negation has scope over the conjunction of \textit{all} alternatives (it is not the case that you can bring along book A and you can bring along book B). This has the effect that the free choice is denied under negation, that is, ‘you cannot freely choose to bring any book’:
\end{styleStandard}

\clearpage\begin{styleStandard}
\label{bkm:Ref43030327}\ \ You cannot bring along just any book. \textit{No puedes traer cualquier libro}. 
\end{styleStandard}

\begin{styleStandard}
\ \ ¬ [${\lozenge}$ you bring along book a1 or ${\lozenge}$ you bring along book a2 or … ], for any book a\textsubscript{i} ${\in}$[F020?]D
\end{styleStandard}

\begin{styleStandard}
\ \ ‘You cannot freely choose which book you bring me.’
\end{styleStandard}

\begin{styleStandard}
The result of the negation is the logical inference that you can bring along some book, but not all books. This logical inference is known in the literature as a \textit{quantitative scalar} implicature (see Chierchia 2006, 2013). I can thus sum up that, according to Aloni (2007), negation has scope over the universal quantification of free choice \textit{cualquiera} when \textit{cualquiera} is used under negation. 
\end{styleStandard}

\begin{styleStandard}
\ \ The prediction of this account is that FCIs under negation always produce the negation of free choice independent of stress. This is indeed the case, as pronominal \textit{cualquiera} produces the same effect (see \ and ). As FC is denied under negation, free choice should not be interpreted as a conversational implicature, that is, only available in certain pragmatic contexts, because it is not cancellable and does not disappear under negation in Modern Spanish. Since negation cannot cancel Free Choice, free choice must then be either a presupposition or a conventional implicature (see Potts 2005 for tests of conventional implicature). Neither of them is cancellable. Both hypotheses exist in the literature: free choice is a presupposition (von Fintel 2000, Choi \& Romero 2007) or free choice is a conventional implicature (Aloni 2007, 2021). In this work, I ignore the status of free choice as a presupposition or conventional implicature and focus instead on whether free choice is the \textit{only} lexical meaning of \textit{cualquiera} and whether other meanings can be derived from this lexical meaning. 
\end{styleStandard}

\begin{styleStandard}
\ \ Aloni’s account is silent about further implications for the evaluation of free choice alternatives under negation, such as the implication that the alternative candidates in \ are interpreted as bad candidates on some pragmatic scale (see Horn 2000). One could derive the following inference using free choice under negation in . The negation of free choice entails careful or determined selection. The speaker in \ thus asks to select a book carefully with determination. Careful selection of a book triggers the inference that the book in \ must have the property of being distinct from other books that are not carefully selected (an interesting book, for instance). It is impossible to carefully select a book that does not stand out among other books. Imagine a scenario where all books are the same. In this scenario, it is impossible to select a book carefully. The hearer could thus interpret \textit{No puedes traer cualquier libro} in \ as the negation of picking out a book with an undistinguished property and hence the imperative or advice to bring a distinctive book in . As I will show, \textit{un libro} \textit{cualquiera} can denote a kind of a book, which has an undistinguished property of a book or a property of a book that all books share (see Chapter 3). As I will argue, it is this kind interpretation that has been lexicalised on \textsc{un} N \textit{cualquiera} in predicative contexts or episodic contexts such as \textit{es }\textsc{un}\textit{ }N\textit{ cualquiera} (see Section 3.3). 
\end{styleStandard}

\subsection[2.4 Frequency of different interpretations of cualquiera]{\textbf{2.4 Frequency of different interpretations of }\textbf{\textit{cualquiera}}}
\begin{styleStandard}
In order to quantify different interpretations of Sp. \textit{cualquiera}, Aloni et al. (2010) have classified corpus data from Corpus del Español (CdE) with respect to the classification system in Figure 1, which uses Haspelmath’s (1997) classification system of indefinites across languages that represents syntactic and semantic properties (see also Section 5.2.3 on the diachronic analysis):
\end{styleStandard}

\begin{center}
\tablefirsthead{}
\tablehead{}
\tabletail{}
\tablelasttail{}
\begin{supertabular}{|m{0.16085985in}m{0.13515985in}m{0.29975986in}m{1.2629598in}m{3.4538598in}|}
\hline
 &
\multicolumn{4}{m{5.3879595in}|}{\textit{Functions on the map}}\\
 &
 &
\textbf{Abbr} &
\textbf{Label} &
\textbf{Example}\\
 &
a. &
SK &
specific known &
\textit{Somebody} called. Guess who?\\
 &
b. &
SU &
specific unknown &
I heard \textit{something}, but I couldn’t tell what it was.\\
 &
c. &
IR &
irrealis &
You must try \textit{somewhere} else.\\
 &
d. &
Q &
question &
Did \textit{anybody} tell you anything about it?\\
 &
e. &
CA &
conditional antecedent &
If you see \textit{anything}, tell me immediately.\\
 &
f. &
CO &
comparative &
In Freiburg, the weather is nicer than \textit{anywhere} else in Germany.\\
 &
g. &
DN &
direct negation &
John didn’t see \textit{anybody}.\\
→ &
h. &
AM &
anti-morphic &
I don’t think that \textit{anybody }knows the answer.\\
→ &
i. &
AA &
anti-additive &
The gravity of such act goes beyond \textit{any }justification.\\
→ &
j. &
FC &
free choice &
\textit{Anybody} can solve this problem.\\
→ &
k. &
UFC &
universal free choice &
John kissed \textit{any} woman with red hair.\\
→ &
l. &
GEN &
generic &
\textit{Any }dog has four legs.\\
→ &
m. &
IND &
indiscriminative &
I do not want to go to bed with \textit{just anyone} anymore.\\\hline
\end{supertabular}
\end{center}
\begin{styletablecaption}
Figure \stepcounter{Figure}{\theFigure}: Classification system of FCIs according to Aloni et al. (2010).
\end{styletablecaption}

\begin{styleStandard}
As can be seen from the Figure 2, the most frequent function of all the possible contexts where \textit{cualquiera} appears in Modern Spanish in the 20th century is clearly the FC(I) function, that is, the distribution of total exhaustification of \textit{cualquiera} over alternatives:
\end{styleStandard}

\begin{styleStandard}
  [Warning: Image ignored] % Unhandled or unsupported graphics:
%\includegraphics[width=4.9602in,height=3.0299in,width=\textwidth]{kellert-img001.png}
 
\end{styleStandard}

\begin{styletablecaption}
Figure \stepcounter{Figure}{\theFigure}: Distribution of \textit{cualquiera} in Modern Spanish (i.e. 20th c.) Aloni et al. (2010: 19).\footnote{\textrm{\ The colours do not have any representative value.}}
\end{styletablecaption}

\begin{styleStandard}
The graph also shows that it can have other functions as well (e.g. the indiscriminative [IND] and ‘no matter’ interpretations). The latter interpretation is not defined in Figure 1. The meaning distribution misses the evaluative reading of \textit{cualquiera}, which I will investigate also with respect to its frequency in Modern Spanish in Chapter 3.
\end{styleStandard}

\subsection[2.5 Summary of different properties of cualquiera in synchrony]{\textbf{2.5 Summary of different properties of }\textbf{\textit{cualquiera}}\textbf{ in synchrony}}
\begin{styleStandard}
The critical review of the state of research in this chapter has shown that \textit{cualquier}a has different syntactic functions (determiner, nominal modifier, pronoun, noun, degree adjective), as well as different semantic functions, as summarised in Table 4.
\end{styleStandard}

\begin{center}
\tablefirsthead{\hline
\textbf{Semantic and syntactic properties} &
\textbf{\textit{algún}} &
\textbf{\textit{cualquier}}\textbf{ NP/}\textbf{\textit{cualquiera}}\textbf{ as pronoun} &
\textbf{NP }\textbf{\textit{cualquiera}} &
\textbf{Source }\\\hline
\textbf{FCI/HI} &
no &
yes &
Yes under modal verbs, subtrigging, comparatives, negation. &
Fălăuş (2018),

A\&M (2011/2018)

Aloni (2007, 2021)\\\hline
\textbf{RCI} &
no &
no &
Yes (in object position 

under agentive verbs)

No (in subject position, adverbs and/or under non-agentive verbs). &
Choi \& Romero (2008), A\&M (2011/2018)\\\hline
\textbf{IGN} &
yes &
 &
Yes in Argentinian Spanish (like Italian \textit{qualsiasi}). &
Chierchia (2006),

Kellert \& Francia, (in press)\\\hline
\textbf{EVAL 1}

\textbf{‘ordinary’}

\textbf{‘bad’/}

\textbf{‘out of place’}

\textbf{EVAL 2} &
no &
no in Sp.

yes in ArgSp &
Yes . for EVAL as ‘ordinary’

Distribution unrestricted in Sp., i.e. \textit{cualquiera} is lexically ambiguous and can appear under modal contexts as well (A\&M 2011/2018) 

Yes in ArgSp for EVAL as ‘ordinary’ and ‘bad/out of place’

Distribution restricted in ArgSp under specific verbs such as copular verbs and 

verbs of saying. &
A\&M (2011/2018), 

Rivero (2011), Fălăuş (2015)

Conde (2011), Rizzo Salierno (2015)

\\\hline
\textbf{NPI } &
yes &
no 

 &
No

 &
Arregui (2006),

Aloni et al. (2010)\\\hline
\textbf{Modifiable by adverbs like }\textbf{\textit{excepto/}}

\textbf{\textit{casi}} &
no &
yes  &
? &
Arregui (2006)

\\\hline
\textbf{Sub-trigging } &
no &
yes  &
? &
Quer (2000), 

Aloni et al. (2010)\\}
\tablehead{\hline
\textbf{Semantic and syntactic properties} &
\textbf{\textit{algún}} &
\textbf{\textit{cualquier}}\textbf{ NP/}\textbf{\textit{cualquiera}}\textbf{ as pronoun} &
\textbf{NP }\textbf{\textit{cualquiera}} &
\textbf{Source }\\\hline
\textbf{FCI/HI} &
no &
yes &
Yes under modal verbs, subtrigging, comparatives, negation. &
Fălăuş (2018),

A\&M (2011/2018)

Aloni (2007, 2021)\\\hline
\textbf{RCI} &
no &
no &
Yes (in object position 

under agentive verbs)

No (in subject position, adverbs and/or under non-agentive verbs). &
Choi \& Romero (2008), A\&M (2011/2018)\\\hline
\textbf{IGN} &
yes &
 &
Yes in Argentinian Spanish (like Italian \textit{qualsiasi}). &
Chierchia (2006),

Kellert \& Francia, (in press)\\\hline
\textbf{EVAL 1}

\textbf{‘ordinary’}

\textbf{‘bad’/}

\textbf{‘out of place’}

\textbf{EVAL 2} &
no &
no in Sp.

yes in ArgSp &
Yes . for EVAL as ‘ordinary’

Distribution unrestricted in Sp., i.e. \textit{cualquiera} is lexically ambiguous and can appear under modal contexts as well (A\&M 2011/2018) 

Yes in ArgSp for EVAL as ‘ordinary’ and ‘bad/out of place’

Distribution restricted in ArgSp under specific verbs such as copular verbs and 

verbs of saying. &
A\&M (2011/2018), 

Rivero (2011), Fălăuş (2015)

Conde (2011), Rizzo Salierno (2015)

\\\hline
\textbf{NPI } &
yes &
no 

 &
No

 &
Arregui (2006),

Aloni et al. (2010)\\\hline
\textbf{Modifiable by adverbs like }\textbf{\textit{excepto/}}

\textbf{\textit{casi}} &
no &
yes  &
? &
Arregui (2006)

\\\hline
\textbf{Sub-trigging } &
no &
yes  &
? &
Quer (2000), 

Aloni et al. (2010)\\}
\tabletail{}
\tablelasttail{}
\begin{supertabular}{|m{0.80385983in}|m{0.51155984in}|m{1.0823599in}|m{2.18586in}|m{1.2011598in}|}

\end{supertabular}
\end{center}
\begin{styletablecaption}
Table \stepcounter{Table}{\theTable}: Different properties of \textit{cualquiera} in Modern Spanish as studied so far in the literature.
\end{styletablecaption}

\begin{styleStandard}
When \textit{cualquiera} acts as a quantificational determiner or pronoun, it is used as a FCI with a restrictive distribution (modal contexts, episodic context with subtrigging, epistemic modality, or negation) in European Spanish. It has a universal interpretation if embedded under a possibility modal. In non-modal environments, such as episodic contexts and predicative contexts in affirmative sentences, N \textit{cualquiera} has usually an evaluative interpretation of ‘unremarkable’ or a ‘random choice’ interpretation in European Spanish. In predicative contexts in negative sentences, it expresses negation of free choice (the so-called ‘not just any’ interpretation). \textit{Cualquiera} has the status of a degree adjective in Argentinian Spanish.
\end{styleStandard}

\subsection[2.6 Diachronic research on cualquiera]{\textbf{2.6 Diachronic research on }\textbf{\textit{cualquiera}}}
\subsubsection[2.6.1 Grammaticalization of cualquiera ]{\textbf{2.6.1 Grammaticalization of }\textbf{\textit{cualquiera}}\textbf{ }}
\begin{styleStandard}
It is still a debate in the literature on the diachrony of \textsc{qualquier} and other indefinites in Romance as to whether they are an innovation of the Romance languages, in particular whether they have undergone a grammaticalisation path that led from separate words into a single compound (henceforth: the Innovation Hypothesis) or whether they are directly related to Latin indefinites, either inherited or remodelled (i.e. by the substitition of elements in the course of language development or even as a conscious imitation of Latin models (henceforth: the Continuation Hypothesis; see C\&P 2009, Pato 2012, Mackenzie 2019, Kellert \& Enrique-Arias, submitted, and references therein). Under the Innovation Hypothesis, most previous diachronic studies concentrate on the question of how the relative pronoun \textit{qual} and the verb \textit{querer} ‘want, wish’ became grammaticalised into a single complex compound (see Palomo 1934, Lombard 1938–39/1947–48, Meier 1950, Rivero 1988, Espinosa \& Sánchez Lancis 2006, C\&P 2009, among others). According to Aloni et al. (2010), the grammaticalisation process cannot be retraced in the diachronic corpus CdE because it took place prior to the written documentation of Old Spanish, or, in other words, the compound \textit{qualquier,\nobreakdash-e,\nobreakdash-a} already existed in the earliest periods in this corpus (see also C\&P 2009, Rivero 1988). I will review both hypotheses in this section.
\end{styleStandard}

\paragraph[2.6.1.1 Continuation hypothesis]{\textit{2.6.1.1 Continuation hypothesis}}
\begin{styleStandard}
According to some linguists, Romance FCIs are based on or and/or influnced by Latin FCIs, including calques from Latin (henceforth: the Continuation Hypothesis). As early as the late 19th century, Meyer-Lübke (1899: 57) assumed that Italian \textit{qualunque} was remodelled on Latin \textit{qualiscumque }‘(of) any (kind) whatever’, composed of \textit{qualis} ‘which’ + \textit{cumque} ‘ever’, where the second component was substituted with \textit{unquam} ‘any, any-‘, hence \textit{qual[is] unqua[m]} ({\textgreater} Old and Modern Italian \textit{qualunque}) (see also Cortalazzo \& Zolli 1979: 335; Rohlfs 1968, vol. II: 222–223). Etymological dictionaries and some linguists also assume that the French indefinite \textit{quiconque} has its source in Old French \textit{qui que} ‘whoever’ + \textit{onques} ‘ever’ ({\textless} Lat. \textit{unquam}), which was influenced by Latin \textit{quicumque} (see Becker 2014 citing Godefroy 1880–1902, vol. 6: 511; Gamillscheg 1969: 737; Greimas 1998: 489; Tobler-Lommatzsch 1925–2018, vol. 8: 91). Menéndez Pidal (1928) assumes that Old Spanish \textit{qual }\textit{quier}, \textit{qui quier}, \textit{qual-se-quiera} represent the Old Spanish equivalents of the Latin indefinites \textit{quilibet, qualislibet}. Meier (1950: 387f.) assumes that Old Sp. \textit{qualquiera} and \textit{quienquiera} ‘whoever’ were already fixed expressions in Old Spanish. However, he admits that the attested morphological variation of -\textit{quiera} (e.g. -\textit{quiere, -quisiere}) poses a problem for the fixed expression hypothesis. Another problem for the Continuation Hypothesis is that Old Spanish \textit{qualquier,\nobreakdash-e,\nobreakdash-a} contains a verb different from the verb inside the Latin indefinites, which have -\textit{libet }(lit. ‘pleases’),\footnote{\ \textrm{Latin has also indefinite forms with the verb form }\textrm{\textit{vis}}\textrm{ ‘you want’: }\textrm{\textsc{quivis}}\textrm{, }\textrm{\textsc{quodvis}}\textrm{ lit. ‘who you want’, ‘what you want’, etc. (see Pato 2012: 304).}} with the consequence that, like in the case of It. \textit{qualunque} mentioned above, one has to assume a substitution with forms of (Old) Sp. \textit{querer} ‘to want, to love’. The Spanish verb \textit{querer} stems from Latin \textsc{quaerere} ’to search, to investigate, try to find’ (C\&P 2009), which has acquired the volitional meaning of ‘want’ (see C\&P 2009, Fn.24, p. 1152). C\&P (2009: 1152) assume that the free choice property can be more directly derived from ‘to try to find’ than from ‘to want, to desire’.\footnote{\textrm{\ This view is problematic for different reasons: First, in Catalan or Italian the Latin verb }\textrm{\textsc{volere}}\textrm{ was chosen (}\textrm{\textit{qualsevol, qualsivoglia}}\textrm{), and, second, typologically, there is a whole range of volitional verbs that can feature in indefinites (Haspelmath 1997), but to my knowledge, verbs meaning ‘to search’ are not among them.}} C\&P also note that Old Spanish \textit{qualquier,-e,-a} can be often found as a graphically separated form in the 13th century, a fact that they interpret as being against the Continuation Hypothesis (see 5.1.2 for details). 
\end{styleStandard}

\begin{styleStandard}
\ \ In most of the accounts that I have subsumed under the label Continuation Hypothesis, it remains somewhat unclear whether the authors believe in an organic development, in calques from written Latin, or in a combination of the two.
\end{styleStandard}

\begin{styleStandard}
\ \ Pato’s (2012) position on the origin of \textit{cualquiera} is not very clear. He seems to suggest that both the Continuation and the Innovation Hypotheses are correct in that—although \textit{qualquier,\nobreakdash-e,\nobreakdash-a} may be an innovation (but he does not explicitely say so)—Old Spanish inherited certain constructions from Latin that can be applied to it. In particular, he is referring to the phenomenon of interposition, in which a noun or an NP is placed between the pronoun and the verb \textit{quier} (e.g. \textit{qual manera quier }lit. ‘which manner one wants’). He explicitly calls it “a structure inherited from Latin” (Pato 2012: 273). Pato seems to believe that \textit{qual quier} was a quasi indefinite pronoun which could be separated by some linguistic material such as nouns, similar to ‘what manner soever’.
\end{styleStandard}

\begin{styleStandard}
\ \ However, Pato argues that the form without interposition, i.e. \textit{qual quier} N(P) ‘which one wants N(P)’ continued the path of grammaticalisation, becoming more fixed after the 13th century, as shown by graphically synthetic form \textit{qualquier} and later \textit{cualquier} (Pato 2012: 297).
\end{styleStandard}

\begin{styleStandard}
\ \ I will first review Pato’s argumentation for the Continuation Hypothesis and then discuss the problems it entails. 
\end{styleStandard}

\begin{styleStandard}
\ \ Pato (2012) assumes that the structure with interposition is a characteristic phenomenon of Spanish in the 13th century, which had already arisen in Latin, where \textit{quilibet} “may be separated by tmesis into \textit{qui} and \textit{libet}” (Gildersleeve \& Lodge 1900: 62, quoted by Pato 2012: 304). The example given by Gildersleeve is Sallust, \textit{De coniuratione Catilinae} 5,4, which reads “quoius rei lubet” ‘of any thing’.\footnote{\ The example itself is not quoted by either Gildersleeve or Pato, but was retrieved from the edition of Sallust’s works by Ahlberg (1919), p. 5. \textit{Quoius} is a variant of \textit{cuius}, genitive of the relative pronoun \textit{qu\=\i}, and \textit{lubet} is a variant of \textit{libet}. \textrm{Pato himself (2012: 305) tries to provide a further example from the late Latin}\textrm{\textbf{ }}\textit{Lex Romana Visigothorum} (which was in use on the Iberian Peninsula): “[…] qui cuius libet rei possessione privati sunt […]”. However, this does not seem to be a case of interposition, if it means, as I think, “who [\textit{qui}] are deprived of the possession of any [\textit{cuius} \textit{libet}] thing [\textit{rei}].”} According to Pato, some Romance languages have inherited this structure. He suggests that the interposition corresponds to the classical poetic structure of tmesis, which is used to separate a syntactic structure for emphatic reasons: “la tmesis o «transposición de palabras» (Institutio Oratoria, Quintiliano) es una licencia poética y retórica clásica (variante del hipérbaton) que consiste en la separación de los dos elementos que forman un compuesto” (Pato 2012: 280). 
\end{styleStandard}

\begin{styleStandard}
\ \ Pato (2012) examines a corpus of works (directed) by Alfonso X of Castile (the Wise), where \textit{qual quier} mostly shows interposition of the following nouns: \textit{estrella} ‘star’, \textit{grado} ‘degree’, \textit{manera} ‘manner’, \textit{logar} ‘place’, and \textit{planeta} ‘planet’ (e.g. \textit{qual manera quier} ‘whichever manner’). The interposition occurs more often than the absence thereof (“sin interposición”), which he analyses as \textit{qual quier} NP (97.5\% vs 74\%, respectively). Pato observes that the form \textit{qual quier} also appears as a single word form in the 13th century. However, he only investigates the separated form in relation with and without interposition (i.e. \textit{qual quier manera} vs \textit{qual manera quier}). The interposition seems mostly characteristic of the scientific texts of the corpus, in particular of the \textit{Los libros del saber de astronomía} and the \textit{Canones de Albateni} (Pato 2012: 301). 
\end{styleStandard}

\begin{flushleft}
\tablefirsthead{}
\tablehead{}
\tabletail{}
\tablelasttail{}
\begin{supertabular}{|m{0.19345984in}|m{1.8531599in}|m{0.43025985in}|m{0.70185983in}|m{0.40385985in}|m{0.71985984in}|m{0.42955986in}|m{0.73025984in}|}
\hhline{~~------}
\multicolumn{1}{m{0.19345984in}}{} &
 &
\multicolumn{2}{m{1.2108599in}|}{\centering \textbf{Con interposición}} &
\multicolumn{2}{m{1.2024599in}|}{\centering \textbf{Sin interposición}} &
\multicolumn{2}{m{1.23856in}|}{\centering \textbf{Totales por texto}}\\\hline
\raggedleft \textbf{1.} &
\textbf{\textit{Libros del saber de astronomía}} &
\raggedleft 58\% &
\raggedleft (334/578) &
\raggedleft 19\% &
\raggedleft (32/168) &
\raggedleft 49\% &
\raggedleft\arraybslash (366/746)\\\hline
\raggedleft \textbf{2.} &
\textbf{\textit{Cánones de Albateni}} &
\raggedleft 15\% &
\raggedleft (86/578) &
\raggedleft {}-{}- &
\raggedleft {}-{}- &
\raggedleft 11.5\% &
\raggedleft\arraybslash (86/746)\\\hline
\raggedleft \textbf{3.} &
\textbf{\textit{Espéculos}} &
\raggedleft 10\% &
\raggedleft (60/578) &
\raggedleft 4.1\% &
\raggedleft (7/168) &
\raggedleft 9\% &
\raggedleft\arraybslash (67/746)\\\hline
\raggedleft \textbf{4.} &
\textbf{\textit{Formas e imágenes}} &
\raggedleft 4.5\% &
\raggedleft (26/578) &
\raggedleft 9\% &
\raggedleft (15/168) &
\raggedleft 5.5\% &
\raggedleft\arraybslash (41/746)\\\hline
\raggedleft \textbf{5.} &
\textbf{\textit{Libro del cuadrante señero}} &
\raggedleft 4\% &
\raggedleft (21/578) &
\raggedleft {}-{}- &
\raggedleft {}-{}- &
\raggedleft 2.8\% &
\raggedleft\arraybslash (21/746)\\\hline
\raggedleft \textbf{6.} &
\textbf{\textit{Judizios de las estrellas}} &
\raggedleft 4\% &
\raggedleft (21/578)) &
\raggedleft 10\% &
\raggedleft (17/168) &
\raggedleft 5\% &
\raggedleft\arraybslash (38/746)\\\hline
\raggedleft \textbf{7.} &
\textbf{\textit{Lapidario}} &
\raggedleft 2.4\% &
\raggedleft (4/578) &
\raggedleft 14\% &
\raggedleft (24/168) &
\raggedleft 5\% &
\raggedleft\arraybslash (38/746)\\\hline
\raggedleft \textbf{8.} &
\textbf{\textit{Libro de las cruces}} &
\raggedleft {}-{}- &
\raggedleft (1/578) &
\raggedleft 13\% &
\raggedleft (22/168) &
\raggedleft 3\% &
\raggedleft\arraybslash (23/746)\\\hline
\raggedleft \textbf{9.} &
\textbf{\textit{General estoria, I}} &
\raggedleft {}-{}- &
\raggedleft (1/578) &
\raggedleft 4\% &
\raggedleft (7/168) &
\raggedleft 1\% &
\raggedleft\arraybslash (8/746)\\\hline
\multicolumn{2}{|m{2.12536in}|}{\textbf{Total number per structure }} &
\raggedleft 97.5\% &
\raggedleft (564/578) &
\raggedleft 74\% &
\raggedleft (124/168) &
\raggedleft 92\% &
\raggedleft\arraybslash (688/746)\\\hline
\end{supertabular}
\end{flushleft}
\begin{styletablecaption}
Table \stepcounter{Table}{\theTable}: Pato (2012: 282).
\end{styletablecaption}

\begin{styleStandard}
Pato assumes that interposition is correlated not only with text genre (scientific texts) but also with a certain information-structural function in which the NP was focused in the position between \textit{qual} and \textit{quier}: “parece sensato pensar que la interposición podía emplearse para enfatizar el sustantivo tmético, por tanto era una variación aprovechable, pero con carácter literario” (Pato 2012: 284).
\end{styleStandard}

\begin{styleStandard}
\ \ There are some problems with Pato’s version of the \textit{Continuation Hypothesis}. Apart from the fact that Pato is not explicit on the exact relation between the Romance and the supposed Latin structure, i.e. is it really a structure inherited from Latin (as he explicitely says) or an effect of Latinising writing, as other parts of his article may suggest? The fact that interposition occurs mostly in scientific texts would rather suggest the latter.\footnote{\ Alternatively, in particular considering the fact that both the \textrm{\textit{Los libros del saber de astronomía}}\textrm{ and the }\textrm{\textit{Canones de Albateni}}\textrm{ depend on Arabic sources, Pato (2012: 306) does not exclude that the interposition phenomenon of }\textrm{\textit{qual quier}}\textrm{ might be a calque from Arabic, as suggested by Fernández González (1992).}} 
\end{styleStandard}

\paragraph[2.6.1.2 Innovation hypothesis/Grammaticalisation Hypothesis]{\textit{2.6.1.2 Innovation hypothesis/Grammaticalisation Hypothesis}}
\begin{styleStandard}
\ \ According to some linguists, wh-based FCIs such as Spanish \textit{cualquiera} developed from wh-relative clauses. 
\end{styleStandard}

\begin{styleStandard}
\ \ An early approach of this type is that by Lombard (1938–39/1947–48). Importantly, the author remarks that the verbal compont of \textit{quienquiera} and similar formations (including our \textsc{qualquier}) cannot be a calque of Latin models, for the simple reason that \textit{quier(e)} and \textit{quiera }do not translate the corresponding verbal components of the Latin items. These are \textit{libet }‘it pleases’ (with no logical subject) and \textit{vis} ‘you want’ (in 2\textsuperscript{nd} pers. sing.), whereas the Spanish formations show the verb \textit{querer} ‘to love, to want’ in the third person (and in the subjunctive). The Spanish structure shows an indeterminate logical subject in the meaning of ‘one’ (a typically Romance usage, but not typical of Latin), which can sometimes be made more precise by using an impersonal reflexive construction (of the type \textit{quien se quier,\nobreakdash-e,\nobreakdash-a). }Lombard (1938–39/1947–48) thus considers not only \textsc{qualquier}\textit{ }but also Sp. \textit{quienquiera, quien se quiera, dondequiera} ‘wherever’,\textit{ }and others, as well as an impressive number of similar formations in many other Romance varieties, as pan-Romance innovations. According to Lombard (1938/39: 187), in the beginning, the verbal forms in these expressions were, in reality, relative clauses, where the choice of the verb was free (e.g. Old Sp. “tenie gomito é lanzaba cual manjar que le daban”\footnote{\ Quoted from Gessner (1895: 159), where the example is identified as stemming from Libro de los exemplos (467 a m).} ‘He was sick, and thrust out whatever food they gave him’), but later the verbs meaning ‘to be’ and ‘to want’ yielded more or less fixed expressions, which could have a relative clause of their own, which could also become fixed (like in \textit{cualquiera que sea, }Port. \textit{como quer que seja}, Catalan, \textit{qualsevulla (qualsevol) que sìa, }It. \textit{qual vuolsi sia}\textbf{\textit{).}}
\end{styleStandard}

\begin{styleStandard}
Ainsi, l’indéfini relatif esp. \textit{cual que }a donné l’indéfini général \textit{cual [que] quiera, }qui a donné le nouvel indéfini relatif \textit{cualquiera que, }qui a donné le nouvel indéfini général \textit{cualquiera que sea, }qui peut même donner le nouvel indéfini relatif \textit{cualquiera que sea que. }De ces trois types d’indéfinis relatifs, seuls les deux premiers, \textit{cual que }et \textit{cualquiera que, }ont réellement une existence.
\end{styleStandard}

\begin{styleStandard}
\ \ \ \ (Lombard 1938/39: 188)
\end{styleStandard}

\begin{styleStandard}
In a similar vein, more recent authors (Company-Company \& Pozas-Loyo 2009, Haspelmath 1997, and others) claim that wh-based FCIs such as Sp. \textit{cualquiera} have emerged as a result of a grammaticalisation process in which free relative clauses were reanalysed as indefinite noun phrases (see also Palomo 1934, Company-Company \& Pozas-Loyo 2009, Girón Alconchel 2012, Brucart 1999: §7.5.7). However, there is no consensus about how exactly the grammaticalisation of Old Sp. \textit{qualquiera} took place. One possible explanation why there is no consensus is that the indefinite \textit{qualquiera} was already present in Old Spanish documents. It is thus difficult to confirm the hypothesis of the grammaticalisation on the basis of empirical proof. 
\end{styleStandard}

\begin{styleStandard}
\ \ The origin of Spanish \textit{cualquiera} according to some authors, is represented in (308) (cf. Cuervo 1886–1994 in C\&P 2009: 1086). Note that this grammaticalisation path is a simplified representation, which takes modern graphical representation as basis (Elvira 2007). 
\end{styleStandard}

\begin{styletablecaption}
Figure \stepcounter{Figure}{\theFigure}: Grammaticalisation path of cualquier in Spanish (C\&P 2009: 1086)
\end{styletablecaption}

\begin{styleStandard}
The most characteristic signs for the grammaticalisation of \textit{qualquier,-e,-a} according to C\&P (2009) are briefly described as follows: The indefinite \textit{qual-quier,\nobreakdash-e,\nobreakdash-a} changed in frequency, in phonology (i.e. the loss of \textit{e} in the form \textit{qual quier}), and in its morphosyntactic status (see also Lehmann 1995 for typical features of grammaticalisation). Its previous semantic interpretation was weakened as well; that is, its etymological interpretation has been lost and it acquired a new syntactic status. I will return to the frequency and especially to the semantic weakening in Section 2.6.2. As for the phonological weakening, the loss of \textit{e} in the form -\textit{quier} is not specific to the item at issue here but due to a general phonological apocope that applied to all kinds of syntactic categories from the 13th to the 14th century. It can thus hardly be counted as evidence for grammaticalisation.
\end{styleStandard}

\begin{styleStandard}
\ \ Given the list of arguments for the Romance \textit{Innovation Hypothesis}, which have been already convincingly discussed in Lombard (1938–39/1947–48), I assume that this hypothesis is stronger than the \textit{Continuation Hypothesis}, i.e. Old Sp.\textbf{\textit{ }}\textit{qualquiera,-e,-a}\textbf{ }is derived from a relative clause structure.
\end{styleStandard}

\subsubsection[2.6.2 Semantic change in time]{\textbf{2.6.2 Semantic change in time}}
\begin{styleStandard}
Not much research exists on the diachrony of the evaluative meaning of \textit{cualquiera} in European or Argentinian Spanish. According to Company Company (2009), the change in meaning of \textit{cualquiera} from FCI to EVAL is due to a \textit{subjectification process}, whereby the meaning of \textit{cualquiera} acquires a more subjective interpretation over time (see also Traugott 1989 for \textit{subjectification} as trigger of meaning change in English). Company Company (2009) suggests that the free choice indefinite \textit{cualquiera} has undergone the change into the evaluative \textit{cualquiera}, triggering speaker-independent alternatives to a more subjective meaning dependent on the speaker’s evaluation. However, it remains to be seen what this subjective interpretation exactly means and how this change towards a more subjective meaning occurred. I will deal with this question in Chapters 4 and 5. With respect to diachrony of \textit{cualquiera} in European Spanish in Chapter 5, I will also study the grammaticalisation of \textit{cualquiera} in order to understand how the occurrence of the evaluative meaning fits into the grammaticalisation process, which according to the literature correlates with a loss and not with enrichment of meanings of the grammaticalised construction (Traugott 1989, among others). With respect to the diachrony of \textit{cualquiera} in Argentinian Spanish, Rizzo Salierno (2013) assumes that the potential linguistic precedent of the gradable predicate \textit{cualquiera} in ArgSp is possibly related to the gradable predicate \textit{cualunque}, an Italianism (see 2.1.4). \textit{Cualunque} appeared in Argentinian Spanish during the Italian immigration (on \textit{qualunque} in Italian, see Chierchia 2006, Aloni et al. 2010, among others) (see 2.1.4). Like \textit{cualquiera} in Argentinian Spanish, \textit{cualunque} also allows degree modification (unlike in Italian, see Kellert \& Francia in press) and can be used predicatively or attributively in a postnominal position, as happens with \textit{cualquiera}, and have an evaluative interpretation like ‘ordinary’:
\end{styleStandard}

\begin{styleStandard}
\ \ El\ \ que\ \ yo\ \ digo\ \ \textit{era\ \ re\ \ cualunque.\ \ }
\end{styleStandard}

\begin{styleStandard}
\ \ the\ \ that\ \ I\ \ say\ \ was\ \ very\ \ \textsc{cualunque}\textit{\ \ }
\end{styleStandard}

\begin{styleStandard}
\ \ ‘The one that I said was very ordinary.’
\end{styleStandard}

\begin{styleStandard}
\ \ (Rizzo Salierno 2013: 28)
\end{styleStandard}

\begin{styleStandard}
\ \ Es\ \ una\ \ flaca\ \ \textit{cualunque\ \ }que\ \ se\ \ cree\ \ que\ \ es\ \ Marilyn Monroe.
\end{styleStandard}

\begin{styleStandard}
\ \ Is\ \ a\ \ woman\ \ \textsc{cualunque}\textit{\ \ }that\ \ her\ \ believes\ \ that\ \ is\ \ Marilyn Monroe
\end{styleStandard}

\begin{styleStandard}
\ \ ‘She is an ordinary woman who thinks she’s Marilyn Monroe.’
\end{styleStandard}

\begin{styleStandard}
\ \ (Rizzo Salierno 2013: 29)
\end{styleStandard}

\begin{styleStandard}
However, the change from an indefinite pronoun/determiner/postnominal modifier \textit{cualunque} to a gradable adjective \textit{cualunque} and for that matter the transfer of this function to \textit{cualquiera} remains unexplained. I will deal with this change in Chapter 4. 
\end{styleStandard}

\begin{styleStandard}
\ \ With respect to the question how different meanings developed, it remains an open question as to how the reanalysis of the indefinite \textit{cualquiera} into an item with the evaluative meaning ‘common/unremarkable’ in European Spanish and into a degree adjective with meaning ‘unremarkable/bad’ or an evaluative noun with the meaning ‘nonsense’ in Argentinian Spanish fits models of linguistic reanalysis in the frameworks of formal semantics (Eckardt 2006, Bar-Asher Siegal 2015, Aloni 2021, among others). According to these authors, it is possible to describe reanalysis on the level of morphosyntactic form (F) and/or on the level of meaning (M) as shown in :
\end{styleStandard}

\begin{styleStandard}
\label{bkm:Ref65487536}\ \ Scenarios of reanalysis (see Aloni 2021):
\end{styleStandard}

\begin{styleStandard}
\ \ Scenario 1:\ \ \ \ \{F1;M1\}\textsuperscript{t1} [F0E0?] \{F2;M2\}\textsuperscript{t2}
\end{styleStandard}

\begin{styleStandard}
\ \ Scenario 2:\ \ \ \ \{F1;M1\}\textsuperscript{t1} [F0E0?] \{F1;M2\}\textsuperscript{t2}
\end{styleStandard}

\begin{styleStandard}
\ \ Scenario 3:\ \ \ \ \{F1;M1\}\textsuperscript{t1} [F0E0?] \{F2;M1\}\textsuperscript{t2}
\end{styleStandard}

\begin{styleStandard}
Reanalysis can take place both at the morphological/syntactical level F and at the semantic level M (Scenario 1), or it can be restricted to one level (Scenarios 2 and 3). Diachronic semantics deals with changes in Scenarios 1 and 2. The important claim of the formal semantic studies is that, at both points of time (t1 and t2), F suits M in a compositional manner—this is necessary for a reanalysis to take place (see Eckardt 2006, Bar-Asher Siegal 2015, Aloni 2021).
\end{styleStandard}

\subsection[2.7 Open questions on cualquiera and outlook to this book]{\textbf{2.7 Open questions on }\textbf{\textit{cualquiera}}\textbf{\textit{ }}\textbf{and outlook to this book}}
\begin{styleStandard}
\ In this section, I discuss open questions and present the outlook to this book.
\end{styleStandard}

\subsubsection[2.7.1 Questions on synchrony]{\textbf{2.7.1 Questions}\textbf{ on synchrony}}
\begin{styleStandard}
There is no detailed discussion or definition of the non-modal meaning of \textit{cualquiera} ‘ordinary/common’ (which I have called the evaluative meaning, EVAL 1) in predicative contexts in Modern Spanish and how this meaning of \textit{cualquiera} is related to other meanings such as the free choice meaning ‘any’ and the random choice meaning ‘random’. The examples discussed in the literature seem to indicate that this evaluative meaning of \textit{cualquiera} depends on the verb class, i.e. it is specifically found with non-volitional verbs (e.g. \textit{copular} verbs, episodic tense, perception verbs) (see 2.3). EVAL 2 is also present with verbs such as \textit{hacer/decir cualquiera} ‘do/say nonsense’ in ArgSp and French \textit{n’importe quoi} (see 2.1). I will therefore look into verb classes more systematically and address the question how the definition of EVAL can capture the variation in meanings under predicative verbs such as the meaning ‘common’ and the deragotary version of it, such as ‘nothing special/unremarkable’ (EVAL 1) and the meaning as ‘nonsense’ under volitional verbs such as \textit{hacer/decir} (EVAL 2). This meaning variation correlates to the variation in varieties of Spanish (Argentinian vs European Spanish) as the meaning ‘nonsense’ (EVAL 2) is only present in Argentinian Spanish (see Chapter 3 for EurSp and Chapter 4 for ArgSp). 
\end{styleStandard}

\begin{styleStandard}
\ \ Another question that needs to be answered is what syntactic and semantic status the postnominal \textit{cualquiera} has in European Spanish. Is it a determiner or something else? And what status does the evaluative \textit{cualquiera} have under volitional verbs in Argentinian Spanish? Is it an indefinite pronoun or a noun?
\end{styleStandard}

\begin{styleStandard}
\ \ It also remains to be seen whether all readings of \textit{cualquiera} or at least a subset of them can be derived from a more general meaning of \textit{cualquiera} or whether \textit{cualquiera} is lexically ambiguous as A\&M (2018) suggest (see Section 2.2.3). My working hypothesis is that the evaluative meaning is the result of a reanalysis of the\textbf{ }\textbf{\textit{modal}}\textit{ cualquiera} (i.e. \textit{cualquiera} that is used under a modal) in certain syntactic and semantic contexts in diachrony. I will address these issues within the diachronic study of \textit{cualquiera} in Chapters 5 and 6.
\end{styleStandard}

\begin{styleStandard}
\ \ The classification system of \textit{cualquiera} chosen by Aloni et al. (2010) in 2.5 is not satisfactory, because it mixes up syntactic and semantic properties of \textit{cualquiera} and does not correlate them. For instance, conditional antecedents and comparative clauses are syntactic structures that can license \textit{cualquiera}. Consequently, conditional antecedents and comparative clauses are syntactic contexts, as both sentence types can license FCIs. From the results in Figure 2 in 2.5, we do not learn whether different functions correlate with different syntactic positions of \textit{cualquiera} (e.g. +/– prenominal position). The classification map in Figure 1 in 2.5 does not contain a definition of the ‘no matter’ Interpretation, making it difficult to judge how it is different from the indiscriminative interpretation (IND). Moreover, the RC and evaluative interpretation are absent from the classification system and the frequency distribution in Figures 1 and 2 in Section 2.5. Given the critical points mentioned with respect to 2.5, I decided to use a more detailed syntactic classification of \textit{cualquiera} for my corpus analysis of the word in synchrony and diachrony (see Chapters 3, 4, and 6). I will describe the evaluative function of \textit{cualquiera} in more detail in Chapters 3 and 4 and relate it to its grammatical function (i.e. N \textit{cualquiera} in Chapter 3 and \textit{decir cualquiera} in Chapter 4). The diachronic study in Chapters 5 and 6 will help me to understand the syntactic and semantic status of \textit{cualquiera} in different positions and different varieties of Spanish. 
\end{styleStandard}

\subsubsection[2.7.2 Questions on diachrony and the diatopic variation of cualquiera]{\textbf{2.7.2 Questions on diachrony and the diatopic variation of }\textbf{\textit{cualquiera}}}
\begin{styleStandard}
One important diachronic question is how the evaluative meaning \textit{cualquiera} came into existence in European and Argentinian Spanish. Does it have the same origin? In order to answer this question I will study the diachrony of \textit{cualquiera} with EVAL 2 in Chapter 4 and the diachrony of EVAL 1 in Chapter 5. According to many linguists (Blank 1997, Traugott 1989, Eckardt 2006, among others), semantic change is often triggered by pragmatic enrichment of the target utterance. For instance, the new meaning can be the result of reanalysing the old meaning with contextual information (some pragmatic implicatures). The open question is whether the evaluative meaning of \textit{cualquiera} in European and Argentinian Spanish is based on some pragmatic reanalysis in a compositional manner. 
\end{styleStandard}

\begin{styleStandard}
\ \ I will argue in Chapters 3, 4, and 5 that the evaluative meaning of \textit{cualquiera} is the result of using free choice indefinites in certain syntactic positions, such as the predicative position. I will show that the appearance of the evaluative meaning of \textit{cualquiera} is an example of the Scenario 1 in . Moreover, I will argue that the evaluative meaning is the result of pragmatic strengthening because the free choice meaning (the original meaning) is enriched by contextual information in certain pragmatic contexts in which some attitude holder x expresses her evaluation towards the proposition ‘every possibility is an option’ (FCI) or towards the proposition ‘every possibility matches agent’s decisions’ (RCI) expressed by \textit{cualquiera}. This contextual interpretation has triggered a reanalysis of \textit{cualquiera} into an evaluative predicate in European and Argentinian Spanish with further reanalysis into a degree predicate in Argentian Spanish. I dub my analysis ‘modal indefinites under attitudes’. My analysis thus combines research from two different camps: modal and polarity sensitive items as discussed in this section and semantics dealing with attitudes and evaluation (e.g. personal taste predicates) (see Heim 1992, Lasersohn 2005, among others), which I will introduce in Chapters 3 and 4.
\end{styleStandard}

\clearpage\section[3 Evaluative meaning of postnominal cualquiera in European Spanish]{\textbf{3 Evaluative meaning of postnominal }\textbf{\textit{cualquiera }}\textbf{in European Spanish}}
\begin{styleStandard}
As I already showed in Chapter 2, the evaluative reading of (\textsc{un)} N \textit{cualquiera} ‘common, ordinary’ can appear after predicative verbs, such as the copular (like) verbs \textit{ser} ‘to be’ and\textit{ parecer} ‘to seem’ (see 2.1.2). In this chapter, I investigate the evaluative meaning of (\textsc{un)} N\textit{ cualquiera} in more detail. In order to do so, I will start with the description of (\textsc{un)} N \textit{cualquiera} in predicative contexts and extend this description to other verbs.
\end{styleStandard}

\begin{styleStandard}
\ \ My observations are primarily based on corpus data from the corpus \textit{Corpus del Español} (CdE), in particular the following subcorpora: CdE\_web dialects\_EurSp. and CdE\_now\_EurSp. On particular occasions, I also use speaker judgements. The crucial result of my observations, which will serve the analysis of \textit{cualquiera}, is that predicative \textsc{un} N \textit{cualquiera} appears extremely often (ca. 98\% of all data), especially under negation (around 80\%), see 3.1.
\end{styleStandard}

\begin{styleStandard}
\ \ As I will argue in the subsequent sections (3.2–3.5), negation and the predicative position of \textsc{un} N \textit{cualquiera} trigger a variety of readings such as ‘an N similar to anyone else’, ‘a common N’, ‘an ordinary N’ which I subsumed under the term EVAL to distinguish it from modal readings such as ‘any’ and ‘random’. In order to explain certain types of data, such as coordination with adjectives, one needs to postulate a reanalysis of \textit{cualquiera} into a (frequency) adjective (e.g. with the meaning ‘common’) in European Spanish (see 3.5). I will argue that negation plays an important role in this reanalysis, as it introduces a contrastive property into the discourse, which strengthens the property interpretation of postnominal \textit{cualquiera} as ‘common’ or ‘unremarkable’.
\end{styleStandard}

\begin{styleStandard}
\ \ The chapter is organised as follows. In 3.1, I first present some data concerning the distribution of postnominal \textit{cualquiera}, which show that it typically occurs in predicative syntactic contexts, and, most frequently, under negation. I will then describe the evaluative function in more detail in 3.2, where I will show that the evaluative function of postominal \textit{cualquiera} comes in the three readings mentioned above. I will also show that what I call the ‘similarity’ reading can be assumed as the basic meaning of postnominal \textit{cualquiera}, from which the ‘common’ and the derogative meaning can be derived. In 3.3, I present and discuss some additional properties of postnominal \textit{cualquiera}, in particular the possibility of an overt experiencer, the extensional meaning of the ‘common’ reading and its relationship to the free choice reading, and the very limited possibility of postnominal \textit{cualquiera} under modal verbs. After a brief intermediate summary in 3.4, I will show in 3.5 that \textit{cualquiera} in its ‘similarity’ reading can be derived from the FCI function and can consequently be analysed as a quantifier under a generic operator in predicative and episodic contexts. 3.6 will offer an analysis according to which \textit{cualquiera}, when it has the ‘common’ reading, is no longer interpreted as a quantifier, but as an item referring to a property of types of individuals. I will argue that both analyses, the quantifier analysis and the adjectival analysis are needed in order to capture the corpus data and the ambiguity between the free choice and the evaluative interpretation of \textit{cualquiera} observed in the literature (see Chapter 2, in particular my discussion of A\&M 2011, 2018).
\end{styleStandard}

\subsection[3.1 Distribution of postnominal cualquiera in the Corpus del Espanol\_web dialects]{\textbf{3.1 Distribution of postnominal }\textbf{\textit{cualquiera }}\textbf{in the Corpus del Espanol\_web dialects}}
\begin{styleStandard}
In this subsection. I only used the subcorpus Corpus del Espanol\_web dialects (henceforth: CdE\_web dialects) because other subcorpora did not allow for the automatic extraction of frequencies in European Spanish by the time I collected data in CdE in 2016.\footnote{\ \textrm{In the meantime, many other subcorpora appeared that are much bigger and also allow for automatic frequency extraction. In the future, I will make a comparison between these subcorpora.}} Henceforth, I will use square brackets [ ] for the different constructions I searched for in this corpus, e.g. [V \textsc{un} N \textit{cualquiera}] (= postnominal \textit{cualquiera} after the verb).
\end{styleStandard}

\begin{styleStandard}
\ \ All examples that I have examined show that 82\% of [V \textsc{un} N \textit{cualquiera}] occur under implicit negation as in \footnote{\textrm{\ I counted as implicit negation only those cases where the context was very clear that the speaker does not believe that the proposition expressed by V }\textsc{un}\textrm{ N }\textrm{\textit{cualquiera}}\textrm{ is true.}}\ or explicit negation as in , whereby explicit negation is a more frequent type (99\% of total cases with negation):
\end{styleStandard}

\begin{styleStandard}
\label{bkm:Ref42868256}\ \ ¿Sigues pensando que fue un día cualquiera? [...] es lunes 5 de diciembre de 2011.
\end{styleStandard}

\begin{styleStandard}
\ \ ‘Are you still thinking that it was a normal day? […] it’s Monday, December 5, 2011.’
\end{styleStandard}

\begin{styleStandard}
\ \ (@ 40)
\end{styleStandard}

\begin{styleStandard}
\label{bkm:Ref42868258}\ \ [...] que esto no es un libro cualquiera, esto es un libro de James Dashner.
\end{styleStandard}

\begin{styleStandard}
\ \ ‘[…] that this is not an ordinary book, this is a book from James Dasner.’
\end{styleStandard}

\begin{styleStandard}
\ \ (@ 41)
\end{styleStandard}

\begin{styleStandard}
Affirmative contexts such as (see \ and ) appear less frequently than negative contexts (see Table 5). Recall that affirmative contexts represent usually a puzzle for the licensing of \textit{cualquiera} as an FCI (see Section 2.2) unless rescued by some operators (see Section 2.3). 
\end{styleStandard}

\begin{styleStandard}
\ \ In affirmative contexts (i.e. contexts without explicit or implicit negation) the verb in [V \textsc{un} N \textit{cualquiera}] can be of different types, such as transitive verbs and predicative verbs, that is, verbs that take predicative complements as their argument (see \ and ). Predicative verbs are the most frequent type of verbs in this corpus as will be shown in Table 6. They are often realised as a copular verb in present indicative, as in \ and :
\end{styleStandard}

\begin{styleStandard}
\label{bkm:Ref42868883}\ \ Parece una chiquilla cualquiera, vestida con una túnica liviana que adorna su cuerpo hermoso y esconde el brazo derecho mutilado.
\end{styleStandard}

\begin{styleStandard}
\ \ ‘She looks like a common little girl, ….’ 
\end{styleStandard}

\begin{styleStandard}
\ \ (@ 42)
\end{styleStandard}

\begin{styleStandard}
\label{bkm:Ref42868888}\ \ De todos modos, tampoco os esperéis la gran historia. Es una historia cualquiera.
\end{styleStandard}

\begin{styleStandard}
\ \ ‘Anyways, don’t expect a big story either. It’s a story like any other.’
\end{styleStandard}

\begin{styleStandard}
\ \ (@ 43) 
\end{styleStandard}

\begin{styleStandard}
There are less occcurrences of other types of verbs that are not copular verbs, e.g. transitive verbs such as \textit{prestó} and \textit{puse} in \ and . Non-copular verbs that appear in [V \textsc{un} N \textit{cualquiera}] are usually realised with perfective aspect, expressed in Spanish with the \textit{pretérito indefinido} (in so-called episodic contexts; see section 2.3 for discussion of episodic contexts):
\end{styleStandard}

\begin{styleStandard}
\label{bkm:Ref42868870}\ \ La NYK \textit{prestó} un contenedor cualquiera a la BBC.
\end{styleStandard}

\begin{styleStandard}
\ \ ‘The NYK borrowed a random container to the BBC.’
\end{styleStandard}

\begin{styleStandard}
\ \ (@ 44)
\end{styleStandard}

\begin{styleStandard}
\label{bkm:Ref42869017}\ \ Le \textit{puse} un nombre cualquiera. 
\end{styleStandard}

\begin{styleStandard}
\ \ ‘I put some random name.’
\end{styleStandard}

\begin{styleStandard}
\ \ (@ 45)
\end{styleStandard}

\begin{styleStandard}
Table 5 summarises the distribution of +/– affirmative mood of [V \textsc{un} N \textit{cualquiera}]. It appears in negative contexts (with implicit or explicit negation) more often than in affirmative ones:\footnote{\ In this and the following tables, the first line shows the percentages and the second line the absolute numbers of occurrences.} 
\end{styleStandard}

\begin{center}
\tablefirsthead{}
\tablehead{}
\tabletail{}
\tablelasttail{}
\begin{supertabular}{|m{1.5906599in}|m{1.9893599in}|m{1.9886599in}|}
\hline
\centering \textbf{V + }\textbf{\textsc{un}}\textbf{ N }\textbf{\textit{cualquiera}} &
\centering \textbf{V + }\textbf{\textsc{un}}\textbf{ N }\textbf{\textit{cualquiera}}\textbf{ in affirmative contexts} &
\centering\arraybslash \textbf{Neg + V + }\textbf{\textsc{un}}\textbf{ N }\textbf{\textit{cualquiera}}\\\hline
\centering 100\% &
\centering 18\% &
\centering\arraybslash 82\%\\\hline
\centering 353 &
\centering 66 &
{\centering 287\par}

{\centering (explicit negation = 283 = 99\% \par}

\centering\arraybslash of total negation)\\\hline
\end{supertabular}
\end{center}
\begin{styletablecaption}
Table \stepcounter{Table}{\theTable}: Distribution of +/– affirmative contexts of V+ \textsc{un}\textit{ }N\textit{ cualquiera }in CdE\_web dialects in European Spanish.
\end{styletablecaption}

\begin{styleStandard}
The most frequent verb type that precedes \textsc{un} N \textit{cualquiera} is the copular verb \textit{ser }‘be’ (98\% of all verbs), whose most frequent verb form (verb token) is the indicative present tense (\textit{es }‘is’, 88\%). The less frequent verb type are transitive verbs such as \textit{poner} ‘put’:
\end{styleStandard}

\begin{center}
\tablefirsthead{}
\tablehead{}
\tabletail{}
\tablelasttail{}
\begin{supertabular}{|m{1.5906599in}|m{1.9893599in}|m{1.9886599in}|}
\hline
\centering \textbf{V + }\textbf{\textsc{un}}\textbf{ N }\textbf{\textit{cualquiera}} &
\centering \textbf{Ser + }\textbf{\textsc{un}}\textbf{ N }\textbf{\textit{cualquiera}} &
\centering\arraybslash \textbf{Other verbs}\\\hline
\centering 100\% &
\centering 98\% &
\centering\arraybslash 2\%\\\hline
\centering 353 &
\centering 345 &
\centering\arraybslash 8\\
 &
\centering \textit{es} (88\% = 308) &
\centering\arraybslash \textit{puse, prestó}\\\hline
\end{supertabular}
\end{center}
\begin{styletablecaption}
Table \stepcounter{Table}{\theTable}: Distribution of verbs in V + \textsc{un}\textit{ }N\textit{ cualquiera }in CdE\_web dialects in European Spanish.\footnote{\textrm{\ Other determiners were disregarded due to low frequency.}}
\end{styletablecaption}

\begin{styleStandard}
The observation that the most frequent verb token is\textit{ es }3\textsuperscript{rd} person singular and not plural suggests that \textsc{un} N \textit{cualquiera }is usually used by speakers to describe a property of a single person or entity and not a property of different entities or people. As we will see in Section 3.2, the most frequent token is \textit{un día cualquiera} ‘an ordinary day’. Thus, \textsc{un} N\textit{ cualquiera} is mostly used to describe the day of speaker’s reference as ordinary or usual. 
\end{styleStandard}

\begin{styleStandard}
\ \ With respect to the distribution of the determiner preceding N \textit{cualquiera}, the most frequent determiner is the indefinite determiner \textsc{un}, as already observed in the literature (see Section 2). However there are few cases of N \textit{cualquiera} being preceded by definite determiners, demonstratives, and quantifiers (see Section 3.6 for data and analysis). In the following sections I will focus on \textsc{un} N \textit{cualquiera}, as it is the most frequent type in European Spanish. As predicative \textsc{un} N \textit{cualquiera} is a very important type in EurSp according to my corpus analysis, I will investigate this type more closely, especially the meanings of \textsc{un} N \textit{cualquiera} under predicative verbs.
\end{styleStandard}

\subsection[3.2 Interpretations of postnominal cualquiera]{\textbf{3.2 Interpretations of postnominal }\textbf{\textit{cualquiera}}}
\begin{styleStandard}
One heuristic for approaching the semantic readings of \textsc{un} N \textit{cualquiera} is to investigate the range and types of appositives and adverbial modifications that co-occur with \textsc{un} N\textit{ cualquiera}, as they can provide us with additional information about the syntactic and semantic status of \textit{cualquiera} (see also Condoravdi 2015: 227 for English). I have looked at the possible appositives and adverbial modifications that \textsc{un} \textit{N cualquiera} can have or appear with in the corpus data CdE\_web dialects and CdE\_now.\footnote{\textrm{\ I have especially investigated the most frequent nouns that have more than five occurences, namely expressions of time (e.g. }\textrm{\textit{día}}\textrm{ ‘day’, 76 occ.) and animate nouns (e.g. }\textrm{\textit{chico}}\textrm{ ‘guy’, 8 occ.). For the sake of simplicity, I ignored nouns that referred to mathematical objects (e.g. variables, points). In mathematical contexts, }\textrm{\textit{es }}\textsc{un}\textrm{ N}\textrm{\textit{ cualquiera}}\textrm{ means that the value of some point, number, or variable is undetermined. This can be paraphrased with a modal meaning that expresses that X can have any possible value or that one can choose any possible value (see the discussion in 2.3.3):}\par 
\setcounter{listWWNumxixleveli}{0}
\begin{listWWNumxixleveli}
\item 
\begin{styleFootnote}
\textrm{X es una variable cualquiera.}
\end{styleFootnote}
\end{listWWNumxixleveli}
\textrm{‘X is an unspecified/undetermined variable’}\par \textrm{‘X can have any possible value.’}\par \textrm{Mathematical objects seem to allow modal paraphrasis, something that needs to be studied more carefully in the future.}} In order to search for adverbial modification automatically, I searched for all strings after \textsc{un} N \textit{cualquiera }that do not start with sentence ending punctuation, such as a question mark (?). This way the search gave me all strings included in the sentence (e.g. \textsc{un} N \textit{cualquiera}\textbf{,} como…).
\end{styleStandard}

\begin{styleStandard}
\ \ I will show that \textsc{un} N \textit{cualquiera }has mostly the evaluative interpretation (EVAL), which has different flavours depending on the context (COM(mon), DER(ogatory), SIM(ilarity)). As will be shown in the following section, the most frequent flavour of EVAL is SIM. This reading has not been observed previously in the literature, in contrast to the ‘unremarkable’ interpretation of \textit{cualquiera} (see 2.3). Even the ‘unremarkable’ interpretation, which has been already observed in the literature in relation with \textit{cualquiera} (see 2.3), appears in two different flavours: COM and DER. These different flavours of the ‘unremarkable’ interpretation have not been studied before in the literature in relation with \textit{cualquiera}. My study of these different flavours of EVAL is, to my knowledge, the first study to investigate EVAL on the basis of corpus data.
\end{styleStandard}

\subsubsection[3.2.1 Readings in affirmative contexts under indicative verbal mood]{\textbf{3.2.1 Readings in affirmative contexts under indicative verbal mood}}
\begin{styleStandard}
The affirmative sentence \textit{Es} \textit{un }N\textsubscript{[animate]}\textit{ cualquiera} often appears with a comparative expression which expresses similarity between the referent to which the property \textsc{un }N\textsubscript{[animate]}\textit{ cualquiera} applies (here: \textit{Juan}) and other individuals of the same comparison class (here: ‘everyone else like you and me’) (hence: SIM):
\end{styleStandard}

\begin{styleStandard}
\ \ Juan es un hombre cualquiera, cargado de errores. Como tú, o como yo.
\end{styleStandard}

\begin{styleStandard}
\ \ ‘Juan is a human being like everyone else, full of flaws. Like you, or like me.’
\end{styleStandard}

\begin{styleStandard}
\ \ (similar to @ 46)
\end{styleStandard}

\begin{styleStandard}
The similarity construction is expressed by the comparison term \textit{como} ‘like’, which expresses similarity between the referent and a high number of individuals that share some (basic or prototypical)\footnote{\textrm{\ The assumption that }\textrm{\textit{common}}\textrm{ properties must be basic follows from the similarity of all individuals denoted by N }\textrm{\textit{cualquiera}}\textrm{. For instance, }\textrm{\textit{un trabajo cualquiera}}\textrm{ is a job that shares with other jobs the most basic properties, whatever it is that defines a job as a job, such as getting up early, working for eight hours, going back home, etc.}} properties or at least one prototypical property with the subject. It is impossible to find a comparison to a set of individuals that contains only one individual or none:
\end{styleStandard}

\begin{styleStandard}
\ \ Juan\textsubscript{j} es un hombre cualquiera, cargado de errores. \# No hay otro hombre que tenga errores como él\textsubscript{j}.
\end{styleStandard}

\begin{styleStandard}
\ \ ‘Juan is a human being like everyone else, full of flaws. \#There is no one else like him.’
\end{styleStandard}

\begin{styleStandard}
In the following example, there is a similarity reading, which in this case expresses that the morning the speaker refers to is a (kind of) morning that is similar to all or many other mornings. Indeed, native speakers deem it possible for \ to have the same interpretation as in . This test shows that even though the comparison is not overtly expressed, speakers can infer an implicit comparison from \textsc{un} N \textit{cualquiera} (see 3.5.1):
\end{styleStandard}

\begin{styleStandard}
\label{bkm:Ref42870625}\ \ Fue una mañana cualquiera, como cada mañana. Había estado leyendo y llamó uno de los muchachos de la cantina.
\end{styleStandard}

\begin{styleStandard}
\ \ ‘It was a morning like any other mornings. I had been reading and one of the guys from the canteen called.’
\end{styleStandard}

\begin{styleStandard}
\ \ (similar to @ 46)
\end{styleStandard}

\begin{styleStandard}
\label{bkm:Ref42870616}\ \ Fue una mañana cualquiera. Había estado leyendo y llamó uno de los muchachos de la cantina.
\end{styleStandard}

\begin{styleStandard}
\ \ ‘It was a morning like any other. I had been reading and one of the guys from the canteen called.’
\end{styleStandard}

\begin{styleStandard}
\ \ (@ 47)
\end{styleStandard}

\begin{styleStandard}
There are also appositives that describe N \textit{cualquiera} as N\textit{ corriente} ‘regular/normal N’ (henceforth: COM) as in \ or N \textit{cualquiera} itself can appear as an appositive clause that describes N \textit{normal} ‘normal N’ as in :
\end{styleStandard}

\begin{styleStandard}
\label{bkm:Ref57554115}\ \ Hoy es un día cualquiera, un día corriente, [...]
\end{styleStandard}

\begin{styleStandard}
\ \ ‘Today is a day like many other days, a normal day, …’
\end{styleStandard}

\begin{styleStandard}
\ \ (@ 48)
\end{styleStandard}

\begin{styleStandard}
\label{bkm:Ref63545479}\ \ Yo soy una persona normal, un tipo cualquiera.
\end{styleStandard}

\begin{styleStandard}
\ \ ‘I am a normal person, a guy like other guys.’
\end{styleStandard}

\begin{styleStandard}
\ \ (@ 49, similar to @ 50)
\end{styleStandard}

\begin{styleStandard}
Finally, some appositives clearly indicate that \textsc{un} N \textit{cualquiera} can have a derogatory meaning (DER). The most frequent contexts in which \textsc{un} N \textit{cualquiera} has this interpretation contain animate nouns. In these contexts, N \textit{cualquiera} is described as not worth mentioning:
\end{styleStandard}

\begin{styleStandard}
\ \ era un tipo cualquiera, sin ninguna particularidad digna de mención.
\end{styleStandard}

\begin{styleStandard}
\ \ ‘he was an ordinary guy, without any particular quality worth mentioning.’
\end{styleStandard}

\begin{styleStandard}
\ \ (@ 51)
\end{styleStandard}

\begin{styleStandard}
In the following example, the deragotory reading is expressed by the use of the adverb \textit{simplemente} ‘just’ or \textit{nada más} ‘nothing but’, which expresses that the person described by N \textit{cualquiera }has no properties other than that described by N \textit{cualquiera}:
\end{styleStandard}

\begin{styleStandard}
\ \ Era simplemente un tipo cualquiera.
\end{styleStandard}

\begin{styleStandard}
\ \ ‘It was just an ordinary guy.’
\end{styleStandard}

\begin{styleStandard}
\ \ (@ 52)
\end{styleStandard}

\begin{styleStandard}
\ \ No soy nada más que un chico cualquiera, que pasaría desapercibido, sin más, normal, que se diría.
\end{styleStandard}

\begin{styleStandard}
\ \ ‘I am nothing but an ordinary boy, that would go unnoticed, nothing more, normal, that one would say.’ 
\end{styleStandard}

\begin{styleStandard}
\ \ (similar to @ 53)
\end{styleStandard}

\begin{styleStandard}
All three flavours are related to each other. SIM implies COM as ‘to be similar to anyone else’ implies ‘to be normal’. The deragotory reading expresses a negative judgement of ‘being normal or similar to anyone else’, which can be paraphrased as ‘being just normal’ and ‘not special’. I will investigate this semantic relation more closely in Section 3.6.
\end{styleStandard}

\begin{styleStandard}
\ \ To sum up so far, in affirmative (predicative) contexts, \textsc{un} N \textit{cualquiera} refers to entities or individuals that share common properties or properties that everyone else has. Thus, \textit{un día} \textit{cualquiera} expresses ‘a day like any other day’. The same nouns can have a ‘common/frequent’ or a ‘deragotory’ reading. A formal analysis of \textit{N cualquiera} as a type/property of individuals will be proposed in 3.6.
\end{styleStandard}

\subsubsection[3.2.2 Readings under explicit and implicit negation]{\textbf{3.2.2 Readings under explicit and implicit negation}}
\begin{styleStandard}
Under implicit negation expressed by the counterfactual interpretation of the verb \textit{sería} in \ or the question context in , \textsc{un} N \textit{cualquiera} can have an evaluative interpretation with all flavours described in 3.2.1, namely COM, SIM, and DER:
\end{styleStandard}

\begin{styleStandard}
\label{bkm:Ref54634548}\ \ Pensaron que sería un día cualquiera. Pero no. La basura en el mar hizo que ese viernes sea diferente.
\end{styleStandard}

\begin{styleStandard}
\ \ ‘They thought that it would have been a common day (COM)/a day like any other day (SIM)/just a normal day (DER). But no. The rubbish in the sea turned this Friday into a special day.’
\end{styleStandard}

\begin{styleStandard}
\ \ (@ 54, similar to @ 55)
\end{styleStandard}

\begin{styleStandard}
\label{bkm:Ref54634539}\ \ ¿Pensabas que este sábado sería un sábado cualquiera?
\end{styleStandard}

\begin{styleStandard}
\ \ ‘You thought that this Saturday would be a normal (COM) Saturday/a Saturday like all other Saturdays (SIM)/just a normal Saturday (DER)?’
\end{styleStandard}

\begin{styleStandard}
\ \ (@ 56, similar to @ 57)
\end{styleStandard}

\begin{styleStandard}
\textit{N cualquiera} under (explicit or implicit) negation is often used to express \textit{a contrast}\textbf{\textit{ }}to a (proper) noun or some other DP that contains a reference to a single individual or type of individual. This individual/entity has the property of being ‘particular’, ‘distinguished’, ‘uncommon’, or having ‘higher value’ than \textit{N cualquiera}, such as Messi (who is known to be an extraordinary football player) in \ and the author John Dashner\footnote{\ The further context of (145) (from a book review) reveals that this author is particularly distinguished for writing books full of suspense and action.} in :
\end{styleStandard}

\begin{styleStandard}
\label{bkm:Ref66615171}\ \ Sin Messi, el Barcelona sería un equipo cualquiera.
\end{styleStandard}

\begin{styleStandard}
\ \ ‘Without Messi, the team Barcelona would be some unimportant football team.’
\end{styleStandard}

\begin{styleStandard}
\ \ (@ 58)
\end{styleStandard}

\begin{styleStandard}
\label{bkm:Ref150779604}\ \ […] que esto no es un libro cualquiera, esto es un libro de James Dashner.
\end{styleStandard}

\begin{styleStandard}
\ \ ‘[…] that this is no ordinary book, this is a book by James Dashner.’
\end{styleStandard}

\begin{styleStandard}
\ \ (@ 41)
\end{styleStandard}

\begin{styleStandard}
To sum up, the negation can trigger a contrastive clause expressed by a proper noun (e.g. Messi) or some DP (\textit{un libro de James Dashner}) which singles out one and only one alternative that has the property of being ‘particular’, ‘uncommon’, or ‘outstanding’. This alternative is opposed to the individuals described by \textsc{un} N \textit{cualquiera}.
\end{styleStandard}

\begin{styleStandard}
\ \ As I will suggest in my analysis, predicative verbs represent a special linguistic context that gives rise to reanalysis of N \textit{cualquiera} into a property type and thus a shift from a quantifier \textit{cualquiera} into a property \textit{cualquiera} (see 3.6). 
\end{styleStandard}

\subsubsection[3.2.3 The semantic/pragmatic status of meanings]{\textbf{3.2.3 The semantic/pragmatic status of }\textbf{meanings}}
\begin{styleStandard}
The question arises as to which of the readings identified in 3.2.1 and 3.2.2 is a lexical meaning of \textit{cualquiera} and which one is just an implicature and then, whether one can cancel this implicature and thus can cancel the meaning of postnominal \textit{cualquiera} (see Aloni 2021, A\&M 2018 on different views).
\end{styleStandard}

\begin{styleStandard}
\ \ As has been argued by Grice (1989), conversational implicatures are usually cancellable, because they highly depend on contextual factors. Conventional implicatures usually resist cancellation because they convey a side comment by the speaker and do not contribute to the at-issue meaning of the sentence (Potts 2005, Simons et al. 2010).
\end{styleStandard}

\begin{styleStandard}
\ \ Let us see whether one of the readings is cancellable. The following examples were presented to five native speakers for acceptability judgements. The word \textit{cualquiera} is not translated into English, because it was up to the native speakers to decide which meaning best fits the context introduced by the \textit{pero} ‘but’ clause:
\end{styleStandard}

\begin{styleStandard}
\label{bkm:Ref42976554}\ \ a.\ \ \#Juan es un hombre cualquiera, pero no es como los demás.
\end{styleStandard}

\begin{styleStandard}
\ \ \ \ ‘Juan is a man, but he is not like others.’
\end{styleStandard}

\begin{styleStandard}
\ \ b.\ \ ?Juan es un hombre cualquiera, pero no es un hombre común.
\end{styleStandard}

\begin{styleStandard}
\ \ \ \ ‘Juan is a man, but he is not ordinary.’
\end{styleStandard}

\begin{styleStandard}
\ \ c.\ \ Juan es un hombre cualquiera, pero no es un hombre sin nivel.
\end{styleStandard}

\begin{styleStandard}
\ \ \ \ ‘Juan is a man, but he is not a man without class.’
\end{styleStandard}

\begin{styleStandard}
None of the speakers accepted the sentence in a), because they found that it triggered a contradiction. The first clause in a) expresses that Juan is a man like others, while the second clause expresses the opposite, namely that Juan is not like others, resulting in a contradiction. The sentence in c) clearly does not lead to a contradiction, because it expresses that Juan is like many others, but he is not a man without class or a simple man. The sentence in b) shows that the COM is also not cancallable, which is demonstrated by the mark “?”. The reason why the meaning ‘ordinary’ triggered different judgements in sentence b) among the speakers is probably that \textit{común} ‘ordinary’ is itself ambiguous between the deragotory meaning ‘ordinary, not special’ and the ‘similarity’ meaning, that is, ‘ordinary, not different from others’. The question is now how to interpret these results where ‘similarity’ and ‘common’ cannot be cancelled and where ‘of low value’ or ‘just ordinary’ can be? I take this to mean that ‘similarity’ and ‘common’ are entailed in the lexical meaning of \textit{cualquiera} and that the deragotory meaning is a contextual meaning derived pragmatically.
\end{styleStandard}

\begin{styleStandard}
\ \ Assuming that the lexical meaning of \textit{cualquiera} is its free choice meaning that triggers alternatives (see Chapter 2), a reasonable assumption would be that SIM and COM are semantically related to the free choice ‘any’. I will suggest such an analysis in 3.5, where I will show that SIM is a special case of ‘any’. I will assume that COM is entailed in SIM (i.e. being similar to anyone else entails having common properties) and DER is a contextual interpretation of both COM and SIM (i.e. having common properties is judged contextually as having not-extraordinary properties).
\end{styleStandard}

\begin{styleStandard}
\ \ In the following section I will investigate further semantic properties of \textsc{un} N \textit{cualquiera} that will help us to define and characterise its meanings more precisely. 
\end{styleStandard}

\subsubsection[3.2.4 +/– Evaluative cualquiera]{\textbf{3.2.4 +/– Evaluative }\textbf{\textit{cualquiera}}}
\begin{styleStandard}
Evaluative\textit{ cualquiera} with readings such as ‘common’ should be distinguished from FC \textit{cualquiera} with the reading ‘any’. Both interpretations can be made explicit in the following example:
\end{styleStandard}

\begin{styleStandard}
\ \ la existencia de un delito cualquiera
\end{styleStandard}

\begin{styleStandard}
\ \ ‘the existence of any criminal act’ \ \  \ \ \ \ \ \ \ \ \ \ \ \ \ \ \ \ \ \ \ \ \ \ \ \ \ \ \ \ \ \ \ \ \ \textit{cualquiera} (FCI)
\end{styleStandard}

\begin{styleStandard}
\ \ ‘the existence of a common criminal act’\ \  \ \ \ \ \ \ \ \ \ \ \ \ \ \ \ \ \ \ \ \ \ \ \ \ \ \ \ \ \ \ \ \ \ \textit{cualquiera} (EVAL)
\end{styleStandard}

\begin{styleStandard}
\ \ (@ 69)
\end{styleStandard}

\begin{styleStandard}
Under the FCI reading, the probability of any sort of criminal act is calculated as a mean of every single criminal act (including a very infrequent and unusual one such as a murder), whereas the probility of a common criminal act is not the probability of any kind of criminal act but the probability of the occurrence of a certain type of criminal act that can be described as a prototypical, usual, or common criminal act (e.g. robbery). The same observation applies to other readings as well. A criminal act that is similar to many other criminal acts is a certain type of criminal act, namely the one that shares properties with many other criminal acts. An unremarkable criminal act is also a certain type of criminal act. Only in the FCI case can \textit{cualquiera} be modified by a subjunctive relative clause:
\end{styleStandard}

\begin{styleStandard}
\ \ la existencia de un delito cualquiera que sea
\end{styleStandard}

\begin{styleStandard}
\ \ ‘the existence of any criminal act that there might be’ FCI only
\end{styleStandard}

\begin{styleStandard}
\ \ \#‘the existence of a common criminal act that there might be’\ \ 
\end{styleStandard}

\begin{styleStandard}
\ \ (similar to @ 70)
\end{styleStandard}

\begin{styleStandard}
Assuming that there are ony three kinds of crimes—A, B, and C—and assuming further that A and B are common, whereas C is very infrequent, \textit{un delito cualquiera} under the reading ‘common’ denotes a more restricted set of crimes, namely the set \{A, B\}. FCI \textit{cualquiera} does not add this information about the distribution, whereas \textit{cualquiera} on the ‘common/frequent’ reading does. Set-theoretically, ‘common’ is thus a subset of a set denoted by FCI \textit{cualquiera. }The extension of criminal acts is thus distinct under the FCI and ‘common’ readings, with the ‘common/frequent’ reading denoting a subset of the FCI denotation with respect to its extension.
\end{styleStandard}

\begin{styleStandard}
\ \ a.\ \ \textit{cualquiera} (FCI) = \{A,B,C\}
\end{styleStandard}

\begin{styleStandard}
\ \ b.\ \ \textit{cualquiera} (EVAL ‘COM/DER/SIM’) = \{A,B\}
\end{styleStandard}

\begin{styletablecaption}
Figure \stepcounter{Figure}{\theFigure}: The set denotation of \textit{cualquiera} as EVAL is a subset of the FCI ‘any’.
\end{styletablecaption}

\begin{styleStandard}
What does \textit{cualquiera} denote by its intension? I suggest analysing the postnominal \textit{cualquiera} in a way similar to adjectives or nominal compounds that modify nouns and describe a property of a property. I call them predicate modifiers (see Section 3.3):
\end{styleStandard}

\begin{styleStandard}
\ \ Juan es un hombre cualquiera/ordinario/simplón.
\end{styleStandard}

\begin{styleStandard}
\ \ ‘Juan is an ordinary/simple man.’
\end{styleStandard}

\begin{styleStandard}
\ \ (@ 3)
\end{styleStandard}

\begin{styleStandard}
I will show later that this predicate modifier analysis of \textit{cualquiera} can serve us as a basis for explaining some aspects of diatopic variation of Spanish \textit{cualquiera}, in particular some striking differences between European and Argentinian Spanish (see 3.3 and Chapter 4 for details). 
\end{styleStandard}

\subsection[3.3 Postnominal cualquiera as a predicate (modifier)]{\textbf{3.3 Postnominal }\textbf{\textit{cualquiera}}\textbf{ as a predicate (modifier}\textbf{)}}
\begin{styleStandard}
I assume that postnominal \textit{cualquiera} should be analysed as a postnominal modifier like the postnominal adjective \textit{común}, \textit{ordinario }‘common, ordinary’:
\end{styleStandard}


\setcounter{listWWNumxleveli}{0}
\begin{listWWNumxleveli}
\item 
\begin{styleStandard}
Un hombre [\textsubscript{Modifier} \textit{cualquiera}] = Un hombre [\textsubscript{Modifier} \textit{común/ordinario}]. 
\end{styleStandard}
\end{listWWNumxleveli}
\begin{styleStandard}
On this analysis, we should expect that postnominal \textit{cualquiera}, can be coordinated, just like postnominal adjectives. And indeed, we find coordination of \textit{cualquiera} with adjectives such as \textit{mediocre }‘mediocre’\textit{, ordinario} ‘ordinary’, and \textit{normal} ‘normal’.
\end{styleStandard}

\begin{styleStandard}
\textbf{\ \ }para un día cualquiera y normal en el Ciaam. 
\end{styleStandard}

\begin{styleStandard}
\ \ ‘for a common and normal day in Ciaam’
\end{styleStandard}

\begin{styleStandard}
\ \ (@ 79)
\end{styleStandard}

\begin{styleStandard}
\ \ y no es para llenar un currículo para una hoja de vida ordinaria y cualquiera.
\end{styleStandard}

\begin{styleStandard}
\ \ ‘and it is not to fill in the CV with a page of an ordinary and normal life.’
\end{styleStandard}

\begin{styleStandard}
\ \ (@ 80)
\end{styleStandard}

\begin{styleStandard}
\ \ comprendí que no era más que un ser mediocre y cualquiera. 
\end{styleStandard}

\begin{styleStandard}
\ \ ‘I understood that I was just a mediocre and normal being’
\end{styleStandard}

\begin{styleStandard}
\ \ (@ 81)
\end{styleStandard}

\begin{styleStandard}
However, this kind of coordination is not very frequent in European Spanish (3 occ. in total out of 357), which suggests that the use of postnominal \textit{cualquiera} as an adjective-like category is not very common in this variety of Spanish.
\end{styleStandard}

\begin{styleStandard}
\ \ Recall from Section 2.2, that N \textit{cualquiera} usually co-occurs with indefinite nouns as observed in the literature (see Section 2.2). However, there are few examples in which NP \textit{cualquiera} is preceded by a quantifier or even by a definite determiner.
\end{styleStandard}

\begin{styleStandard}
\ \ In the following examples, the definite noun refers to a certain kind of man, namely to the common man \textit{as a type}. One could also describe this type as \textit{el hombre de la calle}, lit. ‘the man from the street’ or ‘common/normal man’:
\end{styleStandard}

\begin{styleStandard}
\ \ El hombre cualquiera del autobús no siente vergüenza alguna, aunque debería,...
\end{styleStandard}

\begin{styleStandard}
\ \ ‘[…] The common man from the bus does not feel any shame, although he should,…’
\end{styleStandard}

\begin{styleStandard}
\ \ (@ 82)
\end{styleStandard}

\begin{styleStandard}
\ \ Las armas de este hombre cualquiera, de aspecto normal, se resumen en: prudencia
\end{styleStandard}

\begin{styleStandard}
\ \ ‘The weapons of this common, normal-looking man are summarised in: caution’
\end{styleStandard}

\begin{styleStandard}
\ \ (@ 83)
\end{styleStandard}

\begin{styleStandard}
\textit{Cualquiera} occurs also with quantified nouns (see Section 2.1.3 on the discussion concerning quantified nouns):
\end{styleStandard}

\begin{styleStandard}
\label{bkm:Ref58144031}\ \ Esto lo saben mucha gente cualquiera 
\end{styleStandard}

\begin{styleStandard}
\ \ ‘This is known to many ordinary people’
\end{styleStandard}

\begin{styleStandard}
\ \ (@ 84)
\end{styleStandard}

\begin{styleStandard}
\label{bkm:Ref63546096}\ \ Las cifras del IMC se obtienen tras el estudio de los datos de muchísimas chicas cualquiera.
\end{styleStandard}

\begin{styleStandard}
\ \ ‘The figures in the IMC are obtained after studying many ordinary young women.’
\end{styleStandard}

\begin{styleStandard}
\ \ (Fulton, \textit{Seis semanas}, El índice de masa corporal)
\end{styleStandard}

\begin{styleStandard}
\ \ vamos a morir algún día cualquiera. 
\end{styleStandard}

\begin{styleStandard}
\ \ ‘we will die on some ordinary day’
\end{styleStandard}

\begin{styleStandard}
\ \ (@ 85, similar to @ 86)
\end{styleStandard}

\begin{styleStandard}
The occurrence of definite/quantified nouns with postnominal \textit{cualquiera} finds a straight-forward explanation in my analysis of \textit{cualquiera}. It is because \textit{cualquiera} acts in a way similar to adjectives like ‘common, unremarkable, worthless’, and adjectives can modify different types of noun phrases, including definite noun phrases, demonstrative and quantified noun phrases. 
\end{styleStandard}

\begin{styleStandard}
\ \ I will now present a detailed formal analysis of the postnominal \textit{cualquiera} as a predicate (modifier). 
\end{styleStandard}

\subsubsection[3.3.1 Formalisation of cualquiera as a predicate (modifier) ]{\textbf{3.3.1 Formalisation of }\textbf{\textit{cualquiera}}\textbf{ as a predicate}\textbf{ (modifier)}\textbf{ }}
\begin{styleStandard}
The idea in a nutshell is that postnominal \textit{cualquiera} as a predicate belongs to different types of predicates, which is why it triggers different interpretations such as COM and DER observed in 3.2. In what follows, I will formalize the different meanings of the predicate \textit{cualquiera}. 
\end{styleStandard}

\begin{styleStandard}
\ \ Postnominal \textit{cualquiera} with the meaning \textit{común} ‘common’ as ‘usual’ or ‘frequent’ should be analysed as a frequency adjective that modifies nouns with a reference to kinds, as in \ (see Carlson 1977, Chierchia 1998, Gehrke \& McNally 2015). They do not modify proper nouns, as shown in :
\end{styleStandard}

\begin{styleStandard}
\label{bkm:Ref72354905}\ \ Cats are common in Europe.
\end{styleStandard}

\begin{styleStandard}
\label{bkm:Ref72354921}\ \ *Angela Merkel is common in Germany.
\end{styleStandard}

\begin{styleStandard}
The function of a postnominal frequency adjective such as English \textit{common} as in \textit{a} \textit{common example/student} is to turn kinds into subkinds. It takes a kind as an argument and returns a subkind as an output. 
\end{styleStandard}

\begin{styleStandard}
\ \ At this point, I need to give a short introduction into kind nouns (see Gehrke \& McNally 2015 for this introduction). Nouns denote properties of kinds or individuals, as shown in the representation in \ (see Carlson 1977, Chierchia 1998, Gehrke \& McNally 2015). I use a subscript $\alpha $\textit{\textsubscript{k}} to indicate a variable that ranges over kinds (see Gehrke \& McNally 2015). No subscript $\alpha $ indicates a variable that ranges over atomic individuals:
\end{styleStandard}

\begin{styleStandard}
\label{bkm:Ref42889594}\ \ [27E6?][NP [N car]][27E7?]: $\lambda $\textit{x}\textit{\textsubscript{k}}[car(\textit{x}\textit{\textsubscript{k}})] 
\end{styleStandard}

\begin{styleStandard}
These nouns realise a syntactic Num(ber) projection, which semantically is interpreted as properties of token entities (see Gehrke \& McNally 2015, and references therein). R is Carlson’s (1977) realisation relation. The kind variable is bound off by the existential quantifier in the noun representation in . The numberless noun \textit{car }denotes the set of all kinds of cars (station wagons, sedans, etc.), the (phonologically identical) singular NumP \textit{car }denotes the set of atomic individuals that realise some kind of car, and the plural NumP \textit{cars }denotes the set of token pluralities of some kind of car.
\end{styleStandard}

\begin{styleStandard}
\label{bkm:Ref42889808}\ \ [27E6?][NumP [NP car]][27E7?]: $\lambda $\textit{y}${\exists}$\textit{x}\textit{\textsubscript{k}}[car(\textit{x}\textit{\textsubscript{k}}) \& R(\textit{y}, \textit{x}\textit{\textsubscript{k}})]
\end{styleStandard}

\begin{styleStandard}
I assume that postnominal \textit{cualquiera} can denote a property of kinds, whereas \textit{cualquiera} modifying a noun denotes a property of a property of kinds:
\end{styleStandard}

\begin{styleStandard}
\ \ a.\ \ [27E6?]postnominal cualquiera[27E7?]: $\lambda $\textit{P}$\lambda $\textit{x}\textit{\textsubscript{k}}[(cualquiera (\textit{P}))(\textit{x}\textit{\textsubscript{k}})]
\end{styleStandard}

\begin{styleStandard}
\ \ b.\ \ [27E6?]estudiante cualquiera[27E7?]:\textit{ }$\lambda $\textit{x}\textit{\textsubscript{k}}[(cualquiera (student))(\textit{x}\textit{\textsubscript{k}})]
\end{styleStandard}

\begin{styleStandard}
\ \ \ \ ‘kind of a student.’
\end{styleStandard}

\begin{styleStandard}
The experiencer/evaluator function of \textit{cualquiera}, which can be made salient by \textit{para }x\textit{ }‘for x’ and is represented as a subscript E on the predicate modifier \textit{cualquiera}, which predicates over types of individuals according to some evaluator/experiencer E in a given context C:
\end{styleStandard}

\begin{styleStandard}
\ \ [27E6?]estudiante cualquiera[27E7?]\textsuperscript{c} :\textit{ }$\lambda $\textit{x}\textit{\textsubscript{k }}[(cualquiera\textsubscript{E} (student))(\textit{x}\textit{\textsubscript{k}})] 
\end{styleStandard}

\begin{styleStandard}
\ \ ‘kind of student for the experiencer/evaluator E’
\end{styleStandard}

\begin{styleStandard}
The following satisfaction condition in \ expresses that the kind property represented by N \textit{cualquiera} describing a set of elements denoted by the kind is very frequent. Thus, \textit{a common student }denotes a kind of a student such that all students that represent this kind occur very frequently. What it means to be \textit{un estudiante cualquiera} ‘a common student’ is established pragmatically. Thus, \textit{un estudiante cualquiera} denotes a kind of a student that has usual or frequently occurring grades, say 2.7 (= average grade) in the German system. I adopt the condition of use of the English frequency adjective \textit{common} from Gehrke \& McNally (2015) and adapt it for \textit{cualquiera} on the assumption that it has the same use as the Spanish frequency adjectives \textit{común} ‘common’\textit{ }or\textit{ corriente} ‘frequent’. “(Frequency) distribution in the representation is a function that yields the distribution of a set of entities at the given index, with high values” (see Gehrke \& McNally 2015, my addition of the bracket).
\end{styleStandard}

\begin{styleStandard}
\label{bkm:Ref54641377}\ \ \ \ Condition of use in a context for \textit{cualquiera} with the meaning ‘common/frequent’:
\end{styleStandard}

\begin{styleStandard}
\ \ \ \ ${\forall}$\textit{P}, \textit{x}\textit{\textsubscript{k}}, \textit{i}[(cualquiera(\textit{P}))(\textit{x}\textit{\textsubscript{k}}) at \textit{i }$\leftrightarrow $ [\textit{P}(\textit{x}\textit{\textsubscript{k}}) \& distribution(\{\textit{y }: R( \textit{y}, \textit{x}\textit{\textsubscript{k}}) at \textit{i}\}) = \textit{high}]]
\end{styleStandard}

\begin{styleStandard}
\ \ a.\ \ [27E6?]estudiante cualquiera[27E7?]: $\lambda $\textit{x}\textit{\textsubscript{k}} [(cualquiera(\textit{student}))(\textit{x}\textit{\textsubscript{k}})]
\end{styleStandard}

\begin{styleStandard}
\ \ b.\ \ ${\forall}$\textit{x}\textit{\textsubscript{k}}, \textit{i}[(cualquiera(student))(\textit{x}\textit{\textsubscript{k}}) at \textit{i }$\leftrightarrow $ [student (\textit{x}\textit{\textsubscript{k}}) \& distribution(\{\textit{y }: R( \textit{y}, \textit{xk}) at \textit{i}\}) = \textit{high}]]
\end{styleStandard}

\begin{styleStandard}
The type N \textit{cualquiera} implies a contrast to another type, which denotes a singleton and has a low (frequency) distribution. The satisfaction condition for the meaning N non\textit{{}-cualquiera }or N \textit{distinguido/especial} ‘special/particular N’ are the same as for N \textit{cualquiera} besides the low frequency:
\end{styleStandard}

\begin{styleStandard}
\ \ \ \ ${\forall}$\textit{P}, \textit{x}\textit{\textsubscript{k}}, \textit{i}[(distinguido(\textit{P}))(\textit{x}\textit{\textsubscript{k}}) at \textit{i }$\leftrightarrow $
\end{styleStandard}

\begin{styleStandard}
\ \ \ \ [\textit{P}(\textit{x}\textit{\textsubscript{k}}) \& distribution(\{\textit{y }: R( \textit{y}, \textit{x}\textit{\textsubscript{k}}) at \textit{i}\}) = \textit{low}]]
\end{styleStandard}

\begin{styleStandard}
\ \ a.\ \ [27E6?]estudiante distinguido[27E7?]: $\lambda $\textit{x}\textit{\textsubscript{k}} [(distinguished(\textit{student}))(\textit{x}\textit{\textsubscript{k}})]
\end{styleStandard}

\begin{styleStandard}
\ \ b.\ \ ${\forall}$\textit{x}\textit{\textsubscript{k}}, \textit{i}[(distinguished(student))(\textit{x}\textit{\textsubscript{k}}) at \textit{i }$\leftrightarrow $ [student (\textit{x}\textit{\textsubscript{k}}) \& distribution(\{\textit{y }: R( \textit{y}, \textit{xk}) at \textit{i}\}) = \textit{low}]]
\end{styleStandard}

\begin{styleStandard}
Under this approach, the contribution of \textit{cualquiera} is to restrict the kind denoted by the noun N, that is, to turn a kind into a subkind. For instance, the noun \textit{estudiante} denotes a kind of a student and the NP \textit{estudiante cualquiera} denotes a common kind of student as opposed to a non-common one. On the extension, N denotes a set of students that belong to the kind ‘student’ and \textit{estudiante cualquiera} denotes a restricted set of students, namely a set that is described as a ‘common student’, rather than as a ‘distinguished student’. 
\end{styleStandard}

\begin{styleStandard}
\ \ In case N denotes an event noun (indicated by the subscript $\alpha $\textit{\textsubscript{e}}), \textit{cualquiera} denotes a property of a property of an event (see Gehrke \& McNally 2015 for predicate modifiers of event nouns):
\end{styleStandard}

\begin{styleStandard}
\ \ a.\ \ [27E6?]cualquiera[27E7?]: $\lambda $\textit{P}$\lambda $\textit{x}\textit{\textsubscript{e}}[(cualquiera(\textit{P}))(\textit{x}\textit{\textsubscript{e}})]
\end{styleStandard}

\begin{styleStandard}
\ \ b.\ \ [27E6?]evento/historia cualquiera[27E7?]: $\lambda $\textit{x}\textit{\textsubscript{e}}[(cualquiera(event/story))(\textit{x}\textit{\textsubscript{e}})]
\end{styleStandard}

\begin{styleStandard}
\ \ \ \ ‘frequent event/story. A story/event that frequently occurs’
\end{styleStandard}

\begin{styleStandard}
\ \ Satisfaction condition:
\end{styleStandard}

\begin{styleStandard}
\ \ c.\ \ ${\forall}$\textit{x}\textit{\textsubscript{e}}, \textit{i}[(cualquiera(story))(\textit{x}\textit{\textsubscript{e}}) at \textit{i }$\leftrightarrow $ [story(\textit{x}\textit{\textsubscript{e}}) \& distribution(\{\textit{y }: R( \textit{y}, \textit{x}\textit{\textsubscript{e}}) at \textit{i}\}) = \textit{high}]]
\end{styleStandard}

\begin{styleStandard}
To sum up, \textit{cualquiera} in DET [+/– def.] N \textit{cualquiera}, which can be coordinated with adjectives, is not a determiner, but a nominal predicate that describes the distribution of properties among animate entities (non-temporal frequency) or the distribution of properties of events in time (temporal frequency with the meaning \textit{common} and \textit{frequent in time}).
\end{styleStandard}

\begin{styleStandard}
On this analysis, N \textit{cualquiera} in its COM meaning can be used with determiners and quantifiers as is expected from an adjectival use of \textit{cualquiera} (see , , and – \ for data):
\end{styleStandard}

\begin{styleStandard}
\ \ a.\ \ [27E6?]ese/el estudiante cualquiera[27E7?]: $\iota $ \textit{x}[(cualquiera(student))(\textit{x})]
\end{styleStandard}

\begin{styleStandard}
\ \ b.\ \ [27E6?]algún/otro/un estudiante cualquiera[27E7?]: some \textit{x}[(cualquiera(student))(\textit{x})]
\end{styleStandard}

\begin{styleStandard}
\ \ c.\ \ [27E6?]mucha gente cualquiera[27E7?]: many \textit{x}[(cualquiera(people))(\textit{x})]
\end{styleStandard}

\begin{styleStandard}
As \textit{cualquiera} modifies an NP, which denotes a kind of entity or individual, the predicative N\textit{ cualquiera} can refer to some specific type/kind of individual, such as the common type. This is indeed what we find in the data:\textit{ el hombre cualquiera, la mujer cualquiera} ‘the common type of men or women’ as opposed to \textit{el hombre distinguido/particular} ‘the distinguished type of men’ (see 3.6.1).
\end{styleStandard}

\begin{styleStandard}
\ \ In the following example, repeated from , the definite noun refers to a certain kind of men, namely to the common man \textit{as a type}. One could also describe this type as \textit{el hombre de la calle} lit. ‘the man from the street’ or ‘common/normal man’:
\end{styleStandard}

\begin{styleStandard}
\label{bkm:Ref63545986}\ \ El tipo de el C1, que antes fue el de el [sic] C2, es un timo. Es el prototipo de pícaro español. Un personaje que nuestros más célebres escritores de el [sic ] Siglo de Oro describieron con crudo realismo – cómo olvidar el Lazarillo de Lázaro de Tormes –. \textbf{El hombre cualquiera del autobús no siente vergüenza alguna, aunque debería, por tratar de engañar a otras personas que están dispuestas a ayudar a aquellos extraños que lo pasan mal}.
\end{styleStandard}

\begin{styleStandard}
\ \ ‘[…] The common man from the bus does not feel any shame, although he should,…’
\end{styleStandard}

\begin{styleStandard}
\ \ (@ 82)
\end{styleStandard}

\begin{styleStandard}
The same analysis applies to \textit{el día cualquiera}, namely the most prototypical day of the week:
\end{styleStandard}

\begin{styleStandard}
\ \ diría que es el día más indeterminado, el día cualquiera de la semana, un día metido en la selva de los días
\end{styleStandard}

\begin{styleStandard}
\ \ ‘I would say that it is the most undistinguished day, just some day of the week, a day put in the jungle of the days.’
\end{styleStandard}

\begin{styleStandard}
\ \ (@ 100)
\end{styleStandard}

\begin{styleStandard}
The definite article can thus refer to a specific type of individual, namely a common one, and not some other type, a non-common one. 
\end{styleStandard}

\begin{styleStandard}
\ \ We also find demonstratives, which refer to a specific type of a man of speaker’s reference that has prototypical properties of a man. The example is embedded in a philosophical context, which contains a statement about a certain type of a man:
\end{styleStandard}

\begin{styleStandard}
\label{bkm:Ref63545989}\ \ Las armas de este hombre cualquiera, de aspecto normal, se resumen en: prudencia
\end{styleStandard}

\begin{styleStandard}
\ \ ‘The weapons of this common, normal-looking man are summarised in: caution’
\end{styleStandard}

\begin{styleStandard}
\ \ (@ 83)
\end{styleStandard}

\begin{styleStandard}
Another consequence of the kind analysis of postnominal \textit{cualquiera} is that it explains the lack of occurence with degree modifiers or adverbs like \textit{muy} ‘very’ as in \textit{un hombre *(muy) cualquiera} ‘a common man’. Degree adverbs like \textit{muy} ‘very’ combine with gradable adjectives, not with kind predicates (Kennedy 1999). 
\end{styleStandard}

\begin{styleStandard}
\ \ To recap, postnominal \textit{cualquiera} has been analysed as a frequency adjective ‘common’ in line with Gehrke \& McNally (2015). Its lexical meaning is to describe kinds of individuals or entities that have a high frequency, that is, considering all invidividuals or entities denoted by N, we would observe very often a set of individuals or entities that are described by N\textit{ cualquiera}. 
\end{styleStandard}

\subsubsection[3.3.2 Formalisation of the deragotory meaning of N cualquiera]{\textbf{3.3.2 Formalisation of the deragotory meaning of N }\textbf{\textit{cualquiera}}}
\begin{styleStandard}
The idea is that the deragotory meaning of postnominal \textit{cualquiera} has an additional satisfaction condition on top of describing individuals or types of individuals with a high-frequency distribution. Namely, the second satisfaction condition in \ expresses that N\textit{ cualquiera} in its DER reading, such as \textit{equipo cualquiera} ‘ordinary team’ denotes a type of individuals that is ranked lower on some qualitative scale according to the experiencer E expressed in \textit{para }x ‘for x’ than the type of individuals denoted by \textit{equipo} non-\textit{cualquiera} or a distinguished team (e.g. \textit{equipo con Messi} ‘team with Messi’):
\end{styleStandard}

\begin{styleStandard}
\label{bkm:Ref58094203}\ \ [27E6?]equipo cualquiera\textsubscript{deragotory }[27E7?]\textsuperscript{c}: $\lambda $\textit{x}\textit{\textsubscript{k}}[(cualquiera\textsubscript{E} (\textit{team}))(\textit{x}\textit{\textsubscript{k}})]
\end{styleStandard}

\begin{styleStandard}
\ \ a.\ \ ${\forall}$\textit{x}\textit{\textsubscript{k}}, \textit{i}[(cualquiera(team))(\textit{x}\textit{\textsubscript{k}}) at \textit{i }$\leftrightarrow $ [team(\textit{x}\textit{\textsubscript{k}}) \& distribution(\{\textit{y }: R( \textit{y}, \textit{xk}) at \textit{i}\}) = \textit{high}]]
\end{styleStandard}

\begin{styleStandard}
\ \ b.\ \ ${\forall}$y\textsubscript{k, }${\forall}$x\textsubscript{k, }\textit{i}[(cualquiera[2032?] (y\textsubscript{k}) \& distinguished[2032?](y\textsubscript{k}) at \textit{i} $\leftrightarrow $ [team(\textit{x}\textit{\textsubscript{k}})(\textit{y}\textit{\textsubscript{k}}) \& (\textit{x}\textit{\textsubscript{k}}) {\textless}\textsubscript{lower on some pragmatic scale according to E }(\textit{y}\textit{\textsubscript{k}}) at \textit{i}]]
\end{styleStandard}

\begin{styleStandard}
On this analysis, \textsc{un} N \textit{cualquiera} cannot have an evaluative meaning in contexts that do not include any evaluation on a pragmatic scale according to an experiencer, such as mathematical contexts. This prediction is borne out empirically:
\end{styleStandard}

\begin{styleStandard}
\ \ Es un cuerpo geométrico cualquiera. \# Sin nivel.
\end{styleStandard}

\begin{styleStandard}
\ \ \# ‘It’s an unremarkable geometric body. Without class.’
\end{styleStandard}

\begin{styleStandard}
\ \ (@ 101)
\end{styleStandard}

\begin{styleStandard}
I will show in Chapter 4 that in some languages or language varieties DER is only felicitous under the condition b; that is, the requirement that the set of individuals or entities must be described as usual/frequent withouth the condition b is not given. French \textit{n’importe quoi} (see a) and Argentinian Spanish \textit{cualquiera} (see b) are such a case. They have necessarily a deragotory meaning. According to the deragotory meaning of \textit{n’importe quoi}/\textit{cualquiera, }the evaluator considers that saying something undistinguished is worse than saying something distinguished (see Chapter 4 for detailed analysis):
\end{styleStandard}

\begin{styleStandard}
\label{bkm:Ref72355051}\ \ a.\ \ Marie a dit (du grand) n’importe quoi.\ \ (French)
\end{styleStandard}

\begin{styleStandard}
\ \ \ \ ‘Marie said a big bullshit.’\ \ 
\end{styleStandard}

\begin{styleStandard}
\ \ \ \ \ \ (@ 30)
\end{styleStandard}

\begin{styleStandard}
\ \ b.\ \ Me mandaste cualquiera.\ \ (ArgSp)
\end{styleStandard}

\begin{styleStandard}
\ \ \ \ ‘You told me bullshit’\ \ 
\end{styleStandard}

\begin{styleStandard}
\ \ \ \ (@ 102)
\end{styleStandard}

\subsubsection[3.3.3 The similarity interpretation of postnominal cualquiera]{\textbf{3.3.3 The similarity interpretation of postnominal }\textbf{\textit{cualquiera}}}
\begin{styleStandard}
The goal of this section is to adapt the predicate analysis of postnominal \textit{cualquiera} with a similarity interpretation or SIM as in \textit{es un día cualquiera} ‘it’s a day like any other day’ and to show the connection between SIM and the free choice pronoun \textit{cualquiera} ‘anyone’. \ \ 
\end{styleStandard}

\begin{styleStandard}
\ \ Consider the following two examples:
\end{styleStandard}

\begin{styleStandard}
\label{bkm:Ref72353992}\ \ Juan es un estudiante cualquiera de mi clase.
\end{styleStandard}

\begin{styleStandard}
\ \ ‘Juan is a student \textsc{cualquiera} in my class.’
\end{styleStandard}

\begin{styleStandard}
\ \ (similar to @ 74)
\end{styleStandard}

\begin{styleStandard}
\label{bkm:Ref72354006}\ \ Juan es un estudiante como cualquiera de mi clase.
\end{styleStandard}

\begin{styleStandard}
\ \ ‘Juan is a student like everyone in my class.’
\end{styleStandard}

\begin{styleStandard}
\ \ (similar to @ 75)
\end{styleStandard}

\begin{styleStandard}
\ \ One possible analysis of \ is to interpret it semantically as . Comparative constructions are known to license FCIs (see Aloni 2021 and 2.2). In this view, \textit{un hombre cualquiera }would be interpreted as ‘a man like anyone [=any man]’. The function of postnominal \textit{cualquiera} in \ is thus to express a predicate that has a similar semantic interpretation as the pronoun \textit{cualquiera} in an overt comparative construction as in \ (for comparative clause analyses of cases like in , see Bresnan 1973, Nespor \& Napoli 1986, Sáez 1992, Kennedy 1999, among many others)
\end{styleStandard}

\begin{styleStandard}
\ \ Different analyses have been suggested in the literature on comparative clauses (see Bresnan 1973, Nespor \& Napoli 1986, Sáez 1992, Kennedy 1999, among many others). I will follow the coordination hypothesis of comparatives suggested by Nespor \& Napoli (1986) and Sáez (1992), among others, who assume that comparatives represent coordinations (on either the phrasal or sentence level); that is, \textit{como cualquiera de mi clase} is represented as a coordinated phrase or clause (\&P). I assume that the comparative phrase in \ is clausal because it can be paraphrased as a clause. In order to represent the clausal analysis, I mark the predicate as deleted as shown in the English translation (see Bresnan 1973):
\end{styleStandard}

\begin{styleStandard}
\label{bkm:Ref72354047}\ \ Juan es un estudiante como lo es cualquiera de mi clase.
\end{styleStandard}

\begin{styleStandard}
\ \ Juan is a student like it is any of my class
\end{styleStandard}

\begin{styleStandard}
\ \ ‘Juan is a student like everyone in my class is a student.’
\end{styleStandard}

\begin{styleStandard}
\ \ (similar to @ 76)
\end{styleStandard}

\begin{styleStandard}
Let me suggest a simple syntactic analysis of . I analyse \textit{cualquiera de mi clase} as a subject of a copulative clause, which is inside the comparison clause introduced by \textit{como}. Assuming that \textit{como} conjoins two sentences (see Nespor \& Napoli 1986, Sáez 1992), \ can be interpreted as a coordination of two conjoined copulative sentences on the semantic level:
\end{styleStandard}

\begin{listWWNumxleveli}
\item 
\begin{styleStandard}
‘[Juan is a student] and [[\textsubscript{subj.} everyone from my class] is a student].’
\end{styleStandard}
\end{listWWNumxleveli}
\begin{styleStandard}
\ \ The syntax of \textit{cualquiera} that expresses a comparison is that of a \textit{comparison clause}. This comparison clause is conjoined with the main clause (TP) due to the presence of the comparative particle \textit{como}. In , the first clause has \textit{Juan} as a subject and a predicative DP \textit{un estudiante.} The comparative clause contains \textit{cualquiera} as a subject and a clitic pronoun \textit{lo} which is coreferent with the predicative complement inside the matrix clause:
\end{styleStandard}

\begin{styleStandard}
\label{bkm:Ref72354100}\ \ [\textsubscript{TP} Juan es [\textsubscript{DP }un estudiante]\textsubscript{j}] [\textsubscript{\&P} como lo\textsubscript{j} es cualquiera de mi clase]
\end{styleStandard}

\begin{styleStandard}
\ \ ‘Juan is a student, and everyone (else) is a student in my class.’
\end{styleStandard}

\begin{styleStandard}
\ expresses similarity with respect to some (prototypical) property P* of a student, such as with respect to Juan’s qualifications or grades, which can be made salient in the context such as in . In a less strict interpretation, the sentence in \ with the assumption that the comparative clause represents a conjoined sentence as suggested here could simply mean that Juan and everyone in the speaker’s class are students or have the property of being a student (as opposed to teachers, for example). But this interpretation is too weak and makes false predictions for the use of sentences such as \ in the following context. Imagine a scenario where the speaker wants to say that Juan is a certain type of a student, namely an average student. In this scenario, the meaning ‘Juan is a student and everyone else in my class is a student’ is underinformative. The solution to this problem is pragmatics. What it means to be a student and to be similar to other students is established pragmatically by some property P*; that is, the comparison of students is a comparison with respect to some (prototypical) property P* that holds for students. Thus, the sentence in \ states that Juan and other students in the class of speaker’s reference have the same prototypical property P* as students, which can be usual grades, for instance. From this analysis, we can already pragmatically infer that Juan is an average or unremarkable student. If Juan were a brilliant student, he probably would not be like every other student in the speaker’s class, because it is very unlikely that the whole class consists of brilliant students. 
\end{styleStandard}

\begin{styleStandard}
\ \ I will now formalise the semantic analysis of postnominal \textit{cualquiera} with the similarity reading as the pronoun \textit{cualquiera} in a comparative clause analysis.
\end{styleStandard}

\begin{styleStandard}
\ \ Given the syntax in , the semantic representation entails a universal quantifier \textit{cualquiera} inside the comparative clause \&P. In this case, \textit{cualquiera} is interpreted as a universal quantifier with the meaning ‘everyone’, ranging over individuals in the speaker’s class, and these individuals represent a comparison class of students to whom the subject of the main clause (here: Juan) is compared:
\end{styleStandard}

\begin{styleStandard}
\label{bkm:Ref54682554}\ \ \textit{cualquiera }in a comparison clause
\end{styleStandard}

\begin{styleStandard}
\ \ a.\ \ Sentence: Juan es un estudiante [\textsubscript{ModifP} cualquiera]. Tiene calificaciones normales. ‘Juan is a student like anyone else. He has normal grades.’
\end{styleStandard}

\begin{styleStandard}
\ \ b.\ \ Semantic interpretation: Juan is a student with a property P* and [${\forall}$]([anyone in my class, $\lambda $x. has it (= the property P*)])
\end{styleStandard}

\begin{styleStandard}
\ \ c.\ \ Prose of : Juan is a student with property P* and for each student x in my class, x has \textit{it} [=P*].\footnote{\textrm{\ Another way to interpret the conjoined clause with }\textrm{\textit{cualquiera}}\textrm{ as a subject is to assume that it implies a covert modal or attitude verb that links the expression of similarity to the speaker’s evaluation. In this case, we would have universal quantification over a modal verb that describes speaker’s attitude (see Section 2.2.2). }\textrm{\textit{Juan es un hombre cualquiera}}\textrm{ would be interpreted similar to }\textrm{\textit{Juan es un hombre cualquiera para mí}}\textrm{. ‘Juan is a man cualquiera = According to me, Juan is a man cualquiera.’}\par 
\setcounter{listWWNumxxiiileveli}{0}
\begin{listWWNumxxiiileveli}
\item 
\begin{styleFootnote}
\textrm{It seems to me that Juan is a man with P* and for each student, he/she has P* too.}
\end{styleFootnote}
\end{listWWNumxxiiileveli}
}
\end{styleStandard}

\begin{styleStandard}
The positive outcome of a comparison clause analysis is that it explains the licensing of postnominal \textit{cualquiera} in affirmative predicative contexts. Recall from the puzzle described in Section 2.3 that affirmative predicative and episodic contexts were considered problematic for FCIs such as \textit{cualquiera}. However, in my analysis, the licensing of postnominal \textit{cualquiera} under indicative predicative verbs is explained by its predicate that refers to similarity with all other students from my class. Recall that FCIs are generally licensed in universal contexts in Spanish and other languages (see 2.1 and 2.2).
\end{styleStandard}

\begin{styleStandard}
\ \ Moreover, the suggested analysis explains why the similarity reading cannot be cancelled (see 3.1.3.3). The similarity reading is the result of analysing postnominal \textit{cualquiera} as a modifier that contains a universal quantifier that licenses \textit{cualquiera} in predicative contexts. 
\end{styleStandard}

\begin{styleStandard}
\ \ The suggested analysis also explains why N \textit{cualquiera} is not gradable in European Spanish, as shown in . Postnominal \textit{cualquiera} is a predicate with quantificational force of the same type as ‘every(one), any(one)’ and not a gradable adjective like ‘ordinary’, which is why it cannot be modified by degree adverbs like \textit{muy} ‘very’ (see Chapter 2):
\end{styleStandard}

\begin{styleStandard}
\label{bkm:Ref58074686}\ \ Es un estudiante (*muy) cualquiera. \ \ 
\end{styleStandard}

\begin{styleStandard}
\ \ ‘It’s a student (*very) like anyone else.’
\end{styleStandard}

\begin{styleStandard}
\ \ (@ 77)
\end{styleStandard}

\begin{styleStandard}
Another prediction of the comparative clause analysis is that it should be possible to modify postnominal \textit{cualquiera }by adverbs like \textit{casi} ‘almost’ that restrict the domain of quantification of \textit{cualquiera de mi clase} to ‘almost everyone in my class’ (see 2.1 for such adverbial modification of \textit{cualquiera}). However, the judgements from native speakers of European Spanish and the absence of \textit{un,a N casi cualquiera} in European Spanish books on google books (despite some occurrences in Latin American Spanish), seem to suggest that the postnominal \textit{cualquiera} is not modifiable by the modifier \textit{casi} in European Spanish:
\end{styleStandard}

\begin{styleStandard}
\ \ Fue un día *(casi) cualquiera.\footnote{\textrm{\ There are other cases of adverbial modification of the free choice item }\textrm{\textit{cualquiera}}\textrm{, which shows that }\textrm{\textit{cualquiera}}\textrm{ is a quantifier with the meaning ‘any’. However, }\textrm{\textit{cualquiera}}\textrm{ appears in these examples prenominally as a determiner and not as a postnominal modifier (see Section 4 for further examples of the determiner }\textrm{\textit{cualquiera}}\textrm{ in predicative position):}\par 
\setcounter{listWWNumxiileveli}{0}
\begin{listWWNumxiileveli}
\item 
\begin{styleFootnote}
\textrm{No tiene nada de particular; pero por eso mismo parece Tamayo, porque no tiene nada de particular. }\textrm{\textit{Parece cualquier cosa}}\textrm{ }\textrm{\textit{menos un gran poeta}}\textrm{. (Solos de Clarín, Leopoldo Alas)}
\end{styleFootnote}
\end{listWWNumxiileveli}
\textrm{‘He does not have any particular properties. But for the same reason he resembles Tamayo, because he does not have any particular property. He resembles }\textrm{\textit{anything but a great poet}}\textrm{.’}} \ \ 
\end{styleStandard}

\begin{styleStandard}
\ \ \#‘It was a day like almost any other day.’
\end{styleStandard}

\begin{styleStandard}
\ \ (* according to 2 native speakers of European Spanish, ok according to a Colombian Instagram user and a Mexican Spanish writer Edmée Pardo\footnote{\ \url{https://www.instagram.com/piaguzmanv/}\textstyleInternetlink{\textrm{ and}} \url{https://edmeepardo.com/bibliomemoria-fichas/rfn6txw4j6pg3l2-6kdbj-2e74j-d9r63-zmw5n-az8dg-4e5j4-bamwc-d3fda-c6ygp-efe6w-nr8wf-tnmwk-mwjcd-4wp56}\textrm{ (access, 10 of November 2023)}} and also ok for Michael Robinson, see Prólogo in Pérez Gellida, \textit{Memento mori})
\end{styleStandard}

\begin{styleStandard}
The lack of frequent modification by \textit{casi} in (176) suggests that the postnominal modifier \textit{cualquiera} is not syntactically a freee choice pronoun like the pronoun \textit{cualquiera}, which is modifiable by \textit{casi} very frequently. Consequently, the postnominal \textit{cualquiera} needs to be analysed as a predicate modifier that is semantically related to the pronoun \textit{cualquiera} in the description of the property as ‘like anyone else’.
\end{styleStandard}

\begin{styleStandard}
\ \ I would like to discuss the cancellability of the similarity meaning as shown in the following example. The following example shows that once the postnominal \textit{cualquiera} cooccurs with a an overt comparative clause like \textit{como cualquier mañana} ‘like any other morning’, the interpretation of postnominal \textit{cualquiera} changes into the meaning ‘ordinary’, which corresponds to COM or DER. This suggests that the similarity meaning is not part of the lexical semantic meaning as suggested in the discussion of a), because otherwise the following sentence would express the same information twice and it would be redundant. \ \ 
\end{styleStandard}

\begin{styleStandard}
\ \ Era una mañana cualquiera, como cualquier mañana. 
\end{styleStandard}

\begin{styleStandard}
‘It was an ordinary morning, like any other morning.’ \footnote{\textrm{\ I thank J.Quer for this data point. }\par }\ \ \ 
\end{styleStandard}

\begin{styleStandard}
The issue of what kind status the similarity meaning has (lexical semantic or pragmatic implicature) (Aloni 2021, A\&M 2018) needs to be analysed in the future research.
\end{styleStandard}

\subsection[3.4 Summary and discussion of the predicate modifier analysis]{\textbf{3.4 Summary and discussion of the predicate modifier analysis}}
\begin{styleStandard}
In Section 3.3, I suggested a predicate modifier analysis of postnominal \textit{cualquiera} with the meaning SIM, COM that can explain a certain type of data such as [Determiner/Quantifier N \textit{cualquiera}] and coordination with adjectives. There are several issues that need to be addressed here.
\end{styleStandard}

\begin{styleStandard}
\ \ First, recall from Section Section 2.1.3 (see, in particular, Lapesa 1974, Rivero 2011) that there is the standard assumption in the literature that postnominal \textit{cualquiera} occurs only with indefinite nouns and is hence accompanied by the indefinite article, but not by strong determiners. Linguists have suggested different explanations for this restriction. According to one explanation, postnominal \textit{cualquiera} is itself a strong (quantificational) determiner, and this is why it is incompatible with other strong determiners (Lapesa 1974). Another semantic explanation is that N \textit{cualquiera} has an anti-singleton restriction (see Menéndez-Benito 2010), making it incompatible with definite determiners or demonstratives, as these elements presuppose a uniquely identified referent (singleton restriction) and thus do not co-occur with \textit{cualquiera} as an FCI. However, contrary to these preddictions, I was able to find examples in which postnominal \textit{cualquiera} is preceded by definite determiners, and which my predicate modifier analysis is able to explain.
\end{styleStandard}

\begin{styleStandard}
\ \ Second, there are indeed a lot more data with \textsc{un} N \textit{cualquiera} than with D\textsubscript{def} N \textit{cualquiera} in the corpus (6,938 indefinite occ. vs 171 definite occ.). A simple explanation for this preference for the indefiniteness of N \textit{cualquiera} under copular verbs is that N \textit{cualquiera} is used predicatively. It is a common observation in the literature that predicative NPs are expressed by NPs that are indefinite rather than definite (see Von Geenhoven \& McNally 2005, among others): 
\end{styleStandard}

\begin{styleStandard}
\ \ Era un/*el día normal (para mí).
\end{styleStandard}

\begin{styleStandard}
\ \ ‘It was a/*the usual day (for me).’
\end{styleStandard}

\begin{styleStandard}
\ \ (@ 103)
\end{styleStandard}

\begin{styleStandard}
Third, my analysis predicts that it should be possible to find cross-linguistic evidence for quantifiers reanalysed as nominal predicates. We actually find few cases in Italian and German, where certain quantifiers can be used as nominal predicates in predicative position similar to \textit{un/una cualquiera} in Spanish (see Section 2.1). The only difference to \textsc{un} N \textit{cualquiera} is that the quantifiers in [indefinite + quantifier], such as in \textit{un cualquiera} in Spanish, \textit{un} \textit{niente} in Italian, and \textit{ein Jedermann} in German, have the syntactic status of a noun rather than of a predicate modifier. However, the semantic analysis is the same; that is, the quantifier is interpreted not as a quantifier but as a noun with a property interpretation that describes certain kinds of people. \textit{Ein Jedermann} is a kind of man, as is \textit{un niente} and \textit{un cualquiera.}
\end{styleStandard}

\begin{styleStandard}
\ \ Gianni è un niente. \ \ (Italian)
\end{styleStandard}

\begin{styleStandard}
\ \ ‘Gianni is a nobody.’
\end{styleStandard}

\begin{styleStandard}
\ \ (@ 104)
\end{styleStandard}

\begin{styleStandard}
\ \ Juan ist ein Jedermann. \ \ (German)
\end{styleStandard}

\begin{styleStandard}
\ \ ‘Juan is an everyman.’
\end{styleStandard}

\begin{styleStandard}
Fourth, my analysis allows us to better approach some aspects of the diatopic variation of \textit{cualquiera} in Spanish. The element \textit{cualquiera} in Argentinian Spanish can appear as a (degree) predicate over nouns that refer to atomic individuals such as \textit{Juan} or that refer, for instance, to a certain movie (see 2.4):
\end{styleStandard}

\begin{styleStandard}
\ \ La peli/Juan es (muy) cualquiera. \ \ (ArgSp/*EurSp)
\end{styleStandard}

\begin{styleStandard}
\ \ ‘This movie/Juan is (very) ordinary/bad.’
\end{styleStandard}

\begin{styleStandard}
\ \ (@ 5 / @ 105)
\end{styleStandard}

\begin{styleStandard}
The difference between Argentinian Spanish and its European counterpart is that \textit{cualquiera }is used as a predicate modifier in European Spanish, but not as a predicate (i.e. a predicative adjective). I will explain this difference by the assumption that the analysis of \textit{cualquiera} as an adjective in Argentinian Spanish is at a more advanced stage than in European Spanish (for details, see Chapter 4). Syntactically, it can form a regular copula structure, which can be sketched as follows:
\end{styleStandard}

\begin{styleStandard}
\ \ [\textsubscript{TP} Juan\textsubscript{j} T° es [\textsubscript{PrP} t\textsubscript{j} [\textsubscript{PrP'} Pr° cualquiera]]].
\end{styleStandard}

\begin{styleStandard}
\ \ ‘Juan is ordinary/unremarkable.’
\end{styleStandard}

\begin{styleStandard}
I assume that, in European Spanish \textit{cualquiera} is not used as a predicate that describes properties of individuals. Let me btriefly outline the syntax of \textit{cualquiera} in its COM reading in European Spanish. In European Spanish, postnominal \textit{cualquiera} is a nominal modifier similar to adjectives like ‘common/unremarkable’. I use the labels Modif(ier) for all kinds of modifiers (adjectives, relative clauses, etc.). This type of postnominal \textit{cualquiera} modifies nouns that denote kinds and, as such, N\textit{ cualquiera }can be represented as a KindP (see Zamparelli 2000 for KindPs):
\end{styleStandard}

\begin{styleStandard}
\ \ [\textsubscript{DP} un [\textsubscript{KindP} \textsubscript{Kind°} hombre [\textsubscript{ModifP}\textit{ cualquiera/cumún/ordinario }]]]
\end{styleStandard}

\begin{styleStandard}
\ \ ‘a common/ordinary man’
\end{styleStandard}

\begin{styleStandard}
The whole DP shown in (216) can appear as a predicative complement of the copula verb \textit{be}, but it cannot appear as a predicative complement by itself as in Argentinian Spanish:
\end{styleStandard}

\begin{styleStandard}
\ \ [\textsubscript{TP }Juan\textsubscript{j} [\textsubscript{T'} [T° es [\textsubscript{PrP} t\textsubscript{j} [\textsubscript{PrP' }Pr°\textsubscript{ }[\textsubscript{DP} un [\textsubscript{KindP} hombre [\textsubscript{ModifP}\textit{ cualquiera/cumún/ordinario}]]]]]]]]
\end{styleStandard}

\begin{styleStandard}
\ \ ‘Juan is an ordinary man.’
\end{styleStandard}

\begin{styleStandard}
It can also appear as a nominal modifier of an NP that is the subject of a clause:
\end{styleStandard}

\begin{styleStandard}
\ \ [\textsubscript{DP} \textit{El} [\textsubscript{KindP }\textit{hombre} [\textsubscript{ModifP }\textit{cualquiera}]]] del autobús no siente vergüenza alguna.
\end{styleStandard}

\begin{styleStandard}
\ \ ‘The common man does not feel any shame.’
\end{styleStandard}

\begin{styleStandard}
As we will see in Chapter 4, the Argentinian Spanish \textit{cualquiera} is not restricted to the use of a nominal modifier as in European Spanish. 
\end{styleStandard}

\subsection[3.5 Chapter summary and outlook to further chapters]{\textbf{3.5 Chapter summary and outlook to further chapters}}
\begin{styleStandard}
I started with different observations about \textsc{un} N \textit{cualquiera} in the corpus (see 3.1–3.4) and showed that in affirmative predicative contexts, it has usually an evaluative meaning that manifests itself in three readings, namely ‘similar to anyone else’ (SIM), ‘common (COM)’, and ‘of low value’ (DER) (3.2). 
\end{styleStandard}

\begin{styleStandard}
\ \ I showed in 3.3 that SIM is a special case of interpreting \textit{cualquiera} as a predicate with quantificiational force similar to the quantifier pronoun cualquiera that occurs in a comparison clause and that the COM and DER readings can be derived pragmatically from the SIM reading. 
\end{styleStandard}

\begin{styleStandard}
\ \ In the following chapters, I will examine the diatopic and diachronic variation of \textit{cualquiera} in more detail. Chapter 4 will deal with diatopic variation. Chapters 5 will deal with diachronic variation. 
\end{styleStandard}

\begin{styleStandard}
\ \ The question that will guide my analysis in subsequent sections is why a generalised quantifier with the free choice meaning ‘any’ changes into a predicate modifier as demonstrated in 3 for European Spanish or a gradable predicate as demonstrated for Argentinian Spanish. My short answer to this question is that it is the predicative context, the negation, and the adjectival appositives that trigger the reanalysis of postnominal \textit{cualquiera} as an FCI into a predicate modifier or predicate of an adjective type. Adjectival appositives can elicit different readings of FCIs in predicative position such as ‘similar like anyone else’, ‘common’, and ‘of low value’, and by doing this, they can turn generalised quantifiers of the type {\textless}{\textless}e,t{\textgreater}t{\textgreater} into predicate modifiers of type {\textless}{\textless}et{\textgreater},{\textless}e,t{\textgreater}{\textgreater} and can thus turn the interpretation of ‘any’ into the adjectival interpretation of ‘any’ or ‘common/frequent/normal’. This is how postnominal \textit{cualquiera} gets its adjectival or kind flavour. The adjectival interpretation can be contextually elicited by appositives that specify the property as ‘average/frequent’. The deragotory flavour of ‘nothing special’ is just a negative interpretation of ‘any’ or of ‘average’.
\end{styleStandard}

\begin{styleStandard}
\ \ It is a common observation in the literature that semantic change is often triggered by pragmatic enrichment (see 2.7 and Traugott \& Dasher 2002, Eckardt 2006, among others). In the case of \textit{cualquiera}, the most crucial trigger for semantic change is syntax, in particular the predicative context and negation, which triggers a contrastive reading between the property described by N\textit{ cualquiera} and another property described by N non-\textit{cualquiera}. The function of adjectival appositives that describe the property of postnominal \textit{cualquiera }as ‘common, frequent, low value’\textit{ }elicit the adjectival interpretation. The diachronic path of \textit{cualquiera} can be sketched as follows:
\end{styleStandard}

\begin{styleStandard}
\ \ Diachronic development of \textit{cualquiera}
\end{styleStandard}

\begin{styleStandard}
\ \ 1.\ \ Use of \textit{cualquiera} as FCI in modal constructions.
\end{styleStandard}

\begin{styleStandard}
\ \ 2.\ \ Use of \textit{cualquiera} in predicative position (under negation) with different readings such as similarity or deragotory reading.
\end{styleStandard}

\begin{styleStandard}
\ \ 3.\ \ Reanalysis of FCI \textit{cualquiera} as a property of kind nouns in predicative contexts.
\end{styleStandard}

\begin{styleStandard}
\ \ 4.\ \ Reanalysis of \textit{cualquiera} as a gradable property\textit{ }of proper nouns.
\end{styleStandard}

\begin{styleStandard}
In Section 3.6, I argued that a small set of data of postnominal \textit{cualquiera} in European Spanish shows that postnominal \textit{cualquiera} has reached step 3. In Chapter 4, I will show that Argentinian Spanisch has reached step 4. In Chapters 5 and 6, I will look at the diachronic conditions that have enabled step 3 in European Spanish, namely the grammaticalisation of \textit{cualquiera} from a relative clause into a compound which at some point was possible not only under modal verbs, but also under affirmative predicative verbs.
\end{styleStandard}

\begin{styleStandard}
\ \ In the following chapters, I will study the diatopic and diachronic variation in more detail. According to my analysis, \textit{cualquiera} as FCI should be diachronically prior to readings in affirmative predicative sentences. The lexicalisation of non-modal readings in some varieties of Spanish or other Romance or Non-Romance languages is expected to occur later than encountering FCIs in predicative position. Moreover, I will also look into the development of EVAL ‘nonsense’ in Argentinian Spanish \textit{cualquiera} in Chapter 4.
\end{styleStandard}

\clearpage\section[4 Postnominal and pronominal cualquiera in Argentinian Spanish]{\textbf{4 Postnominal and pronominal }\textbf{\textit{cualquiera }}\textbf{in Argentinian Spanish}}
\begin{styleStandard}
This chapter analyses \textit{cualquiera} in Argentinian Spanish and accounts for the differences from its counterpart in European Spanish (EurSp).\footnote{\textrm{\ Aspects of this topic were presented in my talk “}\textrm{\textit{Cualquiera}}\textrm{ in Argentinian Spanish,” II Encuentro de Dialectos del Español, (Ciudad Real, Universidad de Castilla-La Mancha, 17–18 May 2018). I’m very thankful to the audience of this workshop, especially to Carlos Muñoz.}} I will suggest that \textit{cualquiera} embedded under certain transitive verbs such as ArgSp \textit{mandar/decir cualquiera} ‘say nonsense’ has its origin in \textit{mandar/decir cualquier cosa} ‘say anything possible’ or ‘say something random’. \textit{Cualquiera} under predicative verbs such as \textit{parecer} will be analysed as a degree adjective that has its origin in the FCI \textit{cualquiera} with the interpretation ‘any’.
\end{styleStandard}

\begin{styleStandard}
\ \ The basic data in this section is extracted from the \textit{Corpus del Español web dialects }(CdE\_web dialects, already introduced in Section 1.5 and used in Chapter 3).\footnote{\ \textrm{I also asked about the use of }\textrm{\textit{cualquiera}}\textrm{ in predicative position and under transitive verbs with the possibility of degree modification (e.g. }\textrm{\textit{es muy/re cualquiera}}\textrm{) in the of }\href{https://www.facebook.com/groups/686284488187168/permalink/992126550936292/}{\textstyleInternetlink{\textrm{\textit{Lingüística Argentina}}}}\textrm{ online forum on Facebook}.\textrm{ 18 native speakers responded to my questions. Most speakers accepted }\textrm{\textit{cualquiera}}\textrm{ in predicative position and under transitive verbs. However, not all degree adverbs can modify }\textrm{\textit{cualquiera}}\textrm{, as noted by native speakers. Most speakers accept }\textrm{\textit{re cualquiera}}\textrm{, in contrast to }\textrm{\textit{muy cualquiera}}\textrm{.}}
\end{styleStandard}

\begin{styleStandard}
\ \ This section has the following structure. I will briefly recall the current state of research \textit{cualquiera} in ArgSp from Chapter 2 in Section 4.1. In Section 4.2 I will show that the description in the literature is incomplete and very informal and will object to the suggestion that the source for the special use of \textit{cualquiera} in ArgSp should be derived from an analogy with the Italianism \textit{cualunque} ‘whichever’. I will present new data on \textit{cualquiera} in 4.3 that will serve my synchronic analysis in Section 4.4. And in 4.5, I will suggest that \textit{cualquiera} in ArgSp has lexicalised the speaker’s evaluation of the modal component denoted by \textit{cualquiera} as an FCI, which has triggered a reanalysis into an evaluative adjective or an evaluative noun. I will show further evidence for the analysis in 4.6. 
\end{styleStandard}

\subsection[4.1 Current state of research on cualquiera in Argentinian Spanish]{\textbf{4.1 Current state of research on}\textbf{\textit{ cualquiera}}\textbf{ in Argentinian Spanish}}
\begin{styleStandard}
The item \textit{cualquiera} can have a pejorative or derogatory use (DER, a subtype of the evaluative use or function, EVAL) in Argentinian Spanish with copular predicates such as \textit{parecer} ‘to look like’ and \textit{ser} ‘to be’, similar to EurSp, but, in addition, can be used predicatively and can have degree modification by degree adverbs (Rizzo Salierno 2013: ex. 2, 3.1.1, Chapter 2). The following examples show degree modification of \textit{cualquiera} by the degree adverb \textit{completamente} ‘totally/comletely’, and they are interpreted as ‘(very) bad’ or as ‘(complete) nonsense’. As I will argue in this chapter, the meaning ‘nonsense’ is a pejorative interpretation of ‘anything’ and thus another instance of DER:
\end{styleStandard}

\begin{styleStandard}
\textbf{\ \ }\ \ Es\ \ totalmente\ \ cualquiera\ \ eso\ \ de\ \ haber\ \ puesto\ \ a\ \ Roger\ \ con\ \ la\ \ voz\ \ de\ \ Syd.
\end{styleStandard}

\begin{styleStandard}
\ \ \ \ is\ \ totally\ \ \textsc{cualquiera}\ \ that\ \ of\ \ having\ \ put\ \ to\ \ Roger\ \ with\ \ the\ \ voice\ \ of\ \ Syd
\end{styleStandard}

\begin{styleStandard}
\ \ \ \ ‘Putting Roger with Syd´s voice is total(ly) bullshit/nonsense/out of place/bad.’
\end{styleStandard}

\begin{styleStandard}
\ \ \ \ (Rizzo Salierno 2013: 28)
\end{styleStandard}

\begin{styleStandard}
\textbf{\ \ }\ \ Eso\ \ de\ \ la\ \ tanga\ \ es\ \ completamente\ \ cualquiera.
\end{styleStandard}

\begin{styleStandard}
\ \ \ \ that\ \ of\ \ the\ \ thong\ \ is\ \ completely\ \ \textsc{cualquiera}
\end{styleStandard}

\begin{styleStandard}
\ \ \ \ ‘That thong thing is complete(ly) bullshit/nonsense/out of place/bad.’
\end{styleStandard}

\begin{styleStandard}
\ \ \ \ (Rizzo Salierno 2013: 28)
\end{styleStandard}

\begin{styleStandard}
According to Rizzo Salierno (2013), the potential linguistic precedent of \textit{cualquiera} is possibly related to the Italianism \textit{cualunque}, which seems to have its origin in Lunfardo slang. \textit{Cualunque} in Argentinian Spanish (unlike Italian \textit{qualunque}, Author 2021) can also be used predicatively or attributively in a postnominal position and allows degree modification with \textit{muy} ‘very’ and other degree adverbs such as \textit{completamente}, as happens with ArgSp \textit{cualquiera} (see Kellert \& Francia in press for a historical overview and analysis of \textit{cualunque}):
\end{styleStandard}

\begin{styleStandard}
\textbf{\ \ }\ \ Washizu\ \ se\ \ asustó\ \ porque\ \ sus\ \ zapatos\ \ staban\ \ sucios\ \ y\ \ su\ \ ropa\ \ era
\end{styleStandard}

\begin{styleStandard}
\ \ \ \ Washizu\ \ him\ \ scared\ \ because\ \ his\ \ shoes\ \ were\ \ dirty\ \ and\ \ his\ \ clothing\ \ was
\end{styleStandard}

\begin{styleStandard}
\ \ \ \ muy\ \ \textit{cualunque}.
\end{styleStandard}

\begin{styleStandard}
\ \ \ \ very\ \ \textsc{cualunque}
\end{styleStandard}

\begin{styleStandard}
\ \ \ \ ‘Washizu was scared because his shoes were dirty and his clothes were very \textit{bad/out of place.}’
\end{styleStandard}

\begin{styleStandard}
\ \ \ \ (Rizzo Salierno 2013: 29)
\end{styleStandard}

\begin{styleStandard}
\textbf{\ \ }\ \ Es\ \ una\ \ \textit{flaca}\ \ \textit{cualunque}\ \ que\ \ se\ \ cree\ \ que\ \ es\ \ Marilyn\ \ Monroe.
\end{styleStandard}

\begin{styleStandard}
\ \ \ \ is\ \ a\ \ woman\ \ cualunque\ \ that\ \ her\ \ believes\ \ that\ \ is\ \ Marilyn\ \ Monroe
\end{styleStandard}

\begin{styleStandard}
\ \ \ \ ‘She’s an ordinary woman who thinks she’s Marilyn Monroe.’
\end{styleStandard}

\begin{styleStandard}
\ \ \ \ (Rizzo Salierno 2013: 29)
\end{styleStandard}

\begin{styleStandard}
Rizzo Salierno does not discuss cases of \textit{cualquiera} under transitive verbs, and he does not analyse the meaning ‘nonsense’, probably because he believes that \textit{cualquiera} has its origin in \textit{cualunque} and because he observed \textit{cualunque} only in predicative contexts. As Kellert \& Francia (in press) show \textit{cualunque }is indeed used mostly in predicative contexts or as a nominal modifier. They conclude that the transitive use of \textit{cualquiera} with the meaning ‘nonsense’ cannot have its origin in \textit{cualunque}. The transitive use of \textit{cualquiera} has been previously noted in the literature (see 2.1.4). Conde (2011) assumes that the transitive use of \textit{cualquiera} as ‘nonsense’ is specific to Lunfardo. However, as of yet this use has not been investigated in the literature.
\end{styleStandard}

\subsection[4.2 Open questions and main hypothesis]{\textbf{4.2 Open questions and main hypothesis}}
\begin{styleStandard}
The main question of this section is to determine the origin of the semantic and syntactic change of \textit{cualquiera} and to account for different meanings of \textit{cualquiera} as ‘ordinary’ and ‘nonsense’ and for their different syntactic functions as arguments of transitive verbs and as predicative elements.
\end{styleStandard}

\begin{styleStandard}
\ \ Rizzo Salierno suggests that some analogy with \textit{cualunque} might be the cause of a change. However, despite some similarities between \textit{cualunque} and \textit{cualquiera} in predicative contexts, they have a different distribution under transitive verbs (see Kellert \& Francia). Whereas \textit{cualquiera} occurs as a complement of verbs of saying, \textit{cualunque} does not:
\end{styleStandard}

\begin{styleStandard}
\textbf{\ \ }\ \ Me\ \ contestó/dijo\ \ \textit{cualquiera/*cualunque}
\end{styleStandard}

\begin{styleStandard}
\ \ \ \ Me\ \ answered/said\ \ \textsc{cualquiera/cualunque}
\end{styleStandard}

\begin{styleStandard}
\ \ \ \ ‘(S)he answered me with some bullshit/nonsense.’
\end{styleStandard}

\begin{styleStandard}
\ \ \ \ (Rizzo Salierno 2013: 26)
\end{styleStandard}

\begin{styleStandard}
Even if Rizzo Salierno was right about the assumption that \textit{cualquiera} has acquired its evaluative meaning from \textit{cualunque}, this assumption does not explain the origin of the evaluative meaning, which is absent from the Italian source item \textit{qualunque}. It just shifts the issue of how an indefinite can develop this type of evaluative meaning to \textit{cualunque}, but does not resolve it. Rizzo Salierno does not try to explain the origin of the evaluative use of ArgSp \textit{cualquiera} either historically or in the semantic structure of this item, and thus does not try to relate the evaluative function of \textit{cualquiera/cualunque} to the FCI meaning of these items. 
\end{styleStandard}

\begin{styleStandard}
I will fill this gap. Moreover, I will test Rizzo Salierno’s description of \textit{cualquiera} as a gradable adjective and investigate whether it behaves like a regular gradable adjective that can be modified by different degree modifiers (also including, e.g., items meaning ‘a bit’) and be used in comparative constructions such as (e.g. \textit{más cualquiera que} ‘worse than’). I will test this hypothesis with other degree modifiers and comparative constructions that are usually possible with degree adjectives.
\end{styleStandard}

\begin{styleStandard}
\ \ As I already said, my main hypothesis about the diachronic origin of the different readings of \textit{cualquiera} is that both the evaluative meaning ‘ordinary, bad’ (EVAL 1) and the meaning ‘nonsense’ (EVAL 2) have their origin in the free choice indefinite meaning ‘any’. I already showed in Chapter 3 how this hypothesis applies to EVAL 1, and I now show this for the special ArgSp meaning EVAL 2.
\end{styleStandard}

\begin{styleStandard}
\textbf{\ \ }\ \ Free Choice \textit{cualquiera }‘any’ \textit{{\textgreater} }Adj \textit{cualquiera }‘unremarkable/bad’ (EVAL 1) or Noun \textit{cualquiera} ‘nonsense’ (EVAL 2)
\end{styleStandard}

\begin{styleStandard}
In order to test this hypothesis for EVAL 2, I will look at the distribution of \textit{cualquiera} in synchrony and diachrony in ArgSp in this section and offer an analysis that explains the change in (225). 
\end{styleStandard}

\subsection[4.3 New data on cualquiera in Argentinian Spanish]{\textbf{4.3 New data on }\textbf{\textit{cualquiera }}\textbf{in Argentinian Spanish}}
\begin{styleStandard}
This section presents new data on \textit{cualquiera} in Argentinian Spanish which will be analysed in 4.4.
\end{styleStandard}

\subsubsection[4.3.1 Distribution of special uses of cualquiera across Latin American varieties]{\textbf{4.3.1 Distribution of special uses of }\textbf{\textit{cualquiera }}\textbf{across Latin American varieties}}
\begin{styleStandard}
In this subsection we will see that predicative \textit{cualquiera} and \textit{cualquiera} under transitive verbs such as \textit{mandar} ‘say’ are more frequent in Argentinian Spanish than in other varieties of Spanish.
\end{styleStandard}

\begin{styleStandard}
\ \ The predicative \textit{cualquiera} as in \textit{es cualquiera }(Table 8) and\textit{ parece cualquiera} (Table 9) is most frequent in Argentinian Spanish in contrast to other varieties (decimal numbers represent relative frequency in per million, i.e. how often a word or a word phrase occurs per one million words). The number 0.32 represents the mean of how many times \textit{es cualquiera} appears in all varieties per million words. The crucial observation is that the number of times \textit{es cualquiera} appears in Argentinian Spanish is three times higher than the mean. The number 1.05 represents the average number of instances of \textit{es cualquiera} per million in the Argentinian subcorpus, i.e. a set of one million words in the Argentinian subcorpus contains an average of 1.05 instances of \textit{es cualquiera}. 
\end{styleStandard}

\begin{center}
\tablefirsthead{}
\tablehead{}
\tabletail{}
\tablelasttail{}
\begin{supertabular}{m{0.39275986in}m{0.39345986in}|m{0.39415985in}|m{0.39345986in}|m{0.39275986in}|m{0.39345986in}|m{0.39415985in}|m{0.39345986in}|m{0.39275986in}|m{0.39345986in}|m{0.39415985in}|m{0.39345986in}|m{0.39275986in}|}
\hline
\multicolumn{1}{|m{0.39275986in}|}{\textbf{context}} &
\textbf{all} &
\textbf{per million} &
\textbf{Argentina} &
\textbf{Bolivia} &
\textbf{Chile} &
\textbf{Colombia} &
\textbf{Costa Rica} &
\textbf{Cuba} &
\textbf{Dom. Rep.} &
\textbf{Ecuador} &
\textbf{*Spain} &
\textbf{Guatemala}\\\hline
\multicolumn{1}{|m{0.39275986in}|}{\centering \textit{es cualquiera}} &
\centering 633 &
\centering 0.32 &
{\centering 1.05\par}

\centering 178 &
{\centering 0.20\par}

\centering 8 &
{\centering 0.11\par}

\centering 7 &
{\centering 0.20\par}

\centering 33 &
{\centering 0.78\par}

\centering 23 &
{\centering 0.38\par}

\centering 24 &
{\centering 0.27\par}

\centering 9 &
{\centering 0.21\par}

\centering 11 &
{\centering 0.18\par}

\centering 77 &
{\centering 0.24\par}

\centering\arraybslash 13\\\hline
 &
 &
\textbf{Honduras} &
\textbf{Mexico} &
\textbf{Nicaragua} &
\textbf{Panama} &
\textbf{Peru} &
\textbf{Puerto Rico} &
\textbf{Paraguay} &
\centering \textbf{El Salvador} &
\centering \textbf{USA} &
\centering \textbf{Uruguay} &
\centering\arraybslash \textbf{Venezuela}\\\hhline{~~-----------}
 &
 &
{\centering 0.43\par}

\centering 15 &
{\centering 0.18\par}

\centering 45 &
{\centering 0.25\par}

\centering 8 &
{\centering 0.58\par}

\centering 13 &
{\centering 0.24\par}

\centering 26 &
{\centering 0.40\par}

\centering 13 &
{\centering 0.27\par}

\centering 8 &
{\centering 0.30\par}

\centering 11 &
{\centering 0.36\par}

\centering 59 &
{\centering 0.72\par}

\centering 28 &
{\centering 0.24\par}

\centering\arraybslash 24\\\hhline{~~-----------}
\end{supertabular}
\end{center}
\begin{styletablecaption}
Table \stepcounter{Table}{\theTable}: Distribution of \textit{es cualquiera} in CdE\_web dialects (highlighting my own).
\end{styletablecaption}

\begin{styleStandard}
As we can see, the construction does not seem to be exluded in EurSp, but seems to be extremely marginal there, as it also is in Chile, Bolivia, Colomia, and Mexico (with values ${\leq}$ 0.2 per million). It is more frequent in most other Latin American countries, but the value rises to more than 0.5 per million only in a small part of Central America (Panama, 0.58, and Costa Rica, 0.78) and in South America (in Uruguay, 0.72), whereas in Argentina it rises to 1.05 per million.
\end{styleStandard}

\begin{styleStandard}
\ \ The same observation holds for \textit{parece cualquiera}, which is even more marginal than \textit{es cualquiera}, whose frequency is shown in Table 10. In most Latin American varieties in CdE, \textit{parece cualquiera} is not used at all. It is only used in Argentina, Spain, Mexico, and Peru. In Argentina, it is used with a strikingly higher frequency than in other countries (12 times higher than the average in the three other countries).
\end{styleStandard}

\begin{center}
\tablefirsthead{}
\tablehead{}
\tabletail{}
\tablelasttail{}
\begin{supertabular}{m{0.39275986in}m{0.39345986in}|m{0.39415985in}|m{0.39345986in}|m{0.39275986in}|m{0.39345986in}|m{0.39415985in}|m{0.39345986in}|m{0.39275986in}|m{0.39345986in}|m{0.39415985in}|m{0.39345986in}|m{0.39275986in}|}
\hline
\multicolumn{1}{|m{0.39275986in}|}{\textbf{context}} &
\textbf{all} &
\textbf{per million} &
\textbf{Argentina} &
\textbf{Bolivia} &
\textbf{Chile} &
\textbf{Colombia} &
\textbf{Costa Rica} &
\textbf{Cuba} &
\textbf{Dom. Rep.} &
\textbf{Ecuador} &
\textbf{Spain} &
\textbf{Guatemala}\\\hline
\multicolumn{1}{|m{0.39275986in}|}{\centering \textit{parece cualquiera}} &
\centering 30 &
\centering 0.02 &
{\centering 0.15\par}

\centering 25 &
 &
 &
 &
 &
 &
 &
 &
{\centering 0.01\par}

\centering 3 &
\\\hline
 &
 &
\textbf{Honduras} &
\textbf{Mexico} &
\textbf{Nicaragua} &
\textbf{Panama} &
\textbf{Peru} &
\textbf{Puerto Rico} &
\textbf{Paraguay} &
\centering \textbf{El Salvador} &
\centering \textbf{USA} &
\centering \textbf{Uruguay} &
\centering\arraybslash \textbf{Venezuela}\\\hhline{~~-----------}
 &
 &
 &
{\centering 0.01\par}

\centering 1 &
 &
 &
{\centering 0.01\par}

\centering 1 &
 &
 &
 &
 &
 &
\\\hhline{~~-----------}
\end{supertabular}
\end{center}
\begin{styletablecaption}
Table \stepcounter{Table}{\theTable}: Distribution of \textit{parece cualquiera} in CdE\_web dialects (highlighting my own).
\end{styletablecaption}

\begin{styleStandard}
The use of \textit{cualquiera} under the verb \textit{mandar} ‘say’ is also more frequent in Argentinian Spanish than in the other varieties as shown by the difference in percentages (75\% in Argentinian Spanish vs 25\% in all the other varieties).
\end{styleStandard}

\begin{styleStandard}
\ \ Of the 44 occurrences of this construction in the corpus, 33 (= 75\%) are found for Argentina, as the following table shows. Moreover, the most frequent verb forms are the indicative present and past tense forms, which refer to some past or present speech events, i.e. someone said or is presently saying something \textit{wrong/nonsense}:
\end{styleStandard}

\begin{center}
\tablefirsthead{}
\tablehead{}
\tabletail{}
\tablelasttail{}
\begin{supertabular}{|m{1.2983599in}|m{0.41085985in}|m{0.6101598in}|m{0.70945984in}|m{0.6101598in}|m{0.8073598in}|m{0.6295598in}|}
\hline
\centering \textbf{Context} &
\centering \textbf{All} &
\centering \textbf{Argentina} &
\centering \textbf{Other countries} &
{\centering \textbf{Indicative}\par}

{\centering \textbf{Present}\par}

{\centering \textbf{Imperf}\par}

\centering \textbf{Perfect} &
{\centering \textbf{Subjunctive/}\par}

\centering \textbf{Imperative} &
{\centering \textbf{Gerund}\par}

{\centering \textbf{Infinitive}\par}

\\\hline
\centering \textit{mandan cualquiera} &
\centering 8 &
\centering 6 &
\centering 2 &
\centering + &
 &
\\\hline
\centering \textit{mandaste cualquiera} &
\centering 5 &
\centering 3 &
\centering 2 &
\centering + &
 &
\\\hline
\centering \textit{mandaron cualquiera} &
\centering 5 &
\centering 5 &
\centering 0 &
\centering + &
 &
\\\hline
\centering \textit{mandé cualquiera} &
\centering 4 &
\centering 4 &
\centering 0 &
\centering + &
 &
\\\hline
\centering \textit{mandar cualquiera} &
\centering 3 &
\centering 2 &
\centering 1 &
 &
 &
\centering\arraybslash +\\\hline
\centering \textit{mandando cualquiera} &
\centering 3 &
\centering 2 &
\centering 1 &
 &
 &
\centering\arraybslash +\\\hline
\centering \textit{mando cualquiera} &
\centering 3 &
\centering 3 &
\centering 0 &
\centering + &
 &
\\\hline
\centering \textit{manda cualquiera} &
\centering 3 &
\centering 2 &
\centering 1 &
\centering + &
 &
\\\hline
\centering \textit{mande cualquiera} &
\centering 3 &
\centering 1 &
\centering 2 &
 &
\centering + &
\\\hline
\centering \textit{mandó cualquiera} &
\centering 2 &
\centering 2 &
\centering 0 &
\centering + &
 &
\\\hline
\centering \textit{mandaba cualquiera} &
\centering 2 &
\centering 1 &
\centering 1 &
\centering + &
 &
\\\hline
\centering \textit{mandaban cualquiera} &
\centering 2 &
\centering 1 &
\centering 1 &
\centering + &
 &
\\\hline
\centering \textit{mandamos cualquiera} &
\centering 1 &
\centering 1 &
\centering 1 &
\centering + &
 &
\\\hline
\centering \textbf{Total} &
{\centering 100\%\par}

\centering 44 &
{\centering 75\%\par}

\centering 33 &
{\centering 25\%\par}

\centering 11 &
 &
 &
\\\hline
\end{supertabular}
\end{center}
\begin{styletablecaption}
Table \stepcounter{Table}{\theTable}: Distribution of lemma\textit{ mandar cualquiera} in CdE\_web dialects.
\end{styletablecaption}

\begin{styleStandard}
The difference in frequencies suggests that \textit{cualquiera} has undergone a change in Argentinian Spanish. Despite the clear difference in frequencies of the special uses between Argentinian and other Latin American varieties, the frequency of the special uses is very low compared to other phrases (e.g. \textit{decir algo} ‘say something’, with 408 occ., or \textit{es común} ‘it is common’, with 1,447 occ.). One reasonable explanation of the low frequency of these uses is that they represent informal spoken speech (see also Rizzo Salierno 2013, Conde 2011, among others). 
\end{styleStandard}

\begin{styleStandard}
\ \ The other Latin Spanish varieties will not be considered here, nor will the question be considered as to whether Argentinian Spanish has influenced the use of \textit{cualquiera} in other varieties. 
\end{styleStandard}

\subsubsection[4.3.2 Postnominal cualquiera]{\textbf{4.3.2 }\textbf{Postnominal }\textbf{\textit{cualquiera}}}
\begin{styleStandard}
In this section, I will look into the predicative ArgSp \textit{cualquiera} in detail and see whether it behaves as a gradable adjective in any respect and whether it allows any kind of syntactic combination like usual gradable adjectives do (e.g. different combinations with adverbs, comparatives, etc.). 
\end{styleStandard}

\begin{styleStandard}
\ \ As already observed in Table 4 and in Section 4.1, ArgSp predicative \textit{cualquiera} can be modified by degree adverbs that denote high degrees (i.e. adverbs meaning ‘very’ and the like). The most frequent modifying item is \textit{re}\footnote{\ \textrm{See footnote 4. The detailed linguistic status of }\textrm{\textit{re}}\textrm{, as well as its origin, needs to be addressed in future research. In this thesis, I assume in line with Rizzo Salierno (2013) and Kanwit \& Terán (2020) that it is a degree adverb similar to }\textrm{\textit{muy}}\textrm{ ‘very’ that denotes a high degree. The form }\textrm{\textit{re}}\textrm{ has been described as a derivative prefix from Latin that has grammaticalised as an intensifier (Kanwit \& Terán 2020: 257 and references therein). The intensifier }\textrm{\textit{re}}\textrm{ has developed from a locative or spatial marker to a marker of repetition or iteration, and finally to an intensifier adverb (see Kanwit \& Terán 2020: 257 and references therein).}} (25 occurrences in the corpus). The most important observation is that the predicative \textit{cualquiera} has the meaning ‘nonsense’ if the subject of the predicative complement is a proposition or an event: 
\end{styleStandard}

\begin{styleStandard}
\textbf{\ \ }\ \ ademas\ \ me\ \ parece\ \ \textit{re\ \ cualquiera\ \ }eso\ \ de\ \ q\ \ se\ \ encajo\ \ una\ \ piña\ \ 
\end{styleStandard}

\begin{styleStandard}
\ \ \ \ moreover\ \ to.me\ \ seems\ \ very\ \ \textsc{cualquiera\ \ }that\ \ of\ \ that\ \ (s)he\ \ received\ \ a\ \ punch\ \ 
\end{styleStandard}

\begin{styleStandard}
\ \ \ \ sola\ \ para\ \ hecharle\ \ la\ \ culpa\ \ a\ \ Mar
\end{styleStandard}

\begin{styleStandard}
\ \ \ \ alone\ \ for\ \ throw.her\ \ the\ \ blame\ \ to\ \ Mar
\end{styleStandard}

\begin{styleStandard}
\ \ \ \ ‘Moreover, that receiving a punch simply to blame Mar seems complete nonsense to me.’
\end{styleStandard}

\begin{styleStandard}
\ \ \ \ (@ 106)
\end{styleStandard}

\begin{styleStandard}
All occurrences of predicative \textit{cualquiera} in the Argentinain data of CdE\_web dialects co-occur with adverbial or verbal phrases that express whose negative evaluation it is that \textit{cualquiera} denotes (see the expression of the concept of judge in Lasersohn 2005, among others; and Chapter 3 on the judge parameter of \textit{tasty}{}-predicates). In almost all cases that I found in CdE\_web dialects, the evaluation is linked to the speaker, e.g. \textit{para mí} ‘for me’, \textit{me parece que} ‘It seems to me that…’. Of the 25 occurrences of \textit{parece cualquiera}, 17 are \textit{me parece}, two are \textit{te parece}, and one is \textit{parece}. All cases represent informal speech/dialogue:
\end{styleStandard}

\begin{styleStandard}
\textbf{\ \ }\ \ No\ \ lo\ \ terminé\ \ de\ \ leer\ \ porque\ \ me\ \ pareció\ \ cualquiera.
\end{styleStandard}

\begin{styleStandard}
\ \ \ \ not\ \ it\ \ finish\ \ of\ \ read\ \ because\ \ me\ \ seemed\ \ \textsc{cualquiera}
\end{styleStandard}

\begin{styleStandard}
\ \ \ \ ‘I did not finish reading it because it seemed \textit{nonsense} to me.’
\end{styleStandard}

\begin{styleStandard}
\ \ \ \ (@ 107)
\end{styleStandard}

\begin{styleStandard}
\textbf{\ \ }\ \ Para\ \ mí,\ \ me\ \ parece\ \ cualquiera.
\end{styleStandard}

\begin{styleStandard}
\ \ \ \ for\ \ me\ \ me\ \ seems\ \ \textsc{cualquiera}
\end{styleStandard}

\begin{styleStandard}
\ \ \ \ ‘As for me, I find it \textit{nonsense}.’
\end{styleStandard}

\begin{styleStandard}
\ \ \ \ (@ 108)
\end{styleStandard}

\begin{styleStandard}
\textbf{\ \ }\ \ eso\ \ es\ \ cualquiera,\ \ al\ \ menos\ \ para\ \ mí.
\end{styleStandard}

\begin{styleStandard}
\ \ \ \ that\ \ is\ \ \textsc{cualquiera}\ \ at\ \ least\ \ for\ \ me
\end{styleStandard}

\begin{styleStandard}
\ \ \ \ ‘That is \textit{nonsense}, at least for me.’
\end{styleStandard}

\begin{styleStandard}
\ \ \ \ (@ 109)
\end{styleStandard}

\begin{styleStandard}
However, I also find few cases of evaluative \textit{cualquiera} with 3\textsuperscript{rd} person opinion holder (see also Lasersohn 2005 for the same observation with \textit{tasty}{}-predicates):
\end{styleStandard}

\begin{styleStandard}
\textbf{\ \ }\ \ haría lo mismo pero si hay gente que no le gusta porque piensa que es cualquiera (yo no pienso eso).....
\end{styleStandard}

\begin{styleStandard}
\ \ \ \ ‘I would do the same but there are people that don’t like it because they think it is \textit{nonsense} (I don’t think so)….’
\end{styleStandard}

\begin{styleStandard}
\ \ \ \ (@ 110)
\end{styleStandard}

\begin{styleStandard}
In Section 4.3.3, we will see that \textit{cualquiera} with the meaning ‘nonsense’ also occurs under transitive verbs.
\end{styleStandard}

\begin{styleStandard}
\ \ The meaning of \textit{cualquiera} as ‘nonsense’ can be reinforced by degree adverbs that denote high degrees such as \textit{muy/demasiado/tan}/\textit{re}. Degree modification of \textit{cualquiera} as ‘nonsense’ is unexpected under the assumption that \textit{cualquiera} is a noun when it means ‘nonsense’. I will return to this puzzle in my analysis in 4.4. I translate in this case the modified \textit{cualquiera} as ‘quite nonsensical’, ‘such utter nonsense’, or as ‘completely out of place’:
\end{styleStandard}

\begin{styleStandard}
\textbf{\ \ }\ \ Me parece todo tan cualquiera.
\end{styleStandard}

\begin{styleStandard}
‘Everything seems such utter nonsense to me’
\end{styleStandard}

\begin{styleStandard}
(@ 111)
\end{styleStandard}

\begin{styleStandard}
If the subject of the predicative \textit{cualquiera} refers to a person or an entity, the meaning of \textit{cualquiera} can take on different flavours already discussed in Chapter 3, such as ‘ordinary, unremarkable’ and the pejorative meaning of ‘bad’. 
\end{styleStandard}

\begin{styleStandard}
\textbf{\ \ }\ \ Yo, que soy tan cualquiera como cualquiera.\footnote{\ This \textit{cualquiera} is a pronoun (see Chapter 2).}
\end{styleStandard}

\begin{styleStandard}
‘I who am so ordinary like anyone else.’
\end{styleStandard}

\begin{styleStandard}
\ \ \ \ (@ 112)
\end{styleStandard}

\begin{styleStandard}
Argentinian Spanish also has \textsc{un} N \textit{cualquiera} in predicative position just as in European Spanish (e.g. \textit{una fiesta re cualquiera} ‘a very bad party’). The big difference to European Spanish is that \textit{cualquiera} can also be modified by a degree adjective and be used as a predicative adjective as in (232).
\end{styleStandard}

\begin{styleStandard}
\ \ Note that \textit{cualquiera} can predicate over plural entities as well, as shown in . In the following context, \textit{cualquiera} expresses an evaluative property of entitities such as pictures according to a judge (here: B). In the following context, the evaluation concerns what the pictures represent (i.e. namely that they show a bad example):
\end{styleStandard}

\begin{styleStandard}
\label{bkm:Ref150786385}\textbf{\ \ }A:\ \ Esas\ \ fotos\ \ que\ \ publican\ \ parece\ \ que\ \ dan\ \ mal\ \ ejemplo
\end{styleStandard}

\begin{styleStandard}
\ \ \ \ these\ \ pics\ \ that\ \ publish\ \ seem\ \ that\ \ give\ \ bad\ \ example
\end{styleStandard}

\begin{styleStandard}
\ \ B:\ \ Sí,\ \ son\ \ re\ \ cualquiera.
\end{styleStandard}

\begin{styleStandard}
\ \ \ \ yes\ \ are\ \ very\ \ \textsc{cualquiera}
\end{styleStandard}

\begin{styleStandard}
\ \ A: ‘These pics that they publish it seems that they give the wrong idea.’
\end{styleStandard}

\begin{styleStandard}
B: ‘Yes, they[these photos] are completely out of place. The idea that these pictures represent does not make sense.’
\end{styleStandard}

\begin{styleStandard}
\ \ (similar to @ 113)
\end{styleStandard}

\begin{styleStandard}
The following data show that \textit{cualquiera }need not necessarily combine with adverbs that denote high degrees (against Salierno’s 2013 observation). Actually, we find two occurences with modification by the adverb \textit{medio}, which literally denotes ‘half’ and means ‘a bit/almost’ when combined with gradable adjectives like in \ and \ and \textit{cualquiera }as in \ and :
\end{styleStandard}

\begin{styleStandard}
\label{bkm:Ref150786427}\textbf{\ \ }\ \ \textit{Es\ \ medio\ \ complicado\ \ }este\ \ asunto
\end{styleStandard}

\begin{styleStandard}
\ \ \ \ is\ \ half\ \ complicated\ \ this\ \ case
\end{styleStandard}

\begin{styleStandard}
\ \ \ \ ‘It’s a bit complicated, this case’
\end{styleStandard}

\begin{styleStandard}
\ \ \ \ (@ 114)
\end{styleStandard}

\begin{styleStandard}
\label{bkm:Ref150786436}\textbf{\ \ }\ \ matemáticamente\ \ \textit{es\ \ medio\ \ imposible\ \ }que\ \ todas\ \ las\ \ personas\ \ puedan\ \ pasar
\end{styleStandard}

\begin{styleStandard}
\ \ \ \ mathematically\ \ is\ \ half\ \ impossible\ \ that\ \ all\ \ the\ \ people\ \ can\ \ go
\end{styleStandard}

\begin{styleStandard}
\ \ \ \ por\ \ esa\ \ situación
\end{styleStandard}

\begin{styleStandard}
\ \ \ \ through\ \ that\ \ situation
\end{styleStandard}

\begin{styleStandard}
\ \ \ \ ‘Mathematically it’s almost impossible for all the people to go through that situation’
\end{styleStandard}

\begin{styleStandard}
\ \ \ \ (@ 115)
\end{styleStandard}

\begin{styleStandard}
\label{bkm:Ref150786444}\textbf{\ \ }\ \ es\ \ que\ \ la\ \ defensa\ \ del\ \ rugby\ \ por\ \ los\ \ valores\ \ que\ \ entraña\ \ \textit{es\ \ medio\ \ cualquiera.}
\end{styleStandard}

\begin{styleStandard}
\ \ \ \ is\ \ that\ \ the\ \ defence\ \ of\ \ rugby\ \ of\ \ the\ \ values\ \ that\ \ involves\ \ is\ \ half\ \ \textsc{cualquiera}
\end{styleStandard}

\begin{styleStandard}
\ \ \ \ O sea,\ \ a\ \ mí\ \ es\ \ un\ \ deporte\ \ que\ \ no\ \ me\ \ atrae.\ \ \textbf{\ \ \ \ }
\end{styleStandard}

\begin{styleStandard}
\ \ \ \ that.is\ \ to\ \ me\ \ is\ \ a\ \ sport\ \ that\ \ not\ \ me\ \ attracts\ \ \ \ \ \ 
\end{styleStandard}

\begin{styleStandard}
\ \ \ \ ‘Defending rugby for the values it involves is a bit nonsensical/out of place. That is, it’s a sport that does not appeal me.’
\end{styleStandard}

\begin{styleStandard}
\ \ \ \ (similar to @ 116)
\end{styleStandard}

\begin{styleStandard}
\label{bkm:Ref150786457}\textbf{\ \ }\ \ \textit{es\ \ medio\ \ cualquiera\ \ }eso
\end{styleStandard}

\begin{styleStandard}
\ \ \ \ is\ \ half\ \ \textsc{cualquiera}\ \ that
\end{styleStandard}

\begin{styleStandard}
\ \ \ \ ‘That is a bit\textit{ nonsensical}.’
\end{styleStandard}

\begin{styleStandard}
\ \ \ \ (@ 117)
\end{styleStandard}

\begin{styleStandard}
I even found that the diminuitive\textit{ }suffix\textit{ -ita} can combine with \textit{cualquiera}, and in this case the speaker attenuates the bad quality ‘a bit nonsensical’:
\end{styleStandard}

\begin{styleStandard}
\textbf{\ \ }\ \ Bueno,\ \ cerramos\ \ este\ \ post\ \ que\ \ \textit{es\ \ cualquierita}
\end{styleStandard}

\begin{styleStandard}
\ \ \ \ well\ \ close\ \ this\ \ post\ \ that\ \ is\ \ \textsc{cualquiera}.\textsc{dim}
\end{styleStandard}

\begin{styleStandard}
\ \ \ \ ‘Well, let’s close this post because it’s a bit nonsensical’
\end{styleStandard}

\begin{styleStandard}
\ \ \ \ (@ 118) 
\end{styleStandard}

\begin{styleStandard}
If \textit{cualquiera} has been reanalysed as a gradable adjective, as such data suggest, it should be possible to find it in comparative clauses such as ‘he is worse than anyone else’. Surprisingly, the comparative with \textit{más que} was not found in CdE\_web dialects. However, the following case from the corpus might be interpreted as an elliptical comparative clause:
\end{styleStandard}

\begin{styleStandard}
\textbf{\ \ }\ \ el\ \ reggae\ \ es\ \ \textit{cualquiera}.\ \ Mas\ \ [cualquiera] \ \ el\ \ reggae\ \ nacional.\ \ 
\end{styleStandard}

\begin{styleStandard}
\ \ \ \ the\ \ reggae\ \ is\ \ \textsc{cualquiera}\ \ more\ \ \textsc{[cualquiera]}\ \ the\ \ reggae\ \ national\ \ 
\end{styleStandard}

\begin{styleStandard}
\ \ \ \ ‘Reggae is \textit{bad}. National reggae is (even) worse.’
\end{styleStandard}

\begin{styleStandard}
\ \ \ \ (@ 119, interpolation my own)
\end{styleStandard}

\begin{styleStandard}
There are some examples of comparatives with \textit{cualquiera} on the Internet:
\end{styleStandard}

\begin{styleStandard}
\textbf{\ \ }\ \ Más\ \ cualquiera\ \ que\ \ el\ \ otro\ \ \ \ \ \ \ \ \ \ 
\end{styleStandard}

\begin{styleStandard}
\ \ \ \ more\ \ \textsc{cualquiera}\ \ than\ \ the\ \ other\ \ \ \ \ \ \ \ \ \ 
\end{styleStandard}

\begin{styleStandard}
\ \ \ \ ‘worse than the other one.’
\end{styleStandard}

\begin{styleStandard}
\ \ \ \ (@ 120)
\end{styleStandard}

\begin{styleStandard}
\textbf{\ \ }\ \ Y\ \ con\ \ los\ \ libros\ \ me\ \ pasó\ \ más\ \ o\ \ menos\ \ lo\ \ mismo\ \ aunque
\end{styleStandard}

\begin{styleStandard}
\ \ \ \ and\ \ with\ \ the\ \ books\ \ me\ \ happened\ \ more\ \ or\ \ less\ \ the\ \ same\ \ although
\end{styleStandard}

\begin{styleStandard}
\ \ \ \ no\ \ los\ \ pienso\ \ seguir\ \ porque\ \ son\ \ más\ \ \textit{cualquiera\ \ }que\ \ la\ \ serie\ \ por\ \ ahí
\end{styleStandard}

\begin{styleStandard}
\ \ \ \ not\ \ them\ \ think\ \ follow\ \ because\ \ are\ \ more\ \ \textsc{cualquiera}\ \ than\ \ the\ \ series\ \ for\ \ there
\end{styleStandard}

\begin{styleStandard}
\ \ \ \ ‘And with the books more or less the same thing happened to me, although I will definitely not follow them because they are worse\textit{ }than the series that is out there.’
\end{styleStandard}

\begin{styleStandard}
\ \ \ \ (@ 121)
\end{styleStandard}

\begin{styleStandard}
\textbf{\ \ }\ \ me\ \ parece\ \ mas\ \ cualquiera\ \ que\ \ la\ \ cumbia
\end{styleStandard}

\begin{styleStandard}
\ \ \ \ to.me\ \ seems\ \ more\ \ \textsc{cualquiera}\ \ than\ \ the\ \ cumbia
\end{styleStandard}

\begin{styleStandard}
\ \ \ \ ‘It seems worse to me than the cumbia.’
\end{styleStandard}

\begin{styleStandard}
\ \ \ \ (@ 122)
\end{styleStandard}

\begin{styleStandard}
It thus seems that \textit{cualquiera} does occur in comparative structures, although very rarely, so the corpus is probably too small to detect it or \textit{cualquiera} is still on its way of being reanalysed as a gradable adjective. Another argument that \textit{cualquiera} has been already reanalysed as a gradable adjective is shown by the difference in interpretation of \textit{cualquiera} under negation. Recall from Chapter 2 that, when it is used as an FCI, \textit{cualquiera} has a ‘not just any’ interpretation under negation. However, \textit{cualquiera} has a different interpretation when it occurs in predicative position as an adjective in , in contrast to the indefinite pronoun in : 
\end{styleStandard}

\begin{styleStandard}
\label{bkm:Ref63460902}\textbf{\ \ }\ \ No\ \ es\ \ tan\ \ cualquiera\ \ pero\ \ sí\ \ medio\ \ cualquiera
\end{styleStandard}

\begin{styleStandard}
\ \ \ \ not\ \ is\ \ so\ \ \textsc{cualquiera}\ \ but\ \ yes\ \ half\ \ \textsc{cualquiera}
\end{styleStandard}

\begin{styleStandard}
\ \ \ \ ‘It’s not very \textit{bad }but it is a bit bad.’
\end{styleStandard}

\begin{styleStandard}
\ \ \ \ (@ 123)
\end{styleStandard}

\begin{styleStandard}
\label{bkm:Ref63460922}\textbf{\ \ }\ \ es\ \ un\ \ portal,\ \ y\ \ no\ \ es\ \ cualquiera\ \ es\ \ el\ \ mejor\ \ 
\end{styleStandard}

\begin{styleStandard}
\ \ \ \ is\ \ a\ \ platform\ \ and\ \ no\ \ is\ \ \textsc{cualquiera}\ \ is\ \ the\ \ best
\end{styleStandard}

\begin{styleStandard}
\ \ \ \ portal\ \ de\ \ autos\ \ del\ \ pais
\end{styleStandard}

\begin{styleStandard}
\ \ \ \ platform\ \ of \ \ cars\ \ in.the\ \ country
\end{styleStandard}

\begin{styleStandard}
\ \ \ \ ‘It’s a platform and not just any platform. It’s the best car platform in this country.’
\end{styleStandard}

\begin{styleStandard}
\ \ \ \ (@ 124)
\end{styleStandard}

\begin{styleStandard}
This section has shown that \textit{cualquiera} is used as a gradable category with evaluative meanings auch as ‘nothing special’, ‘bad’, ‘unremarkable’, ‘nonsensical’, or ‘out of place’ in predicative contexts in Argentinian Spanish. It can be modified with adverbs \textit{re/muy/tan/medio} ‘very/a bit’ and to a lesser extent with comparative adverbs such as \textit{más} ‘more’.
\end{styleStandard}

\subsubsection[4.3.3 Pronominal cualquiera under transitive verbs]{\textbf{4.3.3 Pronominal }\textbf{\textit{cualquiera}}\textbf{\textit{ }}\textbf{under transitive verbs}}
\begin{styleStandard}
Until now in this chapter, I have presented data from ArgSp on predicative \textit{cualquiera} in which the item at issue has a series of evaluative meanings such as ‘bad’ or ‘nonsensical’. 
\end{styleStandard}

\begin{styleStandard}
\ \ In sentences with imperative mood, which I also analyse as a modalised constructions (see A\&M 2018 and Section 3.2.3), ArgSp \textit{cualquiera} is interpreted as an indefinite pronoun ‘anything’ (which would correspond to \textit{cualquier cosa }in EurSp). Recall from 2.1.4 that \textit{cualquiera} can be used as a substitute for \textit{cualquier cosa} in EurSp:
\end{styleStandard}

\begin{styleStandard}
\textbf{\ \ }\ \ \textstylehgkelc{\textit{¡}}\textit{Hagamos cualquiera! }(hacer algo, hacer lo que sea, o como que da lo mismo hacer una cosa u otra).
\end{styleStandard}

\begin{styleStandard}
\ \ \ \ ‘Let’s do something! (do something, do anything, or it doesn’t matter whether it’s this thing or the other).’
\end{styleStandard}

\begin{styleStandard}
(@ 125) 
\end{styleStandard}

\begin{styleStandard}
In the following two examples \textit{digo cualquiera} and \textit{hace cualquiera} can be interpreted as ‘I say anything possible’ and ‘she does anything possible’, respectively, i.e. \textit{cualquiera} can be interpreted as a Free Choice Indefinite pronoun ‘anything’ under covert subtrigging. But another possible interpretation is that the speaker uses \textit{cualquiera} to qualify or evaluate the result state of the verbs ‘say’ and ‘do’ as ‘nonsense’ (EVAL 2):
\end{styleStandard}

\begin{styleStandard}
\textbf{\ \ }\ \ al\ \ rato\ \ volvió\ \ porque\ \ sabe\ \ que\ \ yo\ \ dormida\ \ \textit{digo\ \ }cualquiera
\end{styleStandard}

\begin{styleStandard}
\ \ \ \ \ \ \ \ came back\ \ because\ \ know\ \ that\ \ I\ \ asleep\ \ say\ \ \textsc{cualquiera}
\end{styleStandard}

\begin{styleStandard}
\ \ \ \ a.‘She came back because she knows that when I sleep, I say anything possible.’ (FCI)
\end{styleStandard}

\begin{styleStandard}
b.‘She came back because she knows that when I sleep, I talk nonsense/rubbish.’ (EVAL 2)
\end{styleStandard}

\begin{styleStandard}
\ \ \ \ (@ 126)
\end{styleStandard}

\begin{styleStandard}
\textbf{\ \ }\ \ se\ \ va\ \ cuando\ \ quiere,\ \ hace\ \ cualquiera.
\end{styleStandard}

\begin{styleStandard}
\ \ \ \ her/him\ \ goes out\ \ when\ \ wants\ \ does\ \ \textsc{cualquiera}
\end{styleStandard}

\begin{styleStandard}
\ \ \ \ a. ‘(S)he goes out when (s)he wants to do so, (s)he does anything she wants. \ \ \ \ \ \ \ (FCI)
\end{styleStandard}

\begin{styleStandard}
b. ‘(S)he goes out when (s)he wants to do so, (s)he does nonsense.’ \ \ \ \ \ \ \ \ \ \ \ \ \ \ \ \ \ \ \ \ \ \ (EVAL 2)
\end{styleStandard}

\begin{styleStandard}
\ \ \ \ (@ 127)
\end{styleStandard}

\begin{styleStandard}
It thus seems that that the interpretation of \textit{cualquiera} under certain transitive verbs in the indicative as ‘anything’ (FCI) or ‘nonsense’ (EVAL 2) is context-dependent in Argentinian Spanish. In contrast, if \textit{cualquiera} occurs under the verb \textit{mandar} ‘to say’, it has only the evaluative meaning ‘to talk nonsense’, as has already been observed in Section 2.1.4. An example is given in :
\end{styleStandard}

\begin{styleStandard}
\label{bkm:Ref150786556}\textbf{\ \ }\ \ No\ \ supe\ \ qué\ \ decirle,\ \ así\ \ que\ \ le\ \ mandé\ \ cualquiera.
\end{styleStandard}

\begin{styleStandard}
\ \ \ \ not\ \ knew\ \ what\ \ say.her\ \ so\ \ that\ \ her\ \ told\ \ \textsc{cualquiera}
\end{styleStandard}

\begin{styleStandard}
\ \ \ \ ‘I did not know what to tell her, so I told her nonsense.’
\end{styleStandard}

\begin{styleStandard}
\ \ \ \ (@ 128)
\end{styleStandard}

\begin{styleStandard}
A crucial observation is that ArgSp \textit{cualquiera} also has an evaluative meaning like ‘nonsense’ when \textit{mandar} appears under a modal verb:\footnote{\ The context of the example is: “Aca el problema es que estos tipos, como los chantas KK, se abusan porque hay fanáticos como este Juan emilio que no ven la realidad, viven en otro mundo y les pueden mandar cualquiera total no les da la cabeza para darse cuenta que lo estan verseando.”\par ‘…they live in a different world and you can tell them complete bullshit, they don’t have the intelligence to realise that people are mocking them.’}
\end{styleStandard}

\begin{styleStandard}
\textbf{\ \ }\ \ y\ \ les\ \ pueden\ \ mandar\ \ cualquiera total\ \ 
\end{styleStandard}

\begin{styleStandard}
\ \ \ \ and\ \ that\ \ her\ \ sent\ \ \textsc{cualquiera total}
\end{styleStandard}

\begin{styleStandard}
\ \ \ \ ‘and you can tell them complete bullshit.’
\end{styleStandard}

\begin{styleStandard}
\ \ \ \ (@ 129a)
\end{styleStandard}

\begin{styleStandard}
The evaluative reading here is reinforced by the degree adverb \textit{total} ‘complete’. This observation suggests that \textit{cualquiera} under \textit{mandar} ‘to say’, no matter what mood \textit{mandar} has, always has the meaning of ‘to say nonsense/bullshit’ in contrast to other verbs that have different interpretations depending on the verbal mood. The only thinkable explanation is that the evaluative meaning of \textit{cualquiera} as ‘nonsense’ has been lexicalised in the construction of \textit{mandar cualquiera}, whereas with other verbs, the interpretation of \textit{cualquiera} depends on the verbal mood. 
\end{styleStandard}

\subsubsection[4.3.4 Summary of the data]{\textbf{4.3.4 Summary of the data}}
\begin{styleStandard}
Under modal verb or in modalised constructions such as subjunctives and imperatives, ArgSp \textit{cualquiera} has an FCI function in the sense of ‘anything’. Non-modal contexts such as transitive verbs in the indicative past and present tense may also trigger FCI if \textit{cualquiera} is licensed under subtrigging. \textit{Cualquiera} under transitive verbs is something typical for Argentinian Spanish but not for other varieties of Spanish. However, in certain contexts that express evaluation and with certain verbs such as \textit{mandar, cualquiera} has an evaluative interpretation of ‘nonsense/bullshit’. In predicative contexts such as with copular verbs, \textit{cualquiera} has only the evaluative interpretation in Argentinian Spanish. It is in this context that it can be modified by different degree adverbs and occur in degree comparatives, although modification with degree comparative adverbs occurs less frequently than modification with simple degree adverbs such as ‘very/a bit’. The interpretation of \textit{cualquiera} as ‘ordinary’, ‘nonsense’, or ‘anything’ depends strongly on the grammatical category. If it is a pronoun, it has the interpretation as ‘anything’. If it is a noun, it has the interpretation of ‘nonsense’, and if it is an adjective, it has the interpretation of ‘ordinary’ or ‘nonsensical’.
\end{styleStandard}

\begin{center}
\tablefirsthead{}
\tablehead{}
\tabletail{}
\tablelasttail{}
\begin{supertabular}{|m{0.57545984in}|m{2.50456in}|m{0.48305985in}|m{1.1754599in}|m{1.1643599in}|}
\hhline{~----}
\multicolumn{1}{m{0.57545984in}|}{} &
\centering \textbf{Indicative present or past tense} &
\centering \textbf{Modal} &
\centering \textbf{\textit{mandar }}\textbf{‘say’} &
\centering\arraybslash \textbf{Predicative}\\\hline
\centering \textbf{FCI} &
\centering + (subtrigging) &
\centering + &
\centering – &
\centering\arraybslash –\\\hline
\centering \textbf{Eval} &
\centering + &
\centering – &
\centering + &
\centering\arraybslash +\\\hline
\end{supertabular}
\end{center}
\begin{styletablecaption}
Table \stepcounter{Table}{\theTable}: Different interpretations in relation with verbal mood and verb type (+/– predicative verb).
\end{styletablecaption}

\subsection[4.4. Analysis of different uses of cualquiera in Modern Argentinian Spanish]{\textbf{4.4. Analysis of different uses of }\textbf{\textit{cualquiera}}\textbf{\textit{ }}\textbf{in Modern Argentinian Spanish}}
\begin{styleStandard}
In this section I will give an analysis of the different syntactic functions of \textit{cualquiera} in Modern Argentinian Spanish and discuss their origin. I will show that the indefinite use of \textit{cualquiera} ‘any’ is the origin of the evaluative functions of \textit{cualquiera} ‘ordinary’ and ‘nonsense’.
\end{styleStandard}

\begin{styleStandard}
\ \ The data with degree modification suggests that ArgSp \textit{cualquiera} used in predicative contexts is the result of a reanalysis as a gradable adjective denoting a relation between degrees and individuals. It thus behaves in a way similar to English \textit{tall}: \textit{Juan is very tall} denotes that ‘Juan is tall to degree d, and degree d denotes a high degree on a scale of tallness’. \textit{Cualquiera} is thus a predicate denoting a scale of ordinary qualities or the pejorative version of it (i.e. bad qualities; see Section 3 on the pejoration of ‘common’, ‘ordinary’). This kind of predicate has a subjective value, i.e. the badness is not objectively measurable in contrast to tallness. The gradable predicate thus takes an additional argument, which is the experiencer of the bad quality (see Section 4.3):
\end{styleStandard}

\begin{styleStandard}
\textbf{\ \ }\ \ Adj° \textit{cualquiera/cualunque}\footnote{\textrm{\ However, both }\textrm{\textit{cualquiera}}\textrm{ and }\textrm{\textit{cualunque}}\textrm{ are invariable adjectives as they do not inflect for person and gender. Number variation is possible in ArgSp }\textrm{\textit{las chicas son cualquieras/cualunques}}\textrm{ ‘the girls are nothing special’ (see also Rainer 1993 and Section 2.1).}}\textit{ }= \textit{$\lambda $d$\lambda $x}. bad/ordinary\textbf{ }(\textit{d})(\textit{x}) for some experiencer E
\end{styleStandard}

\begin{styleStandard}
e.g. \textit{la peli es re cualquiera}. ‘The movie is very ordinary’
\end{styleStandard}

\begin{styleStandard}
\textbf{\ \ }\ \ [DegP re/medio/más/demasiado/tan [AdjP cualquiera/cualunque]]
\end{styleStandard}

\begin{styleStandard}
As I showed in Section 4.3.2, \textit{cualquiera} can also have the meaning of ‘nonsense’ when the object of evaluation is not an entity or person, as exemplified in (250), but a proposition or an event, as in (252):
\end{styleStandard}

\begin{styleStandard}
\textbf{\ \ }\ \ Lo\ \ que\ \ me\ \ decís\ \ \textit{es\ \ cualquiera}
\end{styleStandard}

\begin{styleStandard}
\ \ \ \ the\ \ that\ \ me\ \ say\ \ is\ \ \textsc{cualquiera}
\end{styleStandard}

\begin{styleStandard}
\ \ \ \ ‘What you are telling me is nonsense.’
\end{styleStandard}

\begin{styleStandard}
\ \ \ \ (@ 130)
\end{styleStandard}

\begin{styleStandard}
The difference in meaning between ‘nonsense’ and ‘ordinary’ or the pejorative version of the latter meaning as ‘bad’ is related to the syntactic and semantic type of the syntactic subject of the predicative \textit{cualquiera.} In case the subject is a sentence or a proposition p of semantic type {\textless}s,t{\textgreater} (i.e. the characteristic function of a set of possible worlds), \textit{cualquiera} means ‘nonsense’. In the following example, \textit{cualquiera} expresses evaluation with respect to what has been said, i.e. the argument of the verb of \textit{decir} ‘to say’. This argument is a proposition. In this example, \textit{cualquiera} is a bare noun, precisely the same one that occurs as a noun under verbs such as verbs of saying\textit{,} such as ArgSp\textit{ decir/mandar} ‘say nonsense’, that usually take propositions as their arguments (e.g. \textit{decir/mandar}) (see 4.5.2 for details). 
\end{styleStandard}

\begin{styleStandard}
\textbf{\ \ }\ \ $\lambda $p cualquiera (p) = ‘nonsense’
\end{styleStandard}

\begin{styleStandard}
\textit{decir/mandar cualquiera} = ‘say nonsense’
\end{styleStandard}

\begin{styleStandard}
One important issue that needs to be addressed by this analysis is the fact that in Spanish and Romance languages in general, bare singular nouns are rare in contrast to indefinite singular nouns with an indefinite article (e.g. \textit{decir una tontería} ‘say a stupid thing’). This restriction does not apply to plurals (e.g. \textit{decir tonterías }‘say stupid things’). I assume that the reason why \textit{cualquiera} appears as a bare noun rather than as an indefinite noun (*\textit{decir un/una cualquiera}) or a plural noun (*\textit{decir cualquieras}) is due to its pronominal origin. The pronoun \textit{cualquiera} ‘anything’ is invariable for plural or singular. Moreover, \textit{cualquiera} that takes a proposition as its argument is described as a kind property; that is, it describes the kind of things that one can say (see 4.5.2 for details). Recall that \textit{cualquiera} that predicates over propositions can also appear with adverbs that reinforce the evaluative property as ‘utter nonsense’ or ‘very nonsensical’. In this case the usual modification is a nominal modifier \textit{total} ‘complete’ as repeated here:
\end{styleStandard}

\begin{styleStandard}
\textbf{\ \ }\ \ y\ \ les\ \ pueden\ \ mandar\ \ cualquiera total\ \ 
\end{styleStandard}

\begin{styleStandard}
\ \ \ \ and\ \ that\ \ her\ \ sent\ \ \textsc{cualquiera total}
\end{styleStandard}

\begin{styleStandard}
\ \ \ \ ‘and you can tell them complete bullshit.’
\end{styleStandard}

\begin{styleStandard}
\ \ \ \ (@ 129b)
\end{styleStandard}

\begin{styleStandard}
However, there are also cases of \textit{cualquiera} with the meaning ‘nonsense’ and degree modification with (\textit{tan, re} ‘very’) when \textit{cualquiera} is used in a predicative position, as already shown in (231) and repeated here.
\end{styleStandard}

\begin{styleStandard}
\textbf{\ \ }\ \ Me parece todo tan cualquiera escribirlo casi me da vergüenza.
\end{styleStandard}

\begin{styleStandard}
‘Everything seems such utter nonsense to me writing it down. It’s almost embarrassing to me.’
\end{styleStandard}

\begin{styleStandard}
(@ 111)
\end{styleStandard}

\begin{styleStandard}
This shows that \textit{cualquiera} can change its syntactic status as an adjective or a bare noun depending on the syntactic context. This flexibility in word categorisation is expected under my semantic analysis, as \textit{cualquiera} denotes a property and whether this property is represented as a noun or as an adjective depends on the syntax.
\end{styleStandard}

\begin{styleStandard}
\ \ The suggested analysis can account for the difference between \textit{cualquiera} and \textit{cualunque} in ArgSp. The element \textit{cualunque} does not have the meaning of ‘nonsense’, as it cannot be used to predicate over propositions and thus be used under transitive verbs such as verbs of saying (see Section 4.2). The difference in meanings between ‘bad/ordinary’ and ‘nonsense’ is a further argument against deriving \textit{cualquiera} from the analogy of \textit{cualunque} (against Rizzo Salierno 2013). I believe that both items \textit{cualquiera} and \textit{cualunque} have developed the gradable property in a parallel way as a result of reanalysis of FCIs as (gradable) predicates. 
\end{styleStandard}

\subsection[4.5 From FC use of cualquiera to its evaluative meanings in ArgSp]{\textbf{4.5}\textbf{\textit{ }}\textbf{From}\textbf{\textit{ }}\textbf{FC use of }\textbf{\textit{cualquiera}}\textbf{ }\textbf{to its evaluative meanings in Arg}\textbf{Sp}}
\begin{styleStandard}
In this section, I will provide an analysis of the indefinite \textit{cualquiera} in ArgSp and show how the indefinite \textit{cualquiera} changed to the evaluative \textit{cualquiera} analysed in 4.5.1. I will suggest different diachronic paths of \textit{cualquiera} as ‘ordinary’ and as ‘nonsense’. This suggestion will be tested on diachronic data of ArgSp in 4.6.3.
\end{styleStandard}

\subsubsection[4.5.1 The change from the postnominal modifier to the adjective cualquiera in ArgSp]{\textbf{4.5.1 The change from the postnominal modifier to the adjective }\textbf{\textit{cualquiera}}\textbf{ in }\textbf{ArgSp}}
\begin{styleStandard}
I assume that the element \textit{cualquiera} with the meaning ‘ordinary’ or ‘bad’ (EVAL 1) is diachronically derived from \textit{cualquiera} in \textsc{un} N \textit{cualquiera} by lexicalising the evaluative property of ‘ordinary, bad’ and by being used as a predicative adjective without a noun:
\end{styleStandard}

\begin{styleStandard}
\textbf{\ \ }\ \ Juan es un hombre cualquiera. (EurSp and ArgSp)
\end{styleStandard}

\begin{styleStandard}
‘Juan is a man like any other man’ SIM
\end{styleStandard}

\begin{styleStandard}
‘Juan is an ordinary man’ COM
\end{styleStandard}

\begin{styleStandard}
STAGE 1= evaluative use of the indefinite \textit{cualquiera} 
\end{styleStandard}

\begin{styleStandard}
\textbf{\ \ }\ \ Juan es un hombre muy cualquiera. (ArgSp)
\end{styleStandard}

\begin{styleStandard}
‘Juan is a (very) ordinary man.’ Lexicalisation of COM and reanalysis into an adjective
\end{styleStandard}

\begin{styleStandard}
STAGE 2 A= reanalysis into an adjective that modifies a noun.
\end{styleStandard}

\begin{styleStandard}
\textbf{\ \ }\ \ Juan es (muy) cualquiera (ArgSp) 
\end{styleStandard}

\begin{styleStandard}
‘Juan is (very) ordinary.’ Separation from the noun
\end{styleStandard}

\begin{styleStandard}
STAGE 2 B= reanalysis into an adjective without nominal modification.
\end{styleStandard}

\begin{styleStandard}
The diachronic steps can be formalised as follows. \textit{Cualquiera} passed from being analysed as a property of individuals or kinds of individuals into a property of degrees that can be used with or without a noun:
\end{styleStandard}

\begin{styleTabellenInhalt}
\textbf{\ \ }\ \  $\lambda $\textit{P}$\lambda $\textit{x}\textit{\textsubscript{ }}[(cualquiera (\textit{P}))(\textit{x})]\ \ 
\end{styleTabellenInhalt}

\begin{styleTabellenInhalt}
‘property (of a kind of) individuals’ \ \ \ (Stage 1)
\end{styleTabellenInhalt}

\begin{styleTabellenInhalt}
\ \ \ \ $\lambda $\textit{d }$\lambda $\textit{P }$\lambda $\textit{x}[(cualquiera (P)(\textit{d}))(\textit{x})]
\end{styleTabellenInhalt}

\begin{styleTabellenInhalt}
‘gradable property of individuals realised as a nominal adjective’\ \ (Stage 2 A)
\end{styleTabellenInhalt}

\begin{styleTabellenInhalt}
\ \ \ \ $\lambda $\textit{d }$\lambda $\textit{x}[(cualquiera\textbf{ }(\textit{d}))(\textit{x})]
\end{styleTabellenInhalt}

\begin{styleTabellenInhalt}
‘gradable property of individuals realised as a predicative adjective’\ \  \ \ (Stage 2 B)
\end{styleTabellenInhalt}

\begin{styleStandard}
The empirical evidence of these diachronic steps will be discussed in 4.6.3 and in Chapter 5.
\end{styleStandard}

\begin{styleStandard}
I assume that the noun \textit{cualquiera} took a different diachronic path. 
\end{styleStandard}

\subsubsection[4.5.2 The change from a pronoun into a noun cualquiera ‘nonsense’ in ArgSp]{\textbf{4.5.2 }\textbf{The change from a pronoun into a noun }\textbf{\textit{cualquiera}}\textbf{ ‘nonsense’ in ArgSp}}
\begin{styleStandard}
My hypothesis is that \textit{cualquiera} under transitive verbs with the evaluative meaning ‘nonsense’ (EVAL 2) (see 4.4.1) has its origin in the indefinite pronoun \textit{cualquiera} with the quantificational force of ‘anything’. Recall from 2.1.4 that \textit{cualquiera} can be used as \textit{cualquier cosa} in ArgSp: 
\end{styleStandard}

\begin{styleStandard}
\textbf{\ \ }\ \ cualquiera\textsubscript{transitive} ‘any’ under ‘say’ = $\lambda $P $\lambda $Q [18E?]x[Q(\textit{x}) \& P(x)] ‘anything’ 
\end{styleStandard}

\begin{styleStandard}
FC interpretation = everything is a possible option.
\end{styleStandard}

\begin{styleStandard}
Recall from Section 2.1 that in European Spanish, inanimate \textit{cualquiera} requires mentioning the domain in the context preceding that in which \textit{cualquiera} refers. As the following example from European Spanish shows, \textit{cualquiera} stands for \textit{cualquier interpretación} (see Section 2.1).
\end{styleStandard}

\begin{styleStandard}
\textbf{\ \ }\ \ Aunque\ \ puedan\ \ admitir\ \ un\ \ alto\ \ número\ \ de\ \ interpretaciones\ \ (pero\ \ no\ \ admiten
\end{styleStandard}

\begin{styleStandard}
\ \ \ \ although\ \ can\ \ admit\ \ a\ \ high\ \ number\ \ of\ \ interpretations\ \ but\ \ not\ \ admit
\end{styleStandard}

\begin{styleStandard}
\ \ \ \ cualquiera).\ \ Hay\ \ también\ \ una\ \ tercera\ \ posibilidad.
\end{styleStandard}

\begin{styleStandard}
\ \ \ \ \textsc{cualquiera}\ \ there.is\ \ too\ \ a\ \ third\ \ possibility
\end{styleStandard}

\begin{styleStandard}
\ \ \ \ ‘Even though they can admit a high number of interpretations (not not just any). There is also a third possibility.’
\end{styleStandard}

\begin{styleStandard}
\ \ \ \ (@ 131; EurSp)
\end{styleStandard}

\begin{styleStandard}
In Argentinian Spanish as well, \textit{cualquiera} can be used to refer back to an item of the a preceding context. In the following examples, the preceding context has \textit{cualquier cosa}, and \textit{casi cualquiera} refers to \textit{cualquier cosa}. Both forms are the argument of the verb \textit{necesitar} ‘to need’ in \ and of \textit{ocurrir} ‘happen’. Note also that \textit{casi} usually modifies FCIs (see Chapter 3 for modification with \textit{casi} ‘almost’):
\end{styleStandard}

\begin{styleStandard}
\label{bkm:Ref150786782}\textbf{\ \ }\ \ \textit{Cualquier\ \ cosa\ \ }que\ \ necesiten\ \ (bueno,\ \ \textit{casi\ \ cualquiera!})\textit{\ \ }nos\ \ avisan!
\end{styleStandard}

\begin{styleStandard}
\ \ \ \ any\ \ thing\ \ that\ \ need\ \ well\ \ almost\ \ anything\ \ us\ \ inform
\end{styleStandard}

\begin{styleStandard}
\ \ \ \ ‘If you need anything (well, almost anything!) tell us!!’
\end{styleStandard}

\begin{styleStandard}
\ \ \ \ (@ 132)
\end{styleStandard}

\begin{styleStandard}
\label{bkm:Ref150786875}\textbf{\ \ }\ \ Hoy\ \ en\ \ día,\ \ uno\ \ se\ \ imagina\ \ que\ \ podia\ \ estar\ \ ocurriendo\ \ \textit{casi\ \ cualquier}
\end{styleStandard}

\begin{styleStandard}
\ \ \ \ today\ \ in\ \ day\ \ one\ \ her/him\ \ imagines\ \ that\ \ could\ \ be\ \ happening\ \ almost\ \ any
\end{styleStandard}

\begin{styleStandard}
\ \ \ \ \textit{cosa.\ \ Casi\ \ cualquiera\ \ }menos,\ \ probablemente,\ \ lo\ \ que\ \ ocurría.
\end{styleStandard}

\begin{styleStandard}
\ \ \ \ thing\ \ almost\ \ anything\ \ minus\ \ probably\ \ the\ \ that\ \ happened
\end{styleStandard}

\begin{styleStandard}
\ \ \ \ ‘Nowadays, one imagines that anything could be happening. Anything but what was actually happening, probably.’
\end{styleStandard}

\begin{styleStandard}
\ \ \ \ (@ 133)
\end{styleStandard}

\begin{styleStandard}
As already noted in Section 2.1.4, \textit{cualquiera} in ArgSp, but not in EurSp, can be used as \textit{cualquier cosa} without previous mention of \textit{cualquier cosa}. \textit{Cualquiera} has lexicalised the meaning of \textit{cualquier cosa} in ArgSp and does not require previous mention of \textit{cualquier cosa}:
\end{styleStandard}

\begin{styleStandard}
\label{bkm:Ref150786811}\textbf{\ \ }\ \ Respondí/Entendí/Hice/Dijo [\textsubscript{DP} \textit{cualquiera}] 
\end{styleStandard}

\begin{styleStandard}
\ \ \ \ = ‘I answered/understood/did/said anything.’
\end{styleStandard}

\begin{styleStandard}
There is a pragmatic reason why \textit{cualquiera} as a verbal argument has the meaning of \textit{cualquier cosa} and not a more specific meaning, such as \textit{cualquier libro} ‘any book’. \textit{Cualquiera} in object position without any previous discourse can only mean \textit{cualquier cosa}, because \textit{cualquier cosa} is more generic than \textit{cualquier libro}. A similar case can be demonstrated with elliptical nouns that appear without previous mention. For example, \textit{me gusta comer} ‘I like eating’ is interpreted as ‘I like eating food or edible things’ and not as ‘I like eating a specific thing, such as chocolate.’
\end{styleStandard}

\begin{styleStandard}
\ \ There are two questions that appear with respect to . The first is what licenses the use of FCI \textit{cualquiera},\textit{ }or\textit{ cualquier cosa }for that matter, under perfective verbs in , given that the licensing of FCIs happens usually under modal operators rather than under perfective verbs (see Chapter 2). The second question is how to account for the evaluative reading of \textit{cualquiera} under transitive verbs, e.g. \textit{decir}/\textit{mandar cualquiera} ‘to say nonsense’ or \textit{hacer cualquiera} ‘to do something out of place/nonsense’ (see 4.3.3).
\end{styleStandard}

\begin{styleStandard}
\ \ To address the first question, I assume that FCI \textit{cualquiera} in episodic contexts such as under transitive perfective verbs (see e.g. \ in Section 2.3 repeated in ) can be licensed under subtrigging where the modal is contained inside a covert relative clause, e.g. \textit{que podría responder} ‘that I could possibly answer’ (see Section 2.2 for subtrigging and the formal representation of it):
\end{styleStandard}

\begin{styleStandard}
\label{bkm:Ref63461360}\textbf{\ \ }\ \ Respondí cualquiera\ \ [que podría decir/responder]
\end{styleStandard}

\begin{styleStandard}
\ \ \ \ ‘I answered anything (I could answer/say).’
\end{styleStandard}

\begin{styleStandard}
Another way in which the ArgSp pronominal FCI \textit{cualquiera} ‘anything’ can be licensed in episodic contexts is to assume a covert modal inside the matrix clause as I represent informally in , e.g. agent’s decision based modal that triggers the random choice interpretation (see Section 2.3.2) or a counterfactual modal (see 2.3.1):
\end{styleStandard}

\begin{styleStandard}
\label{bkm:Ref63461220}\textbf{\ \ }\ \ Respondí cualquiera
\end{styleStandard}

\begin{styleStandard}
\ \ \ \ ‘I answered something random.’
\end{styleStandard}

\begin{styleStandard}
\ \ a.\ \ ‘I said something and this something fulfills my decision to say any possible thing.’ 
\end{styleStandard}

\begin{styleStandard}
(Agent’s decision-based modality)
\end{styleStandard}

\begin{styleStandard}
\ \ b.\ \ ‘The thing I actually said could have been anything else than what I actually said.’ (Counterfactual modal) 
\end{styleStandard}

\begin{styleStandard}
\ \ \ \ (@ 33)
\end{styleStandard}

\begin{styleStandard}
The two licensing conditions (modal operator inside a relative clause, i.e. subtrigging, or inside the matrix clause as random choice modality) create two different truth conditions. Under subtrigging in , \textit{cualquiera} refers to many speech events in the actual world (I said a, b, c, d….) and under random choice in , \textit{cualquiera} refers to a single speech event in the actual world and the actual speech event is contrasted to other speech events that the speaker could have said or that agree with agent’s decision. There is evidence that we need both analyses, i.e. licensing of \textit{cualquiera} under subtrigging and under some modal that triggers random choice. In cases where \textit{cualquiera} can be modified with adverbs like \textit{casi} ‘almost’ or with \textit{except}{}-phrases (see Section 2.1), \textit{cualquiera} is licensed under subtrigging. But when \textit{cualquiera} has an existential inference in a particular context, it is licensed under a modal that triggers random choice. I will first derive the evaluative meaning of \textit{cualquiera} from FCIs under subtrigging, as shown in , and then from random choice modality, as shown in .
\end{styleStandard}

\begin{styleStandard}
\ \ In order to answer the question as to how to account for the evaluative reading under transitive verbs, I assume that the evaluative reading ‘nonsense’ is a pragmatic implicature on \textit{cualquiera} or \textit{cualquier cosa} under transitive verbs that has been lexicalised under certain verbs such as \textit{mandar cualquiera} ‘say nonsense’ in ArgSp. From the sentence in , the hearer can infer that what the speaker said can be of the form p and not(p) at the same time. To utter a contradiction can be evaluated as nonsense. This is the most direct way to account for the meaning ‘nonsense’ of \textit{cualquiera}. Another, indirect, way to account for the meaning ‘nonsense’ is to derive ‘nonsense’ from the random choice meaning of \textit{cualquiera}; that is ‘something random’ can be interpreted as ‘nonsense’ in some pragmatic context. Namely when the agent who said something random did not reflect on what she/he said according to some judge. The result of the lack of reflection can be interpreted as something meaningless or as nonsense according the judge, such as in a situation when the agent was half-asleep. Thus, in this context the evaluator/judge can infer that the agent said something meaningless. Under this explanation, \textit{to say} \textit{nonsense} would be derived from a contextual interpretation or pragmatic strengthening of the meaning of \textit{decir/mandar cualquiera} ‘to say something random’. 
\end{styleStandard}

\begin{styleStandard}
\ \ Under both analyses, the evaluative meaning of \textit{cualquiera} as ‘nonsense’ is a result of a pragmatic reasoning triggered by the use of FCIs under verbs of saying (see Aloni 2005 for indefinites in general, Author 2020a on pragmatic reasoning in Mexican Spanish). 
\end{styleStandard}

\begin{styleStandard}
\ \ I find contexts that support my pragmatic analysis. The following example expresses that what the speaker actually answered does not correspond to what the speaker originally wanted to answer. In other words, the speaker’s intentions do not correspond to his/her actions/speech events. This is why \textit{cualquiera} has the meaning of ‘something unreflected’ in the following example:
\end{styleStandard}

\begin{styleStandard}
\textbf{\ \ }\ \ No\ \ supe\ \ qué\ \ decirle,\ \ Así\ \ que\ \ le\ \ mandé\ \ cualquiera.
\end{styleStandard}

\begin{styleStandard}
\ \ \ \ not\ \ knew\ \ what\ \ say.her\ \ So\ \ that\ \ her\ \ sent\ \ \textsc{cualquiera}
\end{styleStandard}

\begin{styleStandard}
\ \ \ \ ‘I did not know what to tell her, so I told her nonsense.’
\end{styleStandard}

\begin{styleStandard}
\ \ \ \ (@ 134)
\end{styleStandard}

\begin{styleStandard}
Another example, where the contextual interpretation restricts the interpretation of \textit{cualquiera} is given in the following sentence. Here the speaker has done everything except what she aimed to do, namely to fix the heating system. The exception is expressed through the contrast \textit{pero} ‘but’:
\end{styleStandard}

\begin{styleStandard}
\textbf{\ \ }\ \ Intenté\ \ arreglar\ \ la\ \ calefacción,\ \ pero\ \ hice\ \ cualquiera.
\end{styleStandard}

\begin{styleStandard}
\ \ \ \ tried\ \ fixing\ \ the\ \ heating.system\ \ but\ \ did\ \ \textsc{cualquiera}
\end{styleStandard}

\begin{styleStandard}
\ \ \ \ ‘I wanted to fix the heating system, but I did something out of place/nonsense (as a result to my attempt to fix the heating system).’
\end{styleStandard}

\begin{styleStandard}
\ \ \ \ (@ 135)
\end{styleStandard}

\begin{styleStandard}
In the following example \textit{decir} \textit{cualquiera} is interpreted negatively, similar to \textit{mandar la cagada} ‘to say bullshit’. It is interpreted as something unreflected or unreasonable. 
\end{styleStandard}

\begin{styleStandard}
\textbf{\ \ }\ \ Che,\ \ ya\ \ te\ \ mandaste\ \ la\ \ cagada,\ \ dijiste\ \ cualquiera,\ \ tiraste\ \ abajo\ \ a
\end{styleStandard}

\begin{styleStandard}
\ \ \ \ that\ \ ya\ \ you\ \ said\ \ the\ \ bullshit\ \ said\ \ \textsc{cualquiera}\ \ pulled\ \ down\ \ to
\end{styleStandard}

\begin{styleStandard}
\ \ \ \ la\ \ mitad\ \ de\ \ los\ \ hombres\ \ diciendo\ \ que\ \ tienen\ \ que\ \ ser\ \ capaces
\end{styleStandard}

\begin{styleStandard}
\ \ \ \ the\ \ half\ \ of\ \ the\ \ men\ \ saying\ \ that\ \ must\ \ that\ \ be\ \ capable
\end{styleStandard}

\begin{styleStandard}
\ \ \ \ de\ \ embarazar\ \ una \ \ mujer.
\end{styleStandard}

\begin{styleStandard}
\ \ \ \ of\ \ impregnate\ \ a\ \ woman
\end{styleStandard}

\begin{styleStandard}
\ \ \ \ ‘Hey, you said bullshit, you said nonsense, you put half of the men very low by saying that they have to be capable of making a woman pregnant.’
\end{styleStandard}

\begin{styleStandard}
\ \ \ \ (@ 136)
\end{styleStandard}

\begin{styleStandard}
The inclusion of morally questionable alternatives like saying morally questionable things into the quantification domain of \textit{cualquiera} can result in a negative interpretation of \textit{decir} \textit{cualquiera} as saying all possible things including morally bad things:
\end{styleStandard}

\begin{styleStandard}
\textbf{\ \ }\ \ \textit{decir cualquiera} … ‘say anything possible including morally bad things’
\end{styleStandard}

\begin{styleStandard}
This negative meaning of \textit{cualquiera} is thus tantamount to FCI \textit{cualquiera} plus the pragmatic characterisation of the quantificational domain. The negative meaning of \textit{cualquiera }is the result of excluding certain alternatives from the domain of quantification which some evaluator considers as being good alternatives or of characterising the domain of quantification as consisting of all possible alternatives including morally questionable alternatives. I will therefore look into the domain quantification of \textit{cualquiera} in more detail in the next subsection. 
\end{styleStandard}

\subsubsection[4.5.3 Formal analysis of cualquiera as ‘nonsense’]{\textbf{4.5.3 Formal analysis of }\textbf{\textit{cualquiera}}\textbf{ as ‘nonsense’}}
\begin{styleStandard}
I assume that the domain of quantification of ArgSp \textit{cualquiera} can be formally represented as a restriction of the quantification by an \textit{except}{}-phrase as in . Assuming that \textit{except}{}-phrases contribute to the truth conditions of the sentence by restricting the domain of quantification of \textit{cualquiera} to the domain of predicates besides the predicate expressed by the \textit{except}{}-phrase (von Fintel 1994),\footnote{\textrm{\ The assumption that }\textrm{\textit{except}}\textrm{{}-phrases contribute to the truth conditions of the sentence they modify follows from the Restrictiveness Condition (von Fintel 1994): If a sentence with an }\textrm{\textit{except}}\textrm{{}-phrase is true, the same sentence without the }\textrm{\textit{except}}\textrm{{}-phrase is false (without making the exception, the sentence would be false).}} I must assume that \textit{cualquiera} in \ triggers universal quantification, which contributes to the truth conditions of the sentence it appears in:
\end{styleStandard}

\begin{styleStandard}
\label{bkm:Ref63461516}\textbf{\ \ }\ \ ‘I answered \textit{cualquiera} except what I wanted’
\end{styleStandard}

\begin{styleStandard}
every([[\textit{thing}]] – [[\textit{what I wanted to answer}]])(\{x: I answered x\})
\end{styleStandard}

\begin{styleStandard}
\textit{Cualquiera }must be interpreted universally at some level of interpretation, because existential quantification does not allow \textit{except}{}-phrases (see von Fintel 1994, Arregui 2006, Gajewski 2008):
\end{styleStandard}

\begin{styleStandard}
\textbf{\ \ }\#\ \ I answered something except what I wanted.
\end{styleStandard}

\begin{styleStandard}
\textit{Except}{}-phrases of FCIs are a good argument that the universal inference of FCIs must be derived from a wider scope pragmatic operation or on the sentence level and not locally (see Section 2.2 on such an analysis, where the universal interpretation is represented on the propositional level; see also Chierchia 2006, Gajewski 2008, but also Dayal 1998, 2004, who argues that FCIs are in fact universal quantifiers, with the result that \textit{except}{}-phrases operate locally). \textit{Except}{}-phrases should thus operate on wide-scope universal implicature, that is, after universal implicature of FCI has been derived on the sentence or pragmatic level (Gajewski 2008):\footnote{\textrm{\ The focus operator LEAST triggers exhaustivity of the alternative denoted by the exceptive; that is, it triggers the interpretation that the quantification domain of anything includes all alternatives other than the one alternative subtracted from the domain of quantification.}}
\end{styleStandard}

\begin{styleStandard}
\textbf{\ \ }\ \ Bill may have any flavour of ice cream but coffee.
\end{styleStandard}

\begin{styleStandard}
\ \ \ \ LF: [LEAST coffee]\textsubscript{X} Op\textsubscript{FC} [may [Bill have any flavour but X]]
\end{styleStandard}

\begin{styleStandard}
\ \ \ \ (Gajewski 2008, Gajewski 2013: 193)
\end{styleStandard}

\begin{styleStandard}
I assume that in case we have a pragmatically inferred \textit{except}{}-phrase as in \ , repeated from , \textit{cualquiera} is licensed by covert subtrigging (see Quer 2000 and 2.1 for overt subtrigging):
\end{styleStandard}

\begin{styleStandard}
\label{bkm:Ref63461661}\textbf{\ \ }a.\ \ Respondí cualquiera.
\end{styleStandard}

\begin{styleStandard}
‘I answered anything that I could possibly answer except what I wanted to answer.’
\end{styleStandard}

\begin{styleStandard}
\ \ b.\ \ LF: [what I wanted to answer]\textsubscript{X} Op\textsubscript{FC} [I could possibly answer [I answered anything but X]] 
\end{styleStandard}

\begin{styleStandard}
\ \ c.\ \ Prose: For every thing that I could possibly answer, I answered it, except what I wanted to answer.
\end{styleStandard}

\begin{styleStandard}
This analysis explains the universal interpretation of \textit{cualquiera} on the one hand and the licensing of \textit{cualquiera} as FCI via subtrigging on the other.
\end{styleStandard}

\begin{styleStandard}
\ \ Let me just briefly sketch an alternative analysis of \textit{cualquiera} in episodic contexts that relies on random choice modality and not on subtrigging. The analysis is again very similar with the difference that the modal operator that licenses \textit{cualquiera} is some kind of a modal, such as a counterfactual modal or a decision-based modal (see Section 2.3):
\end{styleStandard}

\begin{styleStandard}
\label{bkm:Ref150787301}\textbf{\ \ }a.\ \ Respondí cualquiera. 
\end{styleStandard}

\begin{styleStandard}
‘I answered something random.’
\end{styleStandard}

\begin{styleStandard}
\ \ b.\ \ Existential component
\end{styleStandard}

\begin{styleStandard}
I said something.
\end{styleStandard}

\begin{styleStandard}
\ \ c.\ \ Counterfactual Modal component: And I could have said anything other than what I actually said.
\end{styleStandard}

\begin{styleStandard}
Decision-Based Modal component: And anything else that I did not actually say matches my decision to say something.
\end{styleStandard}

\begin{styleStandard}
The question now is how the evaluative characterisation of \textit{cualquiera} as ‘nonsense’ comes into play. Recall that this is the second challenge that I mentioned at the beginning of 4.5.2. 
\end{styleStandard}

\begin{styleStandard}
\ \ As already described in 4.5.2, I suggested that one could distinguish between contextual and lexicalised evaluative meaning. The contextual evaluative meaning of \textit{cualquiera} is better described as the evaluation of some judge in a given context. This means that a) implies an attitude holder who corresponds to the speaker of the utterance who evaluates the event in which some agent has said anything the agent could possibly answer other than what the agent wanted to answer (henceforth: event\textsubscript{said(agent)(anythingy possible)}). The attitude holder believes that this event is less preferred than a different event e[2032?] in which the agent said something else (henceforth: event\textsubscript{said(agent)(not(anything possible))}). This belief state of the attitude holder introduces an ordering scale of worlds that can be informally described as follows. All the worlds in which an agent said anything possible are less preferred than all worlds in which the agent did not say anything possible but said something specific/distinguished. I represent this ordering scale in line with Stalnaker (1984) and Heim’s (1992) doxastic modal \textit{want}, which orders worlds according to the attitude holder’s preferences. Stalnaker (1984) and Heim (1992) suggest that \textit{want} is treated as a universal quantifier over worlds. “X wants S” is interpreted as “X believes that a world w in which S is true is more desirable than a world in which S is false.” Thus, \textit{want} is interpreted with respect to a doxastic modal base: 
\end{styleStandard}

\begin{styleStandard}
\textbf{\ \ }\ \ [27E6?]want[27E7?]\textsuperscript{g,w} = $\lambda $p ${\in}$ D {\textless}s,t{\textgreater} . $\lambda $x ${\in}$ De. Dox x,w ${\subseteq}$ \{w[2032?]: any world w[2034?] maximally similar to w[2032?] such that w[2033?] ${\in}$ p is more desirable to x in w than any world w[2034?] maximally similar to w[2032?] such that w[2034?] ${\notin}$ p\}\ \ 
\end{styleStandard}

\begin{styleStandard}
\ \ \ \ (Heim 1992: 193)
\end{styleStandard}

\begin{styleStandard}
I follow Heim’s analysis of attitude predicates such as \textit{want} and apply it to \textit{cualquiera}. On this analysis \textit{cualquiera} takes a proposition p, an attitude holder x, and a doxastic modal corresponding to a set of worlds that are ordered according to x’s preferences in the actual world w in such a way as to represent p’s worlds being less preferred than non-p’s worlds:
\end{styleStandard}

\begin{styleStandard}
\label{bkm:Ref150787526}\textbf{\ \ }\ \ [27E6?]cualquiera[27E7?]\textsuperscript{g,w} = $\lambda $p ${\in}$ D {\textless}s,t{\textgreater}. $\lambda $x ${\in}$ De. Dox {\textless}x,w{\textgreater} ${\subseteq}$ \{w[2032?]: any world that represent the negation of p is more preferred than all worlds that represents p, i.e. the proposition that is the argument of \textit{cualquiera}\}\ \ 
\end{styleStandard}

\begin{styleStandard}
\ \ \ \ (my adaptation of Heim’s 1992: 193 doxastic modality to \textit{cualquiera})
\end{styleStandard}

\begin{styleStandard}
If we apply this denotation to , the ordered worlds correspond to any world in which the agent said anything possible is less preferred in w according to the attitude holder than the worlds in which the agent did not say anything possible, that is, in which the agent said something distinguished or particular.
\end{styleStandard}

\begin{styleStandard}
\ \ The evaluative meaning component of \textit{cualquiera} thus corresponds to attitude holder’s preferences in the actual world.
\end{styleStandard}

\begin{styleStandard}
\ \ This analysis also captures the observation that the attitude holder does not have to be the speaker as shown in the following example (see also 4.3.2), by the fact that the attitude holder is represented as a variable in :
\end{styleStandard}

\begin{styleStandard}
\textbf{\ \ }\ \ haría lo mismo pero si hay gente que no le gusta porque piensa que es cualquiera (yo no pienso eso).....
\end{styleStandard}

\begin{styleStandard}
\ \ \ \ ‘I would do the same but there are people that don’t like it because they think it is \textit{nonsense} (I don’t think so)….’
\end{styleStandard}

\begin{styleStandard}
\ \ \ \ (@ 110)
\end{styleStandard}

\begin{styleStandard}
I assume that the lexicalisation of \textit{cualquiera} under \textit{mandar} can be represented as a type-shifting in diachrony. Free Choice indefinite \textit{cualquiera} ‘anything’ of type {\textless}{\textless}e,t{\textgreater}t{\textgreater} shifts into a property of type {\textless}e,t{\textgreater}, which corresponds to a noun that is interpreted as a collective thing that comprises all things that one could possibly say including contradictions. When used as an argument of \textit{mandar}, this collective thing can be paraphrased as ‘nonsense’.
\end{styleStandard}


\setcounter{listWWNumxleveli}{0}
\begin{listWWNumxleveli}
\item 
\begin{styleListParagraph}
\textmd{[D° {\textless}{\textless}e,t{\textgreater}t{\textgreater} }\textmd{\textit{cualquiera}}\textmd{ ‘anything’] [F0E0?] type-shifting [N° {\textless}et{\textgreater} ‘a collective thing that comprises all things that one could possibly say including contradictions’}
\end{styleListParagraph}
\end{listWWNumxleveli}
\begin{styleStandard}
The analysis of \textit{cualquiera} (EVAL 2) is thus the same as the one suggested in Section 3.5 for \textit{cualquiera} in European Spanish (EVAL 1), with the only difference that \textit{cualquiera} under \textit{mandar} in Argentinian Spanish is a property of things that denote propositions and not a property of kinds of people or entities as in EVAL 1. \ 
\end{styleStandard}

\subsubsection[4.5.4 Summary of my analysis and discussion]{\textbf{4.5.4 }\textbf{Summary of my analysis and discussion}}
\begin{styleStandard}
The origin of the transitive \textit{cualquiera }as an FCI is analysed\textit{ }in parallel with\textit{ cualquier cosa},\textit{ }which is interpreted in certain pragmatic contexts as ‘anything possible except …’. I can now suggest the following diachronic steps that have led to the development of the evaluative meaning of \textit{cualquiera} under transitive verbs:
\end{styleStandard}

\begin{styleStandard}
\textbf{\ \ }Diachronic steps of transitive\textit{ cualquiera}
\end{styleStandard}

\begin{styleStandard}
\ \ Step 1: \ \ Use of \textit{cualquiera} as ‘any’ to refer to a domain mentioned in the previous context in EurSp and ArgSp. A special case of this use is the use of \textit{cualquiera} to refer back to previously mentioned \textit{cualquier cosa}.
\end{styleStandard}

\begin{styleStandard}
\ \ Step 2:\ \ Use of \textit{cualquiera} in the sense of \textit{cualquier cosa} without previous mentioning of \textit{cualquier cosa} in ArgSp.
\end{styleStandard}

\begin{styleStandard}
\ \ Step 3:\ \ Use of \textit{cualquiera} as \textit{cualquier cosa} in episodic contexts: 
\end{styleStandard}

\begin{styleStandard}
[\textsubscript{VP} \textit{decir/mandar}[episodic] [\textsubscript{DP} \textit{cualquiera=cualquier cosa}]] [\textsubscript{Adjunct NP} \textit{except} something specific or something the agent wanted to say]]]
\end{styleStandard}

\begin{styleStandard}
\ \ Step 4:\ \ Evaluation of the thing or things that have been said according to preferences of an attitude holder: i.e. to say something randomly implies to say something without thinking are less preferred according to an attitude holder than to say something distinguished or reflected.
\end{styleStandard}

\begin{styleStandard}
\ \ Step 5:\ \ Interpretation of \textit{cualquiera} as ‘nonsense’. 
\end{styleStandard}

\begin{styleStandard}
\ \ Step 6:\ \ Lexicalisation of the pragmatic evaluation as ‘nonsense’ under \textit{mandar. }
\end{styleStandard}

\begin{styleStandard}
Syntactic shift from indefinite quantifier into a noun\textit{.}
\end{styleStandard}

\begin{styleStandard}
As the diachronic development of \textit{cualquiera} under transitive verbs shows, the step towards lexicalisation of the evaluative meaning is the pragmatic interpretation of \textit{cualquiera} as something unreflected or something that can include contradictions according to some evaluator or judge. In contrast to Rizzo Salierno (2013), I do not assume that the transitive \textit{cualquiera} has its origin in an analogy to the element \textit{cualunque}, because \textit{cualquiera} does not share the same distribution and frequency with \textit{cualunque}, but with \textit{cualquier cosa}. There is no occurrence with \textit{cualunque} under transitive verbs. However, we do find \textit{cualquier cosa} under transitive verbs.
\end{styleStandard}

\begin{styleStandard}
\ \ The reanalysis of \textit{cualquiera} as an adjective in predicative position was already discussed in Chapter 3 for European Spanish. The difference from European Spanish is that ArgSp allows degree modification of \textit{cualquiera} and the possibility of using \textit{cualquiera} as a predicative adjective without nominal modification, as shown in Step 2:
\end{styleStandard}

\begin{flushleft}
\tablefirsthead{}
\tablehead{}
\tabletail{}
\tablelasttail{}
\begin{supertabular}{m{0.42885986in}|m{0.41155985in}|m{5.22336in}|}
\hhline{~--}
\liststyleWWNumx
\setcounter{saveenum}{\value{enumi}}
\begin{enumerate}
\setcounter{enumi}{\value{saveenum}}
\item 
\end{enumerate}
 &
\multicolumn{2}{m{5.7136602in}|}{\centering \textbf{Diachronic Steps of Predicative}\textbf{\textit{ Cualquiera}}}\\\hhline{~--}
 &
Step 1: &
Use of \textit{cualquiera} as ‘any’ in postnominal position with the evaluative meaning:

\textit{\ Juan es un hombre cualquiera.} (EurSp and ArgSp)

‘Juan is a man like any other man’ SIM

‘Juan is an ordinary man’ COM\\\hhline{~--}
 &
Step 2: &
Lexicalisation of COM in Argentinian Spanish and reanalysis as an adjective in +/– postnominal position

\textit{Juan es (un hombre) muy cualquiera.} (ArgSp)

‘Juan is a very ordinary man/Juan is very ordinary.’ COM\\\hhline{~--}
\end{supertabular}
\end{flushleft}
\begin{styleStandard}
Chapter 6 will offer a detailed diachronic investigation of the evaluative meaning (EVAL 1) in diachrony. The diachronic analysis of EVAL 2 will be tested in Section 4.6.3. 
\end{styleStandard}

\subsection[4.6 Further evidence for the analysis]{\textbf{4.6 Further evidence for the analysis}}
\begin{styleStandard}
In this section, further evidence for the analysis will be presented on the basis of synchronic and diachronic data.
\end{styleStandard}

\subsubsection[4.6.1 Evidence from other languages]{\textbf{4.6.1 Evidence from other languages}}
\begin{styleStandard}
We find a similar reinterpretation of the free choice indefinite \textit{n’importe quoi} in French. As shown in , the indefinite \textit{n’importe quoi} has different interpretations. It can have a random choice interpretation, as in a), or an evaluative interpretation, as in b) (see Larrivée 2007, Paillard 1998, Pescarini 2009, among others for the ambiguity):
\end{styleStandard}

\begin{styleStandard}
\label{bkm:Ref62641804}\textbf{\ \ }French
\end{styleStandard}

\begin{styleStandard}
\ \ a.\ \ \textit{J’ai dit n’importe quoi (qui m’est passé par la la tête, et ça s’est trouvé à être la bonne réponse !)} 
\end{styleStandard}

\begin{styleStandard}
I have said \textsc{indefinite}
\end{styleStandard}

\begin{styleStandard}
‘I said something random (what came to my mind (first), and it happened to be the right answer).’
\end{styleStandard}

\begin{styleStandard}
\ \ \ \ (Larrivée 2007)
\end{styleStandard}

\begin{styleStandard}
\ \ b.\ \ \textit{Tu dis n’importe quoi. Tu sais pas ce que tu dis.}
\end{styleStandard}

\begin{styleStandard}
‘You’re talking nonsense. You don’t know what you’re saying.’
\end{styleStandard}

\begin{styleStandard}
‘You’re talking bullshit/nonsense.’
\end{styleStandard}

\begin{styleStandard}
\ \ \ \ (@ 137)
\end{styleStandard}

\begin{styleStandard}
\ \ c.\ \ \textit{Tu dis du grand n’importe quoi.}
\end{styleStandard}

\begin{styleStandard}
‘You’re talking bullshit/nonsense.’
\end{styleStandard}

\begin{styleStandard}
\ \ \ \ (@ 138)
\end{styleStandard}

\begin{styleStandard}
If \textit{n’importe quoi} is used with a definite article and adjectival modification as in c), it is no longer ambiguous. It is then used as a noun with the evaluative interpretation of ‘nonsense’. \textit{N’importe quoi} shows similarities to \textit{cualquiera} under transitive verbs, as it can also be used as a noun with an evaluative interpretation. The difference between \textit{cualquiera} and \textit{n’importe quoi} is that \textit{n’importe quoi} cannot appear as a gradual adjective in a predicative position, whereas \textit{cualquiera} can as shown in 4.3.2 (\textit{me parece re cualquiera ‘}it seems a complete bullshit to me’):
\end{styleStandard}

\begin{listWWNumxleveli}
\item 
\begin{styleListParagraph}
\textmd{\ \ \ \ \ \ \ *C’est très n’importe quoi. \ \ \ \  \ \ \ \ \ \ \ \ \ \ \ \ \ \ \ \ \ \ \ \ \ \ \ \ \ \ \ \ \ \ \ (French)}
\end{styleListParagraph}
\end{listWWNumxleveli}
\begin{styleListParagraph}
\textmd{‘This is very bullshit.’}
\end{styleListParagraph}

\begin{styleStandard}
In contrast to \textit{n’importe quoi,} French FCI \textit{quelconque} can be modified with degree adverbs (see Vlachou 2007). It has the meaning of ‘ordinary’:
\end{styleStandard}

\begin{listWWNumxleveli}
\item 
\begin{styleListParagraph}
\textmd{\ \ \ \ \ \ \ \ Ce livre est assez quelconque. \ \ \ \ \ \ \ \ \ \ \ \ \ \ \ \ \ \ \ \ \ \ \ \ \ \ \ \ \ \ \ \ \ \ \ \ \ \ \ \ \ \ \ \ \ \ \ \ \ \ \ \ \ \ \ \ \ \ \ \ \ \ \ \ \ \ \ \ \ \ \ \ \ \ \ \ \ \ \ \ \ \ (French)}
\end{styleListParagraph}
\end{listWWNumxleveli}
\begin{styleListParagraph}
\textmd{‘This book is rather ordinary.’}
\end{styleListParagraph}

\begin{styleListParagraph}
\textmd{(@ 139)}
\end{styleListParagraph}

\begin{styleStandard}
The analysis of \textit{cualquiera} as degree adjective applies to \textit{quelconque} as well.
\end{styleStandard}

\begin{styleStandard}
\ \ A similar pattern was observed in Mexican Spanish for the lexical item \textit{equis} (originally meaning ‘the letter x’), which can be used as an indefinite with a free choice interpretation ‘any’ under conditionals as in , but can also function as a gradable adjective that expresses EVAL1, as shown in \ (see Author 2020a):
\end{styleStandard}

\clearpage\begin{styleStandard}
\label{bkm:Ref63462361}\textbf{\ \ }\ \ Si por equis razón no vas a venir, avisame\ \ (MexSp)
\end{styleStandard}

\begin{styleStandard}
\ \ \ \ ‘If you cannot come for some reason, let me inow.’\ \ 
\end{styleStandard}

\begin{styleStandard}
\ \ \ \ ‘For whatever reason you cannot come, let me now.’\ \ 
\end{styleStandard}

\begin{styleStandard}
\ \ \ \ (Author 2020a: 137)
\end{styleStandard}

\clearpage\begin{styleStandard}
\label{bkm:Ref63462370}\textbf{\ \ }\ \ un hombre bastante/demasiado/muy equis\ \ (MexSp)
\end{styleStandard}

\begin{styleStandard}
\ \ \ \ ‘a very ordinary man (with no special qualities)’\ \ 
\end{styleStandard}

\begin{styleStandard}
\ \ \ \ (Author 2020a: 138)
\end{styleStandard}

\begin{styleStandard}
More research needs to be done in order to testify the change of indefinites into evaluative nouns or adjectives in other languages (see Šimík 2008 for Czech). In the future, I will also compare my results with more recent studies on the Mayan language Chuj (spoken in Guatemala and Mexico) that shows interesting parallels with the contact induced meaning variation of \textit{komon} ‘common, ordinary’ (Alonso-Ovalle \& Royer 2023).
\end{styleStandard}

\begin{styleStandard}
\ \ Another topic that needs further investigation is the licensing of indefinites in episodic contexts without overt subtrigging such as ArgSp \textit{respondí cualquiera} ‘I answered anything possible’ in other languages. Dayal (1998) shows that \textit{any} can be licensed under perfective verbs without overt subtrigging in some English varieties (see Dayal 1998 for English varieties):
\end{styleStandard}

\begin{styleStandard}
\textbf{\ \ }\ \ a. Mary confidently answered any objections.
\end{styleStandard}

\begin{styleStandard}
b. After the dinner, I threw away any leftovers. 
\end{styleStandard}

\begin{styleStandard}
\ \ \ \ (Dayal 1998: 446)
\end{styleStandard}

\begin{styleStandard}
In contrast to English, which rarely allows FCI \textit{any} in episodic contexts, Romance indefinites such as \textit{cualquiera} do allow FCIs in episodic contexts (see 2.3). One possible explanation why \textit{cualquiera} appears more often in episodic contexts than English \textit{any} is that \textit{cualquiera} lexically encodes modality due to its relative clause origin from free choice relative \textit{qual quier} ‘whatever you want’ (see Rivero 2011 and Section 2.2). This explanation also holds for French as \textit{n’importe quoi} lexically encodes a modal ‘it is not important what’. It seems that English \textit{any} does not encode modality lexically and thus does not allow \textit{any} to appear in episodic contexts as in Spanish and French. 
\end{styleStandard}

\subsubsection[4.6.2 Evidence from comparison with cualquier cosa]{\textbf{4.6.2 Evidence from comparison with }\textbf{\textit{cualquier cosa}}}
\begin{styleStandard}
My analysis of ArgSp \textit{cualquiera} is based on the FCI \textit{cualquiera}. It is thus possible that \textit{cualquier cosa} has undergone a similar change. I will therefore compare the two items and their uses. All data is taken from CdE\_web dialects. As I will show \textit{cualquier cosa} is on its way to reanalysis as a Degree Adjective ‘bad/out of place’. 
\end{styleStandard}

\begin{styleStandard}
\textit{\ \ Cualquier cosa} can be used predicatively on par with \textit{cualquiera} and have the same evaluative meaning of ‘bad, incorrect, out of place’. This is shown by the use of \textit{cualquier cosa }and \textit{cualquiera} under predicative verbs such as \textit{ser/parecer} ‘be’/‘behave’.
\end{styleStandard}

\begin{styleStandard}
\textbf{\ \ }\ \ Ese\ \ estudio\ \ es\ \ cualquier\ \ cosa,\ \ realmente\ \ cualquier\ \ cosa.
\end{styleStandard}

\begin{styleStandard}
\ \ \ \ that\ \ study\ \ is\ \ \textsc{cualquier\ \ }thing\ \ really\ \ cualquier\ \ thing
\end{styleStandard}

\begin{styleStandard}
\ \ \ \ ‘This study is bad/out of place, really bad/out of place.’
\end{styleStandard}

\begin{styleStandard}
\ \ \ \ (@ 140)
\end{styleStandard}

\begin{styleStandard}
As with \textit{cualquiera}, \textit{parece cualquier cosa} is often accompanied by some evaluator (e.g. ‘I find’):
\end{styleStandard}

\begin{styleStandard}
\textbf{\ \ }\ \ Me\ \ \textit{parece\ \ cualquier\ \ cosa\ \ }la\ \ designacion\ \ de\ \ este\ \ dia,\ \ francamente
\end{styleStandard}

\begin{styleStandard}
\ \ \ \ me\ \ seems\ \ \textsc{cualquier}\ \ thing\ \ the\ \ appointment\ \ of\ \ this\ \ day\ \ frankly
\end{styleStandard}

\begin{styleStandard}
\ \ \ \ ‘Today’s appointment seems out of place to me, frankly.’
\end{styleStandard}

\begin{styleStandard}
\ \ \ \ (@ 141)
\end{styleStandard}

\begin{styleStandard}
\textbf{\ \ }\ \ Me\ \ \textit{parece\ \ cualquier\ \ cosa\ \ }lo\textbf{\ \ }que\ \ planean\ \ hacer\ \ con\ \ la\ \ línea H.
\end{styleStandard}

\begin{styleStandard}
\ \ \ \ to.me\ \ seems\ \ \textsc{cualquier}\ \ thing\ \ it\ \ that\ \ plan\ \ to.do\ \ with\ \ the\ \ line H
\end{styleStandard}

\begin{styleStandard}
\ \ \ \ No\ \ hay\ \ estudios\ \ que\ \ demuestren\ \ la\ \ factibilida\ \ del\ \ neuvo\ \ recorrido.
\end{styleStandard}

\begin{styleStandard}
\ \ \ \ not\ \ have\ \ studies\ \ that\ \ show\ \ the\ \ feasibility\ \ of.the\ \ new\ \ route
\end{styleStandard}

\begin{styleStandard}
\ \ \ \ ‘I find it nonsense what they are planning to do with the H line.’
\end{styleStandard}

\begin{styleStandard}
\ \ \ \ (@ 142)
\end{styleStandard}

\begin{styleStandard}
A new finding is that \textit{cualquier cosa} can even be used as a degree adjective (only one occurrence has been identfieid in CdE\_web dialects with degree adverb \textit{muy}):
\end{styleStandard}

\begin{styleStandard}
\textbf{\ \ }\ \ No\ \ la\ \ Soporto\ \ en\ \ ninguna.\ \ Verónica\ \ Coincido\ \ con\ \ Maid\ \ en\ \ esa\ \ pareja\ \ es
\end{styleStandard}

\begin{styleStandard}
\ \ \ \ not\ \ her\ \ tolerate\ \ in\ \ no-one\ \ Verónica\ \ Coincido\ \ with\ \ Maid\ \ in\ \ that\ \ couple\ \ is
\end{styleStandard}

\begin{styleStandard}
\ \ \ \ \textit{muy\ \ cualquier\ \ cosa!\ \ }
\end{styleStandard}

\begin{styleStandard}
\ \ \ \ very\ \ \textsc{cualquier}\ \ thing\ \ 
\end{styleStandard}

\begin{styleStandard}
\ \ \ \ ‘I can’t stand any of them. Verónica Coincido with Maid in that couple is completely out of place.’
\end{styleStandard}

\begin{styleStandard}
\ \ \ \ (@ 143) 
\end{styleStandard}

\begin{styleStandard}
Web research has turned up more examples with degree modification by \textit{re/muy} from young bloggers or Facebook users:
\end{styleStandard}

\begin{styleStandard}
\textbf{\ \ }\ \ Mi\ \ pose\ \ es\ \ re\ \ cualquier\ \ cosa\ \ pero\ \ la\ \ edit\ \ es\ \ preciosa!
\end{styleStandard}

\begin{styleStandard}
\ \ \ \ my\ \ pose\ \ is\ \ very\ \ \textsc{cualquier}\ \ thing\ \ but\ \ the\ \ edit\ \ is\ \ gorgeous
\end{styleStandard}

\begin{styleStandard}
\ \ \ \ ‘My pose is completely out of place, but the edit is gorgeous!’
\end{styleStandard}

\begin{styleStandard}
\ \ \ \ (@ 144)
\end{styleStandard}

\begin{styleStandard}
\textbf{\ \ }\ \ Ya\ \ es\ \ muy\ \ cualquier\ \ cosa\ \ ya.
\end{styleStandard}

\begin{styleStandard}
\ \ \ \ now\ \ is\ \ very\ \ \textsc{cualquier}\ \ thing\ \ now
\end{styleStandard}

\begin{styleStandard}
\ \ \ \ ‘It is completely out of place, it definitely is.’
\end{styleStandard}

\begin{styleStandard}
\ \ \ \ (@ 145)
\end{styleStandard}

\begin{styleStandard}
\textbf{\ \ }\ \ Julio Ponce\ \ suena\ \ muy\ \ cualquier\ \ cosa,
\end{styleStandard}

\begin{styleStandard}
\ \ \ \ Julio Ponce\ \ plays\ \ very\ \ \textsc{cualquier}\ \ thing
\end{styleStandard}

\begin{styleStandard}
\ \ \ \ ‘Julio Ponce plays completely out of place.’
\end{styleStandard}

\begin{styleStandard}
\ \ \ \ (@ 146)
\end{styleStandard}

\begin{styleStandard}
I also found \textit{medio} with \textit{cualquier cosa} by young forum users and bloggers:
\end{styleStandard}

\begin{styleStandard}
\textbf{\ \ }\ \ Creo\ \ que\ \ te\ \ podés\ \ presentar\ \ de\ \ la\ \ nada,\ \ pero\ \ conviene\ \ mandarle
\end{styleStandard}

\begin{styleStandard}
\ \ \ \ believe\ \ that\ \ you\ \ can\ \ appear\ \ from\ \ the\ \ nothing\ \ but\ \ is.advisable\ \ send.him
\end{styleStandard}

\begin{styleStandard}
\ \ \ \ un\ \ mail.\ \ Es\ \ \textit{medio\ \ cualquier\ \ cosa\ \ }El\ \ sistema\ \ (y\ \ mejor\ \ ni\ \ hablar\ \ del
\end{styleStandard}

\begin{styleStandard}
\ \ \ \ a\ \ email\ \ is\ \ half\ \ \textsc{cualquier}\ \ thing\ \ the\ \ system\ \ and\ \ better\ \ not.even\ \ talk\ \ of.the
\end{styleStandard}

\begin{styleStandard}
\ \ \ \ adelanto\ \ de\ \ la\ \ segunda\ \ fecha\ \ todavía\ \ estoy\ \ con\ \ la\ \ vena)
\end{styleStandard}

\begin{styleStandard}
\ \ \ \ forward\ \ of\ \ the\ \ second\ \ date\ \ still\ \ am\ \ with\ \ the\ \ mood
\end{styleStandard}

\begin{styleStandard}
\ \ \ \ ‘I think that you can go there directly, but it is advisable to send her/him an email. The system is somewhat out of place (and it’s better if I don’t even talk about bringing the second date forward, I am still in a bad mood).’
\end{styleStandard}

\begin{styleStandard}
\ \ \ \ (@ 147)
\end{styleStandard}

\begin{styleStandard}
\textbf{\ \ }\ \ Sí,\ \ es\ \ cierto,\ \ la\ \ actuación\ \ de\ \ Cecilia\ \ Dopazo\ \ es\ \ \textit{medio\ \ cualquier\ \ cosa\ \ }[…]
\end{styleStandard}

\begin{styleStandard}
\ \ \ \ yes\ \ is\ \ true\ \ the\ \ performance\ \ of\ \ Cecilia\ \ Dopazo\ \ is\ \ half\ \ \textsc{cualquier}\ \ thing\ \ 
\end{styleStandard}

\begin{styleStandard}
\ \ \ \ y\ \ fue\ \ pésima\ \ la\ \ elección\ \ de\ \ ella\ \ para\ \ el\ \ papel.
\end{styleStandard}

\begin{styleStandard}
\ \ \ \ and\ \ was\ \ awful\ \ the\ \ decision\ \ of\ \ her\ \ for\ \ the\ \ role
\end{styleStandard}

\begin{styleStandard}
\ \ \ \ ‘Yes, it’s true, Cecilia Dopazo’s performance was somewhat out of place […] and choosing her for the role was an awful decision.’
\end{styleStandard}

\begin{styleStandard}
\ \ \ \ (@ 148)
\end{styleStandard}

\begin{styleStandard}
These data show that \textit{cualquiera }and \textit{cualquier cosa} can function as gradable adjectives and be modified by different degree modifiers (\textit{muy, re, medio}) that denote +/– high degrees for some speakers. However, more research is needed to identify the sociolinguistic parameters in the acceptability of degree modification of \textit{cualquier cosa}. 
\end{styleStandard}

\begin{styleStandard}
\ \ With respect to degree suffixes, I also found a difference in interpretation between \textit{cualquiera} and \textit{cualquier cosa}. The diminuitive\textit{ }suffix\textit{ -ita} can combine with \textit{cualquiera}, and in this case the speaker attenuates the bad quality ‘a bit bad’ (see 2.1):
\end{styleStandard}

\begin{styleStandard}
\textbf{\ \ }\ \ Bueno,\ \ cerramos\ \ este\ \ post\ \ que\ \ \textit{es\ \ cualquierita}
\end{styleStandard}

\begin{styleStandard}
\ \ \ \ well\ \ close\ \ this\ \ post\ \ that\ \ is\ \ \textsc{cualquier}.\textsc{dim}
\end{styleStandard}

\begin{styleStandard}
\ \ \ \ ‘Well, let’s close this post because it’s a bit nonsensical’
\end{styleStandard}

\begin{styleStandard}
\ \ \ \ (@ 118)
\end{styleStandard}

\begin{styleStandard}
The noun \textit{cosa} can be modified by a diminuitive -\textit{ita} to denote a little tiny thing (a\textit{ minimiser}, see Krifka 1995), and the combination \textit{cualquier cosita} denotes an FCI ‘any little tiny thing’:
\end{styleStandard}

\begin{styleStandard}
\textbf{\ \ }\ \ Dando\ \ educación\ \ de\ \ calidad\ \ y\ \ no\ \ \textit{cualquier\ \ cosita}
\end{styleStandard}

\begin{styleStandard}
\ \ \ \ offering\ \ education\ \ of\ \ quality\ \ and\ \ not\ \ \textsc{cualquier}\ \ thing.\textsc{dim}
\end{styleStandard}

\begin{styleStandard}
\ \ \ \ ‘Offering quality education and not just any little thing.’
\end{styleStandard}

\begin{styleStandard}
\ \ \ \ (@ 149)
\end{styleStandard}

\begin{styleStandard}
\textbf{\ \ }\ \ de\ \ \textit{cualquier\ \ cosita\ \ }llora\ \ peor\ \ que\ \ un\ \ chico
\end{styleStandard}

\begin{styleStandard}
\ \ \ \ of\ \ \textsc{cualquier}\ \ thing.\textsc{dim}\ \ weeps\ \ worse\ \ that\ \ a\ \ boy
\end{styleStandard}

\begin{styleStandard}
\ \ \ \ ‘He weeps over any little thing, even worse than a little boy’
\end{styleStandard}

\begin{styleStandard}
\ \ \ \ (@ 150)
\end{styleStandard}

\begin{styleStandard}
Whereas \textit{cualquier }in \textit{cualquier cosa} is originally a quantificational determiner D° with an FC meaning ‘anything’ (both in European and in Latin American Spanish), as shown in , in ArgSp the whole string \textit{cualquier cosa} has additionally developed a special evaluative meaning and has also been recategoried as an Adj°, see :
\end{styleStandard}

\begin{styleStandard}
\label{bkm:Ref150787837}\textbf{\ \ }\ \ D° cualquier cosa = ‘anything’ $\lambda $P[18E?]x. thing(\textit{x}) \& P(x) + Free Choice Implicature
\end{styleStandard}

\begin{styleStandard}
\ \ \ \ (ArgSp and EurSp) 
\end{styleStandard}

\begin{styleStandard}
\label{bkm:Ref150787813}\textbf{\ \ }\ \ Adj° cualquier cosa = $\lambda $d$\lambda $x. bad (d)(x)
\end{styleStandard}

\begin{styleStandard}
\ \ \ \ (ArgSp) 
\end{styleStandard}

\begin{styleStandard}
The observation that \textit{cualquiera} but not \textit{cualquier cosa} has been found with comparatives and the observation that \textit{cualquier cosa} is less frequent with gradable adjectives seems to suggest that the process of recategorisation into a gradable adjective is more advanced with \textit{cualquiera} than with \textit{cualquier cosa}. If the difference between \textit{cualquiera} and \textit{cualquier cosa} lies in the difference in the ongoing process of lexicalisation, we should observe differences in age distribution or geographical distribution. I will test this hypothesis with native speaker judgements in future research. 
\end{styleStandard}

\subsubsection[4.6.3 Diachronic evidence]{\textbf{4.6.3 Diachronic evidenc}\textbf{e}}
\begin{styleStandard}
For diachronic research of \textit{cualquiera} and \textit{cualquier cosa} in Argentinian Spanish, I first used the diachronic corpus CORDIAM (Corpus Diacrónico y Diatópico del Español, CORDIAM, \href{http://www.cordiam.org}{www.cordiam.org}, which contains a total of 10,889,176 words from 19 Latin-American countries from 1494 to 1905, with documents from archives, journals, and literature. It is annotated for geographical origins. The corpus CORDIAM contains only 104 instances of \textit{cualquiera} in Argentinian Spanish, only one instance of N\textit{ cualquiera}, and zero instances of \textit{cualquier cosa} and \textit{cualquiera} with the meaning of \textit{cualquier cosa}. The remaining instances use \textit{cualquier} NP and \textit{cualquiera} as ‘any NP’ or ‘anyone’: 
\end{styleStandard}

\begin{styleStandard}
\textbf{\ \ }\ \ Ella [...], á este le dirije un cumplimiento, á ese una mirada simpática, á aquel un apóstrofe cualquiera,
\end{styleStandard}

\begin{styleStandard}
\ \ \ \ ‘She sends this one a compliment, to the other one a nice smile, and to that one a random insult.’
\end{styleStandard}

\begin{styleListParagraph}
\textmd{\ \ \ \ (19}\textmd{\textsuperscript{th}}\textmd{, anonymous, journal: }\textmd{\textit{La Camelia}}\textmd{; ArgSp)}
\end{styleListParagraph}

\begin{styleStandard}
I also used a larger diachronic corpus, CORDE (Corpus Diacrónico del Español), which is a diachronic corpus with written texts spanning from Old Spanish to 1974. It contains 250 million written words from different genres, types, and geographical origins. \textit{Cualquiera} appears 535 times in ArgSp starting from the 18th century. Before that, the form \textit{qualquiera} appeared nine times from the 16th to the 18th century (see Chapter 6 on different forms). There is no case of gradable \textit{cualquiera} in predicative position (e.g. \textit{es muy cualquiera}). However, there is one occurrence of predicative \textit{cualquiera} without nominal modification with the meaning ‘nothing important’ or ‘ordinary’: 
\end{styleStandard}

\begin{styleStandard}
\label{bkm:Ref150787886}\textbf{\ \ }\textit{Cualquiera} as a predicative complement
\end{styleStandard}

\begin{styleStandard}
\ \ \ \ \textit{Su ocupación era cualquiera.}
\end{styleStandard}

\begin{styleStandard}
\ \ \ \ his position was cualquiera
\end{styleStandard}

\begin{styleStandard}
\ \ \ \ ‘He had an ordinary position./His position was nothing important’
\end{styleStandard}

\begin{styleListParagraph}
\textmd{\ \ \ \ (1926, Güiraldes, }\textmd{\textit{Don Segundo Sombra}}\textmd{; ArgSp)}
\end{styleListParagraph}

\begin{styleStandard}
One interesting observation is that \textsc{un} N \textit{cualquiera}, such as in , appears more often in affirmative contexts in Argentinian Spanish than it does in European Spanish (56\% in ArgSp vs 18\% in EurSp). This observation of the relatively high frequency of \textsc{un} N \textit{cualquiera} in ArgSp can be interpreted as an indicator that \textit{cualquiera} is undergoing the process of lexicalisation of the meaning ‘ordinary’ and ‘bad’. In Chapter 6, we will look at the development of \textit{cualquiera} in EurSp and investigate the appearance of the evaluative reading (EVAL 1).
\end{styleStandard}

\begin{styleStandard}
\ \ There is no diachronic data that could confirm step 2 in 4.5.3. CORDE does not contain \textit{cualquiera} as a substitute for \textit{cualquier cosa}. CORDE contains 167 examples of \textit{cualquier cosa} in ArgSp among which five examples of \textit{cualquier cosa} are embedded in episodic contexts under verbs of saying (\textit{decir}, \textit{hablar de, murmurar, conversar de}). Note that \textit{cualquier cosa} does not mean ‘nonsense’ in these contexts, but rather has a weaker reading of ‘nothing in particular’ or ‘something unimportant’; that is, the speaker does not care to specify the reference of the object of conversation, as in (\ref{bkm:Ref150787920})a), or the speaker has no plan in mind of what he was going to do, as in (\ref{bkm:Ref150787920})b) (see the indifference or indiscriminative reading in 2.4).
\end{styleStandard}

\begin{listWWNumxleveli}
\item 
\end{listWWNumxleveli}
\begin{styleStandard}
\ \ \textit{Cualquier cosa} in ArgSp
\end{styleStandard}

\begin{styleStandard}
\ \ a.\ \ Y el público, formado por la gente de huella y de estancia, conversaba de cualquier cosa.
\end{styleStandard}

\begin{styleStandard}
\ \ \ \ ‘The general public, herders and ranchers, hummed with talk about everything and nothing.’
\end{styleStandard}

\begin{styleStandard}
\ \ \ \ (1926, Güiraldes, \textit{Don Segundo Sombra}; ArgSp)
\end{styleStandard}

\begin{styleStandard}
\ \ b.\ \ Yo salí muy decidido a cualquier cosa, pero a nada en particular.
\end{styleStandard}

\begin{styleStandard}
\ \ \ \ ‘I was not sure what I was going to do.’
\end{styleStandard}

\begin{styleStandard}
\ \ \ \ (1940, Bioy Casares, \textit{La invención de Morel}; ArgSp)
\end{styleStandard}

\begin{styleStandard}
\textbf{\ \ }c.\ \ Y de pronto, para romper el intolerable silencio, yo decía cualquier cosa que en otro tiempo podía haber tenido interés para el tipógrafo. - La fora ha declarado una huelga de estibadores.
\end{styleStandard}

\begin{styleStandard}
\ \ \ \ ‘And suddenly, in order to break the untolearbale silence, I said just something or other that in another time would have had some interest for the typographer.’
\end{styleStandard}

\begin{styleStandard}
\ \ \ \ (1961,\textit{ }Sábato, \textit{Héroes y tumbas}; ArgSp)
\end{styleStandard}

\begin{styleStandard}
\ \ d.\ \ Estábamos tendidos de espaldas, uno al lado del otro, y ella me acariciaba maquinalmente, yo tenía frío y ella me hablaba de cualquier cosa, de la pelea que acababa de ocurrir en el bar, de las tormentas de marzo, […]
\end{styleStandard}

\begin{styleStandard}
\ \ \ \ ‘We were close together, one against the other, and she caressed automatically, I was cold and she talked to me about something or other, about the fight that had just taken place in the bar, about the storms in March, […] .’
\end{styleStandard}

\begin{styleStandard}
\ \ \ \ (1963, Cortazár, \textit{Rayuela}; ArgSp)
\end{styleStandard}

\begin{styleStandard}
\ \ e.\ \ Suspirando, Traveler se volvió a la cama. A una pregunta soñolienta de Talita, le acarició el pelo y murmuró cualquier cosa. Talita besó el aire, manoteó un poco, se tranquilizó.
\end{styleStandard}

\begin{styleStandard}
\ \ \ \ ‘Sighing, Traveler went back to his room. After Talita asked him a question in a sleepy state, he caressed her hair and whispered something. […]’
\end{styleStandard}

\begin{styleStandard}
\ \ \ \ (1963, Cortazár, \textit{Rayuela}; ArgSp)
\end{styleStandard}

\begin{styleStandard}
To sum up, the diachronic corpora contain only one example of a predicative adjectival use of \textit{cualquiera} and no adjectival use of \textit{cualquier cosa}. They do not show the special ‘nonsense’ reading under verbs of saying. The first examples that appear with degree modification are examples from the modern use of Argentinian Spanish in CdE\_web dialects in the 21st century, as reported in Section 4.3. The earliest occurrences that I could find of \textit{decir/mandar cualquiera} originate from the 21st century, although according to Conde (2011) \textit{decir/mandar cualquiera} is apparently from Lunfardo, which must be older than the 21st century. One possible explanation why there is no diachronic data of special uses of \textit{cualquiera} in ArgSp is that diachronic data is mostly taken from written language, whereas the special uses of \textit{cualquiera} represent spoken language. Another explanation is that the use of \textit{cualquiera} as a gradable adjective with the meaning ‘ordinary’ and ‘bad’ is a recent innovation of native speakers of Argentinian Spanish. Describing diachronic change of spoken language is thus extremely difficult with corpus methods due to the scarcity of corpus data.
\end{styleStandard}

\clearpage\section[5 Diachrony of qualquier revisited]{\textbf{5 Diachrony of }\textbf{\textsc{qualquier }}\textbf{revisited}}
\begin{styleStandard}
The Old Spanish forms of \textit{cualquiera} were rather \textit{qualquier,\nobreakdash-e,\nobreakdash-a} or \textit{qual quier,\nobreakdash-e,\nobreakdash-a} (see C\&P 2009, Pato 2012, Mackenzie 2019, among others). I will use the notation \textsc{qualquier} in reference to both Old and Modern Spanish\textsc{. }By contrast, the notation \textit{qualquier,\nobreakdash-e,\nobreakdash-a} refers exclusively to Old Spanish, and \textit{cualquier,-a} only to Modern Spanish. In this Chapter, I will look at the syntactic and semantic change of the synthetic form of \textsc{qualquier}\textit{,} the form that has been already reanalysed as a compound word, which has its origin in a free relative clause (see Palomo 1934, Lombard 1938–39/1947–48, Meier 1950, Rivero 1988, Espinosa \& Sánchez Lancis 2006, C\&P 2009, among others), illustrated again schematically as follows:
\end{styleStandard}

\begin{styleStandard}
\textbf{\ \ }a.\ \ [\textsubscript{MainClause} ….V\textsubscript{main}... [\textsubscript{CP-Rel.} \textbf{\textit{qualquier }}\textbf{[}\textbf{\textsubscript{C°}}\textbf{ }\textbf{\textit{que}}] [\textsubscript{TP} \textbf{V}\textbf{\textsubscript{[subj.]}}]]] or 
\end{styleStandard}

\begin{styleStandard}
\ \ b.\ \ [\textsubscript{MainClause} [\textsubscript{CP-Rel.} \textbf{\textit{qualquier }}[\textbf{\textsubscript{C°}}\textbf{ }\textbf{\textit{que}}] [\textsubscript{TP} \textbf{V}\textbf{\textsubscript{[subj.]}}]]] ...V\textsubscript{main} ...]
\end{styleStandard}

\begin{styleStandard}
\textbf{\ \ }a.\ \ Sepan quantos esta carta viren como yo [...] \textit{quito e escuso el su portero de la eglesia, qual quier que sea}, por todo tiempo de todo pecho [...] 
\end{styleStandard}

\begin{styleStandard}
\ \ \ \ ‘May all that see this letter know how I […] I exempt and discharge the usher of the church, \textit{whoever (one wants that) it might be}, from all feudal taxes [...]’
\end{styleStandard}

\begin{styleStandard}
\ \ \ \ (13th c., Alfonso X, \textit{Doc. Cast. - León})
\end{styleStandard}

\begin{styleStandard}
\ \ b.\ \ E tan grandes eran las bozes delos hombres \& del ganado que \textit{qualquier que lo viesse} hauria ende gran dolor \& muy gran piedad…, 
\end{styleStandard}

\begin{styleStandard}
\ \ \ \ ‘and the voices of the men and of the cattle were so great that \textit{whoever (one wants that) s/he sees it} would have great pain and very great pity…,’
\end{styleStandard}

\begin{styleStandard}
\ \ \ \ (13th c.,\textit{ }Anonymous,\textit{ Gran Conquista})
\end{styleStandard}

\begin{styleStandard}
272 and 273, respectively, show two kinds of embedding of the relative clause in the 13th century. In 272 a, the relative clause introduced by \textit{qualquier} modifies an element of the main clause (the DP \textit{el su portero de la iglesia}), whereas in (272b) the relative clause itself is the subject of the main clause. As will be shown in 5.2, these are the most frequent types in the 13th century. Other types will be discussed in 5.3. The separated form \textit{qual quier} (the prior stage of the synthetic form \textsc{qualquier}) will not be the focus in this section and the statistical analysis in 5.1 and 5.2. The separated form \textit{qual quier} and the interposition form \textit{qual NP quier} was mostly used in the 13th century and disappeared from that period on (see C\&P 2009, Pato 2012), which is why I do not include this form in this study. 
\end{styleStandard}

\begin{styleStandard}
\ \ I will show in this section that the frequency of the use of the synthetic forms represented by \textsc{qualquier} in sentence types other than relative clauses as illustrated in 274 rose over time. For instance, the frequency of \textsc{qualquier} in main declarative clauses with indicative verbal mood as exemplified in a (schematically analysed in 274) rises significantly (for other instances of sentence types, see 5.2):
\end{styleStandard}

\begin{styleStandard}
\textbf{\ \ }a.\ \ [\textsubscript{Main Declarative Clause} ... Verb\textsubscript{[Indic.] }... [\textsc{qualquier} NP]]\textsubscript{subj} \ \ \ (postverbal subject) or
\end{styleStandard}

\begin{styleStandard}
\ \ b.\ \ [\textsubscript{Main Declarative Clause} ... [\textsc{qualquier}]\textsubscript{subj}. Verb\textsubscript{[Indic.]}]\ \ (preverbal subject)
\end{styleStandard}

\begin{styleStandard}
\textbf{\ \ }a.\ \ Por ende bien \textit{se avise} desto \textit{qualquier} varón. 
\end{styleStandard}

\begin{styleStandard}
\ \ \ \ ‘For this reason any man needs to get informed about this.’
\end{styleStandard}

\begin{styleStandard}
\ \ \ \ (14\textsuperscript{th} c., López de Ayala,\textit{ Libro rimado})
\end{styleStandard}

\begin{styleStandard}
\ \ b.\ \ […], esto \textit{qualquier} \textit{sabría},[...]
\end{styleStandard}

\begin{styleStandard}
‘[…], anybody would know it, […]’
\end{styleStandard}

\begin{styleStandard}
\ \ \ \ (14\textsuperscript{th} c., López de Ayala,\textit{ Libro rimado})
\end{styleStandard}

\begin{styleStandard}
As will be shown in this section, the sentence type and the verbal mood as well as other syntactic factors correlate with differences in the interpretation of \textsc{qualquier} (+/– EVAL). I will test the frequency changes of the type described in 273 vs 274 and other changes quantitatively. The results of these tests will show how \textsc{qualquier} developed further on in time after it changed from a relative clause [\textsubscript{Rel.Cl}.= \textsubscript{D} \textit{qual} \textsubscript{V} \textit{quier} ...], as discussed in Section 5.2.1, into an indefinite determiner [\textsubscript{D} \textit{qualquier/cualquier}] as in \textit{cualquier} \textit{persona} ‘any person’ (see Section 5.2.1 on this grammaticalisation process).
\end{styleStandard}

\begin{styleStandard}
\ \ The ultimate goal of investigating the diachronic change of \textsc{qualquier} is to see how it developed different meanings across time, such as the evaluative meanings ‘ordinary’, ‘common’, and ‘unremarkable’ and whether this development of meanings is related to the (morpho)syntactic structure and change of \textsc{qualquier.} 
\end{styleStandard}

\begin{styleStandard}
\ \ This chapter is organised as follows. 5.1 introduces into the methodology and data selection chosen. The changes will be first described quantitatively in 5.2 and verified statistically. 5.3 suggests a syntactic and semantic analysis of the changes described in 5.2. A summary will be presented in 5.4.
\end{styleStandard}

\subsection[5.1 Methodology and data]{\textbf{5.1 Methodology and data}}
\begin{styleStandard}
\ This Section introduces into the methodology of this Chapter. 
\end{styleStandard}

\subsubsection[5.1.1 Data selection]{\textbf{5.1.1 Data selection}}
\begin{styleStandard}
In order to investigate the development of \textsc{qualquier}, I used the diachronic subcorpus, \textit{Genre/Historical} of Mark Davies’ Corpus del Español (CdE\_G/H): \url{https://www.corpusdelespanol.org/hist-gen/}. This subcorpus contains more than 20,000 Spanish texts from the 1200s to the 1900s.\footnote{\textrm{\ There is another historical corpus, CORDE. This corpus is not annotated in any way and thus does not allow one to search for syntactic configurations. It is only useful for searching for exact phrases. Its main advantage, though, is that CORDE offers the possibility to limit the search to specific countries and registers across periods and varieties. CORDE contains 240 million words from 950–1974 and does not contain any oral data, relying exclusively on written texts from books and the press. It does not provide the relative frequency, which must be calculated manually. For these reasons I decided not to use it.}}
\end{styleStandard}

\begin{styleStandard}
\ \ Different research questions formulated in Chapter 5 have been answered using different data sets of the diachronic corpus CdE\_G/H. The research question concerning the development of \textsc{qualquier} and the correlation between parameters such as time and linguistic features (e.g. verbal position and sentence type) requires manual tagging of every example. For this research question a set of 150 examples per century was created, except for the 14th century, for which only 142 examples could be gathered due to the low number of data in this century in the diachronic corpus CdE\_G/H, which contains only 216 examples of \textsc{qualquier} in total\textsc{. }
\end{styleStandard}

\begin{styleStandard}
\ \ I searched the corpus for the string *ualqu*. The asterisks serve as a wild card character for any letter or sequence of letters. The search was thus sensitive to forms such as: \textit{cualquier, cualquiera, qualquier, qualquiera, qualquiere }(see C\&P 2009, Pato 2012 for graphical variants). The corpus data for answering the aforementioned research question does not include the separated form \textit{qual quier} and thus also excludes cases of interposition such as qual NP quier (see Pato 2012 on interposition). As I already explained, I use the generic notation \textsc{qualquier} to refer to all these variants. 
\end{styleStandard}

\begin{styleStandard}
\ \ The diachronic corpus contains data from the 13th to the 20th century. The corpus represents the centuries as follows: 1200s (1200–1300), 1300s (1300–1400), 1400s (1400–1500), 1500s (1500-1600), 1600s (1600–1700), 1700s (1700–1800), 1800s (1800–1900), 1900s (1900–2000):
\end{styleStandard}

\begin{styleStandard}
Figure \stepcounter{Figure}{\theFigure}: Screenshot of the corpus search\_CdE\_G/H.\footnote{\textrm{\ In Section 5.1 and 5.2 of this Chapter, I use screenshots from the output of the CdE search. In these screenshots, frequencies are indicated either as per million words as absolute numbers or both values are displayed. The information about the relative frequency is especially necessary where I compared frequencies of the whole corpus CdE, because the subcorpora used to represent every century or time period are not equal in CdE and the relative frequency allows us to compare the data between periods (see Kellert 2018 on relative frequency measurements in CdE in Old Spanish).}}
\end{styleStandard}

\begin{styleStandard}
For each century, at least 500 examples (with the exception of the 14th c.) were selected for classification. I annotated the centuries in my data with ordinal numbers (e.g. 13th c.). For the first three centuries (i.e. 13th, 14th, and 15th centuries), I excluded some texts whose edition cannot be considered philologically reliable (see Octavio de Toledo \& Rodríguez Molina 2017 and Chapter 1). Moreover, I excluded cases of coordination and ellipsis (see examples of coordination and ellipsis in Old Spanish in 5.2) so as to be able to annotate verbal properties of the examples with \textsc{qualquier} (see below, Section 5.1.2). After the exclusion of such occurrences of \textsc{qualquier}, I chose at least 150 random examples for each century (142 ex. for the 14th c., as I have explained earlier), which is a sufficient number for a quantitative analysis.\footnote{\ \textrm{In order to make sure that my corpus is big enough to run statistical tests, I contacted} \textrm{Alexander Silbersdorff, Chair of Statistics at the University of Göttingen and Thomas Weskott (Germanistik, University of Göttingen). I especially thank Thomas Weskott for his advice on how to implement the model of logistic regression on “my data”.}} I call the data set acquired with this methodology “my data”. The data as well as the information about the selected texts (represented in percentages) and the cases that have been filtered out (e.g. coordination and ellipsis) can be found in the respective Excel file on GitHub upon the acceptance of this publication.
\end{styleStandard}

\begin{styleStandard}
I use examples from the whole corpus CdE\_G/H in order to make suggestions about previous diachronic stages of the grammaticalisation before \textsc{qualquier} was reanalysed as a compound word (see Section 5.3). 
\end{styleStandard}

\subsubsection[5.1.2 Classification of the examples]{\textbf{5.1.2 Classification of the examples}}
\begin{styleStandard}
As a result of the method described above, my data contain 150 examples with \textsc{qualquier} per century (except the 14\textsuperscript{th} century). I classified these examples using different syntactic and semantic variables that were chosen based on my research questions. These variables can be clustered according to different groups represented as A–F. The graphic variation represents group A; \textsc{qualquier} as the modifier of an adjacent noun vs \textsc{qualquier} alone (i.e. as bare Q without noun modification, or QN/NQ as nominal modification) is Group B; the syntactic function of \textsc{qualquier} (subject, direct object, indirect object, predicative noun, adverb) represents Group C; the position of \textsc{qualquier} with respect to the verb of which \textsc{qualquier} is an argument or which \textsc{qualquier} modifies if it is an adverbial (QV/VQ) represents Group D; the mood, tense, aspect, and finiteness of the verb of which \textsc{qualquier} is an argument or which \textsc{qualquier} modifies if it is an adverbial (indicative/subjunctive, infinitive/gerund/participle, perfective aspect\footnote{\textrm{\ Perfective aspect was chosen to determine whether }\textrm{\textit{cualquiera}}\textrm{ can be used under perfective verbs in diachrony (see Section 2 for the role of perfective verbs in synchrony).}}) constitutes Group E; and the sentence type and sentence mood of the sentence in which \textsc{qualquier} appears represent Group F. The latter group contains the clause type (main clause, subordinated clause) and the sentence mood: declarative or \textsc{qualquier} as antecedent of a relative or complement clause\footnote{\ \textrm{In the descriptive part, I describe a clause introduced by }\textrm{\textit{que}}\textrm{ pretheoretically as a relative or complement clause. The distinction depends strongly on the analysis of }\textrm{\textit{quer}}\textrm{ in }\textrm{\textit{qualquier}}\textrm{.}} introduced by \textit{que} (\textsc{qualquier}\textit{ (XP) que} or\textit{ Q que}), which I analysed as (free) relatives (FRel) or other type.
\end{styleStandard}

\begin{styleStandard}
\ \ Before I present the data, it is important to distinguish between \textsc{qualquier} that appears in a main clause and \textsc{qualquier} that appears in a relative clause or other subordinate clauses, because this difference influences the classification of \textsc{qualquier}. As my interest is to understand the relation of \textsc{qualquier} to the elements that \textsc{qualquier} structurally depends on, I only classified the relation between \textsc{qualquier} and the structurally dependent verb. When \textsc{qualquier} introduced a relative clause, as in 272, I classified \textsc{qualquier} with respect to the variables described above and used in Tables 17–19 below (e.g. its syntactic function or its position to the verb) \textit{inside the relative clause} rather than inside the matrix clause. This classification is a logical consequence of my analysis of relative clauses, that is, \textsc{qualquier} is analysed as part of a relative clause (i.e. as a specifier of a relative clause) and that as such it is structurally connected to the elements inside the relative clause, such as the verb in subjunctive (see for instance Quer 2000 on the importance of the subjunctive in relative clauses).
\end{styleStandard}

\begin{styleStandard}
\ \ I now proceed to go through the overall results concerning the different variables and to illustrate them with examples from the corpus.
\end{styleStandard}

\begin{styleStandard}
\ \ Table 12 shows the absolute frequencies for the variables of group A to D:
\end{styleStandard}

\begin{center}
\tablefirsthead{}
\tablehead{}
\tabletail{}
\tablelasttail{}
\begin{supertabular}{|m{0.48515984in}|m{0.39065984in}|m{0.39065984in}|m{0.39275986in}|m{0.39345986in}|m{0.29625985in}|m{0.29625985in}|m{0.23795986in}|m{0.17885986in}|m{0.17885986in}|m{0.5476598in}|m{0.29765984in}|m{0.27055985in}|m{0.34485984in}|}
\hhline{~~------------}
\multicolumn{1}{m{0.48515984in}}{} &
 &
\multicolumn{2}{m{0.8621598in}|}{\centering \textbf{A}} &
\multicolumn{3}{m{1.1434599in}|}{\centering \textbf{B}} &
\multicolumn{5}{m{1.75596in}|}{\centering \textbf{C}} &
\multicolumn{2}{m{0.6941598in}|}{\centering \textbf{D}}\\\hhline{~~------------}
\multicolumn{1}{m{0.48515984in}}{} &
 &
\multicolumn{2}{m{0.8621598in}|}{\centering \textbf{Orthography}} &
\multicolumn{3}{m{1.1434599in}|}{\centering \textbf{+/- noun modifier}} &
\multicolumn{5}{m{1.75596in}|}{\centering \textbf{Syntactic function}} &
\multicolumn{2}{m{0.6941598in}|}{\centering \textbf{VQ}}\\\hline
\centering \textbf{Century} &
\centering \textbf{Occ} &
\centering \textbf{\textit{qual-}} &
\centering \textbf{\textit{cual-}} &
\centering \textbf{Bare Q} &
\centering \textbf{QN} &
\centering \textbf{NQ} &
\centering \textbf{Subj} &
\centering \textbf{DO} &
\centering \textbf{IO} &
{\centering \textbf{Copula}\par}

\centering \textbf{Predicate} &
\centering \textbf{Adv} &
\centering \textbf{QV} &
\centering\arraybslash \textbf{VQ}\\\hline
\centering 13th &
\centering 150 &
\centering 149 &
\centering 1 &
\centering 102 &
\centering 43 &
\centering 5 &
\centering 94 &
\centering 15 &
\centering 25 &
\centering 5 &
\centering 11 &
\centering 134 &
\centering\arraybslash 16\\\hline
\centering 14th &
\centering 142 &
\centering 139 &
\centering 3 &
\centering 62 &
\centering 78 &
\centering 2 &
\centering 70 &
\centering 15 &
\centering 34 &
\centering 16 &
\centering 7 &
\centering 118 &
\centering\arraybslash 24\\\hline
\centering 15th &
\centering 150 &
\centering 141 &
\centering 9 &
\centering 33 &
\centering 116 &
\centering 1 &
\centering 27 &
\centering 26 &
\centering 65 &
\centering 11 &
\centering 21 &
\centering 102 &
\centering\arraybslash 48\\\hline
\centering 16th &
\centering 150 &
\centering 13 &
\centering 137 &
\centering 37 &
\centering 113 &
\centering 0 &
\centering 35 &
\centering 26 &
\centering 62 &
\centering 13 &
\centering 14 &
\centering 98 &
\centering\arraybslash 52\\\hline
\centering 17th &
\centering 150 &
\centering 12 &
\centering 138 &
\centering 38 &
\centering 111 &
\centering 1 &
\centering 32 &
\centering 42 &
\centering 52 &
\centering 12 &
\centering 12 &
\centering 83 &
\centering\arraybslash 67\\\hline
\centering 18th &
\centering 149 &
\centering 18 &
\centering 131 &
\centering 41 &
\centering 106 &
\centering 2 &
\centering 29 &
\centering 32 &
\centering 66 &
\centering 11 &
\centering 11 &
\centering 91 &
\centering\arraybslash 58\\\hline
\centering 19th &
\centering 150 &
\centering 0 &
\centering 150 &
\centering 55 &
\centering 81 &
\centering 14 &
\centering 24 &
\centering 37 &
\centering 65 &
\centering 14 &
\centering 10 &
\centering 71 &
\centering\arraybslash 79\\\hline
\centering 20th &
\centering 150 &
\centering 0 &
\centering 150 &
\centering 23 &
\centering 123 &
\centering 4 &
\centering 25 &
\centering 38 &
\centering 54 &
\centering 7 &
\centering 26 &
\centering 54 &
\centering\arraybslash 96\\\hline
\end{supertabular}
\end{center}
\begin{styleStandard}
Table \stepcounter{Table}{\theTable}: CdE\_cualquiera\_150 occ. per century, “my data” (1).
\end{styleStandard}

\begin{styleStandard}
In 276, I provide examples for the variables of groups from B to D:
\end{styleStandard}

\begin{styleStandard}
\textbf{\ \ }a.\ \ Bare Q
\end{styleStandard}

\begin{styleStandard}
\ \ \ \ Et mando a~\textit{qualquier}~que touiere estas salinas que les de cadanno estos ocho caffizes de sal bien [...] 
\end{styleStandard}

\begin{styleStandard}
\ \ \ \ ‘And I ordered anyone who has these salterns, that he should give them every year eight \textit{cahices} [measuring unit] of salt […]’
\end{styleStandard}

\begin{styleStandard}
\ \ \ \ (13\textsuperscript{th }c., Alfonso X, \textit{Doc. Cast. - Castilla la Nueva})
\end{styleStandard}

\begin{styleStandard}
\ \ b.\ \ QN
\end{styleStandard}

\begin{styleStandard}
\ \ \ \ Si puede tomar armas~\textit{qualquier~persona}, sea la quistión primera. 
\end{styleStandard}

\begin{styleStandard}
\ \ \ \ ‘[The issue of] whether any person can take weapons, should be the first question.’
\end{styleStandard}

\begin{styleStandard}
\ \ \ \ (15\textsuperscript{th} c., Rodríguez del Padrón,\textit{ Triunfo de las donas})
\end{styleStandard}

\begin{styleStandard}
\ \ c.\ \ NQ
\end{styleStandard}

\begin{styleStandard}
\ \ \ \ y no parecería muy duro que hubiera quien ayudase a llenarle de un \textit{modo}~\textit{cualquiera}. 
\end{styleStandard}

\begin{styleStandard}
\ \ \ \ ‘and it would not seem difficult [to find] that there would be someone who would help to fill it [i.e. the remaining space of an area on a newspaper page] \textit{in any (whichever) way}.’
\end{styleStandard}

\begin{styleStandard}
\ \ \ \ (19\textsuperscript{th} c.,\textit{ }Arenal,\textit{ Beneficencia})
\end{styleStandard}

\begin{styleStandard}
\ \ d.\ \ S=Subject
\end{styleStandard}

\begin{styleStandard}
\ \ \ \ somos del todo enfermos e con mucha amargura; \textit{qualquier tentaçión} nos vençe, aunque non sea muy dura. 
\end{styleStandard}

\begin{styleStandard}
\ \ \ \ ‘we are totally weak and with much bitterness: any temptation wins over us, even if it is not hard [to resist].’
\end{styleStandard}

\begin{styleStandard}
\ \ \ \ (14\textsuperscript{th} c., López de Ayala,\textit{ Libro rimado})
\end{styleStandard}

\begin{styleStandard}
\ \ e.\ \ DO
\end{styleStandard}

\begin{styleStandard}
\ \ \ \ Si puedo~\textit{cualquier~sustancia} doctamente definir, y accidentes describir en la mayor importancia, … 
\end{styleStandard}

\begin{styleStandard}
\ \ \ \ ‘if I can define any kind of substance with erudition and describe any contingent property in its highest importance, ...’
\end{styleStandard}

\begin{styleStandard}
\ \ \ \ (17\textsuperscript{th} c., Villamediana\textit{, Poesía})
\end{styleStandard}

\begin{styleStandard}
\ \ f.\ \ IO
\end{styleStandard}

\begin{styleStandard}
\ \ \ \ [...] E a qualquier delos sobre dichos que mostrare alguna destas escusas bien lo deuen oyr desque el casamiento sea hecho. 
\end{styleStandard}

\begin{styleStandard}
\ \ \ \ ‘and \textit{from anyone} of the aforementioned that will \textit{show} some of these excuses, they well need to hear it before the marriage is made...’
\end{styleStandard}

\begin{styleStandard}
\ \ \ \ (13\textsuperscript{th} c., Alfonso X, \textit{Siete partidas})
\end{styleStandard}

\begin{styleStandard}
\ \ g.\ \ QV
\end{styleStandard}

\begin{styleStandard}
\ \ \ \ Sin dubda Dios perdona el pecado llorado; enpero \textit{qualquier} \textit{se tema}, si puede aver ganado tal graçia; […]
\end{styleStandard}

\begin{styleStandard}
\ \ \ \ ‘Without any doubt, God forgives the regretted sin; but anyone who may be frightened may have obtained such grace…’
\end{styleStandard}

\begin{styleStandard}
\ \ \ \ (14\textsuperscript{th} c., López de Ayala,\textit{ Libro rimado})
\end{styleStandard}

\begin{styleStandard}
\ \ h.\ \ VQ
\end{styleStandard}

\begin{styleStandard}
\ \ \ \ Mi mamá política hacía trampas sin escrúpulos para \textit{ganar~cualquier}~juego en que participara. 
\end{styleStandard}

\begin{styleStandard}
\ \ \ \ ‘My political mum used tricks in an unscrupolous manner to win any kind of game to which she participated.’
\end{styleStandard}

\begin{styleStandard}
\ \ \ \ (20\textsuperscript{th }c., Pisabarro,\textit{ Del agua nacieron los sedientos})
\end{styleStandard}

\begin{styleStandard}
Table 13 shows all numerical values for the variables from Group E:
\end{styleStandard}

\begin{center}
\tablefirsthead{}
\tablehead{}
\tabletail{}
\tablelasttail{}
\begin{supertabular}{|m{0.54205984in}|m{0.44005987in}|m{0.6434598in}|m{0.5504598in}|m{0.6011598in}|m{0.54345983in}|m{0.5483598in}|m{0.5455598in}|m{0.55315983in}|}
\hhline{~~-------}
\multicolumn{1}{m{0.54205984in}}{} &
 &
\multicolumn{7}{m{4.45806in}|}{\centering \textbf{E}}\\\hhline{~~-------}
\multicolumn{1}{m{0.54205984in}}{} &
 &
\multicolumn{7}{m{4.45806in}|}{\centering \textbf{Verbal Aspect/Mood/Finiteness}}\\\hline
 &
 &
\multicolumn{2}{m{1.2726598in}|}{\centering \textbf{Mood}} &
\multicolumn{3}{m{1.8504598in}|}{\centering \textbf{Inf}} &
\multicolumn{2}{m{1.17746in}|}{\centering \textbf{Tense and Aspect}}\\\hline
\centering \textbf{Century} &
\centering \textbf{Occ} &
\centering \textbf{Ind} &
\centering \textbf{Sbjv} &
\centering \textbf{Imperative} &
\centering \textbf{Infinitive} &
\centering \textbf{Gerund} &
\centering \textbf{Participle} &
\centering\arraybslash \textbf{Perfective}\\\hline
\centering 13th &
\centering 150 &
\centering 26 &
\centering 120 &
\centering 3 &
\centering 1 &
\centering 0 &
\centering 0 &
\centering\arraybslash 2\\\hline
\centering 14th &
\centering 142 &
\centering 61 &
\centering 74 &
\centering 0 &
\centering 5 &
\centering 1 &
\centering 1 &
\centering\arraybslash 1\\\hline
\centering 15th &
\centering 150 &
\centering 60 &
\centering 66 &
\centering 0 &
\centering 16 &
\centering 1 &
\centering 7 &
\centering\arraybslash 2\\\hline
\centering 16th &
\centering 150 &
\centering 89 &
\centering 45 &
\centering 1 &
\centering 12 &
\centering 3 &
\centering 0 &
\centering\arraybslash 5\\\hline
\centering 17th &
\centering 150 &
\centering 99 &
\centering 33 &
\centering 2 &
\centering 12 &
\centering 3 &
\centering 1 &
\centering\arraybslash 3\\\hline
\centering 18th &
\centering 149 &
\centering 64 &
\centering 49 &
\centering 1 &
\centering 29 &
\centering 3 &
\centering 3 &
\centering\arraybslash 2\\\hline
\centering 19th &
\centering 150 &
\centering 74 &
\centering 51 &
\centering 2 &
\centering 18 &
\centering 1 &
\centering 4 &
\centering\arraybslash 3\\\hline
\centering 20th &
\centering 150 &
\centering 102 &
\centering 17 &
\centering 0 &
\centering 27 &
\centering 0 &
\centering 4 &
\centering\arraybslash 13\\\hline
\end{supertabular}
\end{center}
\begin{styleStandard}
Table \stepcounter{Table}{\theTable}: CdE\_cualquiera\_150 occ. per century from “my data”.
\end{styleStandard}

\begin{styleStandard}
In 277, I provide examples for the variables of Group E. Recall here again that I only considered the mood of the main verb, iff \textsc{qualquier} was embedded inside a main clause. Otherwise, I considered the mood of the verb inside the relative clause, as in 272.
\end{styleStandard}

\begin{styleStandard}
\textbf{\ \ }a.\ \ Mood: indicative
\end{styleStandard}

\begin{styleStandard}
\ \ \ \ Los actores le restaron elegancia al juego y \textit{desdeñaron~cualquier}~posibilidad de gol. 
\end{styleStandard}

\begin{styleStandard}
\ \ \ \ ‘The actors downplayed the elegance of the game and discounted any possibility for a goal’
\end{styleStandard}

\begin{styleStandard}
\ \ \ \ (20\textsuperscript{th} c\textit{.},\textit{ CR:PrLibre:98May4})
\end{styleStandard}

\begin{styleStandard}
\ \ b.\ \ Mood: subjunctive
\end{styleStandard}

\begin{styleStandard}
\ \ \ \ E tan grandes eran las bozes delos hombres \& del ganado que \textit{qualquier} que lo \textit{viesse} hauria ende gran dolor \& muy gran piedad. 
\end{styleStandard}

\begin{styleStandard}
\ \ \ \ ‘and the voices of the men and of the cattle were so great that whoever (one wants that) s/he sees it would have great pain and very great pity…,’
\end{styleStandard}

\begin{styleStandard}
\ \ \ \ (13\textsuperscript{th} c., Anonymous, \textit{Gran Conquista})
\end{styleStandard}

\begin{styleStandard}
\ \ c.\ \ Imperative
\end{styleStandard}

\begin{styleStandard}
\ \ \ \ Y Él respondióle: -\textit{Demanda qualquier} cosa que todo lo otorgaré. 
\end{styleStandard}

\begin{styleStandard}
\ \ \ \ ‘And He answered him: Ask anything and I will give it to you.’
\end{styleStandard}

\begin{styleStandard}
\ \ \ \ (13th c., Anonymous, \textit{Siete sabios})
\end{styleStandard}

\begin{styleStandard}
\ \ d.\ \ Infinitive
\end{styleStandard}

\begin{styleStandard}
\ \ \ \ … y que hemos dicho poderse \textit{hallar} en~\textit{cualquiera}~de los tres estilos. 
\end{styleStandard}

\begin{styleStandard}
\ \ \ \ ‘… and of which we have said that it is possible to \textit{find} it in \textit{any} of three styles’
\end{styleStandard}

\begin{styleStandard}
\ \ \ \ (18th c., Luzán,\textit{ La poética})
\end{styleStandard}

\begin{styleStandard}
\ \ e.\ \ Gerund
\end{styleStandard}

\begin{styleStandard}
\ \ \ \ […] o en Yahaurpampa, donde hubo la victoria de los Chancas, \textit{siendo~cualquiera}~de aquellos dos puestos más a propósito que el de Cacha […] 
\end{styleStandard}

\begin{styleStandard}
\ \ \ \ ‘or in Yahaurpampa, where there was the victory of the Chancas, any of those two (building) sites being more appropriate than the one in Cacha…’
\end{styleStandard}

\begin{styleStandard}
\ \ \ \ (16th c.,\textit{ }Garcilaso de la Vega,\textit{ Comentarios reales})
\end{styleStandard}

\begin{styleStandard}
\ \ f.\ \ Participle
\end{styleStandard}

\begin{styleStandard}
\ \ \ \ Y si de~\textit{qualquiera}~pasion \textit{enpedidos} se hallan no sentencian en nada fasta ver se libres. 
\end{styleStandard}

\begin{styleStandard}
\ \ \ \ ‘And if they find themselves hindered by any passion, they do not judge until they are free (from this passion).’
\end{styleStandard}

\begin{styleStandard}
\ \ \ \ (15th c., Diego de San Pedro,\textit{ Cárcel})
\end{styleStandard}

\begin{styleStandard}
\ \ g.\ \ Perfect
\end{styleStandard}

\begin{styleStandard}
\ \ \ \ \textit{Han triturado~cualquier}~intento de libertad intelectual de forma totalmente despiadada. 
\end{styleStandard}

\begin{styleStandard}
\ \ \ \ ‘They have destroyed any attempt of intellectual liberty in a totally ruthless way.’
\end{styleStandard}

\begin{styleStandard}
\ \ \ \ (20th c., Mortimer,\textit{ Entrevista (ABC)})
\end{styleStandard}

\begin{styleStandard}
The results for group F are as follows:
\end{styleStandard}

\begin{center}
\tablefirsthead{}
\tablehead{}
\tabletail{}
\tablelasttail{}
\begin{supertabular}{|m{0.61635983in}|m{0.42685983in}|m{0.8163598in}|m{0.8163598in}|m{0.8163598in}|m{0.8163598in}|m{0.8170598in}|}
\hhline{~~-----}
\multicolumn{1}{m{0.61635983in}}{} &
 &
\multicolumn{5}{m{4.39746in}|}{\centering \textbf{F}}\\\hhline{~~-----}
\multicolumn{1}{m{0.61635983in}}{} &
 &
\multicolumn{5}{m{4.39746in}|}{\centering \textbf{Sentence Type/Mood}}\\\hline
\centering \textbf{Century} &
\centering \textbf{Occ} &
\centering \textbf{Subord.} &
\centering \textbf{Main} &
\centering \textbf{Decl} &
\centering \textbf{Q }\textbf{\textit{que}} &
\centering\arraybslash \textbf{Other Types}\\\hline
\centering 13th &
\centering 150 &
\centering 146 &
\centering 4 &
\centering 4 &
\centering 127 &
\centering\arraybslash 19\\\hline
\centering 14th &
\centering 142 &
\centering 111 &
\centering 31 &
\centering 37 &
\centering 75 &
\centering\arraybslash 30\\\hline
\centering 15th &
\centering 150 &
\centering 122 &
\centering 28 &
\centering 37 &
\centering 67 &
\centering\arraybslash 46\\\hline
\centering 16th &
\centering 150 &
\centering 109 &
\centering 41 &
\centering 48 &
\centering 63 &
\centering\arraybslash 39\\\hline
\centering 17th &
\centering 150 &
\centering 96 &
\centering 54 &
\centering 59 &
\centering 33 &
\centering\arraybslash 58\\\hline
\centering 18th &
\centering 149 &
\centering 115 &
\centering 34 &
\centering 39 &
\centering 53 &
\centering\arraybslash 57\\\hline
\centering 19th &
\centering 150 &
\centering 98 &
\centering 52 &
\centering 56 &
\centering 37 &
\centering\arraybslash 57\\\hline
\centering 20th &
\centering 150 &
\centering 89 &
\centering 61 &
\centering 80 &
\centering 15 &
\centering\arraybslash 55\\\hline
\end{supertabular}
\end{center}
\begin{styleStandard}
Table \stepcounter{Table}{\theTable}: CDE\_cualquiera\_150 occ. per century, “my data” (3).
\end{styleStandard}

\begin{styleStandard}
In 278, I provide examples for the variables of Group F:
\end{styleStandard}

\begin{styleStandard}
\textbf{\ \ }a.\ \ Subordinate (Subord.)
\end{styleStandard}

\begin{styleStandard}
\ \ \ \ \& diz que pero que eran tan grandes principes \textit{que qualquier} dellos manternie una huerste yl darie recabdo que en aquella batalla murieron daquella guisa. 
\end{styleStandard}

\begin{styleStandard}
\ \ \ \ ‘and it is said that because they were such important princes that \textit{any} of them could maintain an army, and he would inform (them) that in that battle they died in that manner.’
\end{styleStandard}

\begin{styleStandard}
\ \ \ \ (13th c., Afonso X, \textit{General Estoria IV})
\end{styleStandard}

\begin{styleStandard}
\ \ b.\ \ Main
\end{styleStandard}

\begin{styleStandard}
\ \ \ \ El término pájaro se aplica a \textit{cualquier} ave con capacidad para volar y de pequeño tamaño. 
\end{styleStandard}

\begin{styleStandard}
\ \ \ \ ‘The term bird can be applied to \textit{any bird} able to fly and of small size.’
\end{styleStandard}

\begin{styleStandard}
\ \ \ \ (20th c., \textit{Enc: Ave})
\end{styleStandard}

\begin{styleStandard}
\ \ c.\ \ Declarative (Decl)
\end{styleStandard}

\begin{styleStandard}
\ \ \ \ El término pájaro se aplica a \textit{cualquier} ave con capacidad para volar y de pequeño tamaño. 
\end{styleStandard}

\begin{styleStandard}
\ \ \ \ ‘The term bird is applied to \textit{any bird} able to fly and of small size.’
\end{styleStandard}

\begin{styleStandard}
\ \ \ \ (20th c., \textit{Enc: Ave})
\end{styleStandard}

\begin{styleStandard}
\ \ d.\ \ Q que
\end{styleStandard}

\begin{styleStandard}
\ \ \ \ assi como que pechen omezillo / o otra pena \textit{qualquier que} viere el alcalde que sera guisado. 
\end{styleStandard}

\begin{styleStandard}
\ \ \ \ ‘just as they pay for homicide / or (receive) some other punishment whatever the mayor sees as appropriate.’
\end{styleStandard}

\begin{styleStandard}
\ \ \ \ (14th c., Anonymous, \textit{Leyes de estilo})
\end{styleStandard}

\begin{styleStandard}
\ \ e.\ \ other types
\end{styleStandard}

\begin{styleStandard}
\ \ \ \ … pero quería ser rogado, \textit{para que} no se le imputase jamás por traición \textit{cualquier} siniestro acaecimiento, sino por desgracia. 
\end{styleStandard}

\begin{styleStandard}
\ \ \ \ ‘but he wanted to be asked (to do it), so that any unfortunate occurrence would ever be attributed to him for treason, but for misfortune.’
\end{styleStandard}

\begin{styleStandard}
\ \ \ \ (18th c., Bacallar y Sanna,\textit{ Comentarios})
\end{styleStandard}

\begin{styleStandard}
After this introduction into the data classification, I will present a short introduction of statistical tests I used.
\end{styleStandard}

\subsubsection[5.1.3 Statistical tests and basic notions]{\textbf{5.1.3 Statistical tests and basic notions}}
\begin{styleStandard}
The aim of the statistical tests in this study is to compare frequencies of different constructions with a single period or to compare frequencies of a single construction in different periods. As these frequencies are never numerically identical, statistical tests are required to determine whether the difference in frequency is statistically significant. One possibility for measuring frequency differences is to use the \textit{chi square (}$\chi $²\textit{)} test. This test compares a set of observed frequencies O with a set of expected frequencies E (see Baayen 2008: 77, Author 2018: 64 and references therein): $\chi $² = sum of (O–E)\textsuperscript{2} / E. Expected frequencies (E) are frequencies that we would obtain if the frequency difference was not statistically significant. In order to calculate them, we multiply the sum of observed frequencies and divide them by total number.
\end{styleStandard}

\begin{styleStandard}
\ \ Let us imagine the following linguistic scenario. In the first period (Period 1), Construction 1 occurs 7 times and Construction 2 is observed 12 times. Moreover, in the second period (Period 2), Construction 1is encountered 13 times and Construction 2 occurs 8 times:
\end{styleStandard}

\begin{center}
\tablefirsthead{}
\tablehead{}
\tabletail{}
\tablelasttail{}
\begin{supertabular}{|m{1.8441598in}|m{2.0198598in}|m{2.0191598in}|}
\hhline{~--}
\multicolumn{1}{m{1.8441598in}|}{} &
\centering \textbf{Construction 1} &
\centering\arraybslash \textbf{Construction 2}\\\hline
\centering \textbf{Period 1} &
\centering 7 &
\centering\arraybslash 12\\\hline
\centering \textbf{Period 2} &
\centering 13 &
\centering\arraybslash 8\\\hline
\end{supertabular}
\end{center}
\begin{stylecaption}
\textmd{Table \stepcounter{Table}{\theTable}: Hypothetical data for demonstration of $\chi $².}
\end{stylecaption}

\begin{styleStandard}
We arrive at the expected frequency by calculating expected numbers for each individual cell (i.e. the frequencies we would obtain if there were no association between two variables, that is, between Construction and Period). We do this by multiplying the row sum by the column sum and dividing by the total number. The expected frequencies are given in brackets:
\end{styleStandard}

\begin{center}
\tablefirsthead{}
\tablehead{}
\tabletail{}
\tablelasttail{}
\begin{supertabular}{|m{0.8275598in}|m{1.2990599in}|m{1.2997599in}|m{1.2997599in}|}
\hhline{~---}
\multicolumn{1}{m{0.8275598in}|}{} &
\centering \textbf{Construction 1} &
\centering \textbf{Construction 2} &
\centering\arraybslash \textbf{Total}\\\hline
\textbf{Period 1} &
{\centering 7\par}

\centering (9.5) &
{\centering 12\par}

\centering (9.5) &
\centering\arraybslash 19\\\hline
\textbf{Period 2} &
{\centering 13\par}

\centering (10.5) &
{\centering 8\par}

\centering (10.5) &
\centering\arraybslash 21\\\hline
\textbf{Total} &
\centering 20 &
\centering 20 &
\centering\arraybslash Grand total=40\\\hline
\end{supertabular}
\end{center}
\begin{stylecaption}
\textmd{Table \stepcounter{Table}{\theTable}: Expected and observed frequencies of hypothetical data. Bracketed numbers indicate expected frequencies.}
\end{stylecaption}

\begin{styleStandard}
Now we calculate the formula of $\chi $² : (7-9.5)²/9.5 = 0.6579. We continue to do this for each cell and add all the numbers. The figure we obtain is 2.5062, because 0.6579 + 0.6579 + 0.5952 + 0.5952 = 2.5062. We now calculate the degrees of freedom: (Number of Rows–1) × (Number of Columns–1)(2–1) × (2–1) 1×1=1 (degrees of freedom).
\end{styleStandard}

\begin{styleStandard}
\ \ If the difference between observed and expected frequencies is large enough, we can reject the null hypothesis H\textsubscript{0}, which states that the frequency difference is not relevant (the so-called \textit{Hypothesis of Independence}). The alternative hypothesis to hypothesis H\textsubscript{0} is H\textsubscript{1}, which states that the frequency difference is relevant (see Author 2018: 64 and references therein): 
\end{styleStandard}

\begin{styleStandard}
\ \ H\textsubscript{0 }: $\chi $² = 0 H\textsubscript{1 }: x\textsuperscript{2 }{\textgreater} 0
\end{styleStandard}

\begin{styleStandard}
\ \ The $\chi $² value for 1 degree of freedom must be equal or greater than 3.841, if we assume that it is 5\% likely that H\textsubscript{0} is true (see 5.2 for the application of such a test) (see Baayen 2008: 77, Author 2018: 64 and references therein). If the obtained \textit{chi square} ($\chi $²) is larger than 3.841, it implies that it is unlikely that our observed two frequencies have occurred by chance. The value for our example is 2.5062, which is below the critical value of 3.841. In other words, the two constructions (Construction 1 and Construction 2) are similar with respect to their frequency across time periods. 
\end{styleStandard}

\begin{styleStandard}
\ \ Another important notion of statistics used in this thesis, is the notion of \textit{standard deviation} (SD). The standard deviation measures the variation between a set of values (such as a set of observed frequencies of different constructions or one construction across time). A low standard deviation indicates that the values tend to be close to the mean (also called the expected value) of the set, while a high standard deviation indicates that the values are spread out over a wider range.
\end{styleStandard}

\begin{styleStandard}
\ \ Let us assume we have a normal distribution of data, also called a Gaussian distribution, which has a shape of a bell with a peak at the average value (mean). The standard deviation is a measure of the data that disperse around the mean or the peak. In the following representation, about 68\% of the data points fall within 1 standard deviation of the mean, about 95\% of the data points fall within 2 standard deviations of the mean and about 99\% of the data points fall within 3 standard deviations of the mean. Note that only 3 out of 1,000 points will fall outside the area of 3 standard deviations. The larger the number of the standard deviation, the wider the spread is of the data around the mean. 
\end{styleStandard}

\begin{stylecaption}
\textmd{Figure \stepcounter{Figure}{\theFigure}: Normal distribution in }\url{http://stevegallik.org/cellbiologyolm_statistics.html}\textstyleInternetlink{\textmd{.}}
\end{stylecaption}

\begin{styleStandard}
I also used \textit{Linear regression}, which is a model that can be used to identify the relationship or correlation between different variables by fitting a line to the observed data that best fits the data. 
\end{styleStandard}

\begin{styleStandard}
\ \ Figure 16 shows this for an example from outside linguistics: the relation between cholesterol and body mass index (BMI). It shows that cholesterol and BMI correlate: the higher the BMI, the higher the cholesterol. The graph below shows the estimated regression line superimposed on data points represented as a scatter diagram:
\end{styleStandard}

\begin{stylecaption}
\textmd{Figure \stepcounter{Figure}{\theFigure}: The linear model for the relationship of Cholesterol and body mass index (BMI) in https://sphweb.bumc.bu.edu/otlt/mph-modules/bs/bs704\_correlation-regression/bs704\_correlation-regression\_print.html, accessed November 2020.}
\end{stylecaption}

\begin{styleStandard}
The linear regression can be used in linguistics in order to test whether there is a linear correlation between the frequency of some linguistic construction and time or period, such as to test whether the frequency of postnominal \textsc{qualquier} changes across time or periods (see 5.2.1 for such a test).
\end{styleStandard}

\subsection[Quantitative description]{Quantitative description}
\begin{styleStandard}
We can observe major quantitative differences in time with variables in almost every group, especially if we look at the frequency differences between the 13th century (1200–1300) and the 20th century (1900–2000). The following table contains variables that show a decrease in time. I describe them from left to right: 
\end{styleStandard}

\begin{listWWNumxvleveli}
\item 
\begin{styleStandard}
The subject function of \textsc{qualquier} decreases.
\end{styleStandard}
\item 
\begin{styleStandard}
The use of bare \textsc{qualquier} (bare Q) without nominal modification decreases.
\end{styleStandard}
\item 
\begin{styleStandard}
The preverbal position (QV) decreases. 
\end{styleStandard}
\item 
\begin{styleStandard}
The use of \textsc{qualquier} in clauses with subjunctive (Sbjv) decreases. 
\end{styleStandard}
\item 
\begin{styleStandard}
The use of \textsc{qualquier} in subordinated clauses decreases.
\end{styleStandard}
\item 
\begin{styleStandard}
The use of \textsc{qualquier}\textit{ que} (Q \textit{que}) decreases.
\end{styleStandard}
\end{listWWNumxvleveli}
\begin{center}
\tablefirsthead{}
\tablehead{}
\tabletail{}
\tablelasttail{}
\begin{supertabular}{|m{0.7698598in}|m{0.6108598in}|m{0.6122598in}|m{0.6115598in}|m{0.6108598in}|m{0.6115598in}|m{0.6122598in}|m{0.6080598in}|}
\hline
\centering \textbf{Century} &
\centering \textbf{Occ} &
\centering \textbf{Subject} &
\centering \textbf{bare Q} &
\centering \textbf{QV} &
\centering \textbf{Sbjv} &
{\centering \textbf{Subord.}\par}

\centering \textbf{clause} &
\centering\arraybslash \textbf{Q }\textbf{\textit{que}}\\\hline
\centering 13th &
\centering 150 &
\centering 94 &
\centering 102 &
\centering 134 &
\centering 120 &
\centering 146 &
\centering\arraybslash 127\\\hline
 &
\centering 100\% &
\centering 62.67\% &
\centering 68.00\% &
\centering 89.33\% &
\centering 80.00\% &
\centering 97.33\% &
\centering\arraybslash 84.67\%\\\hline
\centering 14th &
\centering 142 &
\centering 70 &
\centering 62 &
\centering 118 &
\centering 74 &
\centering 111 &
\centering\arraybslash 75\\\hline
 &
\centering 100\% &
\centering 49.30\% &
\centering 43.66\% &
\centering 83.10\% &
\centering 52.11\% &
\centering 78.17\% &
\centering\arraybslash 52.82\%\\\hline
\centering 15th &
\centering 150 &
\centering 27 &
\centering 33 &
\centering 102 &
\centering 66 &
\centering 122 &
\centering\arraybslash 67\\\hline
 &
\centering 100\% &
\centering 18.00\% &
\centering 22.00\% &
\centering 68.00\% &
\centering 44.00\% &
\centering 81.33\% &
\centering\arraybslash 44.67\%\\\hline
\centering 16th &
\centering 150 &
\centering 35 &
\centering 37 &
\centering 98 &
\centering 45 &
\centering 109 &
\centering\arraybslash 63\\\hline
 &
\centering 100\% &
\centering 23.33\% &
\centering 24.67\% &
\centering 65.33\% &
\centering 30.00\% &
\centering 72.67\% &
\centering\arraybslash 42.00\%\\\hline
\centering 17th &
\centering 150 &
\centering 32 &
\centering 38 &
\centering 83 &
\centering 33 &
\centering 96 &
\centering\arraybslash 33\\\hline
 &
\centering 100\% &
\centering 21.33\% &
\centering 25.33\% &
\centering 55.33\% &
\centering 22.00\% &
\centering 64.00\% &
\centering\arraybslash 22.00\%\\\hline
\centering 18th &
\centering 150 &
\centering 29 &
\centering 41 &
\centering 91 &
\centering 49 &
\centering 115 &
\centering\arraybslash 53\\\hline
 &
\centering 100\% &
\centering 19.46\% &
\centering 27.52\% &
\centering 61.07\% &
\centering 32.89\% &
\centering 77.18\% &
\centering\arraybslash 35.57\%\\\hline
\centering 19th &
\centering 149 &
\centering 24 &
\centering 55 &
\centering 71 &
\centering 51 &
\centering 98 &
\centering\arraybslash 37\\\hline
 &
\centering 100\% &
\centering 16.00\% &
\centering 36.67\% &
\centering 47.33\% &
\centering 34.00\% &
\centering 65.33\% &
\centering\arraybslash 24.67\%\\\hline
\centering 20th &
\centering 150 &
\centering 25 &
\centering 23 &
\centering 54 &
\centering 17 &
\centering 89 &
\centering\arraybslash 15\\\hline
 &
\centering 100\% &
\centering 16.67\% &
\centering 15.33\% &
\centering 36.00\% &
\centering 11.33\% &
\centering 59.33\% &
\centering\arraybslash 10.00\%\\\hline
\end{supertabular}
\end{center}
\begin{styleStandard}
Table \stepcounter{Table}{\theTable}: Decrease of frequency in “my data” with respect to variables: subject, bare Q without nominal modification, preverbal, subjunctive, subordinate clause, Q \textit{que}.
\end{styleStandard}

\begin{styleStandard}
Figure 8 shows the decrease in the frequency of these variables across time visually:
\end{styleStandard}

\begin{styleStandard}
Figure \stepcounter{Figure}{\theFigure}: Decrease of frequency in “my data” with respect to the frequency of the variables: subject, bare Q without modification, preverbal, subjunctive, subordinate clause, Q \textit{que}.
\end{styleStandard}

\begin{styleStandard}
Although the decrease is not linear as the frequency of some variables increases in the 15th and 17th centuries (e.g. QV decreases until 17th c., increases in the 18th, and decreases again from the 19th\textsuperscript{ }onwards), we can see an overall decrease from the 13th to the 20th century. This overall decrease is the focus of my statistical and linguistic analysis (see Sections 5.2.1 and 5.3). It should be noted at this point that the interruption of the decrease in the data in the 15th and 17th centuries coincides with the importance of the 15th and 17th centuries in the history of Spanish language, which is subdivided into the 13th, 15th, 17th, and 19th centuries (see 5.2.3). The changes in frequencies in these centuries in Figure 8 may reflect the grammatical changes in these centuries. The separated form \textit{qual quier,}\nobreakdash-\textit{e,}\nobreakdash-\textit{a} disappears from the 16th century onwards (Pato 2012, Chapter 2.6) and that the 15th century can be considered the peak century when \textit{qualquier,-e,-a} grammaticalises completely. The changes in frequency around this peak century might also reflect the grammaticalisation of \textit{qualquier,} \nobreakdash-\textit{e,} \nobreakdash-\textit{a}. I will return to this hypothesis below when I discuss different variables.
\end{styleStandard}

\begin{styleStandard}
\ \ Let us look at some prototypical sentence from the 13th century that contains variables that decrease in their frequency in subsequent centuries. 
\end{styleStandard}

\begin{styleStandard}
\textbf{\ \ }E tan grandes eran las bozes delos hombres \& del ganado que \textit{qualquier que} lo \textit{viesse} hauria ende gran dolor \& muy gran piedad. 
\end{styleStandard}

\begin{styleStandard}
\ \ ‘and the voices of the men and of the cattle were so great that whoever (one wants that) s/he sees it would have great pain and very great pity.’
\end{styleStandard}

\begin{styleStandard}
\ \ (13th c., Anonymous,\textit{ Gran Conquista})
\end{styleStandard}

\begin{styleStandard}
The example in 279 as the following structure. The relative clause \textit{qualquier que} \textit{lo viesse }is the subject of the main verbal phrase \textit{hauria ende gran dolor \& muy gran piedad} and, inside the relative clause headed by the complementiser \textit{que}, \textit{qualquier} is the subject of the verb \textit{viesse}. \textit{Qualquier} does not modify any overt noun, and it is used as a pronoun that introduces free relative clauses similar to \textit{quien} ‘who’ or \textit{el que} ‘(the one) who’; thus, the word order is QV. The verb \textit{viesse} inside the relative clause is used in the subjunctive. The subjunctive has a non-veridical interpretation, that is, one that lacks the presupposition that the proposition of the sentence with subjunctive is true (see Quer 1998, 2000, Becker 2014, and Chapter 2). The subjunctive does not allow a reference to a certain event, but it includes all possible situations in which someone has seen it.
\end{styleStandard}

\begin{styleStandard}
\ \ The following example is similar, with \textit{qualquier} being used as a subject of the complement clause \textit{que le demandaua}. The main difference between 279 and 280 is that the verb of the complement clause is not in the subjunctive but in the indicative. The verb aspect of the complement clause is a durative aspect represented by imperfective form \textit{demandaua}, which has a generic interpretation in this example, due not only to the imperfective but also to the generic adverb \textit{siempre} ‘always’ (see Quer 2000, and Chapter 2 on the generic interpretation of imperfective aspect):
\end{styleStandard}

\begin{styleStandard}
\textbf{\ \ }siempre \ \ daua\ \ a\ \ cualquier\ \ que\ \ le\ \ demandaua.
\end{styleStandard}

\begin{styleStandard}
\ \ always\ \ gave\ \ to\ \ whoever\ \ that\ \ him\ \ asked
\end{styleStandard}

\begin{styleStandard}
\ \ ‘He always gave it to whoever asked him. = Every time someone asked him, he gave it to him.’
\end{styleStandard}

\begin{styleStandard}
\ \ (15th c., Fernando del Pulgar, \textit{Claros varones})
\end{styleStandard}

\begin{styleStandard}
The imperfective verb expressed in the embedded sentence in 280 has a function similar to that in 279; that is, there is no reference to a particular event of someone giving something to someone. Instead, the imperfective verb together with the generic adverb \textit{siempre} ‘always’ expresses all possible events in the past of someone giving something to someone else when asked (see Becker 2014).
\end{styleStandard}

\begin{styleStandard}
\ \ I will argue in Section 5.3 that the decrease of the subjunctive inside the relative clause is related to the degree of grammaticalisation of \textsc{qualquier}. The more transparent \textsc{qualquier} is as a compound containing a wh-word \textit{qual} and the verb \textit{querer}, the more cases we find with \textsc{qualquier} que + V\textsubscript{[subj]}. The less transparent \textit{quier} is as a verb, the fewer cases we find of the type \textit{cualquiera que…} and for that matter, the more cases we find without the subjunctive. The abrupt fall of the subjunctive in the 14th century coincides with the decrease of the separated form \textit{qual quier} in the 14th century and its disappearance in the 15th.
\end{styleStandard}

\begin{styleStandard}
\ \ Let us now discuss other changes. \textsc{qualquier} in relative clauses such as\textit{ }\textsc{qualquier} \textit{que…} can only be realised in a preverbal position with respect to the embedded verb [Main clause … (V\textsubscript{main}) … [\textsc{qualquier}\textit{ que} … V\textsubscript{embed}] … (V\textsubscript{main}) …]. This follows from the analysis of \textsc{qualquier} as a specifier of a relative clause (see, e.g., 272 in this section).
\end{styleStandard}

\begin{styleStandard}
\ \ Another prediction of the Grammaticalisation Hypothesis is that once the string \textit{qual quier} is no longer analysed compositionally as a wh-word and a verb, we should observe an increase of \textsc{qualquier} in sentence types other than relative clauses and the use of \textsc{qualquier} with a verbal mood other than the subjunctive. This is indeed what can be observed in the data in Table 18, which shows that all variables (i.e. postverbal position (VQ), indicative mood and infinitives, perfective aspect, main clause, declaratives, other sentence types, QN) tend to increase in time, especially if we compare the 13th century with the 20th.
\end{styleStandard}

\begin{center}
\tablefirsthead{}
\tablehead{}
\tabletail{}
\tablelasttail{}
\begin{supertabular}{|m{0.6108598in}|m{0.31435984in}|m{0.51155984in}|m{0.41365984in}|m{0.41365984in}|m{0.51155984in}|m{0.51225984in}|m{0.41295984in}|m{0.6108598in}|m{0.41295984in}|}
\hline
\centering \textbf{Century} &
\centering \textbf{Occ} &
\centering \textbf{VQ} &
\centering \textbf{Ind} &
\centering \textbf{Inf} &
\centering \textbf{Perf} &
\centering \textbf{Main} &
\centering \textbf{Decl} &
\centering \textbf{Other} &
\centering\arraybslash \textbf{QN}\\\hline
 &
 &
 &
 &
 &
 &
 &
 &
 &
\\\hline
\centering 13th &
\centering \textbf{150} &
\centering \textbf{16} &
\centering \textbf{26} &
\centering \textbf{1} &
\centering \textbf{2} &
\centering \textbf{4} &
\centering \textbf{4} &
\centering \textbf{19} &
\centering\arraybslash \textbf{43}\\\hline
 &
\centering ~ &
\raggedleft 10.7\% &
\raggedleft 17.3\% &
\raggedleft 0.7\% &
\raggedleft 1.3\% &
\raggedleft 2.7\% &
\raggedleft 2.7\% &
\raggedleft 12.7\% &
\raggedleft\arraybslash 28.7\%\\\hline
\centering 14th &
\centering \textbf{142} &
\centering \textbf{24} &
\centering \textbf{61} &
\centering \textbf{5} &
\centering \textbf{1} &
\centering \textbf{31} &
\centering \textbf{37} &
\centering \textbf{30} &
\centering\arraybslash \textbf{78}\\\hline
 &
\centering ~ &
\raggedleft 16.9\% &
\raggedleft 43.0\% &
\raggedleft 3.5\% &
\raggedleft 0.7\% &
\raggedleft 21.8\% &
\raggedleft 26.1\% &
\raggedleft 21.1\% &
\raggedleft\arraybslash 54.9\%\\\hline
\centering 15th &
\centering \textbf{150} &
\centering \textbf{48} &
\centering \textbf{60} &
\centering \textbf{16} &
\centering \textbf{2} &
\centering \textbf{28} &
\centering \textbf{37} &
\centering \textbf{46} &
\centering\arraybslash \textbf{116}\\\hline
 &
\centering ~ &
\raggedleft 32.0\% &
\raggedleft 40.0\% &
\raggedleft 10.7\% &
\raggedleft 1.3\% &
\raggedleft 18.7\% &
\raggedleft 24.7\% &
\raggedleft 30.7\% &
\raggedleft\arraybslash 77.3\%\\\hline
\centering 16th &
\centering \textbf{150} &
\centering \textbf{52} &
\centering \textbf{89} &
\centering \textbf{12} &
\centering \textbf{5} &
\centering \textbf{41} &
\centering \textbf{48} &
\centering \textbf{39} &
\centering\arraybslash \textbf{113}\\\hline
 &
\centering ~ &
\raggedleft 34.7\% &
\raggedleft 59.3\% &
\raggedleft 8.0\% &
\raggedleft 3.3\% &
\raggedleft 27.3\% &
\raggedleft 32.0\% &
\raggedleft 26.0\% &
\raggedleft\arraybslash 75.3\%\\\hline
\centering 17th &
\centering \textbf{150} &
\centering \textbf{67} &
\centering \textbf{99} &
\centering \textbf{12} &
\centering \textbf{3} &
\centering \textbf{54} &
\centering \textbf{59} &
\centering \textbf{58} &
\centering\arraybslash \textbf{111}\\\hline
 &
\centering ~ &
\raggedleft 44.7\% &
\raggedleft 66.0\% &
\raggedleft 8.0\% &
\raggedleft 2.0\% &
\raggedleft 36.0\% &
\raggedleft 39.3\% &
\raggedleft 38.7\% &
\raggedleft\arraybslash 74.0\%\\\hline
\centering 18th &
\centering \textbf{150} &
\centering \textbf{58} &
\centering \textbf{64} &
\centering \textbf{29} &
\centering \textbf{2} &
\centering \textbf{34} &
\centering \textbf{39} &
\centering \textbf{57} &
\centering\arraybslash \textbf{106}\\\hline
 &
\centering ~ &
\raggedleft 38.7\% &
\raggedleft 42.7\% &
\raggedleft 19.3\% &
\raggedleft 1.3\% &
\raggedleft 22.7\% &
\raggedleft 26.0\% &
\raggedleft 38.0\% &
\raggedleft\arraybslash 70.7\%\\\hline
\centering 19th &
\centering \textbf{149} &
\centering \textbf{79} &
\centering \textbf{74} &
\centering \textbf{18} &
\centering \textbf{3} &
\centering \textbf{52} &
\centering \textbf{56} &
\centering \textbf{57} &
\centering\arraybslash \textbf{81}\\\hline
 &
\centering ~ &
\raggedleft 53.0\% &
\raggedleft 49.7\% &
\raggedleft 12.1\% &
\raggedleft 2.0\% &
\raggedleft 34.9\% &
\raggedleft 37.6\% &
\raggedleft 38.3\% &
\raggedleft\arraybslash 54.4\%\\\hline
\centering 20th &
\centering \textbf{150} &
\centering \textbf{96} &
\centering \textbf{102} &
\centering \textbf{27} &
\centering \textbf{13} &
\centering \textbf{61} &
\centering \textbf{80} &
\centering \textbf{55} &
\centering\arraybslash \textbf{123}\\\hline
 &
~ &
\raggedleft 64.0\% &
\raggedleft 68.0\% &
\raggedleft 18.0\% &
\raggedleft 8.7\% &
\raggedleft 40.7\% &
\raggedleft 53.3\% &
\raggedleft 36.7\% &
\raggedleft\arraybslash 82.0\%\\\hline
\end{supertabular}
\end{center}
\begin{styleStandard}
Table \stepcounter{Table}{\theTable}: Increase of frequency in “my data” with respect to variables: VQ, ind, inf, perf, main, decl, other sentence types, QN.
\end{styleStandard}

\begin{styleStandard}
Table 18 shows the increase in the frequency of these variables across time visually:
\end{styleStandard}

\begin{styleStandard}
Figure \stepcounter{Figure}{\theFigure}: Increase of frequency in “my data” with respect to variables: VQ, ind, inf, perf, main, decl, other sentence types, QN.
\end{styleStandard}

\begin{styleStandard}
As with decreasing frequency (see Figure 9), the 15th and 17th centuries also stand out with respect to the increasing frequency. The frequency of QN increases strongly until the 15th century. The strong increase of QN until the 15th century can be correlated with the observation that the separated form \textit{qual quier,-e,-a} disappears after this century (see Section 5.2.1), which can be interpreted as an indication for completed grammaticalisation. The grammaticalisation of \textit{qual quier,-e,-a} until the 15th century correlates with the strong increase of the determiner \textit{qualquier} from the 13th to the 15th centuries.
\end{styleStandard}

\begin{styleStandard}
\ \ Other variables (e.g. main clause, VQ, and declarative sentence type) increase until the 17th century, with an interruption in the 18th. The grammaticalisation of \textit{qualquier} correlates with the increase of different sentence types as an indefinite form no longer is a relative clause and thus no longer is restricted to a single sentence type.
\end{styleStandard}

\begin{styleStandard}
\ \ Two prototypical examples that illustrate the increased frequency values of these variables are the following sentences with an imperative clause and a declarative clause. In both clauses, \textsc{qualquier} is followed by a noun (\textit{cosa} and \textit{varón}, respectively, with which it forms a single constituent). It is used in postverbal position with indicative verbal mood in a main clause and not introducing a relative or complement clause from which \textsc{qualquier} would be structurally dependent. Note that \textit{qualquier cosa} in 281a. is not an argument of the verb \textit{otorgaré} in the clause introduced by\textit{ que}, but it is the argument of the verb in the main clause:
\end{styleStandard}

\begin{styleStandard}
\textbf{\ \ }a.\ \ Y él respondióle: -\textit{Demanda} \textit{qualquier }cosa que\footnote{\ This \textit{que} corresponds to what Koch \& Oesterreicher (1990: 99) call the ‘polyvalent \textit{que}’, which does not have a clearly subordinating function.} todo lo otorgaré. 
\end{styleStandard}

\begin{styleStandard}
{}-‘Ask anything and I will give it to you.’
\end{styleStandard}

\begin{styleStandard}
\textbf{\ \ }\ \ (13th c., Anonymous, \textit{Siete sabios})
\end{styleStandard}

\begin{styleStandard}
\ \ b.\ \  Por ende bien \textit{se avise} desto \textit{qualquier} varón. 
\end{styleStandard}

\begin{styleStandard}
\ \ \ \ ‘Therefore any man should get well informed about this.’
\end{styleStandard}

\begin{styleStandard}
\ \ \ \ (14th c., López de Ayala, \textit{Libro rimado})
\end{styleStandard}

\begin{styleStandard}
If the frequency change of increasing variables across time in Figure 9 is statistically significant, something that needs still to be shown in the next section, we can conclude that the use of \textit{cualquiera} in relative clauses became less frequent, whereas its use in main clauses became more frequent over time:
\end{styleStandard}

\begin{styleStandard}
\textbf{\ \ }[MainClause ….V... [\textsc{qualquier}\textit{ que} V\textsubscript{[subj.]}]]] = (most frequent in Old Sp., 13th c.)
\end{styleStandard}

\begin{styleStandard}
\ \ [F0E0?] [Main Clause Verb\textsubscript{[Indic.] }\textsc{qualquier} NP] = (most frequent in 20th-century Sp.)
\end{styleStandard}

\begin{styleStandard}
Let us now look at the variables more closely and test the change of frequency statistically.
\end{styleStandard}

\subsubsection[5.2.1 Statistical analysis of relevant changes]{\textbf{5.2.1 Statistical analysis of relevant changes}}
\begin{styleStandard}
This Section introduces into the statistical analysis of relevant changes.
\end{styleStandard}

\paragraph[5.2.1.1 The decrease of Q–V order and Q que]{\textit{5.2.1.1 The decrease of Q–V order and Q que}}
\begin{styleStandard}
Table 19 shows the frequencies associated with different syntactic positions of \textit{cualquiera.} If we look at the overall frequency, it shows a continuous decrease of the preverbal position and a continuous increase of the postverbal position of\textit{ }\textsc{qualquier} in every century. Looking at the frequencies in every single century, we can observe a bigger decrease between the 14th and 15th centuries. This decrease might be motivated with the overall grammatical change of early Old Spanish into late Old Spanish which is considered to be close to the Modern Spanish (see Ariza 1998, Penny 1993, among others). In particular, Old Spanish allowed several constructions of object preposing, which are, however, still attested in the 15th century, as in 285, and even in the 17th century, as in 287. By contrast, postverbal subjects have always been possible in Spanish, as in 283.
\end{styleStandard}

\begin{center}
\tablefirsthead{}
\tablehead{}
\tabletail{}
\tablelasttail{}
\begin{supertabular}{|m{0.8511598in}|m{0.7351598in}|m{0.8080598in}|m{0.90665984in}|m{0.71085984in}|m{1.0045599in}|}
\hhline{~~----}
\multicolumn{1}{m{0.8511598in}}{} &
 &
\multicolumn{2}{m{1.7934599in}|}{\centering \textbf{Q–V}} &
\multicolumn{2}{m{1.7941599in}|}{\centering \textbf{V–Q}}\\\hline
\centering \textbf{Century} &
\centering \textbf{Occ.} &
\centering \textbf{Nr. of occ.} &
\centering \textbf{Rel. frequency} &
\centering \textbf{Nr. of occ.} &
\centering\arraybslash \textbf{Rel. frequency}\\\hline
\centering 13th &
\centering 150 &
\centering 134 &
\centering 89.33\% &
\centering 16 &
\centering\arraybslash 10.67\%\\\hline
\centering 14th &
\centering 142 &
\centering 118 &
\centering 83.10\% &
\centering 24 &
\centering\arraybslash 16.90\%\\\hline
\centering 15th &
\centering 150 &
\centering 102 &
\centering 68.00\% &
\centering 48 &
\centering\arraybslash 32.00\%\\\hline
\centering 16th &
\centering 150 &
\centering 98 &
\centering 65.33\% &
\centering 52 &
\centering\arraybslash 34.67\%\\\hline
\centering 17th &
\centering 150 &
\centering 83 &
\centering 55.33\% &
\centering 67 &
\centering\arraybslash 44.67\%\\\hline
\centering 18th &
\centering 149 &
\centering 91 &
\centering 61.07\% &
\centering 58 &
\centering\arraybslash 38.93\%\\\hline
\centering 19th &
\centering 150 &
\centering 71 &
\centering 47.33\% &
\centering 79 &
\centering\arraybslash 52.67\%\\\hline
\centering 20th &
\centering 150 &
\centering 54 &
\centering 36.00\% &
\centering 96 &
\centering\arraybslash 64.00\%\\\hline
\end{supertabular}
\end{center}
\begin{styleStandard}
Table \stepcounter{Table}{\theTable}: Distribution of \textit{cualquiera} NP in preverbal (Q–V) or postverbal (V–Q) position in “my data”.
\end{styleStandard}

\begin{styleStandard}
The following examples demonstrate minimal pairs (QV and VQ) in “my data”:
\end{styleStandard}

\begin{styleStandard}
\textbf{\ \ }13th century
\end{styleStandard}

\begin{styleStandard}
\ \ a.\ \ QV\ \ Et~\textit{qualquiere~}que les contra esto \textit{fuesse} a el me tornaria por ello. 
\end{styleStandard}

\begin{styleStandard}
\ \ \ \ \ \ ‘And I would go against whoever would be against it.’
\end{styleStandard}

\begin{styleStandard}
\ \ \ \ \ \ (Alfonso X, \textit{Doc. Cast. – Murcia})
\end{styleStandard}

\begin{styleStandard}
\ \ b.\ \ VQ\ \ \& por la altura del polo quando \textit{fuere} sabudo~\textit{qualquier~dellos}. 
\end{styleStandard}

\begin{styleStandard}
\ \ \ \ \ \ ‘and for the height of the pole, as soon as any of them is known’
\end{styleStandard}

\begin{styleStandard}
\ \ \ \ \ \ (Alfonso X, \textit{Albateni})
\end{styleStandard}

\begin{styleStandard}
\textbf{\ \ }14th century
\end{styleStandard}

\begin{styleStandard}
\ \ a.\ \ QV\ \ ca~\textit{qualquier~}que lo faga, çierto sea de ser mal andante ante que muera. 
\end{styleStandard}

\begin{styleStandard}
\ \ \ \ \ \ ‘because anyone that would do it, would be certainly unfortunate before he dies.’
\end{styleStandard}

\begin{styleStandard}
\ \ \ \ \ \ (Anonymous, \textit{Libro del Caballero Zifar})
\end{styleStandard}

\begin{styleStandard}
\ \ b.\ \ VQ\ \ \& el señor \textit{puede} \textit{a~qualquier~}destos mayordomos ante que se despidan del prender los \& tener los presos: 
\end{styleStandard}

\begin{styleStandard}
\ \ \ \ \ \ ‘and the lord can arrest and imprison any of these stewards before they resign from their services to him:’
\end{styleStandard}

\begin{styleStandard}
\ \ \ \ \ \ (Anonymous, \textit{Leyes del estilo})
\end{styleStandard}

\begin{styleStandard}
\textbf{\ \ }15th century
\end{styleStandard}

\begin{styleStandard}
\ \ a.\ \ QV\ \ Y si \textit{de~qualquiera~pasion} enpedidos se \textit{hallan} no sentencian en nada falta ver se libres. 
\end{styleStandard}

\begin{styleStandard}
\ \ \ \ \ \ ‘And if they find themselves hindered by any passion, they do not judge until they are free (from this passion).’
\end{styleStandard}

\begin{styleStandard}
\ \ \ \ \ \ (Diego de San Pedro, \textit{Cárcel})
\end{styleStandard}

\begin{styleStandard}
\ \ b.\ \ VQ\ \ E todos los que en aquellos tiempos vinieron ale \textit{demandar}~\textit{qualquier~cargo}. 
\end{styleStandard}

\begin{styleStandard}
\ \ \ \ \ \ ‘and all those who, at that time, came to ask him from any (kind of) office’
\end{styleStandard}

\begin{styleStandard}
\ \ \ \ \ \ (Fernando del Pulgar, \textit{Claros varones})
\end{styleStandard}

\begin{styleStandard}
\textbf{\ \ }16th century
\end{styleStandard}

\begin{styleStandard}
\ \ a.\ \ QV\ \ Consideraba~\textit{cualquiera}~ocupaçión que la \textit{podía} estorbar, […]
\end{styleStandard}

\begin{styleStandard}
\ \ \ \ \ \ ‘I considered any activity that could bother her, ...’
\end{styleStandard}

\begin{styleStandard}
\ \ \ \ \ \ (Villalón, \textit{El Crotalón})
\end{styleStandard}

\begin{styleStandard}
\ \ b.\ \ VQ\ \ Un desconfiar de sí y un estar siempre temiendo que \textit{podrá }exceder al mío~\textit{cualquiera~mérito} ajeno 
\end{styleStandard}

\begin{styleStandard}
\ \ \ \ \ \ ‘A lack of trust in oneself and a constant fear that any merit of someone else will exceed my own’
\end{styleStandard}

\begin{styleStandard}
\ \ \ \ \ \ (Juana Inés de la Cruz,\textit{ Inundación Castálida})
\end{styleStandard}

\begin{styleStandard}
\textbf{\ \ }17th century
\end{styleStandard}

\begin{styleStandard}
\ \ a.\ \ QV\ \ De don Álvaro, por fe,~\textit{cualquier~}cosa\textit{ creeré}, en razón de la que vi. 
\end{styleStandard}

\begin{styleStandard}
\ \ \ \ \ \ ‘From don Álvaro, in faith, I would \textit{believe} \textit{anything}, based on what I saw\textit{.}’
\end{styleStandard}

\begin{styleStandard}
\ \ \ \ \ \ (Castro, \textit{Malcasados})
\end{styleStandard}

\begin{styleStandard}
\ \ b.\ \ VQ\ \ y no \textit{basta} \textit{cualquiera} para mandar. 
\end{styleStandard}

\begin{styleStandard}
\ \ \ \ \ \ ‘and not anyone is appropriate to be in charge of command.’
\end{styleStandard}

\begin{styleStandard}
\ \ \ \ \ \ (Mira de Amescua, \textit{Examinarse de rey})
\end{styleStandard}

\begin{styleStandard}
\textbf{\ \ }18th century
\end{styleStandard}

\begin{styleStandard}
\ \ a.\ \ QV\ \ pues es sabido y consta a todo el mundo que~\textit{qualquier}~diarista \textit{es} un hombre infalible y no puede engañarse en materia alguna 
\end{styleStandard}

\begin{styleStandard}
\ \ \ \ \ \ ‘because it is known and everybody is aware of the fact that any journalist is an infallible man and cannot be fooled in any subject…’
\end{styleStandard}

\begin{styleStandard}
\ \ \ \ \ \ (Forner, \textit{Gramáticos})
\end{styleStandard}

\begin{styleStandard}
\ \ b.\ \ VQ\ \ pues la semilla de que \textit{se forma}~\textit{cualquiera}~planta, necesariamente vino de otra planta de la especie misma 
\end{styleStandard}

\begin{styleStandard}
\ \ \ \ \ \ ‘because the seed from which any plant grows came necessarily from another plant of the same species’
\end{styleStandard}

\begin{styleStandard}
\ \ \ \ \ \ (Feijoo, \textit{Teatro crítico})
\end{styleStandard}

\begin{styleStandard}
\textbf{\ \ }19th century
\end{styleStandard}

\begin{styleStandard}
\ \ a.\ \ QV\ \ \textit{Cualquiera~me creerá} tu hermanito menor. 
\end{styleStandard}

\begin{styleStandard}
\ \ \ \ \ \ ‘Anybody will believe that I am your younger brother.’
\end{styleStandard}

\begin{styleStandard}
\ \ \ \ \ \ (Tamayo y Baus, \textit{No hay mal})
\end{styleStandard}

\begin{styleStandard}
\ \ b.\ \ VQ\ \ Luego, salía, \textit{y dirigíase} lentamente por las calles hacia\textit{~cualquier~café}, en donde les escribía a su madre y a la novia. 
\end{styleStandard}

\begin{styleStandard}
\ \ \ \ \ \ ‘Then he left, and headed slowly through the streets towards any café, where he would write to his mother and his fiancée.’
\end{styleStandard}

\begin{styleStandard}
\ \ \ \ \ \ (Trigo, \textit{En la carrera})
\end{styleStandard}

\begin{styleStandard}
\textbf{\ \ }20th century
\end{styleStandard}

\begin{styleStandard}
\ \ a.\ \ QV\ \ \textit{Cualquier~persona puede} disfrutar una biblioteca de ensueño.
\end{styleStandard}

\begin{styleStandard}
\ \ \ \ \ \ ‘Any person can enjoy a dream library.’
\end{styleStandard}

\begin{styleStandard}
\ \ \ \ \ \ (Boo Juan Vicente,\textit{ España:ABC: CONTEXT +})
\end{styleStandard}

\begin{styleStandard}
\ \ b.\ \ VQ\ \ Esto lo \textit{puede} hacer\textit{ cualquiera}. 
\end{styleStandard}

\begin{styleStandard}
\ \ \ \ \ \  ‘Anybody can do this.’
\end{styleStandard}

\begin{styleStandard}
\ \ \ \ \ \ (@ 1)
\end{styleStandard}

\begin{styleStandard}
The overall decrease in the order QV is parallel to the overall decrease of the occurrence of \textsc{qualquier} \textit{que} as Table 20 shows. Looking more closely into the frequency distribution of different centuries, the decrease is drastic from the 13th to the 15th century, which might again be correlated with the grammatical change of early and late Old Spanish and the grammaticalisation of \textit{qual quier} as discussed with respect to Figure 9.
\end{styleStandard}

\begin{center}
\tablefirsthead{}
\tablehead{}
\tabletail{}
\tablelasttail{}
\begin{supertabular}{|m{1.1309599in}|m{1.5615599in}|m{1.5622599in}|m{0.9191598in}|}
\hline
\centering \textbf{Century} &
\centering \textbf{Analysed examples/occurences (occ.)} &
\centering \textbf{Nr. of occ.} &
{\centering \textbf{Rel. }\par}

\centering\arraybslash \textbf{frequency}\\\hline
\centering 13th &
\centering 150 &
\centering 127 &
\centering\arraybslash 84.67\%\\\hline
\centering 14th &
\centering 142 &
\centering 75 &
\centering\arraybslash 52.82\%\\\hline
\centering 15th &
\centering 150 &
\centering 67 &
\centering\arraybslash 44.67\%\\\hline
\centering 16th &
\centering 150 &
\centering 63 &
\centering\arraybslash 42.00\%\\\hline
\centering 17th &
\centering 150 &
\centering 33 &
\centering\arraybslash 22.00\%\\\hline
\centering 18th &
\centering 149 &
\centering 53 &
\centering\arraybslash 35.57\%\\\hline
\centering 19th &
\centering 150 &
\centering 37 &
\centering\arraybslash 24.67\%\\\hline
\centering 20th &
\centering 150 &
\centering 15 &
\centering\arraybslash 10.00\%\\\hline
\end{supertabular}
\end{center}
\begin{styleStandard}
Table \stepcounter{Table}{\theTable}: Distribution of QUALQUIER \textit{que} in “my data”.
\end{styleStandard}

\begin{styleStandard}
The following figure visualises the change in the frequency of preverbal \textsc{qualquier} (QV) and \textsc{qualquier} que (\textsc{Q que}) in Table 25:
\end{styleStandard}

\begin{styleStandard}
Figure \stepcounter{Figure}{\theFigure}: Decrease of Q–V and Q \textit{que} (=\textsc{qualquier} \textit{que}) in “my data”.
\end{styleStandard}

\begin{styleStandard}
The question is whether the overall change is statistically significant. On the null hypothesis (see Section 5.1.3), the differences in frequencies per century are not significant; the difference in observed values could be just a random effect of different samples. If the null hypothesis is correct, the frequencies per century should be the same as or very close to the mean. On the intuitive level, the null hypothesis is not correct due to large differences in the data. We know that the maximal frequency is 127 (13th c.) and the minimal frequency is 15 (20th c.). The mean of all frequencies per century is 59. However, the standard deviation (SD) of \textsc{qualquier}\textit{ que} is 32, which is almost half as large as the mean, which is far too high to reflect a normal distribution of the data. Recall from 5.1.3 that the larger the number of the standard deviation, the wider the spread is of the data around the mean. The data in Table 25 cannot be distributed normally. This is also shown by the fact that the data in Figure 18 does not show a bell curve comparable to that of a normal distribution shown in 5.1.3.
\end{styleStandard}

\begin{styleStandard}
\ \ The same difference obtains for the Q–V/V–Q Variable. The maximal frequency is 134 (13th c.) and the minimal frequency is 54 (20th c.). The mean of all frequencies per century is 94. However, the standard deviation of QV is 24, , which is far too high to reflect a normal distribution of the data. Thus, even from a descriptive level, the difference is high enough to conclude that the time or century matters for the variables \textsc{qualquier}\textit{ que} and QV and that these two variables are correlated. The descriptive data is summarised in the following table:
\end{styleStandard}

\begin{center}
\tablefirsthead{}
\tablehead{}
\tabletail{}
\tablelasttail{}
\begin{supertabular}{|m{0.5705598in}m{0.52195984in}|m{0.60315984in}m{0.5031598in}|m{0.72885984in}m{0.8212598in}|}
\hline
\multicolumn{2}{|m{1.1712599in}|}{\centering mean} &
\multicolumn{2}{m{1.1850599in}|}{\centering variation (s²)} &
\multicolumn{2}{m{1.6288599in}|}{\centering standard deviation (Std.)}\\
\centering \textbf{Q–V} &
\centering \textbf{Q }\textbf{\textit{que}} &
\centering \textbf{Q–V} &
\centering \textbf{Q }\textbf{\textit{que}} &
\centering \textbf{Q–V} &
\centering\arraybslash \textbf{Q }\textbf{\textit{que}}\\\hline
\centering 93.875 &
\centering 58.75 &
\centering 564.375 &
\centering 1011.5 &
\centering 23.756578 &
\centering\arraybslash 31.8040878\\\hline
\end{supertabular}
\end{center}
\begin{styleStandard}
Table \stepcounter{Table}{\theTable}: Descriptive statistics of Q–V and of Q \textit{que}.
\end{styleStandard}

\begin{styleStandard}
In order to test the null hypothesis, namely that time (here: Century) has no influence on the frequencies of QV and on \textit{cualquiera} used in relative clauses (i.e. \textsc{qualquier}\textit{ que}), I used a linear regression model as shown, in Figure 11 and 12. These figures show that time plays a role in the frequencies of both constructions; that is, the use of both QV and \textsc{qualquier}\textit{ que} changes statistically significantly over time. The regression output below shows that the Century variable has a significant influence on the syntactic structure (QV) due to the t-value, which correlates to a significance value (t = –4.68, p {\textless} .01) in Figure 11 shows that Century has an influence on the syntactic structure \textsc{qualquier}\textit{ que} due to the high significance value (t = –3.940, p {\textless} .01).\footnote{\textrm{\ To express the result in statistical terminology, I computed a linear regression on the logit-transformed relative frequencies of QV and }\textsc{qualquier}\textrm{\textit{ que}}\textrm{ to test the null hypothesis that the variable Century has no influence on the frequency of QV and/or }\textsc{qualquier }\textrm{\textit{que}}\textrm{. The linear regression showed that Century has a significant influence on QV (t = –4.68, p {\textless} .01) and }\textsc{qualquier}\textrm{ }\textrm{\textit{que}}\textrm{ (t = –3.940, p {\textless} .01).}} The output of the linear regression is given in Figure 12:
\end{styleStandard}

\begin{styleStandard}
Coefficients:
\end{styleStandard}

\begin{styleStandard}
\ \ \ \ \ \ \ \ \ \ \ \ Estimate Std. Error t value Pr({\textgreater}{\textbar}t{\textbar}) \ \ 
\end{styleStandard}

\begin{styleStandard}
(Intercept) \ \ 2.5804 \ \ \ \ \ 0.4449 \ \ \ \ \ 5.800 \ \ \ \ 0.00115 **
\end{styleStandard}

\begin{styleStandard}
Century \ \ \ \ \ \ –0.4122~~~~ 0.0881~ \ \ \ \ –4.678~ \ \ 0.00340 **
\end{styleStandard}

\begin{styleStandard}
Figure \stepcounter{Figure}{\theFigure}: The output of the linear regression from R with respect to QV.
\end{styleStandard}

\begin{flushleft}
\tablefirsthead{}
\tablehead{}
\tabletail{}
\tablelasttail{}
\begin{supertabular}{|m{6.22126in}|}
\hline
Coefficients:

\ \ \ \ \ \ \ \ \ \ \ \ Estimate Std. Error t value Pr({\textgreater}{\textbar}t{\textbar}) \ \ 

(Intercept) \ \ 2.0768 \ \ \ \ 0.6924 \ \ 2.999 \ 0.02403 * 

Century \ \ \ \ \ \ {}-0.5402 \ \ \ \ 0.1371 \ {}-3.940 \ 0.00763 **\\\hline
\end{supertabular}
\end{flushleft}
\begin{styleStandard}
Figure \stepcounter{Figure}{\theFigure}: The output of the linear regression from R with respect to \textsc{qualquier} \textit{que.}
\end{styleStandard}

\begin{styleStandard}
I also ran a\textit{ chi square} test of independence to examine the relation between verbal position and century. The relation between these variables was significant, X\textsuperscript{2} (degree of freedom = 1, 91.1854) p {\textless} 0.00001; that is, VQ is more likely to appear later than QV.
\end{styleStandard}

\begin{styleStandard}
\ \ My results confirm the general observation about the decrease of \textsc{qualquier}\textit{ que} in Old Spanish in the literature (see Chapter 6, C\&P 2009) on the basis of statistical tests. Additionally, I have shown that it is not only \textsc{qualquier}\textit{ que} that decreases but also QV; and for that matter, VQ increases. We will see in Section 5.3 what consequences this change in frequency brings for the semantic interpretation of \textsc{qualquier} and how this change in frequency follows from the Grammaticalisation.
\end{styleStandard}

\paragraph[5.2.1.2 Frequency distribution of QN/NQ]{\textit{5.2.1.2 Frequency distribution of QN/NQ}}
\begin{styleStandard}
My sample from “my data” has low numbers of postnominal \textsc{qualquier} (henceforth NQ) (ranging between 0 and 4 occ. out of 150 occ. in every century). Recall that “my data” do not contain data of coordination and ellipsis. As NQ appears a lot in coordination and elliptical constructions, as will be shown in Section 5.2, NQ is not sufficiently represented in “my data”. In order to test the change of NQ and QN including elliptical sentences and coordination, I used the whole diachronic corpus.
\end{styleStandard}

\begin{styleStandard}
\ \ As we can see just by the observation of the percentage numbers of NQ and QN per every century in the figure and table below, NQ and QN have more or less the same distribution in the 13th century: around 50\% (NQ having 52\% and QN 48\%). This frequency difference changes drastically in the 14th century; as shown below; that is, QN increases, whereas NQ decreases, except for the 18th and 19th centuries, where QN decreases and NQ increases. This difference in frequencies needs an explanation, which will be provided in Section 5.3.
\end{styleStandard}

\begin{stylecaption}
\textmd{Figure \stepcounter{Figure}{\theFigure}: Distribution of percentages of QN/NQ per century in CdE.}
\end{stylecaption}

\begin{center}
\tablefirsthead{\hline
\centering \textbf{Century} &
\centering \textbf{QUALQUIER N in \%} &
\centering \textbf{QUALQUIER N in nr. of Frequency} &
\centering \textbf{N QUALQUIER in \%} &
\centering \textbf{N QUALQUIER in nr. of Frequency} &
\centering\arraybslash \textbf{Total nr. of Frequency of QUALQUIER}\\\hline
\centering 13th &
\centering 48 &
\centering 168 &
\centering 52 &
\centering 185 &
\centering\arraybslash 353\\\hline
\centering 14th &
\centering 83 &
\centering 106 &
\centering 17 &
\centering 21 &
\centering\arraybslash 127\\\hline
\centering 15th &
\centering 90 &
\centering 1,090 &
\centering 10 &
\centering 118 &
\centering\arraybslash 1,208\\\hline
\centering 16th &
\centering 94 &
\centering 3,514 &
\centering 6 &
\centering 222 &
\centering\arraybslash 3,736\\\hline
\centering 17th &
\centering 92 &
\centering 2,026 &
\centering 8 &
\centering 177 &
\centering\arraybslash 2,203\\\hline
\centering 18th &
\centering 94 &
\centering 2,435 &
\centering 6 &
\centering 151 &
\centering\arraybslash 2,586\\\hline
\centering 19th &
\centering 84 &
\centering 2,769 &
\centering 16 &
\centering 544 &
\centering\arraybslash 3,313\\\hline
\centering 20th &
\centering 94 &
\centering 5,542 &
\centering 6 &
\centering 337 &
\centering\arraybslash 5,879\\}
\tablehead{\hline
\centering \textbf{Century} &
\centering \textbf{QUALQUIER N in \%} &
\centering \textbf{QUALQUIER N in nr. of Frequency} &
\centering \textbf{N QUALQUIER in \%} &
\centering \textbf{N QUALQUIER in nr. of Frequency} &
\centering\arraybslash \textbf{Total nr. of Frequency of QUALQUIER}\\\hline
\centering 13th &
\centering 48 &
\centering 168 &
\centering 52 &
\centering 185 &
\centering\arraybslash 353\\\hline
\centering 14th &
\centering 83 &
\centering 106 &
\centering 17 &
\centering 21 &
\centering\arraybslash 127\\\hline
\centering 15th &
\centering 90 &
\centering 1,090 &
\centering 10 &
\centering 118 &
\centering\arraybslash 1,208\\\hline
\centering 16th &
\centering 94 &
\centering 3,514 &
\centering 6 &
\centering 222 &
\centering\arraybslash 3,736\\\hline
\centering 17th &
\centering 92 &
\centering 2,026 &
\centering 8 &
\centering 177 &
\centering\arraybslash 2,203\\\hline
\centering 18th &
\centering 94 &
\centering 2,435 &
\centering 6 &
\centering 151 &
\centering\arraybslash 2,586\\\hline
\centering 19th &
\centering 84 &
\centering 2,769 &
\centering 16 &
\centering 544 &
\centering\arraybslash 3,313\\\hline
\centering 20th &
\centering 94 &
\centering 5,542 &
\centering 6 &
\centering 337 &
\centering\arraybslash 5,879\\}
\tabletail{}
\tablelasttail{}
\begin{supertabular}{|m{0.6087598in}|m{0.9038598in}|m{1.0004599in}|m{1.0011599in}|m{1.0983598in}|m{1.0969598in}|}

\end{supertabular}
\end{center}
\begin{styleStandard}
Table \stepcounter{Table}{\theTable}: Q N (=*ualqu* \_N*) and N Q (\_N* *ualqu*) in the whole CdE (Mar. 2020).
\end{styleStandard}

\begin{styleStandard}
I can thus summarise that, according to the data from CdE, QN increased in time, whereas NQ decreased and there is a big jump in decrease of NQ or increase of QN in the 14th century. This big jump correlates with many other linguistic changes in this time interval of Old Spanish (see Eberenz 2009, among others).
\end{styleStandard}

\begin{styleStandard}
\textbf{\ \ }a.\ \ N \textsc{qualquier} (the most frequent form in the 13th c., the least frequent in the 19th c.)
\end{styleStandard}

\begin{styleStandard}
\ \ b.\ \ \textsc{qualquier} N (the less frequent form in the 13th c., the most frequent in the 19th c.)
\end{styleStandard}

\begin{styleStandard}
In Section 5.3, I will explain this change in frequency of QN/NQ in terms of a linguistic change whereby \textsc{qualquier} changes its syntactic status from the specifier of a relative clause to a noun modifier (NQ) similar to an adjective (N Adj) and then to a quantificational determiner (QN) similar to \textit{algún} N.
\end{styleStandard}

\begin{styleStandard}
\ \ In the next section, I will look into (\textit{otro}) NQ\textit{ }in more detail. 
\end{styleStandard}

\subsubsection[5.2.2 The change from (otro) NQ to un NQ and Def NQ]{\textbf{5.2.2 The change from (}\textbf{\textit{otro}}\textbf{)}\textbf{\textit{ }}\textbf{NQ to }\textbf{\textit{un }}\textbf{NQ}\textbf{\textit{ }}\textbf{and Def NQ}}
\begin{styleStandard}
We will now analyse NQ more closely from the whole corpus (see 5.2.1 for the frequency of NQ in total). Until the 15th century, the word order NQ appears very often in coordination with \textit{o} ‘or’ with nouns preceded by the term \textit{otro }‘other’ as in 292. The majority of the occurrences of the sequence \textit{o otro }NQ \textit{que} (120 occ. with complement clauses introduced by \textit{que} of a total of 196 occurrences of \textit{o otro }NQ, i.e. 61.22\%):
\end{styleStandard}

\begin{styleStandard}
\textbf{\ \ }a.\ \ firiendo en ella con piedra, o con cuchillo o con \textit{otra cosa qualquier} 
\end{styleStandard}

\begin{styleStandard}
\ \ \ \ ‘damaging it [i.e. the Cross] with a stone, or with a knife or with some other thing whatsoever’
\end{styleStandard}

\begin{styleStandard}
\ \ \ \ (13th c., Alfonso X, \textit{Siete Partidas})
\end{styleStandard}

\begin{styleStandard}
\ \ b.\ \ o \textit{otra manera qualquier} que le diese ayuda asabiendas que fuese semeiante o \textit{otra cosa qualquier }que sea 
\end{styleStandard}

\begin{styleStandard}
\ \ \ \ ‘or [providing help to a thief in] some other way whatsoever that would be of help, knowing full well that it is similar [to the above-mentioned things] or any other thing whatever it may be’
\end{styleStandard}

\begin{styleStandard}
\ \ \ \ (13th c., Alfonso X, \textit{Siete Partidas})
\end{styleStandard}

\begin{styleStandard}
\ \ c.\ \ o de \textit{otra razon qualquier} que sea 
\end{styleStandard}

\begin{styleStandard}
\ \ \ \ ‘or for some other reason whatever it may be’
\end{styleStandard}

\begin{styleStandard}
\ \ \ \ (13th c., Alfonso X, \textit{Siete Partidas})
\end{styleStandard}

\begin{styleStandard}
From the 15th century onwards, there are more cases without relative clause modification, which is expected given that relative clause modification decreased in general:
\end{styleStandard}

\begin{styleStandard}
\textbf{\ \ }a.\ \ […] sea debajo de la equinoccial o en \textit{otra parte cualquiera}» A este propósito, trae Celio Rodiginio lo de Arriano, historiador griego 
\end{styleStandard}

\begin{styleStandard}
\ \ \ \ ‘… be it under the equatorial line or in any other place» On this subject, Celio Rodiginio brings up what was said by Arriano, a Greek historian’
\end{styleStandard}

\begin{styleStandard}
\ \ \ \ (16th c., Torquemada, \textit{Jardín})
\end{styleStandard}

\begin{styleStandard}
\ \ b.\ \ secretamente enlazado en el amor de alguna criatura, sea persona, sea \textit{otra cosa cualquiera}. 
\end{styleStandard}

\begin{styleStandard}
\ \ \ \ ‘secretely bound in the love of some creature, be it a person, be it any other thing.’
\end{styleStandard}

\begin{styleStandard}
\ \ \ \ (16th c., Fray Luís de Granada, \textit{Guía de pecadores})
\end{styleStandard}

\begin{styleStandard}
All examples of \textit{otro }N Q\textit{,} independently of the century, have an FCI interpretation (‘other N whatever one wants’). This observation is not surprising given that \textit{otro }N\textit{ }makes reference to the preceding alternatives that all belong to the same class of entities denoted by the noun N (compare English ‘fork, knife, or other utensil’), and the function of postnominal \textsc{qualquier} is to express another alternative that can have any identity or be of any type as long as it belongs to the same class of entities or types denoted by N.
\end{styleStandard}

\begin{styleStandard}
\ \ Looking into the issue of whether elements other than \textit{otro} can precede the sequence NQ, we find that there are also examples of N \textsc{qualquier} with nouns that occur with a determiner. There are very few examples with the indefinite article, of the form of \textit{un/una N qual quier,-e,-a} with the separated spelling \textit{qual quier,-e,-a} (4 occ. in the 13th c., 6 occ. in the 14th, 1 occ. in the 15th):
\end{styleStandard}

\begin{styleStandard}
\textbf{\ \ }a.\ \ … Desi dixo domjngujllo que le diese el Rey \textit{vn omne qual quier} a que
\end{styleStandard}

\begin{styleStandard}
le firiese e quando le oujese ferido que se yrie para el castillo 
\end{styleStandard}

\begin{styleStandard}
\ \ \ \ ‘Then Dominguillo said that the King should give him a man which(ever) one wants whom he could hurt, and after hurting him, he would go to the castle.’
\end{styleStandard}

\begin{styleStandard}
\ \ \ \ (14th c., Anonymous, \textit{Veinte Reyes})
\end{styleStandard}

\begin{styleStandard}
\ \ b.\ \ …en un logar qual quier dela posada. 
\end{styleStandard}

\begin{styleStandard}
\ \ \ \ ‘…at any place in the house.’
\end{styleStandard}

\begin{styleStandard}
\ \ \ \ (13th c., Moamyn,\textit{ Animalias})
\end{styleStandard}

\begin{styleStandard}
The synthetic form \textsc{qualquier }in \textsc{un} N \textsc{qualquier} appears in the 16th century in its modern spelling with \textit{cu}{}- (with total of 2 occ. in this century in the whole corpus). There is no synthetic form of \textsc{un }N \textsc{qualquier }with \textit{qu-} in the whole corpus (only a separated form, as shown in 295). In these examples, \textit{cualquiera} is interpreted as an FCI due to the generic context in which \textit{cualquiera} is used (see Chapter 2 on the licensing of FCIs in generic contexts):
\end{styleStandard}

\begin{styleStandard}
\textbf{\ \ }resulta difícil el destacar en \textit{una ciencia cualquiera}, cuánto más en la filosofía. 
\end{styleStandard}

\begin{styleStandard}
\ \ ‘it turns out to be difficult to stand out in \textit{any science}, even more so in philosophy.’
\end{styleStandard}

\begin{styleStandard}
\ \ (16th c., Sepúlveda, \textit{Epistolario})
\end{styleStandard}

\begin{styleStandard}
\textbf{\ \ }que con tanto más entusiasmo nos dedicamos a \textit{un trabajo cualquiera}, cuanto más nos agrada el hacerlo. 
\end{styleStandard}

\begin{styleStandard}
\ \ ‘We dedicate ourselves to \textit{any kind of job} with the more enthusiasm, the more we like to do it.’
\end{styleStandard}

\begin{styleStandard}
\ \ (16th c., Sepúlveda, \textit{Epistolario})
\end{styleStandard}

\begin{styleStandard}
In the 17th century, the construction \textsc{un}\textit{ }N \textit{cualquiera} is used 6 times in the whole corpus, and it has an ambiguous interpretation between the FCI interpretation (‘any’) and the evaluative interpretation (‘an ordinary’) (see Chapter 3 for the ambiguity of different interpretations of \textit{cualquiera} in 20th/21st-century Spanish). The construction \textsc{un}\textit{ }N\textit{ cualquiera} is used 9 times in the 18th century, 356 times in the 19th, and 169 times in the 20th.
\end{styleStandard}

\begin{styleStandard}
\ \ In the first example, 297, is used under negation with an indicative present tense verb. The examples in 297 show generic contexts with the indicative present and also have an ambiguous interpretation of ‘any’ or ‘ordinary’:
\end{styleStandard}

\begin{styleStandard}
\textbf{\ \ }a.\ \ Para pedir campo a un rey no basta \textit{un hombre cualquiera}, que según la ley del duelo es menester que rey sea. 
\end{styleStandard}

\begin{styleStandard}
\ \ \ \ ‘In order to challenge a king, it is not sufficient to be \textit{an ordinary man} (EVAL), as according to the duelling law he has to be a king himself.’
\end{styleStandard}

\begin{styleStandard}
‘In order to challenge a king, it is not sufficient to be\textit{ just any man} (FCI), as according to the duelling law he has to be a king himself.’
\end{styleStandard}

\begin{styleStandard}
\ \ \ \ (17th c., Anonymous, \textit{El rey don Alfonso})
\end{styleStandard}

\begin{styleStandard}
\ \ b.\ \ Quien se huelga de parecerle discreto a \textit{un hombre cualquiera}, ¿cómo piensa que trata a Dios cuando no se le da nada de no parecerle discreto? 
\end{styleStandard}

\begin{styleStandard}
\ \ \ \ ‘Someone who likes to be seen as discreet by \textit{any man} (FCI), how does he think he treats God when he doesn’t receive anything for not being seen by him as discreet?’
\end{styleStandard}

\begin{styleStandard}
\ \ \ \ ‘Someone who likes to be seen as discreet by by \textit{an ordinary man} (EVAL), how does he think he treats God when he doesn’t receive anything for not being seen by him as discreet?’
\end{styleStandard}

\begin{styleStandard}
\ \ \ \ (17th c., Zabaleta, \textit{El día de fiesta, primera parte})
\end{styleStandard}

\begin{styleStandard}
\ \ c.\ \ Para instruirse en el idioma de la medicina y comer sus aforismos basta \textit{un curso cualquiera}. 
\end{styleStandard}

\begin{styleStandard}
\ \ \ \ ‘in order to learn the language of medicine and to eat its aphorisms, any course (FCI)/an ordinary course (EVAL) is sufficient.’ (1732, \textit{Vida}, Torres Villarroel)
\end{styleStandard}

\begin{styleStandard}
\ \ \ \ (17th c., Zabaleta,\textit{ El día de fiesta, primera parte})
\end{styleStandard}

\begin{styleStandard}
From this time onwards, \textsc{un}\textit{ }N \textit{cualquiera} is used as EVAL or FCI depending on the context as in present-day Spanish (see Chapter 3). In the following example, \textit{un alojamento cualquiera} can be interpreted as FCI ‘any accomodation’ or as ‘a normal accomodation’ excluding ‘abnormal’ or ‘unlivable’ accomodations:
\end{styleStandard}

\begin{styleStandard}
\textbf{\ \ }\ \ En esta perplejidad escribo al administrador del conde de Peñalba, encargándole que me busque \textit{un alojamiento cualquiera} de acuerdo con usted, .... 
\end{styleStandard}

\begin{styleStandard}
\ \ \ \ ‘I’m writing in this perplexity to the administrator of the Count of Peñalba, charging him to find me some accomodation in accordance with you, …’
\end{styleStandard}

\begin{styleStandard}
\ \ \ \ (18th c., Jovellanos, \textit{Correspondencia})
\end{styleStandard}

\begin{styleStandard}
The first examples of \textsc{un} N \textit{cualquiera} in affirmative predicative contexts appeared in the 19th century:
\end{styleStandard}

\begin{styleStandard}
\textbf{\ \ }\ \ Es \textit{un contrato cualquiera}. vos necesitáis dinero, yo os necesito, ....
\end{styleStandard}

\begin{styleStandard}
\ \ \ \ ‘It’s a usual contract: you need money, I need you, …’
\end{styleStandard}

\begin{styleStandard}
\ \ \ \ (19th c., Zorrilla, \textit{Entre clérigos y diablos})
\end{styleStandard}

\begin{styleStandard}
The use of\textit{ N cualquiera }with definite articles or demonstratives such \textit{la, el, esta, este, aquel} N (Def N \textit{cualquiera}) is less frequent than the use of \textit{un/otro} N \textit{cualquiera} (less than 5\% of all occ.) and also appears not earlier than the 19th century (see Chapter 3 for present-day Spanish). An explanation of the late appearance of Def N \textit{cualquiera} will be given in Section 5.3. All examples of Def \textit{N cualquiera }have an evaluative interpretation of ‘common/ordinary’ (see also Chapter 3 for Modern Spanish):
\end{styleStandard}

\begin{styleStandard}
\textbf{\ \ }La cuestión empezó \textit{ese día cualquiera} del siglo XVI, en que don Felipe II decretó que los españoles no podrían estudiar en Universidades extranjeras.
\end{styleStandard}

\begin{styleStandard}
\ \ ‘the problem began on this ordinary day of the 16th century, when Don Felipe decreed, that Spaniards could not study at foreign universities.’ 
\end{styleStandard}

\begin{styleStandard}
\ \ (20th c., Ballester, \textit{Torre Del Aire})
\end{styleStandard}

\begin{styleStandard}
Given these observations, it is now possible to explain the rise of NQ\textit{ }in the 18th and 19th centuries observed in Section 5.2.1.2. The introduction of the new semantic meaning of \textsc{un} NQ ‘an ordinary N’ correlates with the introduction of new syntactic constructions: the use of \textsc{un} NQ\textit{ }under affirmative predicative contexts and the use of NQ under definite articles. These new forms and the new meaning have triggered the increase in frequency of NQ\textit{ }in the 18th and 19th centuries. As the construction (\textit{otro}) NQ practically disappeared in the 19th century and Q(\textit{otro})N was used instead (Table 19), this explains the trend of the general fall of NQ after the 14th century (see Section 5.2.1.2).
\end{styleStandard}

\begin{styleStandard}
\ \ To sum up the development of N\textit{ cualquiera} (bold letters represent higher frequency). It represents a change in different constructions with N\textit{ cualquiera. }It shows that the construction \textit{otro/otra} \textsc{qualquier} (\textit{que }Verb\textsubscript{Subj.})\textbf{ }was first used in the 13th century until the 15th. The construction \textsc{un}\textit{ }N\textit{ cualquiera} appeared in the 15th century and the older construction disappeared. From the 19th century onwards, \textsc{un}\textit{ }N \textit{cualquiera} was used in predicative position. Fewer data of N \textit{cualquiera} contain a definite or demonstrative \textit{la,el,este,esta} N \textit{cualquiera}:
\end{styleStandard}

\begin{styleStandard}
\ \ \& \textit{otro,a} \textsc{qualquier} \textit{que }Verb\textsubscript{Subj.} (13th–15th c.) 
\end{styleStandard}

\begin{styleStandard}
\ \ \ \ (Verb\textsubscript{Ind.}) \textsc{un} N \textit{cualquiera} (Verb\textsubscript{Ind.}) (from 15th c. until today) {\textgreater}{\textgreater}
\end{styleStandard}

\begin{styleStandard}
\ \ \textit{\ \ }\textbf{pedicative verb}\textbf{\textsubscript{Ind.}}\textbf{ }\textbf{\textsc{un}}\textbf{ N }\textbf{\textit{cualquiera}} /\textit{la,el,este,esta} N \textit{cualquiera} (from 19th c. until today)
\end{styleStandard}

\begin{styleStandard}
The distribution in 301 co-occurs with a semantic distribution given the observation that \textit{otro} NQ can only have FCI due to preceding alternatives, whereas \textit{un} NQ is ambiguous between two interpretations (FCI/EVAL) depending on the verbal mood (+/– modal) and Def NQ has only the interpretation of EVAL:
\end{styleStandard}

\begin{styleStandard}
\textbf{\ \ }\textit{otro,a} N \textsc{qualquier que}\textbf{ }….FCI (13th–15th c.) 
\end{styleStandard}

\begin{styleStandard}
\ \ \ \ \ \ \textsc{un} N \textit{cualquiera} FCI/\textbf{EVAL} (from the 16th c. until today) 
\end{styleStandard}

\begin{styleStandard}
\ \ \textit{\ \ la,el,este,esta} N \textit{cualquiera} \textbf{EVAL} (from the 19th c. until today)
\end{styleStandard}

\begin{styleStandard}
As I will argue in Section 5.3, the preference for \textsc{un}\textit{ }N\textit{ cualquiera} over \textit{(otro) }N\textit{ qual quiera }rising\textit{ }and later the appearance of Def N\textit{ cualquiera} is the result of the reanalysis of \textit{cualquiera} from an item introducing a relative clause with FCI into a synthetic nominal modifier with different lexical meanings (FCI and EVAL). 
\end{styleStandard}

\subsubsection[Summing up and evaluating the results ]{Summing up and evaluating the results }
\begin{styleStandard}
It has been shown that the frequency of some variables increased , such as postverbal \textit{cualquiera} (VQ), sentence types other than relative clauses, the verb mood indicative and QN. Moreover, UN/Def NQ came into existence in later centuries:
\end{styleStandard}

\begin{flushleft}
\tablefirsthead{}
\tablehead{}
\tabletail{}
\tablelasttail{}
\begin{supertabular}{|m{6.19836in}|}
\hline
\liststyleWWNumx
\begin{enumerate}
\item {\bfseries \textmd{In early centuries (12th–13th c.) }}
\end{enumerate}
\ \ NQ is more frequent than QN.

\ \ Q \textit{que} is more frequent than other sentence types.

\ \ Subjunctive is more frequent than indicative.

\ \ \textsc{un} N \textsc{qualquier} with indicative did not exist yet.

\ \ \textsc{qualquier} is only interpretable as FCI.

\liststyleWWNumx
\setcounter{saveenum}{\value{enumi}}
\begin{enumerate}
\setcounter{enumi}{\value{saveenum}}
\item {\bfseries \textmd{In later centuries (19th–20th c.)}}
\end{enumerate}
\ \ QN is more frequent than NQ.

\ \ Q\textit{ que} is less frequent than other sentence types.

\ \ Indicative is more frequent than subjunctive.

\ \ \textsc{un} N \textsc{qualquier} with indicative verbal mood appears.

\ \ \textsc{qualquier} has different meanings (FCI, EVAL).\\\hline
\end{supertabular}
\end{flushleft}
\begin{styleStandard}
The main results suggest that \textsc{qualquier} has undergone some changes which correlate with syntatic position and verbal mood/type and (in)definiteness of the NP.
\end{styleStandard}

\begin{styleStandard}
\ \ As was shown in Section 2.6, no quantitative analysis had been done so far on the variables described in Section 5.2, namely the syntactic function of \textsc{qualquier}, its relation to the verb (position and mood), and its relation to the sentence mood. I also have shown when exactly new constructions \textsc{un}/Def N \textsc{qualquier} appear for the first time, and I have quantified them. The construction \textsc{un} N \textsc{qualquier} appears already in the 13th century, but it is used very infrequently in contrast to \textit{otro} N \textsc{qualquier}. In later centuries, \textsc{un} N \textsc{qualquier} starts being used more frequently (from the 18th c. onwards). The construction Def N \textsc{qualquier }is a recent phenomenon (from the 19th c. onwards) which is still less frequent than \textsc{un} N \textsc{qualquier. }I have identified the time interval when the meanings of \textsc{qualquier} started to change, namely from the 16th century.
\end{styleStandard}

\begin{styleStandard}
\ \ In the next section, I will use these results to investigate the Grammaticalisation Hypothesis that was described pretheoretically in Section 5.2. I will especially focus on the question of how the meaning changes of \textsc{qualquier} fit into the Grammaticalisation Hypothesis.
\end{styleStandard}

\subsection[The semantic side of the grammaticalisation]{The semantic side of the grammaticalisation}
\begin{styleStandard}
In this section, I will spell out the change from a relative clause into a quantificational pronoun and determiner from the semantic point of view. The change into a quantificational modifier will be discussed as well.
\end{styleStandard}

\subsubsection[5.3.1 The prenominal determiner]{\textbf{5.3.1 The prenominal determiner}}
\begin{styleStandard}
\ \ I assume that the determiner/pronoun \textit{qualquier} has its origin in the relative clause and that it is the quantificational interpretation of a free relative clause such as \textit{qual quier(e) que} that triggers the reanalysis into a quantificational determiner. I will first illustrate this assumption by the interpretation of a relative clause with a free choice interpretation in Modern Spanish and then apply the semantic interpretation to Old Spanish. As 305 shows, free relative clauses can be replaced by a universal quantificational pronoun ‘everyone’:
\end{styleStandard}

\begin{styleStandard}
\ \ Puedes invitar a quien tú quieras. = Puedes invitar a todos (los que tú quieras invitar).
\end{styleStandard}

\begin{styleStandard}
\ \ \ \ ‘You can invite whom you want (to invite).= You can invite everyone (you want to invite).’
\end{styleStandard}

\begin{styleStandard}
I adopt the denotation of English Free Relatives (FRs) in English (see Šimík 2018) to (FRs) in Spanish. The truth-conditions of 305 are provided in 306.
\end{styleStandard}

\begin{styleAuthor}
\ \ a.\ \ [[quien tu quieras invitar]] = $\lambda $P.${\forall}$x [want[2032?] (x)(you[2032?]) → P(x)]
\end{styleAuthor}

\begin{styleStandard}
(the set of properties that all people that you want to invite)
\end{styleStandard}

\begin{styleStandard}
\ \ b.\ \ [[305)]] = 1 iff
\end{styleStandard}

\begin{styleListParagraph}
\textmd{${\forall}$x [want[2032?] (x)(you[2032?])(invite[2032?])] → can (invite[2032?](x)(you))]}
\end{styleListParagraph}

\begin{styleListParagraph}
\textmd{For all people you want to invite, you can invite them.}
\end{styleListParagraph}

\begin{styleStandard}
Free relatives can thus be analysed as DPs which can be modelled by the insertion of a covert DP with the definite interpretation ‘the NP that’ or universal interpretation (‘everyone’) depending on the verb mood above the relative clause structure; free relatives embedded under modal verbs usually have a universal interpretation:
\end{styleStandard}

\begin{styleListParagraph}
\textmd{\ \ Puedes invitar [DP a todos [CP a}\footnote{\textrm{\ I assume that }\textrm{\textit{a}}\textrm{ is not a preposition but a case marker that belongs to the DP (see Lopez 2000, among others).}}\textmd{ quien tú quieras]]}
\end{styleListParagraph}

\begin{styleListParagraph}
\textmd{\ \ ‘you can invite everyone you want’}
\end{styleListParagraph}

\begin{styleStandard}
It is thus possible to analyse free relatives in Old Spanish as DPs under the assumption that empty syntactic categories correlate with certain semantic interpretations (see also Rivero 1988):
\end{styleStandard}

\begin{styleListParagraph}
\textmd{\ \ [Spec, TP [DP D° ‘everyone’] [CP1 qualj TP1 pro T° quier [CP2 tj que TP2 pro T° lo viesse tj]]] T° hauria ende gran dolor].}
\end{styleListParagraph}

\begin{styleStandard}
\ \ ‘…who one wants that had seen it, would have had great pain and a lot of pity.’
\end{styleStandard}

\begin{styleListParagraph}
\textmd{\ \ (13th c.,}\textmd{\textit{ }}\textmd{Anonymous, }\textmd{\textit{Gran Conquista}}\textmd{)}
\end{styleListParagraph}

\begin{styleStandard}
Once \textit{qual quier} lost its transparent relative clause structure, it was reanalysed as a determiner D° or a pronoun DP and filled the empty DP in an earlier stage (see 5.3.1):
\end{styleStandard}

\begin{styleAuthor}
\ \ [\textsubscript{TP} [\textsubscript{DP} [\textsubscript{CP} qual quier que lo viesse]] \textsubscript{T° }hauria ende gran dolor]. 
\end{styleAuthor}

\begin{styleAuthor}
\ \ [\textsubscript{TP} [\textsubscript{DP} qualquier (NP)] \textsubscript{T° }hauria ende gran dolor]. 
\end{styleAuthor}

\begin{styleListParagraph}
\textmd{\ \ (my example for demonstration)}
\end{styleListParagraph}

\begin{styleStandard}
Let us now see how the syntactic change was interpreted semantically. A compositional analysis of Old Spanish example in (309) is presented in 311. I assume that in Old Sp. \textit{qual} \textit{quier} means something like ‘which(ever) one wants (to choose)’. I underline the interpretation of the verb \textit{quier} and discuss it below in detail. The interpretation is derived step by step. First the relative clause is interpreted as an open proposition due to the wh-feature of \textit{qual} and the clause structure (CP) in (i), then the relative clause is interpreted as a nominal modifier, that is, as a property of the noun N in (ii), then the DP is interpreted as universal quantification in (iii). The truth conditions of the sentence in which the DP occurs are spelled out in 311.
\end{styleStandard}

\begin{styleAuthor}
\ \ a.\ \ [[ [CP qualquier que lo viesse] ]] = \ \ 
\end{styleAuthor}

\begin{styleAuthor}
\ \ \ \ (i) would see it (x) ${\wedge}$ one wants to choose (x)\ \ open proposition
\end{styleAuthor}

\begin{styleStandard}
\ \ b.\ \ [[ [\textsubscript{NP} Ø ‘person’ [\textsubscript{CP }qualquier que lo viesse]] ]] =\textbf{\ \ }
\end{styleStandard}

\begin{styleStandard}
\ \ \ \ (ii) $\lambda $x[person (x) ${\wedge}$ would see it (x) ${\wedge}$ one wants to choose (x)]\ \ property
\end{styleStandard}

\begin{styleStandard}
\ \ c.\ \ [[ [\textsubscript{DP} [\textsubscript{NP} [\textsubscript{CP }qualquier que lo viesse]]] ]] =\ \ 
\end{styleStandard}

\begin{styleStandard}
\ \ \ \ (iii) $\lambda $P. All x[person (x) ${\wedge}$ one wants to choose (x) ${\wedge}$ would see it (x) ${\wedge}$ P(x)]\ \ universal quantifier
\end{styleStandard}

\begin{styleStandard}
\ \ d.\ \ [[qualquier que lo viesse hauria ende gran dolor]] = 1 if would have had[2032?] (All x [person(x) ${\wedge}$ one wants to choose (x) ${\wedge}$ would see him (x)]) (great pain)\ \ 
\end{styleStandard}

\begin{styleStandard}
\ \ \ \ prose: For every person x that one wants to choose and that would see it, x would have great pain.
\end{styleStandard}

\begin{styleStandard}
I will now discuss the semantics of \textit{qual quier} under the transparent meaning of the verb \textit{quiere} and the subject \textit{pro} at Stage 1 of the relative clause analysis. The most important question is whether this meaning is part of truth conditions of the statement in (311). I assume that it is not; that is, the open proposition x would have great pain, is independently true, whether one chooses this value or that value for x. In other words, facts in the real world are independent of belief states. Thus, the statement \textit{any even number (wichever one wants to choose) can be divided by 2} is true independent of whether someone wants to choose 4, 6, 8, etc. We can analyse the literal meaning of \textit{qual quier} rather as a side comment on the part of the speaker about the variable x, similar to appositive relative clauses as in \textit{My colleague (who is a nice guy) moved to Germany}. This side comment can be interpreted as a different illocution from the truth-conditional statement, namely as an invitation of the speaker to the addressee to choose any value for x. In this case, \textit{pro} is interpreted as a pronominal form of the addressee, as in \textit{pro} \textit{quisieres} ‘you (informal) would like’ in 312. Note that in the example 312 \textit{qual quisieres }interrupts the sentence \textit{pon una estrella fixa por su longura}, which confirms the side comment analysis of \textit{qual quisieres}: 
\end{styleStandard}

\begin{styleStandard}
\textbf{\ \ }damos te exiemplo a esto. pon una estrella fixa \textit{qual quisieres} por su longura del eguador con el grado con que passa medio cielo.
\end{styleStandard}

\begin{styleStandard}
\textbf{\ \ }‘We are giving you an example of this. Put a fixed star (which you want) along the equator in an angle such that it passes through the middle of the sky. 
\end{styleStandard}

\begin{styleListParagraph}
\textmd{\ \ (13th c., Alfonso X, }\textmd{\textit{Astronomía}}\textmd{)}
\end{styleListParagraph}

\begin{styleStandard}
The form \textit{qual quisier*} becomes less frequent in the 14th and 15th centuries and disappears completely in the 16th century. This is expected with the decrease of transparency of \textit{qual quisier*}. 
\end{styleStandard}

\begin{styleStandard}
\ \ In more formal texts, such as legal texts, where the addressee represents the general public or the reader of the law such as the context of 312, the form \textit{qual quisieres} ‘which you (informal) want’ cannot be used. The subject \textit{pro} in 312 refers to the general public and the legal statement is true no matter which value the general public wants to choose for x, where x represents a person such that if x saw it, x would have great pain. According to the impersonal interpretation of \textit{pro}, it is guaranteed that the set of x is defined as large as possible (i.e. it contains all possible people the public can think of choosing for x). One direct consequence of this analysis of \textit{qual quier} as a side comment is that one should be able to find two verbs of \textit{querer}, one inside the proposition that has an effect on the truth conditions, and another inside a side comment that does not have any effect on the truth conditions. According to this analysis, only the embedded verb \textit{quisiere} in 313 has an effect on the truth conditions and \textit{qual quier} is interpreted as a side comment ‘which value one wants to choose for x’, where x is represented by \textit{qual}:
\end{styleStandard}

\begin{styleStandard}
\textbf{\ \ }Degrado deuen ser dadas las cosas espirituales \& no por preçio. \textit{onde qual quier que quisiere entrar en orden de religion no deue dar precio ninguno} 
\end{styleStandard}

\begin{styleStandard}
\ \ ‘Spiritual things must be given not bought. \textit{Whoever would want to enter in the realm of} \textit{religion should not give any price}.’
\end{styleStandard}

\begin{styleListParagraph}
\textmd{\ \ (13th c.,}\textmd{\textit{ }}\textmd{Alfonso X, }\textmd{\textit{Siete partidas}}\textmd{)}
\end{styleListParagraph}

\begin{styleStandard}
The crucial assumption is that at the stage, when \textit{querer} was analysed as a lexical verb (i.e. before the 13th c. and partially also in 13th c.), the side comment expressed by \textit{qual quier} as ‘which one wants to choose’ or by \textit{qual quisieres} ‘which you, informal, would like to choose’ is literally expressed, whereas in later periods, when \textit{quier} is no longer interpreted as a verb, free choice is not literally encoded in the relative clause, but in the indefinite itself as part of its conventional meaning (see Chapter 1 and example (4) repeated here in 314):
\end{styleStandard}

\begin{styleStandard}
\textbf{\ \ }\textit{Puedes}\ \ \textit{traer}\ \ \textit{cualquier}\ \ \textit{libro.}\ \ \ \ 
\end{styleStandard}

\begin{styleStandard}
\ \ can\ \ bring\ \ \textsc{cualquier}\ \ book\ \ \ \ 
\end{styleStandard}

\begin{styleStandard}
\ \ ‘You can bring any book.’
\end{styleStandard}

\begin{styleStandard}
\ \ Conventional meaning: ‘You can bring me a book, and each book is a possible option.’
\end{styleStandard}

\begin{styleStandard}
\ \ (Aloni et al. 2011: 2)
\end{styleStandard}

\begin{styleStandard}
Another consequence of the change is that the choice to pick a value for x no longer depends on the generic subject \textit{pro} of \textit{quier }as in Old Spanish, but on the agent of the clause in which \textit{cualquiera} is used. In the Modern Spanish example in 314, the subject of free choice is expressed by the grammatical subject \textit{pro}, which refers to the addressee, who can pick any book he/she wants, not the generic \textit{pro} inside \textit{cualquier}, which is no longer transparent in Modern Spanish. Another consequence of the grammaticalisation is that the indicative mood and perfective aspect increased (see 5.2). As was demonstrated in Section 2.3, \textit{cualquiera} has a ‘random’ interpretation or an evaluative interpretation under volitional verbs with perfective aspect (see Sections 2.3 and 3), repeated here in 315. Under both interpretations (RCI and EVAL), the generic \textit{pro} is no longer visible as the agent of choice. In 315 the random choice depends on the explicit agent Juan and 315 does not express any choice at all, but speaker’s evaluation.
\end{styleStandard}

\begin{styleStandard}
\ \ Juan compró un libro cualquiera. ‘Juan bought a book.’
\end{styleStandard}

\begin{styleStandard}
\ \ a.\ \ John bought a book, and he choose it indiscriminately. (Random Choice)
\end{styleStandard}

\begin{styleStandard}
\ \ b.\ \ John bought a book, and the book is not special or remarkable (according to the speaker). (Unremarkable Reading, EVAL)
\end{styleStandard}

\begin{styleStandard}
\ \ \ \ (A\&M 2018: 2)
\end{styleStandard}

\begin{styleStandard}
According to my analysis, the semantic change can be summarised in the following steps. The grammaticalisation of \textsc{qualquier} has induced a change in the interpretation of \textsc{qualquier} with respect to the suject of ‘want to choose a value for \textit{qual}’. When \textit{qualquier} was still transparent (before the 13th c. and partially still in 13th c.), free choice depended on the generic \textit{pro} and was expressed explicitly by the verbal phrase \textit{pro}\textsubscript{[generic] }\textit{quier}. After the loss of transparency, free choice was encoded in the indefinite and the agent of free choice depended on the grammatical subject/agent of the clause. Random Choice and EVAL appeared with the rise of indefinite nouns, indicative mood, and the volitional/predicative verb class:
\end{styleStandard}

\begin{styleStandard}
\ \ 1.\ \ \textit{qual} \textit{pro} \textsubscript{V°/T°} \textit{quiere que} V-subj. (before the 13th c.) [F0E0?]
\end{styleStandard}

\begin{styleStandard}
\ \ \ \ Free Choice expressed explicitly and the agent of Free Choice is \textit{pro}\textsubscript{[generic]}
\end{styleStandard}

\begin{styleListParagraph}
\textmd{\ \ }2.\textmd{\ \ Indefinite pronoun/determiner }\textmd{\textit{cualquier,-a}}\textmd{ (after the 13th c. till today) [F0E0?]}
\end{styleListParagraph}

\begin{styleListParagraph}
\textmd{\ \ \ \ Free Choice expressed }\textmd{\textit{implicitly}}\textmd{ as a conventional implicature and the agent of Free Choice is the agent of the matrix verb.}
\end{styleListParagraph}

\begin{styleStandard}
\ \ 3.\ \ \textsc{un} N\textit{ cualquiera }under volitional verbs in indicative [F0E0?]
\end{styleStandard}

\begin{styleListParagraph}
\textmd{\ \ \ \ Random Choice Modality (from the 18th c. till today)}
\end{styleListParagraph}

\begin{styleStandard}
\ \ 4.\ \ \textsc{un} N\textit{ cualquiera} under predicative verbs in indicative (from the 18th c. till today)
\end{styleStandard}

\begin{styleStandard}
EVAL
\end{styleStandard}

\begin{styleStandard}
Given the syntactic and semantic change suggested in 316, the grammaticalisation of the indefinite does not correlate with a loss or weakening in meaning as C\&P (2009) suggest—the grammaticalised form \textsc{qualquier} is not void of meaning—but it correlates with a change from an \textit{explicit }to an\textit{ implicit }encoding of Free Choice, once \textit{qual quier} is no longer analysed literally (already in 13th c.) and it also correlates with the change of the agent of Free Choice (from generic to grammatical subject/agent of the matrix clause). From the 13th century until today, Free Choice meaning remains, but it is encoded differently. Only much later from the 18th century onwards, we observe that a \textit{new meaning} (EVAL) appears due to the increase of indicative mood, predicative verbs, and the increase of \textsc{un} N \textit{cualquiera}. I will discuss the meaning change to EVAL in 5.3.2.
\end{styleStandard}

\subsubsection[5.3.2 Analysis of qualquier as a postnominal modifier]{\textbf{5.3.2 Analysis of }\textbf{\textsc{qualquier}}\textbf{ as a postnominal modifier}}
\begin{styleStandard}
The relative clause \textsc{qualquier}\textit{ que…} could also appear postnominally with an overt noun as antecedent in the 13th century.
\end{styleStandard}

\begin{styleStandard}
\ \ Let us now look into this type more closely. The antecedent of the relative clause introduced by \textit{qual(quier) }could not only appear with bare nouns, as shown in 5.3.1 (which was the most frequent construction in the 13th century, as shown in 5.2), but also with overt quantificational determiners as in 317, the frequency of which is very low in the 13th c.):
\end{styleStandard}

\begin{styleStandard}
\textbf{\ \ }a.\ \ Mas si por \textit{alguna otra razon qualquier} que no fuese delas sobre dichas enestas leyes deseredase el padre a su fijo […]
\end{styleStandard}

\begin{styleStandard}
\ \ \ \ ‘But if for some other (whatsoever) reason, that is not among the ones mentioned before in these laws, the father disinherits his son […]’
\end{styleStandard}

\begin{styleStandard}
\ \ \ \ (13th c., Alfonso X, \textit{Siete Partidas})
\end{styleStandard}

\begin{styleStandard}
\ \ b.\ \ [...] que ayan poder de tomar casas vinnas huertas oliuares molinos tiendas atahonas aldeas e \textit{todo otro heredamiento qualquiere} que sea de todo omne [....]. 
\end{styleStandard}

\begin{styleStandard}
\ \ \ \ ‘[…] that they have the power to take houses, vineyards, vegetable gardens, olive groves, tents, horse mills, hamlets, and any other inheritance, whichever it might be, from every person […].’
\end{styleStandard}

\begin{styleStandard}
\ \ \ \ (13th c., Alfonso X,\textit{ Doc. Cast. - Andalucía})
\end{styleStandard}

\begin{styleStandard}
\ \ c.\ \ o por abenençia o por aluedrio o por conposicion o por \textit{toda otra manera qualquiere} que el touiere por bien 
\end{styleStandard}

\begin{styleStandard}
\ \ \ \ ‘by the agreement or by arbitral award or by a deal or by any other manner that he approves of’
\end{styleStandard}

\begin{styleStandard}
\ \ \ \ (13th c., Alfonso X, \textit{Doc. Cast. - Castilla la Nueva})
\end{styleStandard}

\begin{styleStandard}
\ \ d.\ \ […] Desi dixo domjngujllo que le diese el Rey \textit{vn omne qual quier} a que le firiese e quando le oujese ferido que se yrie para el castillo 
\end{styleStandard}

\begin{styleStandard}
\ \ \ \ ‘Then Dominguillo said thet the King should give him a man which(ever) one wants whom he could hurt, and after hurting him, he would go to the castle.’
\end{styleStandard}

\begin{styleStandard}
\ \ \ \ (14th c., Anonymous, \textit{Veinte Reyes})
\end{styleStandard}

\begin{styleStandard}
This shows that \textit{qualquier} cannot be interpreted as a quantificational determiner at this stage yet. Otherwise we would end up with two different quantifiers that quantify over the same NP or over the same domain. This is impossible in natural language; consider, for instance, English *‘some whichever reason’ or *‘every whichever inheritance’ or *‘every whichever manner’ (see also Chapter 2 on the exclusion of double quantification). Let us analyse the example of \textit{todo otro heredamiento que} in 317, of which I only present the relevant part in 318:
\end{styleStandard}

\begin{styleStandard}
\textbf{\ \ }...que ayan poder de tomar casas, .... e \textit{todo otro heredamiento qualquiere que sea}
\end{styleStandard}

\begin{styleStandard}
\ \ ‘…that they can take every inheritance that one wants it to be’ 
\end{styleStandard}

\begin{styleStandard}
The interpretation is derived step by step in 319. First, the relative clause is interpreted as an open proposition due to the wh-feature of \textit{qual} and the clause structure (CP), then the relative clause is interpreted as a nominal modifier, that is, as a property, then the DP is interpreted as universal quantification. The truth conditions of the sentence in which the DP occurs are spelled out in 319d.\footnote{\textrm{\ Another possibility would be to interpret the postnominal }\textrm{\textsc{qualquier}}\textrm{ as a free relative clause without an overt antecedent. On this analysis, one could interpret the free relative clause as a proposition with a discourse variable and not as a property that modifies an antecedent NP:}\par 
\setcounter{listWWNumxxiileveli}{0}
\begin{listWWNumxxiileveli}
\item 
\begin{styleFootnote}
\textrm{Some other thing. Whatever it [=this thing] might be.}
\end{styleFootnote}
\end{listWWNumxxiileveli}
\par \textrm{I leave it for future research, to distinguish between the two options. }}
\end{styleStandard}

\begin{styleAuthor}
\ \ a.\ \ [[qualquier que sea]] =\ \ 
\end{styleAuthor}

\begin{styleAuthor}
\ \ \ \ (i) there might be(x)\ \ open proposition
\end{styleAuthor}

\begin{styleStandard}
\ \ b.\ \ [[heredamiento qualquier que sea]] =\textbf{\ \ }
\end{styleStandard}

\begin{styleStandard}
\ \ \ \ (ii) $\lambda $x[inheritance (x) ${\wedge}$ there might be (x)]\ \ property
\end{styleStandard}

\begin{styleStandard}
\ \ c.\ \ [[todo otro heredamiento qualquier que sea]] =\ \ 
\end{styleStandard}

\begin{styleStandard}
\ \ \ \ (iii) $\lambda $P. All x[other (x) ${\wedge}$ inheritance(x) ${\wedge}$ there might be ${\wedge}$ P(x)]\ \ universal quantifier
\end{styleStandard}

\begin{styleStandard}
\ \ d.\ \ [[e.g. (318) = que ayan poder de tomar todo otro heredamiento qualquier que sea]] = 1 iff\ \ 
\end{styleStandard}

\begin{styleStandard}
\ \ \ \ can take (All x [inheritance(x) ${\wedge}$ there might be (x)]) (they)\ \ 
\end{styleStandard}

\begin{styleStandard}
\ \ \ \ (prose: according to the law they can take every inheritance that there might be)
\end{styleStandard}

\begin{styleStandard}
The only difference between headed relatives as in 319 and free relatives suggested in 5.3.2 is that the nominal antecedent is overtly present in headed relatives in contrast to free relatives with a covert nominal antecedent in 5.3.2. The function of the relative clause CP \textit{qual quier que sea }is the same. It modifies an (overt) nominal antecedent and is thus of a property type {\textless}e,t{\textgreater} ,similar to adjectives or other nominal modifiers: 
\end{styleStandard}

\begin{styleStandard}
\textbf{\ \ }\textit{Qualquier} with an overt antecedent:
\end{styleStandard}

\begin{styleStandard}
\ \ a.[\textsubscript{DP}\textit{\textsubscript{{\textless}{\textless}e,t{\textgreater}t{\textgreater}}} alguna/toda [\textsubscript{Spec,NP} otra \textsubscript{{\textless}e,t{\textgreater}} [\textsubscript{NP {\textless}e,t{\textgreater}} razon/manera [\textsubscript{CP{\textless}e,t{\textgreater}} [\textsubscript{Spec,CP} [\textsubscript{D°} \textit{qualquier}\textsubscript{j}] \textsubscript{C[2032?] C°} \textit{que} \textsubscript{TP2} \textit{pro}\textit{\textsubscript{ }}\textsubscript{T°} \textit{sea} t\textsubscript{j}]]]]
\end{styleStandard}

\begin{styleStandard}
[Warning: Draw object ignored]
\end{styleStandard}

\begin{styleStandard}
I will show now what happened after \textit{qual quier que… }was no longer analysed as a relative clause in headed relatives. The idea in a nutshell is that the string \textit{qual quier }lost its transparent status as consisting of the wh-pronoun \textit{qual} and the verbal form \textit{quier} (see 5.3.1); that is, the wh-feature disappeared, the verb was no longer analysed as such, with the consequence that the embedded clause introduced by \textit{que} was no longer necessary. This is why the frequency of \textit{qualquier,-e,-a que}… decreases, as already shown in 5.2. The postnominal \textit{qualquier,-e,-a} thus can appear without the embedded clause (see, e.g., 321):
\end{styleStandard}

\begin{styleStandard}
\textbf{\ \ }\ \ o con cuchillo o con \textit{otra cosa qualquier} 
\end{styleStandard}

\begin{styleStandard}
\ \ \ \ ‘or with a knife or with some other thing whichever one wants’ 
\end{styleStandard}

\begin{styleStandard}
\ \ \ \ (13th c., Alfonso X, \textit{Siete Partidas})
\end{styleStandard}

\begin{styleStandard}
The crucial question is what syntactic category and semantic interpretation the postnominal \textit{qualquier,-e,-a} acquired in structures like 321, once it was no longer analysed transparently. 
\end{styleStandard}

\begin{styleStandard}
I assume that the postnominal \textit{cualquiera} has preserved its syntactic status as a nominal modifier (like relative clauses) and its semantic status as a property of type {\textless}e,t{\textgreater}, as shown in . The only difference to relative clauses like those in 320 is that the new postnominal \textsc{cualquier} is no longer clausal, but a complex compound word (for details, see Chapter 3):
\end{styleStandard}

\begin{styleStandard}
\ \ [DP una [NP cosa/razon [ModifP \textit{cualquiera}]]] 
\end{styleStandard}

\begin{styleStandard}
[Warning: Draw object ignored]
\end{styleStandard}

\begin{styleStandard}
As I already pointed out, postnominal \textit{cualquiera} does not have the evaluative meaning before the 18th century (see 5.1 and 5.2). The free choice interpretation of the postnominal modifier is triggered by some modal verb, subjunctive, or generic operator (see Chapter 2):
\end{styleStandard}

\begin{styleStandard}
\textbf{\ \ }postnominal \textit{cualquiera}
\end{styleStandard}

\begin{styleStandard}
\ \ …sea debajo de la equinoccial o (sea) en \textit{otra parte cualquiera}» A este propósito, trae Celio Rodiginio lo de Arriano, historiador griego
\end{styleStandard}

\begin{styleStandard}
\ \ ‘… be it under the equinoctial or (be it) in any other place» For this purpose, Celio Rodiginio says of Arriano, Greek historian’ (brackets my own)
\end{styleStandard}

\begin{styleStandard}
\ \ (16th c., Torquemada,\textit{ Jardín})
\end{styleStandard}

\begin{styleStandard}
The free choice meaning of the postnominal modifier \textit{cualquiera} in 324 is similar to the meaning of a postnominal relative clause \textit{qual quiera} at earlier stages (see 5.3.1). It can be roughly paraphrased as some place x and x can have any value (one wants). The semantic shift from FCI to Eval in the 18th century in 324 can be analysed as a shift from the free choice \textit{cualquiera} as ‘any (one wants)’ to a property \textit{cualquiera} that ‘anyone has’ and thus to a property with the meaning ‘ordinary’, because this property does not distinguish individuals from each other (see Chapter 3 for details): 
\end{styleStandard}

\begin{styleStandard}
\textbf{\ \ }Juan es un hombre cualquiera (para mí)
\end{styleStandard}

\begin{styleStandard}
\ \ ‘Juan is a man with a property that anyone has’ (literal) ={\textgreater}
\end{styleStandard}

\begin{styleStandard}
‘Juan is an ordinary man (to me).’
\end{styleStandard}

\begin{styleStandard}
\ \ (@ 151)
\end{styleStandard}

\begin{styleStandard}
The evaluative interpretation has triggered a further shift of postnominal modifying \textit{cualquiera} into nominal \textit{cualquiera} as in \textit{es una cualquiera, }which appears late in the CdE (from the 19th\textsuperscript{ }c.). Note that, semantically, this nominal \textit{cualquiera} is still interpreted as an evaluative predicate of the type {\textless}e,t{\textgreater}, but behaves differently from the postnominal function in that it also encodes the feature [+animate] and female gender:
\end{styleStandard}

\begin{styleStandard}
\textbf{\ \ }\textit{¡ya se ve!, es una cualquiera};\newline
‘One can already see! She is a promiscuous woman;’ 
\end{styleStandard}

\begin{styleStandard}
\ \ (19th c., Rizal,\textit{ Noli me tangere} )
\end{styleStandard}

\begin{styleStandard}
Recall that the semantic shift into the evaluative predicate has triggered a further syntactic reanalysis in Argentinian Spanish as a gradable adjective (see Chapter 4):
\end{styleStandard}

\begin{styleStandard}
\ \ Juan es muy cualquiera (para mí)\ \ (ArgSp)
\end{styleStandard}

\begin{styleStandard}
\ \ [ordinary[2032?] (Juan) (to me) (to a high degree)]\ \ 
\end{styleStandard}

\begin{styleStandard}
\ \ ‘Juan is very common/ordinary to me.’\ \ 
\end{styleStandard}

\begin{styleStandard}
\ \ \ \ (@ 152)
\end{styleStandard}

\begin{styleStandard}
I assume that it is the semantic interpretation of relative clauses as properties that enables one to interpret \textit{cualquiera} as a property with an evaluative meaning ‘ordinary’ in predicative position such as \textit{es un hombre cualquiera} ‘he is an ordinary man’ or as an evaluative property of women that has been lexicalised in \textit{una cualquiera} (Chapter 3 on Eval). The rise of the indicative mood and predicative verbs has further enhanced the property interpretation of \textit{cualquiera} as in N \textit{cualquiera} (see Chapter 3). 
\end{styleStandard}

\begin{styleStandard}
\ \ The different analyses of prenominal and postnominal \textit{cualquiera} explains why only the postnominal \textit{cualquiera} has different lexical meanings (FCI, EVAL) in European Spanish, but not the prenominal \textit{cualquiera}. The reason for this is that only the postnominal \textit{cualquiera} has a predicative status and is interpreted as a nominal modifier similar to adjectives, whereas the quantificational determiner is a functional category that is more restricted in meanings than a lexical category.
\end{styleStandard}

\subsection[5.4 Summary and discussion]{\textbf{5.4 Summary and discussion}}
\begin{styleStandard}
In this section, I have shown on the basis of numerically quantified and statistically analysed data how \textsc{qualquier} has changed on the syntactic and semantic level. I have suggested an analysis that shows how the reanalysis of \textsc{qualquier} as a relative clause into a quantificational determiner (which C\&P 2009 and others call grammaticalisation of \textsc{qualquier}) has taken place, namely from a relative clause to lexicalisation as the single word \textsc{qualquier }(without the ending \textit{{}-a}), which was subsequently reanalysed as a quantificational determiner D° with FC interpretation ‘any(one/thing)’:
\end{styleStandard}

\begin{styleStandard}
\textbf{\ \ }Change from relative clause to a determiner (so-called Grammaticalization)
\end{styleStandard}

\begin{styleStandard}
\ \ \ \ Free Rel. Clause [F0E0?] Lexicalisation as a single word [F0E0?] Determiner as in \textit{cualquier} N
\end{styleStandard}

\begin{styleStandard}
(Interpretation of Determiner \textit{cualquier}: FCI only)
\end{styleStandard}

\begin{styleStandard}
I have suggested that the change of postnominal \textsc{qualquier}\textit{ }(with the ending \textit{{}-a}) was a change from a headed relative clause into a postnominal modifier with the subsequent change into a noun and a degree adjective.
\end{styleStandard}

\begin{styleStandard}
\textbf{\ \ }Headed Rel. Clause [N \textit{qual quiera}] [F0E0?] Lexicalisation of \textit{qualquiera} as a single word [F0E0?]
\end{styleStandard}

\begin{styleStandard}
\ \ \ \ Interpretation of \textit{cualquiera} as a property (EVAL) (e.g. \textsc{un}\textit{ N cualquiera}) [F0E0?] degree variable of the property \textit{cualquiera} in ArgSp (e.g. \textit{una pelí muy cualquiera} EVAL)
\end{styleStandard}

\begin{styleStandard}
Recall from Chapter 2.6 that C\&P (2009) assume that \textsc{qualquier} has lost gradually its semantic interpretation under the assumption that grammaticalisation correlates with a loss or weakening of semantic content (Lehmann 2002 [1982], C\&P 2009). If C\&P (2009) mean that \textsc{qualquier} lost the literal meaning of \textit{querer} ‘want’, I agree with this observation. The crucial change according to my analysis is the change in the agent of Free Choice Modality (i.e. from the generic \textit{pro} of verb \textit{quier} to the agent of the matrix verb). Unlike what is usually claimed for grammaticalisation, indefinites such as \textsc{qualquier} do not lose semantic content. The reason may be that this is not a prototypical case of grammaticalisation. \textsc{qualquier} certainly has grammaticalised to some degree, as it lost the transparent relative clause interpretation of the wh-word \textit{qual} + modal verb \textit{querer }and was reanalysed as a single word (see Lehmann 2002 [1982] on prototypical features of grammaticalisation). However, further semantic and syntactic changes of (postnominal) \textit{cualquiera} in European and of \textit{cualquiera} in Argentinian Spanish appeared. Strictly speaking, the postnominal \textit{cualquiera} did not grammaticalise into a functional category like the prenominal \textit{qual quier }did in European Spanish in the 13th century. This is probably why postnominal \textit{cualquiera} and prenominal \textit{cualquier} differ in syntax and semantic interpretations in European Spanish. 
\end{styleStandard}

\clearpage\section[6 Overall summary and conclusions]{6 Overall summary and conclusions}
\begin{styleStandard}
This thesis had two major aims. The first aim was to probe a compositional analysis of meaning change of modal indefinites. In order to achieve this goal, it was necessary to characterise and define different meanings of \textit{cualquiera} in Modern Spanish varieties. The reason why \textit{cualquiera} is particularly interesting for a linguistic study is that this word has different meanings that appear to be unrelated such as the Free Choice (FC) meaning ‘any’ and evaluative meanings such as ‘ordinary’ and ‘low value’ (which I called EVAL1) and ‘nonsense’ (EVAL 2). I concentrated on European and Argentinian Spanish, because these varieties show different degrees of lexicalisation of these meanings. One research question was whether the different meanings of \textit{cualquiera} in Spanish varieties represent a case of polysemy (Hypothesis 1) or whether \textit{cualquiera} has only one basic meaning (assumedly the FC meaning ‘any’) and other meanings can be inferred from the general meaning pragmatically (Hypothesis 2). Moreover, I investigated the relation between syntax and semantics, i.e. whether certain meanings appear under certain morphosyntactic conditions, which is the core assumption of the compositional framework of meaning change. The second aim of this thesis was to investigate the diachrony of these meanings, i.e. when these meanings appeared in time and under which linguistic conditions they occurred. The major scope of this thesis is thus a synchronic and diachronic description and analysis of different meanings of \textit{cualquiera} in European and Argentinian Spanish. In the first half of this thesis, I investigated the synchrony and in the second half the diachrony of \textit{cualquiera}. I will summarize how I proceeded in my analysis and what are the most important results and new findings to the study of \textit{cualquiera}.
\end{styleStandard}

\begin{styleStandard}
\ \ In order to describe the synchronic meanings of \textit{cualquiera}, I first reviewed the data and various analyses as discussed in the literature concerning \textit{cualquiera} in Modern European and Argentinian Spanish. I have shown that semantic and syntactic analyses are missing in the literature on \textit{cualquiera} with EVAL 1 ‘unremarkable’ present in European and Argentinian Spanish and EVAL 2 ‘nonsense’ present only in Argentinian Spanish.
\end{styleStandard}

\begin{styleStandard}
\ \ In Chapter 3, I presented new data and a new analysis of N \textit{cualquiera} with EVAL 1. I analysed new data from a subcorpus of Corpus del Español (CdE\_web) and have identified various flavors of the evaluative meaning ‘unremarkable’, namely COM (common: ‘regular, normal’, DER (derogative: ‘ordinary, without any particular quality worth mentioning’). I argued that EVAL 1 in European Spanish is a special case of interpreting \textit{cualquiera} as a quantifier with a free choice interpretation ‘any’. I have shown that by deriving the evaluative meaning from the quantifier meaning, \textit{cualquiera} is not a case of polysemy (i.e. Hypothesis 2 of Section 1.1 is verified and Hypothesis 2 is falsified, contra Alonso Ovalle \& Menendez Benito 2018). In order to derive the evaluative meaning ‘ordinary’ from ‘any’, I have analysed \textit{cualquiera} as a quantificational predicate (modifier) similar to adjectives, which has the interpretation of a comparative clause expressing similarity (SIM(ilarity)), i.e. \textit{Juan es un hombre cualquiera} is interpreted as ‘Juan is a man and he is like anyone else’ (whereby ‘anyone else’ is a set of comparable male individuals, see Section 3.3). All other flavors mentioned above are pragmatic inferences of the SIM meaning ‘like anyone else’. The free choice meaning ‘any’ is the basic meaning of \textit{cualquiera} and EVAL 1 ‘ordinary’ follows from this basic meaning. This analysis is the first analysis that provides an explanation how different meanings and functions of \textit{cualquiera} are interconnected and it provides further support for the compositionality view on meaning change.
\end{styleStandard}

\begin{styleStandard}
\ \ I discussed a small data set in the corpus CdE\_web (e.g., definite determiner N \textit{cualquiera}) that shows signs of lexicalization. According to this data set, the SIM meaning (i.e. \textit{cualquiera} as ‘like anyone else’) has been reanalyzed as a property of kinds of individuals with the meaning ‘common/typical/ordinary’. This lexical meaning co-occurs with definite nouns and quantified nouns (e.g. \textit{el hombre cualquiera} ‘the typical person/man’ or \textit{algún día cualquiera} ‘on some ordinary day’) as well as with coordination (e.g. \textit{un ser mediocre y cualquiera} ‘a mediocre and ordinary human being’).
\end{styleStandard}

\begin{styleStandard}
\ \ In Chapter 4, I analysed \textit{cualquiera} with EVAL 1 and EVAL 2 in Argentinian Spanish and suggested an analysis for the interpretations and licensing of \textit{cualquiera} in episodic contexts in this variety. Episodic contexts usually do not license modal indefinites like \textit{cualquier} ‘any’. I suggested that \textit{cualquiera} with EVAL 1 should be analysed as a degree adjective in Argentinian Spanish. I have argued that the difference between EurSp and ArgSp \textit{cualquiera} with EVAL 1 is primarily syntactic: In ArgSp but not in EurSp, \textit{cualquiera} is used as a predicative adjective with degree modification (e.g. \textit{es muy cualquiera} ‘he/she/it is very ordinary’). I analysed the meaning of the adjective \textit{cualquiera} as ‘bad’, as a pejorative interpretation of \textit{cualquiera} as ‘common’ or ‘ordinary’ under the assumption that speakers who use the pejorative version ‘bad’ assume that ‘being ordinary’ counts as a negative property in contrast to ‘being special’ or ‘having a distinctive property’. I have proposed analysing EVAL 2 ‘nonsense’ in parallel with the free choice indefinite\textit{ cualquier cosa }‘anything’\textit{, }which is interpreted in episodic contexts as ‘anything possible’ or as ‘something random’. The evaluative meaning ‘nonsense’ is a pragmatic inference, which comes about as follows. The judge (usually the speaker) evaluates the modal component of the free choice \textit{cualquiera} (‘every possibility is an option’) on a scale of preferences, namely as the least preferred option. Under verbs of saying like \textit{decir} or \textit{mandar} ‘to say’ in ArgSp, this pragmatic inference works as follows. The expression ArgSp. \textit{decir/mandar cualquiera} ‘saying anything possible’ implies saying contradictory things or saying something random or saying unreflected things according to the judge (usually the speaker). This pragmatic evaluation has been lexicalised in Argentinian Spanish under transitive verbs of saying and this evaluative meaning further generalized to other verbs, e.g. predicative verbs (e.g. ArgSp. \textit{lo que me decís es cualquiera} ‘what you are telling me is bullshit/nonsensical’). The diachronic data of Argentinian Spanish \textit{cualquiera} is unfortunately too scarce to draw strong conclusions. Nevertheless, there are some hints that suggest that the origin of the meaning ‘nonsense’ is its occurrence under verbs of \textit{decir} ‘to say’. Concerning the research question of this thesis whether EVAL 1 and EVAL 2 are lexicalized meanings\textit{, }I have argued that the answer is positive for Argentinian Spanish (or, in other words, Hypothesis 1 holds) as these meanings correlate with different syntactic categories than the indefinite \textit{cualquiera} with FC ‘anything’, namely degree adjectives and evaluative nouns. However, what looks like a case of polysemy in Modern Argentinian Spanish is not a polysemy from the diachronic perspective. I argued that, in diachrony EVAL 1 and EVAL 2, are the result of pragmatic inferencing that has been lexicalized and has driven syntactic reanalysis into different categories in Modern Argentinian Spanish, namely from a pronoun into gradable adjective with the meanings ‘ordinary’ and ‘nonsensical’ and into an evaluative noun with the meaning ‘nonsense’. I have shown that this path of reanalysis from FC \textit{cualquiera} into evaluative \textit{cualquiera }is not a special property of \textit{cualquiera} in Argentinian Spanish, but is also present in French (e.g. \textit{n’importe quoi}), Mexican Spanish \textit{equis} (‘any’ in modal contexts and ’ordinary’ under predicative verbs) (Author 2020) and other, non-Romance languages, as for instance in Czech (see Šimík 2008). Chapter 4 presents the first in-depth contrastive analysis of \textit{cualquiera} in Argentinian and European Spanish, based on corpus data and quantitative analysis. 
\end{styleStandard}

\begin{styleStandard}
\ \ After the synchronic analysis of evaluative meanings and the diachronic analysis of EVAL 2 in Argentinian Spanish, I focused on the diachronic analysis of \textit{cualquiera} in European Spanish in order to find evidence for my synchronic analysis, namely that the free choice meaning is the basic meaning and the evaluative meanings are secondary or following from the free choice meaning. My assumption was that the free choice meaning occurred diachronically first and evaluative meanings later. Moreover, I looked at the change of linguistic conditions (e.g. change of episodic contexts, syntactic position, etc.) to see how the meanings arise diachronically. 
\end{styleStandard}

\begin{styleStandard}
\ \ In Chapter 5, I have shown on the basis of quantitative data and their statistical analysis on the basis of historical data from CdE when and how \textit{cualquiera} has developed its functions and meanings in the history of Old and Middle Spanish as well as Modern European Spanish. I analysed different syntactic properties such as the position of \textit{cualquiera} with respect to the verb and the noun and have shown on the basis of statistical tests that \textit{cualquiera} has been used increasingly in different syntactic structures such as declaratives, in postverbal position, etc., which I interpreted as a hint for the grammaticalisation hypothesis. I have modeled the path of grammaticalisation of \textit{cualquiera} in a generative framework and have analysed the meaning components that have changed from the relative clause with free choice interpretation \textit{qual quier,e,a} ‘which one wants’ to \textit{qualquier,e,a} ‘whichever’. The most important semantic component that changed from the meaning of relative clause \textit{qual quier,e,a} to the free choice indefinite \textit{qualquier,e,a} is related to the judge parameter. The free choice interpretation of the indefinite (in contrast to the relative clause) is no longer dependent on the generic \textit{pro} as in \textit{qual pro}\textit{\textsubscript{generic}}\textit{ quier,e,a,} but can be associated with the agent of the matrix clause, (e.g. \textit{Pablo} \textit{puede invitar a cualquiera }‘Pablo can invite anyone’). What both linguistic expressions (relative clause Old Spanish \textit{qual quier,e,a} and the modern indefinite \textit{cualquiera}) have in common is that they do not contribute to the truth conditional semantic meaning of the clause ‘Pablo can invite someone’. They have rather the function of a side-comment, e.g. ‘Pablo has free choice of inviting whoever he likes’. I have shown diachronic evidence for the side-comment analysis as relative clauses such as \textit{qual quier,e,a} behave very much like adjuncts and they can even interrupt a sentence (e.g. \textit{pon una estrella fixa qual quisieres por su longura del eguador} \ ‘put a star [which you want] along the equator’).
\end{styleStandard}

\begin{styleStandard}
\ \ Another important finding of my diachronic study is that the postnominal \textit{cualquiera} is the result of a lexicalisation process by which the relative clause \textit{qual} \textit{quier,e,a} has been lexicalised into a single compound and has persevered its semantic function as a nominal property (Noun+ qualquier,e,a). The evaluative meaning of postnominal \textit{cualquiera}, which appeared in the 17\textsuperscript{th} century, is thus directly related to the property function of relative clauses in a postnominal position in contrast to the prenominal quantifier \textit{cualquier}. I have concluded that the postnominal \textit{cualquiera} did not undergo grammaticalisation in the strict sense as it preserved the same function as in Old Spanish. The diachronic analysis in this chapter is the first in-depth diachronic analysis that combines statistical analysis with a theoretical explanation of meaning change of \textit{cualquiera} and offers a new perspective on the grammaticalisation of \textit{cualquiera}. A more general contribution to the research on grammaticalisation is that I have shown that certain categories like indefinites do not necessarily correlate with a loss in meaning as is standardly assumed with respect to grammaticalisation. Moreover, I have provided evidence that meaning change is not unidirectional, that is, functional categories can change into lexical categories, as is the case of the evaluative meaning of \textit{cualquiera. }
\end{styleStandard}

\begin{styleStandard}
\ \ In sum, this book provides further support for the transparency and compositionality principle involved in meaning change and variation in the domain of modal indefinites like Spanish \textit{cualquiera} and it shows that the grammaticalization does not need to co-occur with meaning loss as standardly assumed in the literature on grammaticalization. 
\end{styleStandard}

\clearpage\section[7 References]{\textbf{7} \textbf{References}}
\subsection[7.1 Corpora and electronic dictionaries]{7.1 Corpora and electronic dictionaries}
\begin{styleStandard}
CdE \textit{= Corpus del Español} (Ed. Mark Davies, 2001-) [URL: \url{https://www.corpusdelespanol.org/}]\textit{.}
\end{styleStandard}

\begin{styleStandard}
CdE\_G/H\textit{ = Corpus del Español} (Genre / Historical Ed. Mark Davies, 2001-) [URL: \url{https://www.corpusdelespanol.org/hist-gen/}]\textit{.}
\end{styleStandard}

\begin{styleStandard}
\textit{CdE web dialects = Corpus del Español: Two billion words, 21 countries }(Web / Dialects. Ed. Mark Davies, 2016-) [URL: \url{http://www.corpusdelespanol.org/web-dial/}].
\end{styleStandard}

\begin{styleStandard}
\textit{CdE now }=\textbf{ }\textit{Corpus of News on the Web (NOW): 10 billion words from 20 countries, updated every day}. (NOW. Ed. Mark Davies 2016-) [URL: \url{https://www.corpusdelespanol.org/now/}]
\end{styleStandard}

\begin{styleStandard}
CORDE = \textit{Corpus Diacrónico del Español} [URL: \url{http://corpus.rae.es/cordenet.html}]. 
\end{styleStandard}

\begin{styleStandard}
CORDIAM = \textit{Corpus Diacrónico y Diatópico del Español de América}. [URL: \url{https://www.cordiam.org/}].
\end{styleStandard}

\begin{styleStandard}
DLE, s.v. \textit{cualquiera}. \textit{Diccionario de la lengua española}, 23.4 online edn. 
\end{styleStandard}

\begin{styleStandard}
\ \ [URL: \url{https://dle.rae.es/cualquiera?m=form}\textstyleInternetlink{]}
\end{styleStandard}

\begin{styleStandard}
DILe Diccionario Latinoamericano de la Lengua Española 
\end{styleStandard}

\begin{styleStandard}
\ \ [URL: \url{http://untref.edu.ar/diccionario/uso.php}]
\end{styleStandard}

\subsection[7.2 Primary sources]{7.2 Primary sources}
\subsubsection[7.2.1 Internet links of examples]{7.2.1 Internet links of examples}
\begin{flushleft}
\tablefirsthead{}
\tablehead{}
\tabletail{}
\tablelasttail{}
\begin{supertabular}{m{0.7955598in}m{5.35876in}}
@ 1  &
\url{http://salomonlaverdad.blogspot.com/2013/08/en-este-pais-el-banco-puede-llegar.html}, as accessed in Februar 2021.\\
@ 2  &
\url{https://www.ipadizate.es/2013/05/18/jailbreak-ios-613-ipad-mini-70915/}, as accessed in Februar 2020.\\
@ 3  &
https://es.ign.com/lost-planet-3/65870/preview/avance-de-lost-planet-3-para-ps3-360-y-pc, as accessed in June 2021.\\
@ 4  &
https://www.fotogramas.es/peliculas-criticas/a524660/cuando-todo-esta-perdido/, as accessed in June 2021.\\
@ 5  &
\url{https://twitter.com/FachinelliR/status/407370344672923648}, as accessed in May 2021.\\
@ 6  &
\url{https://twitter.com/Hernan_Merenzon/status/1258651692809191426}, as accessed in May 2021.\\
@ 7  &
\url{https://nomeolvidedevos.com/category/uncategorized/page/6/}, as accessed in May 2021.\\
@ 8  &
\url{https://es.blastingnews.com/salud-belleza/2014/12/ocho-de-cada-diez-veganos-vuelven-a-comer-carne-por-problemas-fisicos-y-psicologicos-00196381.html}, as accessed in Februar 2020.\\
@ 9  &
\url{https://elsumario.com/los-gatos-no-comen-cualquier-cosa/}, as accessed in May 2021.\\
@ 10  &
\url{https://orlandosuarez.net/2007/01/25/a-modo-de-recordatorio/}, as accessed in May 2021.\\
@ 11  &
\url{https://www.periodistadigital.com/cultura/religion/20121107/victoria-molins-monja-enamorada-jesus-noticia-689401361348/}, as accessed in May 2021.\\
@ 12  &
\url{https://elpais.com/diario/1998/07/26/cvalenciana/901480698_850215.html}, as accessed in December 2020. \\
@ 13  &
\url{https://historiadospodcast.wordpress.com/2019/09/17/magazine-12-cualquiera-lo-hubiera-hecho-mejor/}, as accessed in December 2020.\\
@ 14  &
\url{http://gramaticaespanola.com/publ/sintaxis/los_cuantificadores_ii_caracteristicas_gramaticales_de_los_principales_cuantificadores/el_cuantificador_cualquiera/23-1-0-388}, as accessed in December 2020.\\
@ 15  &
\url{https://www.elotrolado.net/hilo_ohio-proyecto-de-ley-antiaborto-obligaria-a-medicos-hacer-algo-cientificamente-imposible_2357962}, as accessed in December 2020.\\
@ 16  &
\url{http://www.galiciaartabradigital.com/archivos/179321}, as accessed in December 2020.\\
@ 17  &
\url{http://migrantesdelbalon.com/marcos-gondra-subir-de-categoria-fue-un-aliciente-para-seguir-aqui/}, as accessed in May 2021.\\
@ 18  &
\url{https://sede.csn.gob.es/Sede20/registro-general-rec?tipoPagina=p3&idItem=141&lang=es&usu=null}, as accessed in June 2021.\\
@ 19  &
\url{https://www.fnac.es/Ignatius-Farray-La-comedia-como-salvavidas/cp6068/w-4}, as accessed in June 2021.\\
@ 20  &
\url{https://www.elmundo.es/loc/celebrities/2017/09/01/59a8484546163ff6758b45e5.html}, as accessed in June 2021.\\
@ 21  &
\url{https://www.guitarristas.info/foros/como-adquirir-mas-velocidad-problemas-cansancio/220836}, as accessed in June 2021.\\
@ 22  &
\url{http://nosolocoaching.com/tag/sesgos-cognitivos/}, as accessed in May 2021.\\
@ 23  &
https://parclick.es/parking-benidorm/parking-benidorm-city-center-transfer, as accessed in June 2021.\\
@ 24  &
\url{http://www.agarzon.net/los-gritos-de-las-victimas/}\textstyleInternetlink{, as accessed in March 2017.}\\
@ 25  &
\url{https://www.uv.es/~ivorra/Filosofia/TC/14.htm}, as accessed in December 2020.\\
@ 26  &
\url{https://www.asivaespana.com/otros/ortega-smith-no-esta-preparado-para-ciertos-temas}, as accessed in June 2021.\\
@ 27  &
\href{https://www.apologeticamariana.com/por-que-honrar-y-amar-a-la-virgen-maria/,%20}{\textstyleInternetlink{https://www.apologeticamariana.com/por-que-honrar-y-amar-a-la-virgen-maria/,}} as accessed in December 2020.\\
@ 28  &
\url{https://cvc.cervantes.es/foros/leer_asunto1.asp?vCodigo=44034}, as accessed in June 2021.\\
@ 29  &
\url{https://www.eleconomista.es/empresas-finanzas/noticias/6948017/08/15/Wang-Jianlin-dueno-del-Edificio-Espana-nombrado-el-chino-mas-rico-del-mundo.html}, as accessed in June 2021.\\
@ 30  &
\url{https://www.senscritique.com/film/Avalanche_Sharks_les_dents_de_la_neige/critique/233566206}, as accessed in May 2021.\\
@ 31  &
\url{https://www.ouest-france.fr/pays-de-la-loire/le-mans-72000/le-mans-reouverture-des-terrasses-c-est-du-grand-n-importe-quoi-pour-ce-restaurateur-233ee604-b4c6-11eb-b617-49ac21059b35}, as accessed in May 2021.\\
@ 32  &
\url{https://es-es.facebook.com/twx02/posts/524279494781343}, as accessed in June 2021.\\
@ 33  &
\url{http://www.foros-fiuba.com.ar/viewtopic.php?p=201073}\textstyleInternetlink{, as accessed in May 2021.}\\
@ 34  &
\url{http://ivegotbrokenwings.blogspot.com/2009/}, as accessed in May 2021.\\
@ 35  &
\url{https://brainly.lat/tarea/185905}, as accessed in Fabruar 2021.\\
@ 36  &
\url{https://malosfics.foroes.org/t14929-una-maid-a-tu-servicio-por-dianis-mar-vocaloid}, as accessed in May 2021.\\
@ 37  &
\url{http://passerinaversicolor.blogspot.com/2010/12/el-pajaro-dulce-encanto.html}, as accessed in Februar 2021.\\
@ 38  &
https://cadizcarnaval.es/mas-carnaval/baul-de-letras/150-comparsa-calle-de-la-mar-conservo-de-ti, as accessed in Februar 2021.\\
@ 39  &
\url{https://www.llibrerialacultural.com/es/libros/el-dia-de-hoy_4419980634}, as accessed in June 2021.\\
@ 40  &
http://allaenelfarodelfindelmundo.blogspot.com, \ as accessed in December 2020.\\
@ 41  &
\url{http://lavidasecretadeloslibros.blogspot.com/2012/04/resena-las-pruebas-james-dashner.html}, as accessed in December 2020.\\
@ 42  &
\url{https://www.agapea.com/Gustavo-Martin-Garzo/Y-que-se-duerma-el-mar-9788490321942-i.htm}, as accessed in December 2020.\\
@ 43  &
\url{https://www.lesbiana.es/2013/08/23/historias-de-♀-sayonara-baby-parte-12-y-final/}, as accessed in December 2020.\\
@ 44  &
\url{http://www.cabovolo.com/2010/04/container-bbc-box-viaje.html}, as accessed in Februar 2021.\\
@ 45  &
\url{https://www.diariovasco.com/alto-urola/yasmina-bianca-fleschin-20190524002023-ntvo.html}, as accessed in June 2021.\\
@ 45  &
\url{https://diariorombe.es/los-jovenes-de-bata-deambulan-por-las-calles-en-busca-de-electricidad-para-cargar-los-telefonos/}, as accessed in June 2021.\\
@ 46  &
\url{https://aminoapps.com/c/k-pop-es/page/blog/como-conoci-a-bts/d3Lh_buLr8wqDVaQE8Pnpzdq5VWZgJEtP}, as accessed in June 2021.\\
@ 47  &
\url{https://home.cc.umanitoba.ca/~fernand4/empareja.html}, as accessed in Februar 2021.\\
@ 48  &
\url{https://espidofreire.com/index.php?option=com_zoo&task=item&item_id=36&Itemid=304}, as accessed in Februar 2021.\\
@ 49  &
\url{https://www.ole.com.ar/boxeo/preso-destino_0_B19dtm1she.html}, as accessed in Februar 2021.\\
@ 50  &
\url{http://www.teknlife.com/reportaje/el-peligro-de-no-pensar-por-uno-mismo/}, as accessed in June 2021.\\
@ 51  &
\url{https://www.wattpad.com/23260961-regreso-al-pasado-relato}, as accessed in Februar 2021.\\
@ 52  &
\url{https://www.lavanguardia.com/internacional/20171002/431752309557/autor-tiroteo-las-vegas-stephen-paddock.html}, as accessed in Februar 2021.\\
@ 53  &
\url{http://www.destinorpg.es/2015/11/analisis-divinity-original-sin-enhanced.html}, as accessed in June 2021.\\
@ 54  &
\url{https://www.eluniverso.com/noticias/2016/07/09/nota/5680188/pescador-salinas-libero-redes-ballenato/}, as accessed in Februar 2021.\\
@ 55  &
\url{https://igualdadanimal.org/noticia/2016/05/26/asi-seria-un-dia-cualquiera-si-fueses-un-animal-de-granja/}, as accessed in June 2021.\\
@ 56  &
\url{https://www.rockandpop.cl/2017/01/ganate-una-experiencia-chevymusic-in-the-city/}, as accessed in Februar 2012.\\
@ 57  &
\url{https://twitter.com/giantsgaming/status/1253686474769342464?lang=es}, as accessed in June 2021.\\
@ 58  &
\url{https://www.mundodeportivo.com/20150512/2076215618/beckenbauer-sin-messi-el-barcelona-seria-un-equipo-cualquiera.html}, as accessed in Februar 2021.\\
@ 59  &
\url{https://nuevoperiodico.com/para-la-mayoria-de-los-encuestados-el-9-de-mayo-es-el-dia-de-la-victoria/}, as accessed in June 2021.\\
@ 60  &
\url{https://lahierbaroja.com/category/literatura/libros-no-memorables/}, as accessed in June 2021.\\
@ 61  &
\url{https://www.tripadvisor.es/ShowUserReviews-g642221-d9704826-r640124480-Bar_Pizzeria_Catania-Mataro_Catalonia.html}, as accessed in June 2021.\\
@ 62  &
\url{https://sevilla.abc.es/alfinaldelapalmera/noticias-betis/sevi-gordillo-cede-numero-3-manu-carrasco-179716-1561404405-201906242126_noticia.html?ref=https:%2F%2Fwww.corpusdelespanol.org%2Fnow%2Fx3.asp%3Fxx%3D1}, as accessed in December 2020.\\
@ 63  &
http://www.menshealth.es/articulo/La-mejor-dieta-para-tus-musculos, as accessed in March 2017.\\
@ 64  &
\url{http://sonicando.com/?p=2687}, as accessed in December 2020.\\
{\bfseries \textmd{@ 66 }} &
{\bfseries \url{https://www.elperiodicodearagon.com/noticias/aragon/filosofia-llama-puerta_1319416.html}\textstyleInternetlink{\textmd{, }}\textmd{as accessed in December 2020.}}\\
{\bfseries \textmd{@ 67 }} &
{\bfseries \url{https://twitter.com/papel_xd}\textmd{, as accessed in May 2021.}}\\
{\bfseries \textmd{@ 68 }} &
{\bfseries \url{http://blog.webershandwick.es/page/4/?_=1476835200000}\textmd{, as accessed in June 2021.}}\\
{\bfseries \textmd{@ 69 }} &
{\bfseries \url{http://www.asociacionapadevi.com/blog/wp-content/uploads/2017/12/Bases-protocolo-policia-maltrato-animal-M%C2%BA-interior.pdf}\textmd{, as accessed in June 2021.}}\\
{\bfseries \textmd{@ 70 }} &
{\bfseries \textmd{https://www.lexdixit.es/la-detencion-en-espana/, as accessed in June 2021.}}\\
{\bfseries \textmd{@ 71 }} &
{\bfseries \url{https://www.collinsdictionary.com/dictionary/english-spanish/any}\textmd{, as accessed in Februar 2021.}}\\
{\bfseries \textmd{@ 72 }} &
{\bfseries \url{https://www.filmaffinity.com/es/userreviews/1/293307.html}\textmd{, as accessed in June 2021.}}\\
{\bfseries \textmd{@ 73 }} &
{\bfseries \url{https://www.pinterest.de/pin/7388786875893171/?autologin=true}\textmd{, as accessed in June 2021.}}\\
{\bfseries \textmd{@ 74 }} &
{\bfseries \url{https://www.elespanol.com/deportes/juegos-olimpicos/20170709/229977406_0.html}\textmd{, as accessed in June 2021.}}\\
{\bfseries \textmd{@ 75 }} &
{\bfseries \url{http://www.izquierdadiario.es/Julio-Cortazar-los-caminos-del-escritor}\textmd{, as accessed in June 2021.}}\\
{\bfseries \textmd{@ 76 }} &
{\bfseries \url{https://www.laverdad.es/culturas/incapaz-mirarles-movil-20190322004714-ntvo.html}\textmd{, as accessed in June 2021.}}\\
@ 77  &
\url{https://comics.jotdown.es/esenciales/buenas-noches-punpun/}, as accessed in June 2021.\\
@ 78  &
\url{https://www.goodreads.com/book/show/48949315-este-es-un-libro-cualquiera}, as accessed in June 2021.\\
{\bfseries \textmd{@ 79 }} &
{\bfseries \url{https://www.cronicanorte.es/rescatan-31-perros-de-raza-que-se-vendian-ilegalmente/4816}\textmd{, as accessed in December 2020.}}\\
{\bfseries \textmd{@ 80 }} &
{\bfseries \url{http://www.radiored.com.co/soldado-denuncia-a-esteban-santos-hijo-del-presidente-santos/}\textmd{, as accessed in March 2017.}}\\
{\bfseries \textmd{@ 81 }} &
{\bfseries \url{https://escritorasunidas.blogspot.com/2010/09/ingrid-betancourt-no-hay-silencio-que.html}\textmd{, as accessed in December 2020.}}\\
{\bfseries \textmd{@ 82 }} &
{\bfseries \url{http://ytodocomenzoasurgir.blogspot.com/2013/08/el-hombre-que-no-sentia-verguenza.html}\textmd{, as accessed in December 2020.}}\\
{\bfseries \textmd{@ 83 }} &
{\bfseries \url{https://www.elmundo.es/cronica/2001/CR291/CR291-06.html}\textmd{, as accessed in December 2020.}}\\
@ 84  &
\url{https://debates.coches.net/discussion/37556/articulo-aprende-a-ahorrar-gasolina/p2}, as accessed in June 2021.\\
{\bfseries \textmd{@ 85 }} &
{\bfseries \url{https://segundacita.blogspot.com/2012/06/violencia-y-otras-cuestiones.html}\textmd{, as accessed in December 2020.}}\\
{\bfseries \textmd{@ 86 }} &
{\bfseries \url{https://www.lavozdeasturias.es/noticia/viral/2020/02/08/carla-vigo-homenajea-madre-erika-ortiz-aniversario-fallecimiento/00031581164068183156338.htm}\textmd{, as accessed in June 2021.}}\\
{\bfseries \textmd{@ 87 }} &
{\bfseries \url{http://elcomercio.pe/tvmas/farandula/patricio-parodi-anuncio-medidas-legales-contra-antonio-pavon-noticia-1784294}\textmd{, as accessed in March 2017.}}\\
{\bfseries \textmd{@ 88 }} &
{\bfseries \url{http://infinitosdepapel.blogspot.com/2012/07/te-voy-echar-de-menos.html}\textmd{, as accessed in June 2021.}}\\
{\bfseries \textmd{@ 89 }} &
{\bfseries \url{https://www.mundodeportivo.com/20150512/2076215618/beckenbauer-sin-messi-el-barcelona-seria-un-equipo-cualquiera.html}\textmd{, as accessed in December 2020.}}\\
{\bfseries \textmd{@ 90 }} &
{\bfseries \url{http://blogs.elcomercio.pe/busconovio/2009/12/tenia-novio-ya-no.html}\textmd{, as accessed in March 2017.}}\\
{\bfseries \textmd{@ 100 }} &
{\bfseries \url{http://www.letras.mysite.com/vallejos4.htm}\textmd{, as accessed in December 2020.}}\\
{\bfseries \textmd{@ 101 }} &
{\bfseries \url{https://www.slideshare.net/yamipaar/cuerpos-geometricos-11539005}\textmd{, as accessed in June 2021.}}\\
{\bfseries \textmd{@ 102 }} &
{\bfseries \url{https://www.wattpad.com/761059418-%E2%9C%A7%C3%A1ngel-%E2%86%A6dillom%E2%9C%A7-%2715-te-quiero}\textmd{, as accessed in Februar 2021.}}\\
{\bfseries \textmd{@ 103 }} &
{\bfseries \url{https://docplayer.es/87325278-Experiencias-de-des-encuentro-de-las-mujeres-victimas-de-violencia-de-genero-con-el-sistema-penal-fatima-perez-jimenez-ainoa-santander-romera.html}\textmd{, as accessed in June 2021.}}\\
{\bfseries \textmd{@ 104 }} &
{\bfseries \url{https://winenews.it/it/4-scodelle-di-porcellana-bianca-di-lorenzo-franchini_320880/}\textmd{, as accessed in June 2021.}}\\
{\bfseries \textmd{@ 105 }} &
{\bfseries \url{https://autoblog.com.ar/2013/01/18/critica-toyota-86-gt/comment-page-27/}\textmd{, as accessed in June 2021.}}\\
@ 106  &
\url{http://www.espaciocris.com/2009/10/avance-del-capitulo-116.html}, as accessed in November 2020.\\
@ 107  &
\url{https://relatodelpresente.com.ar/el-desafio-de-la-blancura/}, as accessed in November 2020.\\
@ 108  &
\url{http://analizandoelimperialismojuntos.blogspot.com/2016/04/blog-post.html}, as accessed in November 2020.\\
@ 109  &
\url{http://focoeconomico.org/2012/04/01/que-sabemos-sobre-la-emision-y-la-inflacion/}, as accessed in November 2020.\\
@ 110  &
\url{http://www.mundotkm.com/hot-news-96313-alarma-en-telefe-por-el-bajo-rating-de-aliados}, as accessed in March 2017.\\
@ 111  &
\url{http://mexicomemata.blogspot.com/2011/03/todo-lo-que-pensas-se-desvanece-en-el.html}, as accessed in November 2020.\\
@ 112  &
\url{https://aquelarreforos.com.ar/index.php?topic=2180.30}, as accessed in December 2020.\\
@ 113  &
\url{https://www.taringa.net/+humor/jaja-son-re-cualquiera_1cpgt7}, as accessed in June 2021.\\
@ 114  &
\url{http://abelfer.wordpress.com/2012/12/26/que-hacer-con-la-inflacion/}, as accessed in November 2020.\\
@ 115  &
\url{http://www.unidiversidad.com.ar/hormigas-negras-el-arte-hecho-compromiso-popular}, as accessed in November 2020.\\
@ 116  &
\url{https://twitter.com/flortundis/status/1030498827294441473}, as accessed in June 2021.\\
@ 117  &
\url{https://www.fabio.com.ar/4633}, as accessed in November 2020.\\
@ 118  &
\url{https://asitalmundobotija.wordpress.com/2009/01/19/entre-la-moda-y-el-verano/}, as accessed in November 2020.\\
@ 119  &
\url{http://frasesrockeras.blogspot.com/2010/07/hay-que-persistir-e-hinchar-las-pelotas.html}, as accessed in November 2020.\\
@ 120  &
\url{https://paju07.th-cam.com/}, as accessed in November 2020.\\
@ 121  &
\url{http://www.leyendo-vuelo.com/2016/10/12-noches-de-halloween-10-pretty-little.html}, as accessed in November 2020.\\
@ 122  &
\url{http://www.friki.net/informes/56117-pablo-lezcano-hara-un-cover-de-clapton.html}, as accessed in November 2020.\\
@ 123  &
\url{https://elblogdeluisdavid.wordpress.com/2016/05/01/los-limites-del-yo/}, as accessed in November 2020.\\
@ 124  &
\url{https://autoblog.com.ar/2012/05/13/llega-autoblog-radio/}, as accessed in November 2020.\\
@ 125  &
\url{https://forum.wordreference.com/threads/hacen-cualquiera.1497124/}, as accessed in November 2020.\\
@ 126  &
http://blogsdelagente.com/soy-gordita-y-que/, as accessed on March 2017.\\
@ 127  &
https://blogs-fcpolit.unr.edu.ar/redaccion1-margarit/actos-de-habla-en-escenas-de-la-vida-cotidiana/, as accessed in November 2020.\\
@ 128  &
http://www.elortiba.org/old/dichabon.html, as accessed in November 2020.\\
@ 129  &
http://www.eduardocastillopaez.com.ar/2009/05/los-sonidos-de-una-mediocre-campana.html, as accessed in November 2020.\\
@ 130  &
http://untref.edu.ar/diccionario/notas-detalles.php?nota=8, as accessed in November 2020.\\
@ 131  &
https://lamajadescalza.wordpress.com/literatura/, as accessed in November 2020.\\
@ 132  &
\url{https://marcandoelpolo.com/como-cuando-y-donde-encontrar-trabajo-en-nueva-zelanda/}, as accessed in November 2020.\\
@ 133  &
\url{https://www.noticiasdealava.eus/opinion/mesa-de-redaccion/2016/09/02/jole-michele/336274.html}, as accessed in November 2020.\\
@ 134  &
\url{http://www.elortiba.org/old/dichabon.html}, as accessed in November 2020.\\
@ 135  &
\url{https://forum.wordreference.com/threads/hacen-cualquiera.1497124/}, as accessed in November 2020.\\
@ 136  &
\url{https://twitter.com/search?q=%22dijiste%20cualquiera%22&src=typed_query}, as accessed in November 2020.\\
@ 137  &
\url{https://www.jeuxvideo.com/forums/1-7683-12063418-2-0-1-0-le-plus-rapide-pour-les-h-space.htm}, as accessed in May 2021.\\
@ 138  &
\url{https://twitter.com/Zawa75/status/1215536068658966528}, as accessed in May 2021.\\
@ 139  &
\url{https://www.babelio.com/livres/Stangerup-Lhomme-qui-voulait-etre-coupable/397057}, as accessed in May 2021.\\
@ 140  &
\url{http://blogs.lanacion.com.ar/vaso-medio-lleno/uncategorized/el-dinero-si-compra-la-felicidad/}\textstyleInternetlink{, as accessed in March 2017.}\\
@ 141  &
\url{http://blogs.lanacion.com.ar/boquitas-pintadas/discriminacion-y-homofobia/ahora-las-mujeres-lesbianas-tambien-tendran-su-dia/}, as accessed in November 2020.\\
@ 142  &
\url{http://www.sergiobergman.com/el-subte-h-suma-de-a-tres/}, as accessed in November 2020.\\
@ 143  &
\url{http://blogs.lanacion.com.ar/cine/porque-si/no-les-veo-futuro/}, as accessed in November 2020.\\
@ 144  &
\url{https://www.facebook.com/CeresuCosplay/posts/550023348477014}, as accessed in November 2020.\\
@ 145  &
\url{https://twitter.com/EloisaAranda_/status/879385115260063746}, as accessed in November 2020.\\
@ 146  &
\url{https://twitter.com/Puchilla/status/591043944370298880}, as accessed in November 2020.\\
@ 147  &
\url{http://foro.gustfront.com.ar/viewtopic.php?t=1148&start=180}, as accessed in November 2020.\\
@ 148  &
\url{https://www.filmaffinity.com/co/user/rating/583761/539230.html}, as accessed in November 2020.\\
@ 149  &
\url{https://elnosoyloquedeberia.wordpress.com/2013/08/16/que-cosa-fuera-la-massa-sin-cantera/}, as accessed in November 2020.\\
@ 150  &
\url{http://maldealzheimer.wordpress.com/2005/03/13/un-intento-de-explicacion/}, as accessed in November 2020.\\
@ 151 &
\url{https://es-es.facebook.com/pg/Miley-Cyrus-Para-Vos-Una-Chica-Cualquiera-Para-Mi-Un-Modelo-A-Seguir-261016447350740/photos/?tab=album&album_id=261016797350705}, as accessed in June 2021.\\
@ 152 &
\url{https://twitter.com/clarincom/status/759113822305157120}, as accessed in June 2021.\\
\end{supertabular}
\end{flushleft}
\subsubsection[7.2.2 Primary sources of CdE\_G/H, CORDE and CORDIAM]{7.2.2 Primary sources of CdE\_G/H, CORDE and CORDIAM}
\begin{flushleft}
\tablefirsthead{}
\tablehead{}
\tabletail{}
\tablelasttail{}
\begin{supertabular}{m{1.7955599in}m{4.42126in}}
Albateni  &
Alfonso X, \textit{Cánones de Albateni}, in \textit{Electronic Texts and Concordances of the Madison Corpus of Early Spanish Manuscripts and Printings},\textit{ }prepared by John O’Neill, Madison / New York: Hispanic Seminary of Medieval Studies, 1999. CD-ROM (=París Arsenal 8322).\\
Animalias  &
Moamyn,\textit{ Libro de las animalias que cazan}, in \textit{Electronic Texts and Concordances of the Madison Corpus of Early Spanish Manuscripts and Printings},\textit{ }prepared by John O’Neill, Madison/New York: Hispanic Seminary of Medieval Studies, 1999. CD-ROM.\textit{ }(=Madrid Nacional Res. 270). \\
Astronomía  &
Alfonso X, \textit{Libros del saber de astronomía}, in \textit{Electronic Texts and Concordances of the Madison Corpus of Early Spanish Manuscripts and Printings},\textit{ }prepared by John O’Neill, Madison/New York: Hispanic Seminary of Medieval Studies, 1999. CD-ROM (=Madrid Biblioteca Universitaria Complutense 156)\textit{.}\\
Beneficencia  &
Concepción Arenal,\textit{ La beneficencia, la filantropía y la caridad}. Alicante: Biblioteca Virtual Miguel de Cervantes, 1999. \url{http://www.cervantesvirtual.com/nd/ark:/59851/bmcmk682}. \\
Cárcel\textit{ } &
Diego de San Pedro.\textit{ Cárcel de Amor}, in \textit{Electronic Texts and Concordances of the Madison Corpus of Early Spanish Manuscripts and Printings},\textit{ }prepared by John O’Neill, Madison/New York: Hispanic Seminary of Medieval Studies, 1999. CD-ROM (=Sevilla Cuatro compañeros alemanes 1492-03-03)\textit{.}\\
Claros varones\textit{ } &
Fernando del Pulgar, \textit{Claros varones de Castilla},\textit{ }in \textit{Electronic Texts and Concordances of the Madison Corpus of Early Spanish Manuscripts and Printings},\textit{ }prepared by John O’Neill, Madison/New York: Hispanic Seminary of Medieval Studies, 1999. CD-ROM. (=Sevilla Polono 1500-04-24)\textit{.}\\
Comentarios\textit{ } &
Bacallar y Sanna, Vicente, Marqués de San Felipe,\textit{ Comentarios de la guerra de España e historia de su rey Felipe V, El Animoso}, Alicante: Biblioteca Virtual Miguel de Cervantes, 1999. \href{http://www.cervantesvirtual.com/obra/comentarios-de-la-guerra-de-espana-e-historia-de-su-rey-felipe-v-el-animoso--0/}{\textstyleInternetlink{http://www.cervantesvirtual.com/obra/comentarios-de-la-guerra-de-espana-e-historia-de-su-rey-felipe-v-el-animoso-{}-0}}/.\\
Comentarios reales  &
Inca Garcilaso De la Vega, \emph{\textmd{\textit{Comentarios reales,}}\textmd{ }\textmd{e}}studio preliminar y notas de José Durand, Tomos I, II y III. Lima: Universidad Nacional Mayor de San Marcos, 1959.\\
Correspondencia  &
Gaspar Melchor de Jovellanos, \textit{Obras completas. Correspondencia}, tomo II, III y IV, edición crítica, introducción y notas de José Caso González, Oviedo: Universidad de Oviedo, Centro de Estudios del siglo XVIII, 1986.\\
Crónica de Aragon  &
Gauberto Fabricio de Vagad,\textit{ Crónica de Aragon}, in Francisco Marcos-Marín, Angel Gómez Moreno, Charles B. Faulhaber and John J. Nitti (eds.), \textit{ADMYTE: Archivo Digital de Manuscritos y Textos Españoles}. 1999. Madrid: Micronet. \\
Crónica de España  &
Diego de Valera, \textit{Crónica de España}, in Francisco Marcos-Marín, Angel Gómez Moreno, Charles B. Faulhaber and John J. Nitti (eds.), \textit{ADMYTE: Archivo Digital de Manuscritos y Textos Españoles},\textit{ }Madrid: Micronet. 1999 (=Sevilla Alfonso del Puerto 1482).\\
CR:PrLibre:98May4.  &
\url{http://www.prensalibre.co.cr} as accessed via \url{https://www.corpusdelespanol.org/hist-gen/} in May 2021.\\
Del agua &
Victor Pisabarro,\textit{ Del agua nacieron los sedientos}, Badosa.com, 1993. https://badosa.com/n025.\\
Doc. Cast. - Andalucía  &
Alfonso X, \textit{Documentos castellanos de Alfonso X – Andalucía}, in \textit{Textos y Concordancias Electrónicos de Documentos }\textit{Castellanos de Alfonso}, prepared by María Teresa Herrera, María Nieves Sánchez, María Estela González de Fauve and María Purificación Zabía. Madison: Hispanic Society of America. 1999. CD-ROM. \\
Doc. Cast. – 

Castilla la Nueva  &
Alfonso X, \textit{Documentos castellanos de Alfonso X - Castilla la Nueva}, in \textit{Textos y Concordancias Electrónicos de Documentos Castellanos de Alfonso},\textit{ }prepared by María Teresa Herrera, María Nieves Sánchez, María Estela González de Fauve and María Purificación Zabía, Madison: Hispanic Society of America, 1999. CD-ROM.\\
Doc. Cast. – León  &
Alfonso X, \textit{Documentos castellanos de Alfonso X – León}, in \textit{Textos y Concordancias Electrónicos de Documentos Castellanos de Alfonso}, prepared by María Teresa Herrera, María Nieves Sánchez, María Estela González de Fauve and María Purificación Zabía, Madison: Hispanic Society of America, 1999. CD-ROM.\\
Doc. Cast. – Murcia  &
Alfonso X, \textit{Documentos castellanos de Alfonso X – Murcia}, in \textit{Textos y Concordancias Electrónicos de Documentos Castellanos de Alfonso},\textit{ }prepared by María Teresa Herrera, María Nieves Sánchez, María Estela González de Fauve and María Purificación Zabía, Madison: Hispanic Society of America, 1999. CD-ROM.\\
Don Segundo Sombra  &
Ricardo Güiraldes, \textit{Don Segundo Sombra}, Buenos Aires: Editorial Proa, 1926.\\
El Crotalón  &
Cristóbal de Villalón. \textit{El Crótalon de Cristophoro Gnophoso}, Alicante[202F?]: Biblioteca Virtual Miguel de Cervantes, 1999. \url{http://www.cervantesvirtual.com/nd/ark:/59851/bmctx391}.\\
El día de fiesta  &
Juan de Zabaleta,\textit{ El día de fiesta en Madrid y sucesos que en él pasan}, Enrique Suárez Figaredo (ed.) Lemir 20 (2016) - Textos: 145-344.

\url{https://parnaseo.uv.es/lemir/Revista/Revista20/textos/02_Zabaleta_Dia_de_Fiesta.pdf}\textit{.}\\
El rey don Alfonso  &
Anonymous, \textit{El rey don Alfonso, el de la mano horadada}, in Carlos Mata Induráin (ed.) \textit{Comedias burlescas del Siglo de Oro}, tomo I. Madrid/Frankfurt: Iberoamericana/Vervuert. 1998.\\
En la carrera  &
Felipe Trigo,\textit{ En la carrera}. Alicante[202F?]: Biblioteca Virtual Miguel de Cervantes, 1999.

http://www.cervantesvirtual.com/nd/ark:/59851/bmcnk391.\\
Enc: Ave\textit{ } &
\ \url{http://es.encarta.msn.com/artcenter_/browse.html} as accessed via \url{https://www.corpusdelespanol.org/hist-gen/} in May 2021.\\
Entre clérigos y diablos  &
José Zorrilla. \textit{Entre clérigos y diablos ó El encapuchado}, Alicante[202F?]: Biblioteca Virtual Miguel de Cervantes, 2000. \url{http://www.cervantesvirtual.com/nd/ark:/59851/bmcsn043}.\\
Entrevista (ABC) &
Mortimer,\textit{ Entrevista (ABC)}, \url{http://www.abc.es}, as accessed via \url{https://www.corpusdelespanol.org/hist-gen/} in May 2021.\\
Epistolario  &
Juan Ginés de Sepúlveda, \textit{Epistolario de Juan Ginés de Sepúlveda: Selección}, ed. Angel Losada, Madrid: Ediciones Cultura Hispánica del Centro Iberoamericano de Cooperación, 1979.\\
España:ABC  &
Boo Juan Vicente,\textit{ España:ABC: CONTEXT +}, \url{http://www.abc.es}, as accessed via \url{https://www.corpusdelespanol.org/hist-gen/} in May 2021.\\
Espéculo  &
Alfonso X, \textit{Espéculo}, in \textit{Electronic Texts and Concordances of the Madison Corpus of Early Spanish Manuscripts and Printings},\textit{ }prepared by John O’Neill, Madison/New York: Hispanic Seminary of Medieval Studies, 1999. CD-ROM (=Madrid Nacional ms. 10123).\\
Examinarse de rey  &
Antonio Mira de Amescua, \textit{Examinarse de rey}, ed. María Celeste Martínez Calvo, Alicante: Biblioteca Virtual Miguel de Cervantes, 2014. 

http://www.cervantesvirtual.com/obra/examinarse-de-rey/.\\
Exemplario por ABC\textit{ \ } &
\textit{Exemplario por ABC}, in \textit{Electronic Texts and Concordances of the Madison Corpus of Early Spanish Manuscripts and Printings},\textit{ }prepared by John O’Neill, Madison/New York: Hispanic Seminary of Medieval Studies, 1999. CD-ROM. (New York: Hispanic Society. Transcribed by Anthony Cárdenas and John O'Neill).\\
Flujo del espíritu humano  &
Eduard Genís Sol, \textit{El flujo del espíritu humano. La teoría del yin y del yang y la de las cinco fases}. Medicina China Hoy, 2014.

\\
Libro de los exemplos &
Gayangos, Pascual de (ed.), \textit{Sánchez de Vercial, Clemente, Libro de los exemplos por a.b.c.} in Escritores en prosa anteriores al siglo xv, recogidos é ilustrados por Pascual de Gayangos, Madrid: M. Rivadeneyra, 1860, pp. 443-542 (= Biblioteca de Autores Españoles, 51).\\
General Estoria IV  &
Alfonso X, \textit{General Estoria IV}, in \textit{Electronic Texts and Concordances of the Madison Corpus of Early Spanish Manuscripts and Printings},\textit{ }prepared by John O’Neill, Madison/New York: Hispanic Seminary of Medieval Studies, 1999. CD-ROM (=Roma Vaticana Urb lat 539).\\
Gramáticos  &
Juan Pablo Forner, \textit{Los gramáticos: historia chinesca}, Alicante: Biblioteca Virtual Miguel de Cervantes, 2000. \url{http://www.cervantesvirtual.com/nd/ark:/59851/bmck9354}.\\
Gran Conquista  &
Anonymous,\textit{ Gran Conquista de Ultramar}, in \textit{Electronic Texts and Concordances of the Madison Corpus of Early Spanish Manuscripts and Printings}, prepared by John O’Neill, Madison / New York: Hispanic Seminary of Medieval Studies. 1999. CD-ROM (=Salamanca Giesser 1503-06-21).\\
Guía de pecadores  &
Fray Luís de Granada, \textit{Guía de pecadores en la qual se contiene una larga y copiosa exhortación a la virtud, y guarda de los mandamientos divinos}, Madrid: Antonio de Sancha. 1781. \url{http://cdigital.dgb.uanl.mx/la/1080046599_C/1080046599_T1/1080046599_MA.PDF}\\
Héroes y tumbas  &
Ernesto Sábato (1961), \textit{Sobre héroes y tumbas}. Caracas: Ayacucho, 1986.\\
Hist. Ecl. Ind.  &
Fray Gerónimo de Mendieta, \textit{Historia eclesiástica indiana}, editada por Joaquín García Icazbalceta, Alicante[202F?]: Biblioteca Virtual Miguel de Cervantes, 1999. \url{http://www.cervantesvirtual.com/nd/ark:/59851/bmczs2p6}.\\
Inundación Castálida  &
Sor Juana Inés de la Cruz, \textit{Inundación castálida},\textit{ }Alicante[202F?]: Biblioteca Virtual Miguel de Cervantes, 2000. \url{http://www.cervantesvirtual.com/nd/ark:/59851/bmchm544}.\\
Jardín  &
\ Antonio de Torquemada, \textit{Jardín de flores curiosas}, ed. Enrique Suárez Figaredo. Lemir 16 (2012), Textos: 605-834. https://parnaseo.uv.es/Lemir/Revista/Revista16/Textos/07\_Jardin\_Flores\_Torquemada.pdf.\\
La invención de Morel  &
Adolfo Bioy Casares, \textit{La invención de Morel} (1940), ed. Daniel Martino, \textit{La invención de Morel. Plan de evasión. La trama celeste}, Caracas: Ayacucho, 2002.\\
La poética &
Ignacio de Luzán, \textit{La poética o reglas de la poesía en general y de sus principales especies}, Alicante: Biblioteca Virtual Miguel de Cervantes, 1999. http://www.cervantesvirtual.com/obra/la-poetica-o-reglas-de-la-poesia-en-general-y-de-sus-principales-especies-{}-0/.\\
Leyes del estilo  &
Anonymous, \textit{Leyes del estilo}, in \textit{Electronic Texts and Concordances of the Madison Corpus of Early Spanish Manuscripts and Printings},\textit{ }prepared by John O’Neill, Madison / New York: Hispanic Seminary of Medieval Studies. 1999. CD-ROM (=Salamanca Impresor de la Gramática de Nebrija 1500-04-12).\\
Libro de las cruces  &
Alfonso X, \textit{Libro de las cruces}, in \textit{Electronic Texts and Concordances of the Madison Corpus of Early Spanish Manuscripts and Printings},\textit{ }prepared by John O’Neill, Madison/New York: Hispanic Seminary of Medieval Studies, 1999. CD-ROM. (=Madrid Nacional ms. 9294).\\
Libro rimado  &
Pedro López de Ayala,\textit{ Libro rimado de palacio}, Alicante: Biblioteca Virtual Miguel de Cervantes,\textit{ }2004.\textit{ }\url{http://www.cervantesvirtual.com/nd/ark:/59851/bmc0z727}.\\
Malcasados  &
Guillén de Castro, \textit{Los malcasados de Valencia}, critical edition by Luciano García Lorenzo, Madrid: Castalia. 1976.\\
Memento mori  &
César Pérez Gellida, \textit{Memento mori}. In \textit{Trilogía «Versos, canciones y trocitos de carne»}. Penguin Random House Grupo Editorial España, 2015.\\
No hay mal &
Manuel Tamayo y Baus, \textit{No hay mal que por bien no venga}, Alicante[202F?]: Biblioteca Virtual Miguel de Cervantes, 1999. http://www.cervantesvirtual.com/nd/ark:/59851/bmcgx477.\\
Noli me tangere  &
Rizal, José. \textit{Noli me tángere: (El país de los frailes): novela tagala}. Alicante[202F?]: Biblioteca Virtual Miguel de Cervantes. \url{http://www.cervantesvirtual.com/nd/ark:/59851/bmcfq9s5}. 2004.\\
Poesía  &
Juan de Tassis, Conde de Villamediana\textit{, Poesía}, ed. Maria Teresa Ruestes, Barcelona: Planeta [Clásicos Universales Planeta], 1991.\\
Rayuela  &
Julio Cortazár, \textit{Rayuela }(1936), ed. Julio Ortega and Saúl Yurkievich. Madrid: Archivos, 1991.\\
Secretos y cenizas  &
Mercedes Santos, \textit{Secretos y cenizas: Una historia de amor ante el peligro de una guerra que pudo cambiar el destino de América.} Cute Ediciones SRL, 2012.\\
Seis semanas\textit{ } &
Venice A.\textit{ }Fulton, \textit{Seis semanas para ser un pibón}, Barcelona: Libros Cúpula. 2013.\\
Siete partidas  &
Alfonso X, \textit{Siete partidas}, in \textit{Electronic Texts and Concordances of the Madison Corpus of Early Spanish Manuscripts and Printings},\textit{ }prepared by John O’Neill, Madison/New York: Hispanic Seminary of Medieval Studies. 1999. CD-ROM (=Sevilla Meinardo Ungut y Estanislao Polono 10-25-1491)\textit{.}\\
Siete sabios &
Anonymous. \textit{Siete sabios de Roma}, ed. Ventura de la Torre, Madrid: Miraguano, 1993.\\
Teatro crítico\textit{ } &
Benito Jerónimo Feijoo, \textit{Teatro crítico universal}, vol. 5, Alicante: Biblioteca Virtual Miguel de Cervantes, 1999. \url{http://www.cervantesvirtual.com/nd/ark:/59851/bmcmp4z4}.\\
Torre del aire  &
Gonzalo Torrente Ballester, \textit{Torre del aire}, ed. César Antonio Molina. La Coruña: Excma. Diputación Provincial de La Coruña, 1992.\\
Triunfo de las donas  &
Juan Rodríguez del Padrón, \textit{Triunfo de las donas y cadira de onor}, Alicante: Biblioteca Virtual Miguel de Cervantes. 1999. \url{http://www.cervantesvirtual.com/nd/ark:/59851/bmcvd6s7}.\\
Veinte Reyes  &
Anonymous, \textit{Crónica de Veinte Reyes}, in \textit{Electronic Texts and Concordances of the Madison Corpus of Early Spanish Manuscripts and Printings}, prepared by John O’Neill. Madison / New York: Hispanic Seminary of Medieval Studies. 1999. CD-ROM (=Escorial Monasterio Y-I-12).\\
\end{supertabular}
\end{flushleft}
\subsection[7.3 Secondary sources]{7.3 Secondary sources}
\begin{styleStandard}
Ahlberg, Axel. 1919.(ed.) \textit{C. Sallusti Crispi. Catilina; Iugurtha; Orationes et epistulae; Excerptae de historiis}. Editio Maior. Lipsia: Teubner. \url{https://archive.org/details/catilinaiugurtha00salluoft/page/n3/mode/2up}.
\end{styleStandard}

\begin{styleStandard}
Aloni, Maria. 2007. Free choice and exhaustification: an account of subtrigging effects. In Luise McNally \& Estela Puig-Waldmueller (eds.), \textit{Proceedings of Sinn und Bedeutung 11}, 16-30. Barcelona: Universitat Pompeu Fabra.
\end{styleStandard}

\begin{styleDLiteratur}
Aloni, Maria. 2021. Indefinites as fossils: the case of wh-based free choice. In Chiara Gianollo, Klaus von Heusinger \& Maria Napoli (eds.), \textit{Determiners and quantifiers: functions, variation, and change}. Leiden: Brill.
\end{styleDLiteratur}

\begin{styleStandard}
Aloni, Maria, Ana Aguilar Guevara, Angelika Port, Katrin Schulz \& Radek Šimík. 2010. Free choice items as fossils. Paper presented at the Workshop on Indefiniteness Crosslinguistically, Berlin, 25-26 Februar. 
\end{styleStandard}

\begin{styleStandard}
\ \ \textstyleInternetlink{https://www.researchgate.net/publication/48325717\_Free\_choice\_items\_as\_fossils}
\end{styleStandard}

\begin{styleStandard}
Aloni, Maria, Ana Aguilar-Guevara, Machteld de Vos, Angelika Port, Radek Šimík \& Hedde Zeijlstra. 2011. Semantics and pragmatics of indefinites: methodology for a synchronic and diachronic corpus study. In Stefanie Dipper \& Heike Zinsmeister (eds.), \textit{Beyond semantics: corpus-based investigations of pragmatic and discourse phenomena} (Bochumer Linguistische Arbeitsberichte), vol. 3, 1–16. Ruhr Universität Bochum: BLD.
\end{styleStandard}

\begin{styleStandard}
Aloni, Maria \& Angelika Port. 2013. Epistemic indefinites crosslinguistically. In Yelena Fainleib, Nicolas LaCara \& Yangsook Park (eds.), \textit{Proceedings of the 41st Annual Meeting of the North East Linguistic Society}, vol. 1, 29–42. Amherst, MA: GLSA.
\end{styleStandard}

\begin{styleStandard}
Alonso-Ovalle, Luis \& Royer 2021. Random Choice from Likelihood: The Case of Chuj (Mayan), Journal of Semantics, Volume 38, Issue 4, November 2021, Pages 483–529, \url{https://doi.org/10.1093/jos/ffab009}
\end{styleStandard}

\begin{styleStandard}
Alonso-Ovalle, Luis \& Paula Menéndez-Benito. 2010. Modal Indefinites, \emph{\textmd{\textit{Natural Language Semantics}}}, 18(1). 1–31.
\end{styleStandard}

\begin{styleStandard}
Alonso-Ovalle, Luis \& Paula Menendez-Benito. 2011. Expressing Indifference: Spanish Un NP Cualquiera. In Neil Ashton, Anca Chereches \& David Lutz (eds.), \textit{Proceedings of the 21st Semantics and Linguistic Theory Conference}. 333-52. New Brunswick, NJ: Rutgers University Press.
\end{styleStandard}

\begin{styleStandard}
Alonso-Ovalle, Luis \& Paula Menéndez-Benito. 2013. A Note on the Derivation of the Epistemic Effect of Spanish Algún as an Implicature. In Anamaria F\textcyrillic{{\U\cyra}l{\U\cyra}uş (ed.), }\textit{Alternatives in Semantics: Palgrave Studies in Pragmatics, Language and Cognition}, 36–49. London: Palgrave Macmillan UK.
\end{styleStandard}

\begin{styleStandard}
Alonso-Ovalle, Luis \& Paula Menéndez-Benito. 2015. \textit{Epistemic Indefinites: Exploring Modality Beyond the Verbal Domain}. Oxford: Oxford University Press.
\end{styleStandard}

\begin{styleStandard}
Alonso-Ovalle, Luis \& Paula Menéndez-Benito. 2018. Projecting possibilities in the nominal domain: Spanish uno cualquiera. \textit{Journal of Semantics} 35(1). 1\nobreakdash-41.
\end{styleStandard}

\begin{styleStandard}
Ariza Viguera, Manuel. 1998a. \textit{El comentario filológico de textos}. Madrid: Arco/Libros.
\end{styleStandard}

\begin{styleStandard}
Ariza Viguera, Manuel. 1998b. \textit{Antología de la prosa medieval española}. Madrid: Biblioteca Nueva.
\end{styleStandard}

\begin{styleStandard}
Arregui, Ana. 2006. \textit{Cualquier}, exception phrases and negation. In Jenny Doetjes \& Paz González (eds.), \textit{Romance Languages and Linguistic Theory 2004: Selected Papers from }Selected papers from ‘Going Romance’, Leiden, 9–11 December 2004, 1-22. Amsterdam: John Benjamins.
\end{styleStandard}

\begin{styleStandard}
Baayen, R. Harald. 2008. \textit{Analyzing Linguistic Data. A Practical Introduction to Statistics Using R}. Cambridge: Cambridge University Press. 
\end{styleStandard}

\begin{styleStandard}
Bar-Asher Siegal, Elitzur A. 2020. A formal approach to reanalysis: The case of a negative counterfactual marker. \textit{Proceedings of the Linguistic Society of America} 5(2). 34–50. \url{https://doi.org/10.3765/plsa.v5i2.4792}.
\end{styleStandard}

\begin{styleStandard}
Becker, Martin. 2014. \textit{Welten in Sprache. Zur Entwicklung der Kategorie «Modus» in romanischen Sprachen}. Berlin: De Gruyter.
\end{styleStandard}

\begin{styleStandard}
Bello, Andrés. 1945 [1847]. \textit{Gramatica de la lengua castellana destinada al uso de los americanos}. Buenos Aires: Sopena.
\end{styleStandard}

\begin{styleStandard}
Benincà, Paola \& Guglielmo Cinque. 2010. La frase relativa. In Giampaolo Salvi \& Lorenzo Renzi (eds.), \textit{Grammatica dell’Italiano Antico}, 469–507. Bologna: il Mulino.
\end{styleStandard}

\begin{styleStandard}
Blank, Andreas. 1997. \textit{Prinzipien des lexikalischen Bedeutungswandels am Beispiel der romanischen Sprachen}. Tübingen: Niemeyer.
\end{styleStandard}

\begin{styleStandard}
Bloch, Oscar \& Walter von Wartburg. 1975. \textit{Dictionnaire étymologique de la langue française}. Paris: P.U.F.
\end{styleStandard}

\begin{styleStandard}
Borkowska, Paulina \& Grzegorz A. Kleparski. 2007. It befalls words to fall down: pejoration as a type of semantic change. \textit{Studia Anglica Resoviensia} 4. 33\nobreakdash-50.
\end{styleStandard}

\begin{styleStandard}
Bosque, Ignacio. 2001. Adjective position and the interpretation of indefinites. \textit{Current issues in Spanish syntax and semantics} 53. 17-38.
\end{styleStandard}

\begin{styleStandard}
Bresnan, John. 1973. Syntax of the Comparative Clause Construction in English, \textit{Linguistic Inquiry} 4. 275–343.
\end{styleStandard}

\begin{styleStandard}
Brucart, Joseph M. 1999. La estructura del sintagma nominal: las oraciones de relativo. In Ignacio Bosque \& Violeta Demonte (eds.), \textit{Gramática descriptiva de la lengua española}, 395-522. Madrid: Espasa Calpe.
\end{styleStandard}

\begin{styleStandard}
Bruyne, Jacques de. 2011. \textit{Spanische Grammatik}. Berlin, Boston: De Gruyter. https://doi.org/10.1515/9783110919189
\end{styleStandard}

\begin{styleStandard}
Buridant, Claude. 2000. \textit{Grammaire nouvelle de l'ancien français. }Paris: Sedes. 
\end{styleStandard}

\begin{styleStandard}
Bybee, Joan. 2006. From Usage to Grammar: The Mind’s Response to Repetition. \textit{Language}. 82(4). 711–733.
\end{styleStandard}

\begin{styleStandard}
Caponigro, Ivano. 2004. The semantic contribution of wh-words and type shifts: Evidence from free relatives crosslinguistically. In Robert B. Young (ed.), \textit{Proceedings from Semantics and Linguistic Theory (SALT)} \textit{XIV}, 38-55. Ithaca, NY: CLC Publications, Cornell University.
\end{styleStandard}

\begin{styleStandard}
Carlson, Gregory N. 1977. \textit{Reference to kinds in English}. Amherst: University of Massachusetts, Amherst dissertation.
\end{styleStandard}

\begin{styleStandard}
Charnavel, Isabelle. 2015. Same, different and other as comparative adjectives – A uniform analysis based on French. \textit{Lingua} 156. 129-174.
\end{styleStandard}

\begin{styleStandard}
Chierchia, Gennaro. 2006. Broaden your views: Implicatures of domain widening and the '{}'logicality'{}' of language. \textit{Linguistic Inquiry} 37. 535–590.
\end{styleStandard}

\begin{styleStandard}
Chierchia, Gennaro. 1998. \textit{Reference to Kinds across Language}. Natural Language Semantics 6 (4). 339–405. \url{https://doi.org/10.1023/A:1008324218506}.
\end{styleStandard}

\begin{styleStandard}
Choi, Jinyoung \& Maribel Romero. 2008. Rescuing existential free choice items in episodic sentences. In Oliver Bonami \& Patricia Cabredo-Hofherr (eds.), \textit{Empirical Issues in Syntax and }Semantics, vol. 7, 77–98.
\end{styleStandard}

\begin{styleStandard}
Cinque, Guglielmo. 2015. Three phenomena discriminating between “raising” and “matching” relative clauses.\textit{ Semantics-Syntax Interface} 2(1). 1-27.
\end{styleStandard}

\begin{styleStandard}
Company-Company, Concepción. 2003. Grammaticalización y cambio sintáctico en la historia del español. \textit{Medievalia} 34. 3-61.
\end{styleStandard}

\begin{styleStandard}
Company-Company, Concepción. 2009. Parámetros de gramaticalización en los indefinidos compuestos en el español. In Fernando Sánchez Miret (ed.), \textit{Romanística sin complejos. Homenaje a Carmen Pensado}, 71-104. Bern: Peter Lang.
\end{styleStandard}

\begin{styleStandard}
Company-Company, Concepción \& Julia Pozas-Loyo. 2009. Los indefinidos compuestos y los pronombres generic impersonales omne y uno. In Concepción Company-Company (ed.), \textit{Sintaxis historica de le lengua espanola (segunda parte: La frase nominal)}, 1073-1219. Mexico City: Fondo de Cultura Económica-Universidad Nacional Autónoma de Mexico.
\end{styleStandard}

\begin{styleStandard}
Conde, Oscar. 2011. \textit{Diccionario etimológico del lunfardo}. Kindle Edition. Buenos Aires: Taurus.\href{https://www.amazon.com/-/de/dp/B006GFT6DK/ref=sr_1_1?__mk_de_DE=%20ÅMÅŽÕÑ&dchild=1&keywords=conde+diccionario+etimologico&qid=1614335360&s=digital-text&sr=1-1}{\textstyleInternetlink{https://www.amazon.com/-/de/dp/B006GFT6DK/ref=sr\_1\_1?\_\_mk\_de\_DE= ÅMÅŽÕÑ\&dchild=1\&keywords=conde+diccionario+etimologico\&qid=1614335360\&s=digital-text\&sr=1-1}}.
\end{styleStandard}

\begin{styleStandard}
Condoravdi, Cleo. 2015. Ignorance, indifference, and individuation with wh-ever. In Luis Alonso-Ovalle \& Paula Menéndez-Benito (eds.), \textit{Epistemic Indefinites: Exploring Modality Beyond the Verbal Domain}, 213-243. Oxford: Oxford University Press.
\end{styleStandard}

\begin{styleStandard}
Cortelazzo, Manlio \& Paolo Zolli. 1979. \emph{\textmd{\textit{Dizionario etimologico}}}\textstylest{ \textit{della lingua italiana}. Bologna: Zanichelli.}
\end{styleStandard}

\begin{styleStandard}
Croft, William \& Alan Cruse. 2004. Cognitive linguistics. Cambridge: Cambridge University Press.10.1017/CBO9780511803864Search in Google Scholar
\end{styleStandard}

\begin{styleStandard}
Cuervo, Rufino José. 1953-1959. \textit{Diccionario de construcción y régimen de la lengua caste-llana, continuado y editado por el Instituto Caro y Cuervo}, 2nd edn. Santa Fe de Bogotá: Instituto Caro y Cuervo.
\end{styleStandard}

\begin{styleStandard}
Dayal, Veneeta. 1998. “Any” as Inherently Modal. \textit{Linguistics and Philosophy}. 21(5). 433–476.
\end{styleStandard}

\begin{styleStandard}
Dayal, Veneeta. 2004. The universal force of free choice any. In Pierre Pica, Johan Rooryck \& Jeroen van Craenenbroeck (eds.), \textit{The Linguistic Variation Yearbook 4}, 5-40. Amsterdam: John Benjamins.
\end{styleStandard}

\begin{styleStandard}
Dayal, Veneeta. 2013. A Viability Constraint on Alternatives for Free Choice. In Anamaria F\textcyrillic{{\U\cyra}l{\U\cyra}uş (ed.), }\textit{Alternatives in Semantics: Palgrave Studies in Pragmatics, Language and Cognition}, 88–122. London: Palgrave Macmillan. 
\end{styleStandard}

\begin{styleStandard}
\ \ \url{https://doi.org/10.1057/9781137317247_4}.
\end{styleStandard}

\begin{styleStandard}
Den Dikken, Marcel \& Anastasia Giannakidou. 2002. From hell to polarity. \textit{Linguistic Inquiry }33(1). 31–61.
\end{styleStandard}

\begin{styleStandard}
Detges, Ulrich et al. 2021, ‘Positioning reanalysis and reanalysis research’. Journal of Historical Linguistics 5: 1-49.
\end{styleStandard}

\begin{styleStandard}
De Smet, Hendrik \& Olga Fischer. 2017. The role of analogy in language change: Supporting constructions. In Marianne Hundt, Sandra Mollin \& Simone Pfenninger (eds.), The changing English language: Psycholinguistic perspectives, 240–268. Cambridge: Cambridge University Press.
\end{styleStandard}

\begin{styleStandard}
Di Tullio, Ángela. 2013. Aspectos morfosintácticos del español argentino resultantes del contacto con el italiano. In \textit{Perspectivas teóricas y experimentales sobre el español de la Argentina.-(Lingüística iberoamericana; v. 56)}. Madrid: Iberoamericana. 
\end{styleStandard}

\begin{styleStandard}
\ \ \url{http://digital.casalini.it/9783954871971}.
\end{styleStandard}

\begin{styleStandard}
Di Tullio, Ángela Lucía. 2015. Cualquiera. \textit{Diccionario Latinoamericano de la Lengua Española}. Buenos Aires: Universidad nacional de tres de febrero. 
\end{styleStandard}

\begin{styleStandard}
\ \ \url{http://untref.edu.ar/diccionario/notas-detalles.php?nota=8}.
\end{styleStandard}

\begin{styleStandard}
Eberenz, Rolf. 2009. La periodización de la historia morfosintáctica del español: propuestas y aportaciones recientes. \textit{Cahiers d'Études Hispaniques Médiévales} 32. 181-201. 
\end{styleStandard}

\begin{styleStandard}
Eberenz, Rolf. 2000. Los regimientos de peste a fin es de la Edad Media: configuración de un nuevo género textual. In Daniel Jacob \& Johannes Kabatek (eds.), \textit{Lengua medieval y tradiciones discursivas en la Península Ibérica: descripción gramatical: pragmática histórica: metodología}, 79-96. Madrid: Iberoamericana Vervuert. 
\end{styleStandard}

\begin{styleStandard}
\ \ \url{http://digital.casalini.it/9783865278432}.
\end{styleStandard}

\begin{styleStandard}
Eckardt, Regine. 2006. \textit{Meaning Change in Grammaticalization: An Enquiry into Semantic Reanalysis}. Oxford: Oxford University Press.
\end{styleStandard}

\begin{styleStandard}
Elvira, Javier. 2007. Algo más que palabras: Uso y significado en las locuciones del español. \textit{Verba Hispanica} XV(b). 109-125.
\end{styleStandard}

\begin{styleStandard}
Elvira, Javier. 2020. Lexicalización primero, gramaticalización después. Aproximación a la génesis de cualquiera. In José Martínez Alcalde, Juan Pedro Sánchez Méndez, Francisco Javier Satorre Grau, Amparo Ricós Vidal, Adela García Valle, Francisco Pedro Pla Colomer \& Santiago Vicente Llavata (eds.). \textit{El español y las lenguas peninsulares en su }\textit{diacronía: miradas sobre una historia compartida}, 153-168. Valencia: Tirant Humanidades. 
\end{styleStandard}

\begin{styleStandard}
Espinosa Elorza, Rosa María \& Carlos Sánchez Lancis. 2006. Cuantificadores indefinidos en la General Estoria (Tercera parte). \textit{Revista española de Lingüístca} 36(1). 127-156.
\end{styleStandard}

\begin{styleStandard}
Etxeberria, Urtzi \& Anastasia Giannakidou. 2014. Anti-specificity and the role of number: The case of Spanish ‘algún’/’algunos’. Presented at the 24th Colloquium on Generative Grammar, Universidad Autónoma de Madrid, Madrid.
\end{styleStandard}

\begin{styleStandard}
Fanselow, Gisbert. 1988. Aufspaltung von NP und das Problem der ‘freien’ Wortstellung.\textit{ Linguistiche Berichte} (114), 91-113.
\end{styleStandard}

\begin{styleStandard}
Fernández González, Etelvina. 1992. San Isidoro de León. \textit{Cuadernos de Arte Español} 53. Madrid: Historia 16.
\end{styleStandard}

\begin{styleStandard}
Fintel, Kai-Uwe von. 1994. Restrictions on quantifier domains. PhD diss., University of Massachusetts AAI9434544. \href{https://scholarworks.umass.edu/dissertations/%20AAI9434544}{\textstyleInternetlink{https://scholarworks.umass.edu/dissertations/ AAI9434544}}
\end{styleStandard}

\begin{styleStandard}
Fintel, Kai-Uwe von. 2000. Whatever. \textit{Proceedings of the 10th Semantics and Linguistic Theory Conference}, held June 2-4, 2000 at Cornell University, edited by Brendan Jackson and Tanya Matthews. Ithaca, NY: Cornell University, 27–40.
\end{styleStandard}

\begin{styleStandard}
Fontana, Josep. 1993. \textit{Phrase structure and the syntax of clitics in the history of Spanish}. Ph.D. dissertation, University of Pennsylvania.
\end{styleStandard}

\begin{styleStandard}
Frege, G. 1997. “On Sinn and Bedeutung (1892).” In M. Beaney (ed.), The Frege Reader: pp. 151-171.
\end{styleStandard}

\begin{styleStandard}
Fălăuş, Anamaria. 2013. Broaden your views, but try to stay focused: A missing piece in the polarity system. In Ivano Caponigro \& Carlo Ceccchetto (eds.),~\textit{From grammar to meaning: The spontaneous logicality of language},~81-107. Cambridge: Cambridge University Press. 
\end{styleStandard}

\begin{styleStandard}
Fălăuş, Anamaria. 2014. Pick some but not all alternatives!, In Marie-Hélène Côté \& Eric Mathieu (eds.), \textit{Variation within and across Romance Languages}, 155-172. Amsterdam: John Benjamins.
\end{styleStandard}

\begin{styleStandard}
Fălăuş, Anamaria. 2015. Romanian epistemic indefinites. In Luis Alonso-Ovalle \& Paula Menéndez-Benito 2015 (eds.), \textit{Epistemic Indefinites: Exploring Modality Beyond the Verbal Domain}, 60-81. Oxford: Oxford University Press.
\end{styleStandard}

\begin{styleStandard}
Gajewski, Jon. 2008. Licensing strong NPIs. In Lukasz Abramowicz, Stefanie Brody, Toni Cook, Ariel Diertani, Aviad Eilam, Keelan Evanini, Kyle Gorman, Laurel MacKenzie, and Joshua Tauberer (eds.) \textit{Pennsylvania Working Papers in Linguistics} 14(1). \url{https://repository.upenn.edu/pwpl/vol14/iss1/13}.
\end{styleStandard}

\begin{styleStandard}
Gajewski, Jon 2013. An analogy between a connected exceptive phrase and polarity items. In Regina Eckardt, Manfred Sailer, Eva Csipak, Mingya Liu. \textit{Beyond ‘Any‘ and ‘Ever’: New Explorations in Negative Polarity Sensitivity}. Berlin: De Gruyter Mouton. 183-213.
\end{styleStandard}

\begin{styleStandard}
Gamillscheg, Ernst. 1969. \textit{Etymologisches Wörterbuch der französischen Sprache}. Heidelberg: Winter.
\end{styleStandard}

\begin{styleStandard}
Geehoven, Veerle van \& Louise McNally. 2005. On the property analysis of opaque complements. \textit{Lingua} 115 (6). 885-914. 
\end{styleStandard}

\begin{styleStandard}
Gehrke, Berit \& Louise McNally. 2015. Distributional modification: The case of frequency adjectives. \textit{Language} 91(4). 837–870.
\end{styleStandard}

\begin{styleStandard}
Gessner, Emil. 1895. Das spanische indefinite Pronomen. \textit{Zeitschrift für romanische Philologie} 19(2). 153–169. \url{https://doi.org/10.1515/zrph.1895.19.2.153}.
\end{styleStandard}

\begin{styleStandard}
Giannakidou, Anastasia. 2001. The Meaning of Free Choice. \textit{Linguistics and Philosophy} 24(6). 659–735. \url{https://doi.org/10.1023/A:1012758115458}.
\end{styleStandard}

\begin{styleStandard}
Giannakidou, Anastasia \& Josep Quer. 2013. Exhaustive and non-exhaustive variation with anti-specific indefinites: free choice versus referential vagueness. \textit{Lingua} 126. 120\nobreakdash-149.
\end{styleStandard}

\begin{styleStandard}
Gianollo, Chiara. 2018. \textit{Indefinites between Latin and Romance}. Oxford: Oxford University Press.
\end{styleStandard}

\begin{styleStandard}
Girón Alconchel, José Luis. 2012. Gramaticalización como creación de lengua a partir del habla. Relativos e indefinidos compuestos en los \textit{Fueros de Aragón }y en el \textit{Fuero de Teruel}. \textit{Archivio de Filología Aragonesa} 68. 15-38.
\end{styleStandard}

\begin{styleStandard}
Giusti, Giuliana \& Rossella Iovino. 2011. Evidence for a Split DP in Latin\textit{. University of Venice Working Papers in Linguistics} 21. 111-129.
\end{styleStandard}

\begin{styleStandard}
Godefroy, Frédéric. 2006 [1880–1902]. \textit{Dictionnaire de l‘ancienne langue française et de tous ses dialectes du 9e au 15e siècle}, Paris, Editions Champion Electronique.
\end{styleStandard}

\begin{styleStandard}
Greimas, Algirdas Julien. 1998. \textit{Dictionnaire de l’ancien français: le Moyen Âge}. Paris: Larousse.
\end{styleStandard}

\begin{styleStandard}
Grice, H. Paul. 1975. Logic and Conversation. In Peter Cole \& Jerry L. Morgan (eds.), \textit{Syntax and Semantics Volume 3: Speech Acts}, 41-48. New York: Academic Press.
\end{styleStandard}

\begin{styleStandard}
Hacquard, Valentina. 2010. On the event relativity of modal auxiliaries. \textit{Natural Language Semantics} 18(1). 79–114.
\end{styleStandard}

\begin{styleStandard}
Hamblin, Charles. L. 1973. Questions in Montague English. \textit{Foundations of Language }10. 41-53.
\end{styleStandard}

\begin{styleLiteraturStandard}
Hanssen, Federico. 1910. \textit{Spanische Grammatik auf historischer Grundlage}. Halle: Max Niemeyer.
\end{styleLiteraturStandard}

\begin{styleLiteraturStandard}
Hanssen, Federico. 1913. \textit{Gramática histórica de la lengua castellana}. Halle: Max Niemeyer.
\end{styleLiteraturStandard}

\begin{styleStandard}
Haspelmath, Martin. 1997. \textit{Indefinite pronouns}. Oxford: Oxford University Press.
\end{styleStandard}

\begin{styleLiteraturStandard}
Heim, Irene. 1992. Presupposition projection and the semantics of attitude verbs. \textit{Journal of Semantics} 9. 183-221.
\end{styleLiteraturStandard}

\begin{styleStandard}
Heim, Irene \& Angelika Kratzer. 1998. \textstyleacopre{\textit{Semantics in Generative Grammar}. Oxford: Blackwell.}
\end{styleStandard}

\begin{styleStandard}
Heusinger, Klaus von. 2002. Specificity and definiteness in sentence and discourse structure. \textit{Journal of Semantics} 19 (3). 245–274. \url{https://doi.org/10.1093/jos/19.3.245}.
\end{styleStandard}

\begin{styleStandard}
Heusinger, Klaus von. 2007. Referentially Anchored Indefinites. In Ilena Comorovski \& Klaus von Heusinger (eds.), \textit{Existence: Semantics and Syntax} \textit{(Studies in Linguistics and Philosophy)}, 273–292. Dordrecht: Springer Netherlands. \url{https://doi.org/10.1007/978-1-4020-6197-4_10}. 
\end{styleStandard}

\begin{styleStandard}
Heusinger, Klaus von. 2011. Specificity. In Klaus von Heusinger, Claudia Maienborn \& Paul Portner (eds.), \textit{Semantics. An International Handbook of Natural Language Meaning},\textit{ }vol. 2, 1025-1058. Berlin: De Gruyter.
\end{styleStandard}

\begin{styleStandard}
Horn, Laurence R. 2000. Pick a theory (not just \textit{any} theory): Indiscriminatives and the free-choice indefinite. In Laurence R. Horn \& Yasuhiko Kato (eds.), \textit{Negation and Polarity: Syntactic and Semantic Perspectives}, 147–192. Oxford: Oxford University Press.
\end{styleStandard}

\begin{styleStandard}
Jacobson, Pauline. 1995. On the Quantificational Force of English Free Relatives. Emmon Bach, Eloise Jelinek, Angelika Kratzer \& Barbara H. Partee (eds.), \textit{Quantification in Natural Languages}, 451-486. Dordrecht: Springer.
\end{styleStandard}

\begin{styleStandard}
Jayez, Jacques \& Lucia M. Tovena. 2005. Free-Choiceness and Non-Individuation. \textit{Linguistics and Philosoph}y 28. 1-71. 
\end{styleStandard}

\begin{styleStandard}
Kanwit, Matthew \& Virginia, Terán. 2020. Proposing a tripartite intensifier system: re, muy, and bien in Buenos Aires and Tucumán, Argentina. Diego Pascual y Cabo \& Idoia Elola (eds.) \textit{Current Theoretical and Applied Perspectives on Hispanic and Lusophone Linguistics}. 254-272. John Benjamins Publishing Company, Amsterdam/Philadelphia.
\end{styleStandard}

\begin{styleStandard}
\emph{\textmd{Kayne}}\textstyleacopre{, Richard S. }\emph{\textmd{1994}}\textstyleacopre{. \textit{The }}\emph{\textmd{Antisymmetry}}\textstyleacopre{\textit{ of Syntax}. Cambridge, MA: MIT Press.}
\end{styleStandard}

\begin{styleStandard}
Kennedy, Christopher. 1999. \textit{Projecting the Adjective: The Syntax and Semantics of Gradability and Comparison.} New York: Routledge.
\end{styleStandard}

\begin{styleStandard}
Kellert, Olga. 2018. Questions with definite markers in (Old) Romance, with focus on Old Spanish. \textit{Isogloss} 4.1, 55–83. DOI: 10.5565/rev/isogloss.56
\end{styleStandard}

\begin{styleStandard}
Kellert, Olga. 2020. Semantic and syntactic change of the indefinite \textit{equis} in Mexican Spanish, in Remus Gergel \& Jonathan Watkins (Eds.) \textit{Quantification and scales in change}. Berlin: Language Science Press. 131–159
\end{styleStandard}

\begin{styleStandard}
Kellert, Olga. 2020. La diacronía del indefinido \textit{cualquiera}, in Santos Rovira \& José María (Eds.) \textit{Raíces y horizontes del Español}. Perspectivas dialectales, históricas y sociolingüísticas. Lugo: Axac, 57–67 
\end{styleStandard}

\begin{styleStandard}
Kellert, Olga. 2021. Indefinite qualunque from Old to Modern Italian. Journal of Historical Syntax 5.9 (Special Issue: Paola Crisma \& Giuseppe Longobardi (Eds.) Proceedings of the 20th Diachronic Generative Syntax (DiGS) Conference. (https://ojs.ub.uni-konstanz.de/hs/index.php/hs/article/view/82/39)
\end{styleStandard}

\begin{styleStandard}
Kellert, Olga. 2021. The diachronic development of the evaluative meaning of qualunque in Italian, in Chiara Gianollo et al. (Eds.) Determiners and quantifiers: functions, variation, and change.(= Syntax and Semantics). Brill.
\end{styleStandard}

\begin{styleStandard}
Kellert, Olga \& Andrés Enrique Arias (in press) The diachrony of Catalan qualsevol. In Olga Kellert, Malte Rosemeyer, Sebastian Lauschus (Eds.), Romance Indefinites: From Semantics to Pragmatics. Berlin: Language Science Press.
\end{styleStandard}

\begin{styleStandard}
Kellert, Olga \& \ Marika Francia (in press) Cualunque in Argentinian Spanish and qualunque in Italian. In Olga Kellert, Malte Rosemeyer \& Sebastian Lauschus (Eds.), Romance Indefinites: From Semantics to Pragmatics. Berlin: Language Science Press.
\end{styleStandard}

\begin{styleStandard}
Koch, Peter \& Wulf Oesterreicher. 1990. \textit{Gesprochene Sprache in der Romania: Französisch, Italienisch, Spanisch}. \textit{Gesprochene Sprache in der Romania}. T\"{u}bingen: Niemeyer.
\end{styleStandard}

\begin{styleStandard}
Kratzer, Angelika \& Shimoyama, Junko. 2002. Indeterminate Phrases: the View from Japanese. In Yukio Otsu (ed.), \textit{The Proceedings of the Third Tokyo Conference on Psycholinguistics}, 1-25. Tokyo: Hituzi Syobo.
\end{styleStandard}

\begin{styleStandard}
\ \ \url{http://semanticsarchive.net/Archive/WewNjc4Z/}\textstyleInternetlink{ (accessed 9 April 2021).}
\end{styleStandard}

\begin{styleStandard}
Krifka, Manfred, Francis Jeffry Pelletier, Gregory Carlson, Alice ter Meulen, Gennaro Chierchia \& Godehard Link. 1995. Genericity: An Introduction. In Gregory N. Carlson \& Francis Jeffrey Pelletier (eds.), \textit{The Generic Book}, 1-124. University of Chicago Press.
\end{styleStandard}

\begin{styleStandard}
Lapesa, Rafael. 1964. \textit{Historia de la lengua española}. Madrid: Gredos.
\end{styleStandard}

\begin{styleStandard}
Lapesa, Rafael. 1996 [1975]. El sustantivo sin actualizador en español: la ausencia del determinante en la lengua española. In Ignacio Bosque (ed.), \textit{Presencia y ausencia de determinante en la lengua española}, 121-140. Madrid: Visor.
\end{styleStandard}

\begin{styleStandard}
Lapesa, Rafael. 2000. \textit{Estudios de morfosintaxis histórica del español}. Madrid: Gredos.
\end{styleStandard}

\begin{styleStandard}
Larrivée, Pierre. 2007. Comment déterminer la valeur qualitative des indéfinis de sélection arbitraire?. \textit{Zeitschrift für französische Sprache und Literatur}. Franz Steiner Verlag 117(1). 1–13.
\end{styleStandard}

\begin{styleStandard}
Lasersohn, Peter. 2005. Context dependence, disagreement, and predicates of personal taste. \textit{Linguistics and Philosophy} 28. 643–686.
\end{styleStandard}

\begin{styleannotationtext}
\textstylehighlight{Lehmann}, Christian. 2002 [1982]. \textit{Thoughts on grammaticalization}, 2nd edn. Erfurt: Universität Erfurt [=\textit{Arbeitspapiere für Sprachwissenschaft der Universität Erfurt} 9].
\end{styleannotationtext}

\begin{styleStandard}
Lombard, Alf. 1938. Une classe spéciale de termes indéfinis dans les langues romanes. \textit{Studia Neophilologica} 11(1-2). 186-209.
\end{styleStandard}

\begin{styleStandard}
Lombard, Alf. 1947. A propos de \textit{quienquiera}. \textit{Studia Neophilologica} 20(1-2). \textstylest{21-36.}
\end{styleStandard}

\begin{styleStandard}
López, Luis. 2009. \textit{A Derivational Syntax for Information Structure} (Oxford Studies in Theoretical Linguistics). Oxford: Oxford University Press. 
\end{styleStandard}

\begin{styleStandard}
\ \ \url{https://doi.org/10.1093/acprof:oso/9780199557400.001.0001}. 
\end{styleStandard}

\begin{styleStandard}
Macías Villalobos, Cristóal. 1997. \textit{Estructuras y funciones del demostrativo en el español moderno}. Malaga: Universidad de Malaga.
\end{styleStandard}

\begin{styleStandard}
\textstylest{Mackenzie, Ian. 2019. \textit{Language Structure Variation and Change. The case of Old Spanish Syntax.} Palgrave: McMillan. }
\end{styleStandard}

\begin{styleStandard}
Martínez García, Hortensia. 2005. \textit{Construir bien en español: la corrección sintáctica}. Oviedo: Universidad de Oviedo.
\end{styleStandard}

\begin{styleStandard}
Meier, Harri. 1950. Indefinita vom Typus span. \textit{cualquiera}, it. \textit{qualsivoglia}. \textit{Romanische Forschungen}, 62(4). 385-401.
\end{styleStandard}

\begin{styleStandard}
Melis, Chantal, Marcela Flores \& Sergio Bogard. 2004. La historia del español: Propuesta de un tercer periodo evolutivo. \textit{Nueva revista de filología hispanica} 51. 1-56.
\end{styleStandard}

\begin{styleStandard}
Menéndez-Benito, Paula. 2010. On universal Free Choice items. \textit{Nat Lang Semantics} 18. 33-64.
\end{styleStandard}

\begin{styleStandard}
Menéndez Pidal, Ramón. 1928. \textit{Manual de gramática histórica española}. Madrid: Espasa-Calpe.
\end{styleStandard}

\begin{styleStandard}
Menéndez Pidal, Ramón. 1966. \textit{Manual de gramática histórica española}. Madrid: Espasa-Calpe.
\end{styleStandard}

\begin{styleStandard}
Meyer-Lübke, Wilhelm. 1899. \textit{Grammatik der romanischen Sprachen: vol. Romanische Syntax}. Leipzig: Reisland.
\end{styleStandard}

\begin{styleStandard}
Muller, Claude. 2007. Les indéfinis free choice confrontés aux explications scalaires\textit{. Travaux de linguistique} 54. 83-96.
\end{styleStandard}

\begin{styleStandard}
Napoli, Donna Jo \& Marina Nespor. 1986. Comparative structures in Italian. \textit{Language}, 62(3). 622-653. \href{http://doi.org/10.1353/lan.1986.0060}{\textstyleInternetlink{doi:10.1353/lan.1986.0060}}.
\end{styleStandard}

\begin{styleStandard}
Octavio de Toledo, Álvaro \& Javier Rodríguez Molina. 2017. La imprescindible distinción entre texto y testimonio: el CORDE y los criterios de fiabilidad lingüística. \textit{Scriptum digital.} 6. 5-68. https://raco.cat/index.php/scriptumdigital/article/view/329258
\end{styleStandard}

\begin{styleStandard}
Paillard, Denis. 1997. N’importe qui, n’importe quoi, n’importe quel N. \textit{Langue française} 116. 100-114.
\end{styleStandard}

\begin{styleStandard}
Palomo, José R. 1934. The Relative Combined with Querer in Old Spanish. \textit{Hispanic Review} 2(1). 51-64. 
\end{styleStandard}

\begin{styleStandard}
Palomo, José R. 1936. ‘Siquiere’ y sus variantes. \textit{Hispanic Review} 4(1). 66-68.
\end{styleStandard}

\begin{styleStandard}
Pato, Enrique. 2012. \textit{Qual manera quier}: la «interposición» en los indefinidos compuestos del español medieval. \textit{Revista de filología española} 92(2). 273-310.
\end{styleStandard}

\begin{styleStandard}
Penny, Ralph. 1993. \textit{Gramática histórica del español}. Barcelona: Ariel
\end{styleStandard}

\begin{styleStandard}
Pescarini, Sandrine. 2009. \textit{Analyse synchronique et diachronique de l'item à choix libre n'importe quel: comparaison avec ‘tout’}. PhD diss, Université de Nancy.
\end{styleStandard}

\begin{styleLiteraturStandard}
Potts, Christopher. 2005. \textit{The logic of conventional implicatures}. Oxford: Oxford University Press.
\end{styleLiteraturStandard}

\begin{styleStandard}
Sánchez López, Cristina. 1999. Los cuantificadores: Clases de cuantificadores y estructuras cuantificativas. In Ignacio Bosque \& Violeta Demonte (eds.), \textit{Gramática descriptiva de la lengua española}. 1025–1128. Madrid: Espasa-Calpe.
\end{styleStandard}

\begin{styleStandard}
Quer, Josep. 1998. \textit{Mood at the Interface}. The Hague: Holland Academic Graphics.
\end{styleStandard}

\begin{styleStandard}
Quer, Josep. 2000. Licensing free choice items in hostile environments: The role of aspect and mood. \textit{SKY Journal of Linguistics}, 13. 251-268.
\end{styleStandard}

\begin{styleStandard}
Radford, Andrew. 2004. \textit{Minimalist Syntax: Exploring the Structure of English}. Cambridge: Cambridge University Press.
\end{styleStandard}

\begin{styleStandard}
Rainer, Franz. 1993. \textit{Spanische Wortbildungslehre.} Tübingen: Niemeyer.
\end{styleStandard}

\begin{styleStandard}
Rivero, María Luisa. 1988. La sintaxis de \textit{qual} \textit{quiere} y sus variantes en el español antiguo. \textit{Nueva Revista de Filología} \textit{Hispánic}a 36(1). 47-73.
\end{styleStandard}

\begin{styleStandard}
Rivero, María Luisa. 2011. Cualquiera posnominal: Un desconocido cualquiera. \textit{Cuadernos de la Asociacicion de Linguistica y Filología de la America Latina} 3. 60–80.
\end{styleStandard}

\begin{styleStandard}
Rizzo Salierno, Melisa. 2013. El uso adjetivo de \textit{cualquiera}. In Ángela L. Di Tullio (ed.), \textit{Aproximaciones al estudio del espa\~{n}ol de la Argentina}. Universidad Nacional del Comahue. 23-31. \url{http://rdi.uncoma.edu.ar/bitstream/handle/123456789/1569/}
\end{styleStandard}

\begin{styleStandard}
\ \ Aproximaciones.pdf?sequence=1\&isAllowed=y
\end{styleStandard}

\begin{styleStandard}
Rohlfs, Gerhard. 1966-1968. \textit{Grammatica storica della lingua italiana e dei suoi dialetti }(3 vols.). Torino: Einaudi. 
\end{styleStandard}

\begin{styleStandard}
Rooth, Mats Edward. 1985. Association with focus (Montague Grammar, semantics, only, even). University of Massachusetts, Amherst dissertation.~
\end{styleStandard}

\begin{styleStandard}
\ \ https://scholarworks.umass.edu/dissertations/AAI8509599.
\end{styleStandard}

\begin{styleStandard}
Sáez, Luis A. 1992. Comparison and Coordination. In Joseba A. Lakarra \& Jon Ortiz de Urbina (eds.), \textit{Syntactic Theory and Basque Syntax}, 387-415. San Sebastian: Universidad del País Vasco. https://ojs.ehu.eus/index.php/ASJU/article/view/9408.
\end{styleStandard}

\begin{styleStandard}
Sanchez Lopez, Cristina. 1999. Los cuantificadores: clases de cuantificadores y estructuras cuantificativas. In Ignacio Bosque \& Violeta Demonte (eds.), \textit{Gramática descriptiva de la lengua española}, 3 vols. 1025-1128. Madrid: Espasa Calpe.
\end{styleStandard}

\begin{styleStandard}
\emph{\textmd{Seco}}\textstyleacopre{, Manuel. }\emph{\textmd{1986}}\textstyleacopre{. \textit{Diccionario de dudas y dificultades de la lengua española}, 9 edn. Madrid: Espasa Calpe.}
\end{styleStandard}

\begin{styleStandard}
Šimík, Radek. 2008. \textit{On free-choice-like wh-based expressions in Czech}. MS, University of Groningen.
\end{styleStandard}

\begin{styleStandard}
Šimik, Radek. 2018. Free relatives. In Daniel Gutzmann, Lisa Matthewson, Cécile Meier, Hotze Rullmann, \& Thomas Ede Zimmermann (eds.), \textit{The Wiley Blackwell Companion to Semantics}. Oxford: Wiley. 
\end{styleStandard}

\begin{styleStandard}
\ \ https://onlinelibrary.wiley.com/doi/10.1002/9781118788516.sem093
\end{styleStandard}

\begin{styleannotationtext}
Sitaridou, Ioanna. 2011. Word order and information structure in Old \textstylehighlight{Spanish}. \textit{Catalan Journal of Linguistics} 10. 159–184.
\end{styleannotationtext}

\begin{styleStandard}
Stalnaker, Robert. 1984. \textit{Inquiry}. Cambridge: MIT Press.
\end{styleStandard}

\begin{styleStandard}
Stark, Elisabeth. 2006. \textit{Indefinitheit und Textkohärenz: Entstehung und semantische Strukturierung indefiniter Nominaldetermination im Altitalienischen}. [=Beihefte zur Zeitschrift für romanische Philologie, 336]. Berlin: De Gruyter.
\end{styleStandard}

\begin{styleStandard}
Stephenson, Tamina. 2007. Indicative conditionals have relative truth conditions. \textit{Proceedings from the Annual Meeting of the Chicago Linguistic Society} 43(1). 231–242.
\end{styleStandard}

\begin{styleStandard}
Stockwell, Robert P., Hall-Partee, Barbara \& Schacter, Paul. 1973. \textit{The major syntactic structures of English.} New York: Holt, Rinehart and Winston.
\end{styleStandard}

\begin{styleStandard}
Tobler, Adolf \& Lommatzsch, Erhard. 1925–2002. \textit{Altfranzösisches Wörterbuch}. 12 vols. Stuttgart: Franz Steiner Verlag.
\end{styleStandard}

\begin{styleStandard}
Traugott, Elizabeth Closs. 1989. On the Rise of Epistemic Meanings in English: An Example of Subjectification in Semantic Change. \textit{Language} 65 (1). 31-55. 
\end{styleStandard}

\begin{styleStandard}
Traugott, Elizabeth C. \& Richard B. Dasher. 2002. \textit{Regularities in Semantic Change}. \ Cambridge: Cambridge University Press. 
\end{styleStandard}

\begin{styleannotationtext}
\textstyleacopre{Trousdale, Graeme. 2014. On the }\emph{\textmd{relationship between grammaticalization}}\textstyleacopre{\textit{ }and\textit{ }}\emph{\textmd{constructionalization}}\textstyleacopre{. \textit{Folia Linguistica} 48 (2). 557–578. \ \ \ \ \ \ \ \ \ \ \ \ \ \ \ \ \ \ \ \ \ \ \ \ \ \ \ \ \ \ \ \ \ \ \ \ \ \ \ \ \ \ \ \ \ \ \ \ \ \ \ \ \ \ \ \ \ \ \ \ \ \ \ \ \ \ \ \ \ \ \ \ \ \ \ \ \ \ \ \ \ \ \ \ \ \ \ \ \ \ \ \ \ \ \ \ \ \ \ \ \ \ \ \ \ \ \ \ }
\end{styleannotationtext}

\begin{styleStandard}
Truswell, Robert \& Nikolas Gisborne. 2015. Quantificational variability and the genesis of English headed wh-relatives. In Eva Csipak \& Hedde Zeijlstra (eds.), \textit{Proceedings of Sinn und Bedeutung }19. 639–656.
\end{styleStandard}

\begin{styleStandard}
Vlachou, Evangelia. 2007. \textit{Free choice in and out of context: semantics and distribution of French, Greek and English Free Choice Items}. Utrecht: LOT. 
\end{styleStandard}

\begin{styleStandard}
Vlachou, Evangelia. 2012. Delimiting the class of free choice items in a comparative perspective: Evidence from the database of French and Greek free choice items. \textit{Lingua}, 122. 1523-1568.
\end{styleStandard}

\begin{styleStandard}
Wedgwood, Hensleigh. 1872. \emph{\textmd{\textit{A Dictionary of English Etymology}}\textmd{. }}\textstyleacopre{London: Trübner \& Co.} 
\end{styleStandard}

\begin{styleannotationtext}
Wolfe, Sam. 2015. The nature of Old \textstylehighlight{Spanish} verb second reconsidered. \textit{Lingua} 10. 595–640.
\end{styleannotationtext}

\begin{styleStandard}
Zamparelli, Roberto. 2000. \textit{Layers in the determiner phrase}. Rochester: University of Rochester dissertation. New York: Garland.
\end{styleStandard}

\begin{styleStandard}
\emph{\textmd{Zimmermann, Thomas Ede. 1993. On the proper treatment of opacity in certain verbs.}}\textstyleacopre{\textit{ }}\emph{\textmd{\textit{Natural Language Semantics}}}\textstyleacopre{\textit{ }2\textit{ }(1). 149-179.}
\end{styleStandard}

\end{document}
